% 本文件由 LaTeX 排版 Agent 根据有效 translation.json 自动生成。
% author=Lactantius; work_id=0701; title=神圣的制度

\chapter{神圣的制度}

% structure: p0001 source=h1 target=omitted reason=page-title-already-rendered-by-parent

\Emph{第一卷} 论对诸神的虚假敬拜\newline{}\Emph{第二卷} 论错误的起源\newline{}\Emph{第三卷} 论哲学家的虚假智慧\newline{}\Emph{第四卷} 论真智慧与宗教\newline{}\Emph{第五卷} 论正义\newline{}\Emph{第六卷} 论真正的敬拜\newline{}\Emph{第七卷} 论幸福生活\par

\SourceRecord{Translated by William Fletcher.；Ante-Nicene Fathers；Vol. 7.；Edited by Alexander Roberts, James Donaldson, and A. Cleveland Coxe.；Buffalo, NY: Christian Literature Publishing Co.,；1886.；<http://www.newadvent.org/fathers/0701.htm>.}

% structure: p0001 source=h1 target=section reason=page-title

\section{神圣诏书，卷一（论假神崇拜）}

% structure: p0002 source=h2 target=subsection reason=relative-heading-rank; base=h2

\subsection{序言——认识真理的价值，从古至今何其重大。}

凡有卓越天赋之人，当全心投入学问时，便蔑视一切公私事务，将所能付出的辛劳全部用于探究真理；他们以为，探究并认识人神之道的法则，远胜于终日忙于聚敛财富或积累荣誉。因为财富与荣誉皆脆弱而属尘世，仅关乎身体的装饰，不能使人变得更善或更义。那些人实在最配得真理的知识——他们如此渴望认识真理，甚至将真理置于万有之上。显然，有人舍弃财产，完全放弃享乐，为的是无牵无挂、不受阻碍地追随那单纯的真理，且惟独追随真理。真理的名声与权威在他们心中如此有力，以至他们宣称至善的赏报即包含于真理之中。然而，他们并未得到所愿的对象，同时白费了辛劳与努力；因为真理——即创造万有的至高天主的奥秘——不能凭我们自身的能力与知觉获得。否则，天主与人之间便无区别，倘若人的思想竟能达到那永恒威严的计谋与安排。既然人凭己力无法知晓神圣的运作方式，天主便不容许人继续在寻求智慧之光中迷失，在无法穿透的黑暗里徒然徘徊；祂最终开启了人的眼睛，将真理的探究作为自己的恩赐，为的是显明人之智慧的虚无，并给在错误中徘徊的人指示获得永生的道路。\par

然而，很少有人善用这天上的恩惠与礼物，因为真理隐藏在幽暗之中，模糊难见；有学问的人因真理缺少合适的辩护者而轻视它，无学问的人则因真理本性的严峻——倾向于恶习的人性无法承受——而憎恨它；因为美德中混合着苦味，而恶习却以快乐调味，人们被前一种所伤，被后一种所慰，便一头栽倒，被美好事物的假象欺骗，以恶为善。——因此，我以为应当迎战这些错误，好让有学问的人被导向真智慧，无学问的人被导向真宗教。这一职业远比修辞学更美好、更有益、更光荣；我们曾长期从事修辞学，训练青年并非为了美德，而是全然为了狡诈的恶行。如今，我们必当更合宜地讨论天上的训诫，借以教导人的心灵敬拜真威严。凡传授善言之道者，不如教导人虔敬与无罪生活者对人类事业更有功劳；正因如此，哲学家在希腊人中比演说家享有更大的荣耀。因为哲学家被视为正确生活的导师，这远为卓越：善言仅属少数人，而善生属所有人。然而，那虚拟诉讼的实践对我们大有裨益，使我们如今能以更丰富的言辞和能力为真理辩护；因为真理虽可无辞令而受辩护——正如许多人常做的那样——但仍需以清晰雅致的言语加以阐述与讨论，使之既凭自身力量，又借言辞之光彩，更有效地流入人心。\par

% structure: p0005 source=h2 target=subsection reason=relative-heading-rank; base=h2

\subsection{第一章——论宗教与智慧。}

因此，我们着手讨论宗教与神圣之事。若某些最伟大的演说家——犹如他们职业中的老手——在完成辩护的著作后，最终投身于哲学，视之为对劳苦最正当的休息，并在探求那些无法寻得的事物中苦思冥想，以致他们为自己寻求的似乎不是闲暇而是辛劳，且其辛劳远超过先前的工作；那么，我岂不更应如投奔最安全的港湾，投奔那虔敬、真实、神圣的智慧？在此智慧中，一切皆备于言说，悦耳听闻，易于理解，值得担当！若某些精通法律与正义的智者编纂并颁布了民法典籍，用以平息纷争公民的争执与冲突，那么我们以文字遵循神圣诏书，岂不更好、更合宜？我们不再谈论雨水滴落、水流改道或权益主张，而是谈论希望、生命、救恩、永生及天主，为要终结致命的迷信与最可耻的错误。\par

我们现在在您尊贵的名下开始这部作品，啊，伟大的皇帝\Person{君士坦丁}，您是罗马诸王公中首位弃绝谬误、承认并尊崇独一真天主之威严的。因为当那最幸福的一天照耀世界时，最崇高天主将您提升到繁荣的权力顶峰，您以卓越的开端步入那对众人有益且可向往的统治，您恢复了被推翻和夺走的正义，赎回了他人最可耻的行为。为回报此行为，天主将赐予您幸福、德行和长寿，使您即使在年老时也能以年轻时开始的正义治理国家，并将罗马之名的守护传给您子女，正如您从您父亲那里所领受的。因为在世界其他地区仍对义人施暴的恶人，全能者也将以其迟延相应的严厉报应其恶行；因为祂对待虔诚之人是最慈爱的父亲，对待不虔之人则是最正直的审判者。在我要捍卫祂的宗教和神圣崇拜时，我还能向谁呼求，对谁发言，除了那为人类事务恢复了正义和智慧的人呢？\par

因此，让我们抛开那些提出任何确定之事的地上哲学的作者，接近正道；因为如果我认为他们是通往美好生活的充分合适的向导，我本人就会跟随他们，并劝勉别人也跟随他们。但由于他们彼此之间激烈争执，且大多自相矛盾，显然他们的道路绝非笔直；因为他们各自根据自己的意志为自己划出了不同的道路，给寻求真理的人留下了巨大混乱。然而，既然真理从天上启示给我们这些接受了真正宗教奥秘的人，并且我们跟随天主——智慧的导师和真理的向导——我们将毫无性别或年龄区分地召集所有人到天上的牧场。因为对灵魂来说，没有比真理的知识更令人愉快的食物了；为维护和阐释这一真理，我们奉献了七本书，尽管这一主题几乎是无限且难以衡量的劳苦；因此，如果有人想要详细阐述并彻底追溯这些内容，他将有如此丰富的题材，以至于书籍找不到边界，言语没有尽头。但因此我们将简要地综括一切，因为我们即将提出的内容如此清晰明了，以至于更令人惊奇的是，真理在人们看来如此晦暗，尤其是在那些通常被视为有智慧的人眼中，或者因为人们只需受我们训练——即从他们所纠缠的谬误中被召回，走向更好的生活道路。\par

如果如我所希望的，我们将达到这一点，我们将把他们送到那极其丰富和丰沛的学问源泉，他们可从其中大量汲取，平息内心产生的渴望，并熄灭他们的热情。一切对他们而言将是容易的、可成就的和清晰的，只要他们不厌烦以耐心阅读或聆听来领会智慧之训导。因为许多人固执地依附于虚妄的迷信，硬起心肠反对显而易见的真理，与其说他们善待了他们错误坚持的宗教，不如说他们虐待了自己；他们本有直路，却寻求迂回弯曲；他们离开平地，却滑向悬崖；他们离开光明，却盲目而脆弱地躺在黑暗中。我们必须为这些人提供帮助，使他们不与自己争斗，并愿意最终从根深蒂固的谬误中被释放出来。如果他们一旦明白自己为何而生，他们必会如此；因为他们乖僻的原因就是自我无知：如果有人获得了真理的知识并摆脱了这种无知，他将知道他的生活应导向什么目标，以及应如何度过。我这样简要地定义这知识的要旨：不可不凭智慧奉行任何宗教，也不可不凭宗教认可任何智慧。\par

% structure: p0010 source=h2 target=subsection reason=relative-heading-rank; base=h2

\subsection{第二章 论世事之中有天主眷顾}

因此，既已承担起解释真理的职责，我认为不必从那个自然显得首要的探究开始，即是否有一个眷顾万物的天主眷顾，或者万物是偶然被创造或被治理；这一观点由\Person{德谟克利特}引入，并被\Person{伊壁鸠鲁}所确认。但在他们之前，\Person{普罗泰戈拉}做了什么——他对诸神提出疑问；以及后来的\Person{狄亚戈拉斯}——他否定诸神；还有一些其他人，他们不认为诸神存在，除非假定没有眷顾？然而，这些观点受到其他哲学家，尤其是\OriginalTerm{斯多葛学派}{Stoics}的最有力反对，他们教导说，宇宙既无法在没有神圣智慧的情况下被创造，也无法存在，除非由最高智慧治理。但即使\Person{西塞罗}，虽他是学园派的辩护者，也多次详细讨论过治理万物的眷顾，确认了斯多葛派的论据，并自己提出许多新论据；他在自己的所有哲学著作中，尤其是在那些论诸神本性的著作中，都如此做了。\par

的确，要驳斥那些持有乖谬看法之少数人的虚妄，藉着各族和支派的见证（他们在这一点上毫无分歧）并非难事。因为没有人是如此野蛮，且具有如此未开化的禀性，当他举目望天时，虽然他不知道这整个可见宇宙是由哪一位天主的上智安排所治理，却能不从万物的宏大、运动、秩序、恒常、效用、美丽和融合中，认识到确有某种上智安排，并且那以奇妙方式存在的事物，必是由某种更高的智慧所预备。而对于我们，无疑可以十分容易地尽情发挥这一点。但因为这个主题在哲学家当中已被激烈讨论，而那些取消上智安排的人，似乎已被明智而雄辩的人充分解答；又因为在我们所承担的这部著作中的不同地方，需要讨论天主上智安排的智慧，所以我们暂且搁置这一探讨——它与其它问题如此紧密相连，以至于我们若不同时讨论上智安排的主题，似乎便无法讨论任何主题。\par

% structure: p0013 source=h2 target=subsection reason=relative-heading-rank; base=h2

\subsection{宇宙是由一位天主还是多位神明的权能所治理？}

因此，让我们的工作从那个紧随并接续第一个问题的探讨开始：宇宙是由一位天主还是多位神明的权能所治理？凡拥有理智且善于反省的人，没有谁会不明白：是同一存在者创造了万物，并以创造它们的那同一能力治理它们。因为何需多位来维持宇宙的治理？除非我们碰巧以为，如果不止一位，每一位都会拥有较小的能力和力量。而那些主张有许多神明的人，确实造成了这一点；因为那些神明必定是软弱的，因为单独而言，没有其他神明的帮助，他们无法维持如此巨大整体的治理。但是天主，永恒的心智，无疑在每一部分都是圆满、完全和完美的。如果这是真的，那么祂必须是一。因为完全的能力或完美，保持其自身特有的稳固。而那无法从中取走任何东西的，可被视为坚固；那无法增添任何东西的，可被视为完美。\par

谁能怀疑，一位拥有全世界治理权的君王会是最有权能的呢？而且这不无道理，因为凡处存在的万物都属于他，所有来自四面八方的资源都集中在他一人身上。但如果不止一人分割世界的治理权，无疑每人都会拥有较少的权力和能力，因为每个人都必须把自己限制在指定的份额内。同样，如果有不止一位神明，他们的分量也会较轻，因为其他神明也拥有同样的能力。但完美的本性允许在拥有整体者身上比只拥有整体之一小部分者身上存在更大的完美。但天主，如果祂是完美的（正如祂应该是的），那么祂不能不只是一，因为祂是完美的，以致万有都可以在祂内。因此，神明的优越和权能必定更弱，因为每一神明所欠缺的，正是在其他神明内所拥有的；所以神明越多，他们的能力就越小。我为何要提及这至高的能力和神能是全然不可分割的呢？因为凡可分割的东西，也必定难免毁灭。但如果毁灭远离天主，因为祂是不朽坏的、永恒的，那么神能就是不可分割的。因此，天主是一，如果能够容纳如此大能力的不能是别的：然而那些认为有许多神明的人，说他们之间划分了职能；但我们将适时地讨论所有这些事情。同时，我断言这一点，这属于当前的主题。如果他们之间划分了职能，问题就回到同一点：他们中任何一个都无法代替所有。那么，如果别人都闲着，他却不能治理万有，他就不能是完美的。因此，对于宇宙的治理，更需要的是一位之完美卓越，而非许多之不完美能力。但想象如此庞大的规模不能被一位存在者所统治的人，是受了欺骗。因为如果他认为，那有能力创造宇宙的唯一天主，也不能治理祂所创造的，那么他就没有明白天主尊威的伟大和能力。但如果他在心中领会，那神圣的工程是何等浩大——从前它什么也不是，却由天主的权力和智慧从无中造出——这工程只能由一位开始并完成，那么他将明白，那由一位所建立的事物，更容易由一位来治理。\par

有人或许会说，像宇宙这样浩大的工程，不可能由单一个体制造，而必须由许多个体合作。但无论他认为这许多个体有多少、有多大——无论他赋予这许多个体何等伟大、权能、卓越与威严——我将这一切尽归于一位，并宣称这一切在于这一位：在他内，这些属性的量度是人所不能构思，亦无法言表的。既然我们在此主题上，无论是领悟还是言辞皆有所欠缺——因为人的心胸无法容纳如此伟大理解的光明，有死的舌也无法解释如此伟大的事理——我们理当领悟并说出这一事实。我再次看到，另一方面可能有人主张：那许多神明正是我们所以为的唯一天主。但这不可能如此，因为那些神明各自的能力将无法进一步发展，因其他神明的能力相互阻扰。因为要么每位必不能超越其自身的界限，要么倘若其超越了界限，则必会将其余神明逐出其疆界。相信有多神的人，未能看到可能出现的情况：有些神明在意愿上彼此对立，由此在他们之间产生争执与冲突；正如\Person{荷马}所描述的神明们彼此交战，有些希望\Place{特洛伊}被攻陷，有些则反对。因此，宇宙必须由一位的意志所统治。因为除非各部分的大权归属于同一眷顾，整体本身将无法存在；因为每个只关心专属自身之事，就如同没有一位将军与统帅就无法进行战争一般。如果在一支军队中，有与军团、大队、中队、骑兵队一样多的将军，首先，军队无法列阵，因为每人都将拒绝危险；也很难被统治或控制，因为所有人都将运用自己独特的谋略，其多样性所造成的伤害将大于带来的利益。同样，在管理自然事务中，除非有一位万物整体所关心的对象，一切都会解体而陷入衰败。\par

但说宇宙由多位的意志所治理，无异于宣告在同一身体内有许多心灵，因为肢体有众多不同的职分，以致可以设想有各自的心灵分别统辖各自的感官；还有众多情感，我们常被它们所激动，或愤怒，或欲望，或喜乐，或恐惧，或怜悯，以致可以设想所有这些情感中有许多心灵在运作；若有人这样说，他本身似乎连那唯一的心灵也丧失了。但如果在一个身体内，一个心灵统管如此众多的事物，同时又关注整体，为何有人竟认为宇宙不能被一位所统治，而能被多于一位所统治？因为那些主张多神的人知道这一点，他们说这些神明各自管辖不同的职责和部分，但仍有一位至高的统治者。因此，按照这一原则，其他神明将不是神明，而是侍从和臣仆，是那位最伟大、全能的所派定执行这些职责的，并且他们本身将臣服于其权威与命令。因此，如果他们彼此并非平等，他们便不全是神明；因为服侍者与统治者的身份不可能相同。如果天主是至高权能的称谓，那么他必是不朽的、完美的、不能受苦的、不受制于任何其他存在；因此，那些不得不服从于这位唯一至大天主的，就不是神。但持此意见的人并非无故受骗，我们将在稍后揭示这错误的缘由。现在，让我们以见证来证明天主性权力的唯一性。\par

% structure: p0018 source=h2 target=subsection reason=relative-heading-rank; base=h2

\subsection{第四章  先知们也曾预告唯一的天主}

众多先知宣扬并宣告唯一的天主；因为他们充满唯一天主的灵感，以一致和谐的声音预言了未来的事。但那些不认识真理的人认为这些先知不可信；因为他们说那些声音不是神圣的，而是出自人的。诚然，因为他们宣告唯一的天主，他们若不是疯子，便是骗子。但我们看到，他们的预言已经成就，并且每日仍在成就中；他们一致同意的预见教导我们：他们并非出于疯狂。因为被疯狂冲动的人，莫说预言未来，就连连贯地说话也不可能。那么，说这些事的人是骗子吗？欺骗的体系与他们本性何其相违——毕竟他们自己阻止他人行一切欺诈？因为他们被天主派遣，就是要成为祂威严的传报者，并纠正人的邪恶。\par

此外，虚构说谎的倾向属于那些贪图财富、渴求利益的人——这种品性远离了那些圣洁的人。因为他们如此执行托付于他们的职务，以致漠视一切生活所需，不仅不为将来积蓄，甚至不为当日劳作，满足于天主所供给的现时食物；这些人不仅没有收益，反而忍受折磨与死亡。因为正义的规诫令邪恶及度不洁生活者憎恶。故此，那些罪行被揭露并被禁止的人，极其残酷地折磨并杀害了这些先知。因此，那些不贪图利益的人，既无倾向也无动机去行骗。为何我还要说他们中的一些是王子，甚至是君王，贪欲与欺诈的嫌疑完全不可能落在他们身上，而他们却与其他人一样，以同样的先知预见宣告唯一的天主呢？\par

% structure: p0021 source=h2 target=subsection reason=relative-heading-rank; base=h2

\subsection{第五章  诗人和哲学家的见证}

然而，让我们放下先知们的见证，免得引用那些普遍不被相信之人的证据显得不够充分。让我们转向作家们，并为证明真理而引用那些他们惯常用来反对我们的人——即诗人和哲学家——作为证人。于证明天主的独一性，我们不会失败；并非因为他们已确知真理，而是真理本身的力量如此之大，以致无人能盲目到看不见眼前的神性光辉。因此，诗人们虽在其诗作中极力美化诸神，并以最高赞誉夸大其功绩，却仍十分频繁地承认万物由一位神灵或心智统摄和管理。最古老的诗人\OriginalTerm{俄耳甫斯}{Orpheus}，与诸神本身同时代——据记载他曾与\OriginalTerm{廷达瑞俄斯}{Tyndarus}诸子和\OriginalTerm{赫拉克勒斯}{Hercules}一同随\OriginalTerm{阿尔戈英雄}{Argonauts}航行——他论及那位真实而伟大的天主为首生者，因为无物先于祂产生，万物皆由祂而生。他又称祂为\OriginalTerm{法涅斯}{Phanes}，因为当时尚无任何事物时，祂最先显现并从无限中出来。因他无法在心智中理解此存在的起源和本性，便说祂生于无垠之气：\Quote{首生者，法厄同，广袤气之子；}因他再无更多可说。他断言此存在是诸神之父，为祂的缘故造了天，并为祂的儿女预备一个共同的居所：\Quote{祂为不朽者建造了不朽的家园。}于是，在本性与理性引导下，他理解到有一种超越的伟大力量创造了天地。因他不能说\OriginalTerm{朱庇特}{Jupiter}是万有之作者，因其生于\OriginalTerm{萨图尔努斯}{Saturn}；也不能说萨图尔努斯本人是作者，因据传他由天所生；但他不敢将天立为原始之神，因他看出那是宇宙的一个元素，本身也必有一位作者。这一思考引导他认出了那位首生之神，他将其置于首位并赋予首位。\par

\OriginalTerm{荷马}{Homer}未能给我们提供任何关于真理的信息，因他所写的是人事而非神事。\OriginalTerm{赫西俄德}{Hesiod}本可提供，因他将诸神的谱系合为一卷书；但他仍未能给我们提供信息，因他不是从天主创造主开始，而是从混沌开始——那是一团粗糙无序的混乱质料——而他本应先说明混沌本身于何时、以何种方式、从何而来开始存在或具有一致性。无疑，正如万物由某位工匠排列、安排和制作，质料本身也必须由某位存在形成。那么，非天主莫属，万物皆服于其权能下。但他回避承认这一点，因他畏惧未知的真理。因如他所愿，他是受缪斯灵感在赫利孔山上唱出那首歌；但他是在先前的默想和准备之后才来的。\par

我们的诗人中，\OriginalTerm{维吉尔}{Maro}最先接近真理，他这样论及至高的天主，称祂为心智和神灵：——\par

\Quote{须知，天、地、大海，\newline{}月亮苍白之球，星辰之列，\newline{}都由一个灵魂滋养，\newline{}一位神灵，其天火\newline{}在形体的每一部分燃烧，\newline{}激动那伟大的整体。}\par

且为避免有人不知那有如此大能的神灵是什么，他在另一处点明说：\Quote{因为神性遍及所有陆地、海的区域和天的深处；羊群、牛群、人类和一切兽族，于其出生时，都从祂获得微弱的生命。}\par

\OriginalTerm{奥维德}{Ovid}也在他著名作品的开头，毫不掩饰名称地承认宇宙由天主安排，他称祂为世界的塑造者、万物的工匠。然而，若俄耳甫斯或我们国家的这些诗人总是持守他们在本性引导下所察觉的，他们便会理解真理，并获得我们所追随的同一种学问。\par

但是，关于诗人的话就到此为止。现在让我们来看哲学家，他们的权威更为重要，他们的判断也更值得信赖，因为人们相信他们所关注的不是虚构之事，而是对真理的探求。米利都的泰勒斯是七贤之一，据说是最早探究自然原因的人。他说水是生成万物的元素，而天主是那从水中塑造万物的理智。因此，他将万物的质料归之于水汽；他将万物产生的起始和原因归于天主。毕达哥拉斯如此定义天主的本质：\Quote{如同一个灵魂，往来运行，弥漫于宇宙的所有部分和整个自然界，一切被造的生物都从它获得生命。}阿那克萨戈拉说，天主是无限的理智，凭自身的能力运行。安提西尼主张，人们所信奉的神有很多，但有一位自然界的天主是唯一的那位，即整个宇宙的造物主。克莱安西斯和阿那克西美尼断言，气是首要的神祇；我们的诗人也赞同这一观点：\Quote{然后全能的天父\OriginalTerm{以太}{Æther}降下丰沛的雨水，进入他喜悦配偶的怀抱；他自身伟大，与她伟大的身体结合，滋养她所有的后代。}克里西波斯称天主为赋有神圣理智的自然力量，有时又称之为神圣必然性。芝诺也称祂为神圣的自然律。所有这一切观点，尽管不确定，都指向一点——他们都同意有一个天主眷顾存在。因为无论它是自然、以太、理性、理智、命运必然性、神圣律法，还是你称之为其他任何东西，都是我们所称的同一的天主。名称的多样性并不构成障碍，因为就其本身含义而言，它们都指向同一个对象。亚里士多德虽然自相矛盾，说出并持有相互对立的见解，但总体上仍证实有一位理智主宰着宇宙。柏拉图，被公认为最智慧的人，明白而公开地主张一位天主的统治；他不称祂为以太、理性或自然，而是按祂的真实所是称为天主，并且这如此完美奇妙的宇宙是由祂所构造。西塞罗在许多方面追随并模仿他，在他所写的关于法律的书中，经常承认天主，并称祂为至高的；当他以这种方式论及诸神的本质时，他提出证据证明宇宙由祂治理：\Quote{没有比天主更高的事物：因此世界必须由祂治理。因此天主不受任何自然界的服从或支配；因此祂亲自治理整个自然界。}但他在他的\WorkTitle{慰藉}中定义了天主本身是什么：\Quote{天主本身，按我们所能领悟的，也只能被领悟为一个自由而不受束缚的理智，远离一切必朽的物质性，感知并推动万物。}\par

然而，罗马人中最敏锐的斯多葛派塞内加，又是多么频繁地以应得的赞美来追随至高的天主啊！因为当他讨论夭折的主题时，他说：\Quote{你不明白你的审判者的权威和威严，祂是世界的统治者、天上和众神的天主，那些我们分别崇拜和尊崇的神祇都依赖祂。}同样在他的\WorkTitle{劝诫}中：\Quote{这存在者，当祂为那最美丽的构造打下最初的基础，并开始这项自然所知道的无以复加的作品时，为使万有服事各自的统治者，虽然祂已将自己扩展至整个身体，祂仍将诸神造为祂王国的仆役。}他在论及天主的话题上，说了多少与我们自己的作者类似的话啊！但这些事我暂且搁置，因为它们更适合主题的其他部分。目前，足以证明那些最具天赋的人触及了真理，并几乎把握了它，若不是习俗——被错误见解所迷惑——将他们拉回；由于这习俗，他们既认为有其他神祇存在，又相信天主为人使用所造的那些事物，好像它们赋有感知，竟当作神祇来持有和崇拜。\par

% structure: p0030 source=h2 target=subsection reason=relative-heading-rank; base=h2

\subsection{第六章 论天主的见证，以及西比尔和她们的预言。}

现在让我们转向神圣的见证；但我要先提出一个类似神圣见证的见证，一则因其极为古老，二则因我将要提及的那人曾被从人间取去而置于众神之中。据西塞罗记载，大祭司卡尤斯·科塔在与斯多葛派就迷信及关于众神的各种流行意见进行辩论时，为遵循学园派的惯例使一切成为不确定，他说有五位墨丘利；在依次列举了四位之后，他说第五位就是杀死阿尔戈斯的那位，并因此逃往埃及，将法律和文字赐予埃及人。埃及人称他为透特；埃及人一年的第一个月，即九月，便从他得名。他还建造了一座城镇，希腊人至今称之为赫尔墨波利斯（墨丘利之城），费奈的居民以宗教敬礼尊崇他。他虽是人，却极为古老，且渊博地通晓各类学问，以致对许多学科和技艺的知识为他赢得了\OriginalTerm{三倍伟大的}{Trismegistus}之名。他写了许多关于认识神圣事物的书籍，在其中他宣示了至高唯一天主的威严，并用我们所使用的同样名称提及祂——天主和圣父。为使无人探询祂的名字，他说祂是无名的，且因其唯一性，祂不需要名字的专有性。以下是他的原话：\Quote{天主是唯一的，但唯一者不需要名字；因为自有者是无名的。}因此，天主没有名字，因为祂是唯一的；除非在需要区分众人的情况下，才需要专有名字，以便你可以用各自的标记和称呼来指称每个人。但天主，因祂永远是一，没有专有的名字。\par

我还需提出关于神圣神谕和预言的见证，这些更为可靠。因为我们所反驳的人或许会认为，不应相信诗人，仿佛他们编造了虚构，也不应相信哲学家，因为他们作为人也会犯错。马尔库斯·瓦罗，无论在希腊人中还是拉丁人中，都无人比他更有学问，在他写给大祭司卡尤斯·凯撒的关于神圣事物的书中，当谈论到十五人团时，他说西彼拉神谕集并非仅一位西彼拉的作品，而是被统称为西彼拉神谕集，因为古代的预言女先知都被称为西彼拉，或来自德尔斐女祭司的名字，或来自她们宣告众神的旨意。因为在艾欧利亚方言中，他们用\Emph{Sioi}一词称呼众神，而非\Emph{Theoi}；对于旨意，他们用\Emph{bule}一词，而非\Emph{boule}；——于是西彼拉之名便得自\Emph{Siobule}。但他说西彼拉共有十位，他逐一列举，并注明记述每位西彼拉的作家：第一位来自波斯，尼卡诺尔曾提及她，他撰写了马其顿的亚历山大的事迹；——第二位来自利比亚，欧里庇得斯在\WorkTitle{拉弥亚}的序言中提及她；——第三位来自德尔斐，克吕西波在他所著的\WorkTitle{论占卜}中提及她；——第四位是意大利的基梅里亚人，奈维乌斯在他的\WorkTitle{布匿战争}书中以及皮索在他的编年史中提及她；——第五位来自埃律特赖亚，埃律特赖亚的阿波罗多洛斯断言她是他的同乡，并预言了当希腊人出发前往伊利昂时，特洛伊注定毁灭，而荷马将写谎言；——第六位来自萨摩斯，埃拉托色尼写道，他在萨摩斯人的古老编年史中发现了关于她的书面记录。第七位来自库迈，名叫阿玛尔忒亚，有些人称她为赫罗菲勒或德莫菲勒，据说她将九卷书带给国王塔奎尼乌斯·普里斯库斯，并索要三百腓力金币，国王拒绝如此高价，嘲笑这女人的疯狂；她当着国王的面烧了三卷，仍为所剩的书索要同样的价格；塔奎尼乌斯更认为这女人疯了；当她再次烧了三卷后仍坚持索要同样的价格时，国王被打动，用三百金币买下了剩余的书卷。后来在重建卡皮托利山之后，这些书卷的数量增加了；因为它们被从意大利和希腊的所有城市，尤其是从埃律特赖亚收集而来，以所归属的西彼拉的名义被带到罗马。此外，第八位来自赫勒斯滂，出生在特洛伊领土上的马尔佩苏斯村，靠近格尔吉图斯镇；本都的赫拉克利德斯写道，她生活在梭伦和居鲁士的时代；——第九位来自弗里吉亚，在安基拉发出神谕；——第十位来自提布尔，名叫阿尔布内亚，她在提布尔被当作女神崇拜，靠近阿尼奥河畔，据说她的雕像在河深处被发现，手中拿着一本书。元老院将她的神谕移至卡皮托利山。\par

所有这些巫女的预言都被提出并受到如此珍视，只有库迈巫女的预言除外，罗马人将其书卷隐藏起来，他们认为任何除了\Emph{十五人团}之外的人检查这些书卷都是不合法的。而且每个巫女都有各自独立成卷的作品，但由于这些作品都署有巫女之名，它们被认为出自同一人之手；它们被混淆了，每个作品都不能被区分并归于其作者，只有埃律特莱巫女例外，因为她既在诗句中插入了她自己真实的名字，又预言了她将被称作埃律特莱人，尽管她出生于巴比伦。但我们在需要使用她们的见证时，也将不加区分地谈论巫女。所有这些巫女都宣告一位天主，尤其是埃律特莱巫女，她在其他巫女中被视为更著名和高贵；因为最勤勉的作家费内斯特拉在谈到\Emph{十五人团}时说，在重建卡皮托利之后，执政官盖乌斯·库里奥向元老院提议派遣使节前往埃律特莱，搜寻并带回罗马的巫女著作；于是，便派遣了普布利乌斯·加比尼乌斯、马尔库斯·奥塔基利乌斯和卢基乌斯·瓦莱里乌斯，他们带回罗马大约一千行由私人抄写的诗句。我们之前已指出瓦罗也说了同样的话。现在，在这些使节带回罗马的诗句中，有关于唯一天主的这些见证：—\par

1. \Quote{独一的天主，祂是独自的、最全能的、非受造的。}\par

这是唯一至高的天主，祂创造了天，并以众光装饰它。\par

2. \Quote{但有一位独一无比尊能的天主，祂创造了天、太阳、星辰、月亮、肥沃的大地、以及海水的波涛。}\par

既然祂独一是宇宙的缔造者，是构成宇宙或包含于其中的万物的工匠，那么它就见证祂唯独应受敬拜：—\par

3. \Quote{敬拜那独一是世界统治者的那位，祂独自从永远到永远存在。}\par

另一位巫女，无论她是谁，当她声称她将天主的声音传递给人类时，如此说：—\par

4. \Quote{我是独一的天主，除我以外没有别的天主。}\par

我现在本应接着引用其他巫女的见证，只是这些已足够，且我要为更适宜的时机保留其他见证。但既然我们在那些偏离真理、事奉虚假宗教的人面前辩护真理的事业，我们应当提出什么样的证据来反驳他们，而非用他们自己神祇的见证来驳斥他们呢？\par

% structure: p0042 source=h2 target=subsection reason=relative-heading-rank; base=h2

\subsection{第七章  论阿波罗与众神的见证。}

的确，阿波罗——他们以为他超越众神，尤其有预言能力，在科洛丰发出神谕（我想是因为亚洲的宜人使他从德尔斐迁离）——当有人问他祂是谁，或者究竟什么是一位天主时，他以二十一行的诗句回答，其开头如下：—\par

\Quote{自生、无师、无母、不动摇，\newline{}一个甚至不能用言语穷尽的名字，居于火中，\newline{}这就是天主；而我们，祂的使者，是天主的一小部分。}\par

有谁会怀疑这是关于朱庇特说的呢？朱庇特既有母亲，又有名字。我为何要说墨丘利——那至大的三次方者，我上文已提及过——不仅像阿波罗一样说天主是\Quote{无母的}，而且也是\Quote{无父的}，因为祂无任何其他来源，只从自己而出？因为祂不能由任何人产生，而祂自己产生了万物。我想，我已通过论证充分教导，并通过见证确认了那本身足够明显的事实：即宇宙只有一位君王、一位圣父、一位天主。\par

但或许有人会问我们西塞罗书中霍尔田西乌斯所问的同一个问题：如果天主是独一的，那么怎样的独处才是幸福的？仿佛我们主张祂是独一的，就是说祂是荒凉和孤独的。毫无疑问，祂有臣仆，我们称之为使者。这是真实的，我之前说过塞内卡在他的\WorkTitle{劝勉}中说过：天主产生了祂王国的臣仆。但这些既不是神，也不愿被称为神或受敬拜，因为他们只执行天主的命令和旨意。然而，那些被普遍敬拜的神也并非神，其数目是少而固定的。但如果崇拜众神的人认为他们敬拜的是我们所称为至高天主的臣仆的那些存在，那么他们没有理由嫉妒我们这些说只有一位天主、否认有许多神的人。如果众多神祇令他们喜悦，我们说的不是十二位，也不是像俄耳甫斯所说的三百六十五位；但我们却从另一方面指出他们因认为神如此之少而犯了无数错误。不过，让他们知道他们应该被称为什么名字，免得他们冒犯真天主，因为他们把祂的名字赋予不止一位，却显扬了祂的名字。让他们相信自己的阿波罗吧，他在那同一答复中，正如他从朱庇特那里夺去统治权一样，也从其他神那里夺去了他们的名字。因为第三节显示，天主的臣仆不该被称为神，而该被称为天使。他确实关于自己说了谎话；因为他虽属于魔鬼之列，却将自己算作天主的使者，然后在其他答复中他又承认自己是一个魔鬼。因为当他被问及他希望如何被祈求时，他这样回答：—\par

\Quote{全智、全学、精通技艺者，请听，啊，魔鬼。}\par

同样，当应某人之请，他发出诅咒反对斯敏提阿的阿波罗时，他以这样一节开始：—\par

\Quote{啊，世界的和谐，携带光明，全智的魔鬼。}\par

那么，除了他亲口承认自己处于真天主的鞭挞和永罚之下，还有什么可说的呢？因为在另一答复中他也说：—\par

\Quote{那些无倦地周游大地和海上的魔鬼，\newline{}都屈服于天主的鞭挞之下。}\par

我们在第二卷中谈论两者的主题。在此期间，对我们来说足够了的是：当他想要荣誉并把自己置于天上时，他如事情的本质所是地承认了那些永远站在天主身边者应当如何被称呼。\par

因此，愿人远离谬误，抛弃腐化的迷信，承认他们的圣父与主——祂的卓越无可估量，祂的伟大无从感知，祂的始源无法理解。当人灵的热切关注及其锐利的敏锐与记忆达至祂时，所有路径仿佛已穷尽并耗尽，于是它停顿、迷惘、失效；没有任何超越之处可供继续前行。但因存在者必然曾有起始，既然在祂之前无物存在，所以祂是在万有之前自祂自身产生的。因此，\Person{阿波罗}称祂为\Quote{自生者}，\Person{西彼尔}称祂为\Quote{自造者}、\Quote{非受造者}与\Quote{非制作而成者}。而敏锐的\Person{塞内卡}在其\WorkTitle{劝诫}中看到并表达了这一点。他说：\Quote{我们依存于他人。}因此我们仰望某位我们亏欠那在我们内最卓越之物者。他人带我们进入存在，他人塑造了我们；但天主凭其自身的大能创造了自己。\par

% structure: p0054 source=h2 target=subsection reason=relative-heading-rank; base=h2

\subsection{第八章  论天主没有形体，并且为生育不需性别差异}

因此，由这些如此众多且具权威的见证，证明了宇宙是由唯一天主的权能与眷顾所治理；\Person{柏拉图}在\WorkTitle{蒂迈欧篇}中断言祂的能量与威严如此巨大，以致无人能在心中构想，或言语表达，因祂那超越且不可计量的能力。那么，有人能怀疑对天主而言有任何艰难或不可能之事吗？祂凭其眷顾设计，凭其能量建立，凭其判断完成那些如此伟大奇妙的工程，甚至现在仍以祂的神维系它们，以祂的权能治理它们——祂是不可理解的、不可言说的，且除祂自身外无人能完全认识。因此，当我频繁思考这如此伟大威严的主题时，那些敬拜诸神者有时显得如此盲目、如此缺乏反省、如此麻木、如此与哑默动物无几，竟相信那些从自然两性交合出生者能拥有任何威严与神性的影响；因为\Person{厄立特里亚的西彼尔}说：\Quote{神不可能从男人的腰和\Emph{女人的子宫}被塑造出来。}如果这是真的，正如它实际是的，那么显然赫丘利、阿波罗、巴库斯、墨丘利和朱庇特以及其余的都不过是人，因为他们是从两性出生的。但有什么如此远离神的本质，如同那祂自身为延续人类种族而赋予必朽者的行为，且没有实体物质就无法进行呢？\par

因此，如果诸神是不死不灭的，那另一个性别的需要有什么意义呢？他们自身并不需要延续，因为他们永远存在。因为在人类和其他动物中，性别差异与生育及产子的唯一理由，无非是所有种类的生物，因受其必朽性条件注定死亡，可通过相互延续而得以保存。但天主，祂是不死不灭的，无需性别差异，也无需延续。有人或说：这样的安排是必要的，为使他们能有服侍祂者，或受祂管辖者。既然天主是全能的，无需女性中介就能生出儿子，那么女性性别有何必要呢？因为若祂已将能力赐予某些微小生物，使它们\par

\Quote{能以口从叶子和甜草中为自己聚集后代，}\par

为何有人会认为天主自身除非与另一性别结合，否则不能有后代呢？因此，无人如此轻率而不明白，那些被愚昧无知者视为并敬拜为神的，不过仅仅是凡人。那么，有人会说：为何他们曾被相信是神？无疑，因为他们曾是极伟大且有权能的君王；并且因他们德性、功绩或所发现技艺的功绩，被曾受其治理者所爱戴，从而被奉献于永久的纪念。若有人对此怀疑，让他思量他们的功业与事迹，古代诗人和史学家已将这一切传述下来。\par

% structure: p0059 source=h2 target=subsection reason=relative-heading-rank; base=h2

\subsection{第九章  论\Person{赫丘利}及其生与死}

\Person{赫丘利}——那个以其勇猛最为著名，并被视作诸神中的阿非利加努斯者——岂非以其放荡、情欲和奸淫玷污了那据称他曾遍行并净化过的世界吗？这不足为奇，因他是由与\Person{阿尔克墨涅}的奸淫之交所生。\par

在他身上能有什么神性呢？他竟被自己的恶习奴役，违反一切法律，以耻辱、不名誉和暴行对待男女双方。的确，他所做的那些伟大而奇妙之事，不应被视为值得归因于神圣卓越的事。因为！征服一头狮子和一头野猪，以箭射落飞鸟，清扫王室的马厩，战胜一名悍妇并夺其腰带，连同其主人杀死野马，这有何等了不起呢？这些乃有勇气且英雄之人——但仍然是人之事；因为他所战胜的那些是可朽坏的。因为没有什么力量如此强大，正如那位演说家所言，是不能被铁和强力削弱和击破的。但征服心神、抑制愤怒，乃最勇敢之人的本分；而这些他从未做过也不能做到：因为做这些事的人，我不与卓越之人相比，反判断他最似神明。\par

我希望他还曾就情欲、奢侈、贪欲和傲慢添加些什么，以便完善他所判断为与神相似之人的优越。因为不要认为战胜狮子的人比战胜那囚于自身之内的狂暴野兽（即愤怒）更为勇敢；也不要认为击落最贪婪飞禽的人比抑制最贪婪欲望的人更强；更不要认为征服好战的阿玛宗人胜过征服情欲（那摧毁贞洁与名声的征服者）；或者清除马厩粪便的人比清除内心恶习的人更了不起——后者是更为毁灭性的邪恶，因为它们是自身特有的，而那些（外在的）则是可以避免和防范的。由此而来，唯有节制、适度、公正的人才应被判定为勇敢的人。但若有人思考天主的作为，他会立刻判断所有那些最琐碎之人都羡慕的东西为可笑。因为他们不是以他们所不认识的\OriginalTerm{神圣能力}{divine power}来衡量，而是以他们自身力量的软弱来衡量。因为没有人会否认这一点：即\Person{赫丘利}不仅是国王\Person{欧律斯透斯}的仆人——这在某种程度上可能显得光荣——而且还是不贞妇人\Person{翁法勒}的仆人，她曾命令他穿着她的衣服坐在她脚边，执行指定的任务。可耻的卑劣！但这就是享乐被估价的价格。怎么！有人会说，你认为诗人是可信的吗？为什么我不该这样认为？因为叙述这些事的并非\Person{卢基里乌斯}，也非\Person{卢齐安}（他对人还是神都不宽容），而是那些尤其赞美神的美德的人。\par

那么，如果我们不相信那些赞美他们的人，我们该相信谁呢？让那些认为这些人是说谎者的人提供其他我们可以信赖的作者，来教我们这些神是谁，他们如何以及源于何处，他们的力量是什么，他们的数目、权能，在他们之中有什么值得钦佩和崇拜的——简而言之，有什么更可靠、更真实的奥秘。他不会提供这样的权威。那么，让我们相信那些不是为了指责而是为了宣扬赞美而说话的人吧。于是，他与阿尔戈英雄们航海，并洗劫了\Place{特洛伊}，因\Person{拉俄墨冬}拒绝付给他（为拯救其女儿而应得的）报酬而愤怒，由此可知他生活的年代。他还因愤怒和疯狂杀死了他的妻子和子女。这就是人们认为是神的那位吗？但他的继承人\Person{菲罗克忒忒斯}却不这么认为，他对他即将被焚烧时举起了火把，目睹了他的四肢和筋肉的燃烧与消耗，并将他的骨灰葬在\Place{厄塔山}，作为回报，他得到了他的箭。\par

% structure: p0064 source=h2 target=subsection reason=relative-heading-rank; base=h2

\subsection{第十章：关于\Person{阿斯克勒庇俄斯}、\Person{阿波罗}、\Person{尼普顿}、\Person{马尔斯}、\Person{卡斯托尔}和\Person{波鲁克斯}、\Person{墨丘利}和\Person{巴克斯}的生平与事迹。}

\Person{阿斯克勒庇俄斯}除了治愈\Person{希波吕托斯}之外，还行了什么值得神圣尊荣的事？他的出身对\Person{阿波罗}来说也并非没有耻辱。他的死无疑更为著名，因为他赢得了被神用雷电击中的殊荣。\Person{塔尔奎提乌斯}在一篇关于著名人物的论文中说，他出身父母不详，被遗弃，被一些猎人发现；由一只狗哺养，后被交给\Person{喀戎}，学习了医术。他还说他是麦西尼亚人，但曾住在\Place{厄庇道洛斯}。\Person{图利乌斯}也说，他葬在\Place{库诺苏拉}。他的父亲\Person{阿波罗}做了什么？难道他不是因为炽热的爱情，极其可耻地放牧他人的羊群，并为\Person{拉俄墨冬}建造城墙——他与\Person{尼普顿}一起受雇却可被有罪不罚地拒绝报酬？正是从他开始，这个背信弃义的国王学会了拒绝履行他与神订立的任何契约。同时，他在爱上一个美少年时，竟施暴于他，并在嬉戏中杀死了他。\par

\Person{马尔斯}犯了杀人罪，经雅典人出于偏袒而免于谋杀指控，唯恐他显得过于凶猛和野蛮，又与\Person{维纳斯}通奸。\Person{卡斯托尔}和\Person{波鲁克斯}在抢夺他人妻子时，不再是一对孪生兄弟。因为\Person{伊达斯}出于对伤害的嫉妒，用剑刺穿了\Emph{兄弟之一}。诗人们叙述他们交替地生与死：所以他们不仅是最不幸的神，也是所有凡人中最不幸的，因为他们不允许只死一次。然而\Person{荷马}与其他诗人不同，他只记录他们都死了。因为当他在\Place{特洛伊}城墙上让\Person{海伦}坐在\Person{普里阿摩斯}旁边，认出所有希腊首领，却只找不到她的兄弟时，他在演讲中加了这样一行诗：——\par

\Quote{她这样说着，却不知在\Place{斯巴达}，\newline{}他们的故土，他们已躺在黄土之下。}\par

\Person{墨丘利}，这个窃贼和挥霍无度者，除了他的欺诈的记忆之外，还留下了什么有助于他的名声？无疑他配得上天堂，因为他教授了角力场的训练，并且是第一个发明七弦琴的人。\Person{父神利柏}理应拥有最高权威，并在神明的元老院中居首位，因为他是他们中惟一（除了\Person{朱庇特}之外）得胜、率领军队并征服印度人的。但是那位非常伟大且不可战胜的印度统帅却被爱情和欲望极其可耻地征服了。因为他带着他柔弱的随从来到\Place{克里特}，在海滩上遇到一个不贞的女人；由于印度胜利的自信，他想要证明他的男子气概，以免显得过于柔弱。于是他要了那个出卖父亲、杀害兄弟、被另一个丈夫抛弃和休弃的女人为妻，他封她为\Person{利柏拉}，并与她一同升入天堂。\par

这一切神祇之父\Person{朱庇特}的所作所为如何？在习惯的祈祷中，他被称为\Quote{至善至大}。难道不是从他童年早期起，就被证明为不敬，几乎是弑亲者，因为他将父亲逐出王国，流放他，甚至不等他年老力竭而死——他如此渴望统治？当他以暴力和武力夺取父亲的王位后，遭到泰坦们的战争攻击，这成为人类祸患的开端；待战胜他们并获致持久和平后，他余生便沉溺于淫乱和奸情。我不愿提那些被他玷污的贞女，因为那通常被认为尚可容忍。但我不能忽略\Person{安菲特律翁}和\Person{廷达瑞俄斯}的案例，他将他们的家充满了耻辱和恶名。但他达到了不敬和罪愆的顶峰，在于劫持那位王室少年。因为似乎还不够，在强暴妇女后使自己蒙羞，还要玷污自己的性别。这才是真正的奸淫，是违反本性的。犯下这些罪行的人能否被称为\Quote{至大}，是个问题；无疑他不是\Quote{至善}；腐败者、奸淫者和乱伦者无权享有此名；除非我们人类犯了错误，将行如此之事者称为邪恶和堕落，并判断他们最应受各种惩罚。但\Person{马尔库斯·图利乌斯}指责\Person{盖乌斯·维勒斯}犯奸淫是愚蠢的，因为他所崇拜的\Person{朱庇特}也犯了同样的罪；他指责\Person{普布利乌斯·克劳狄乌斯}与妹妹乱伦，而那位\Quote{至善至大}者自己的妹妹也正是他的妻子。\newline{}\par

% structure: p0070 source=h2 target=subsection reason=relative-heading-rank; base=h2

\subsection{第十一章 论\Person{朱庇特}以及\Person{萨图尔努斯}和\Person{乌拉诺斯}的起源、生平、统治、名号和死亡}

那么，谁如此无知，竟想象那甚至不应在地上为王的人在天上统治？一位诗人描写\Person{丘比特}的凯旋时并非没有幽默：在那本书中，他不仅将\Person{丘比特}描绘为众神中最有力量的，而且描绘为他们的征服者。因为，数算每位神祇的恋情——他们正是通过这些恋情进入\Person{丘比特}的权力和统治之下——他布置了一场游行，其中\Person{朱庇特}连同其他神祇，被锁链牵着走在庆祝凯旋的战车前。诗人优美地描绘了这一幕，但它离真相并不远。因为那没有德行、被欲望和邪恶情欲所制伏的人，并非如诗人所虚构的那样臣服于\Person{丘比特}，而是臣服于永死。但且让我们停止谈论道德；让我们审视事实，为使人明白他们可怜地陷于何种错误。平民百姓想象\Person{朱庇特}在天上统治；学者和文盲都同样被说服。因为宗教本身、祈祷、颂歌、神殿和偶像都证明了这一点。然而他们也承认他是由\Person{萨图尔努斯}和\Person{瑞亚}所生。他如何能显现为神，或如诗人所说，被相信为人和万物的创造者，而在他出生之前已有无数千万的人存在——比如那些在\Person{萨图尔努斯}统治时期生活、比\Person{朱庇特}更早享受阳光的人？我看到一个神在最早期为王，另一个在后来的时代为王。因此，以后可能还有另一个。因为如果先前的王国被改变，为何我们不期望后来的也可能被改变呢？除非是\Person{萨图尔努斯}可能产生一个比自己更有力量者，而\Person{朱庇特}却不可能？然而神的治理总是不变的；或者，如果它可以改变——这是不可能的——那么它无疑在任何时候都是可变的。\newline{}\par

那么，朱庇特有可能失去他的王国，就像他的父亲失去它一样吗？毫无疑问，是这样的。因为当那位神祇既不饶恕处女，也不饶恕有夫之妇时，他唯独没有侵犯特提斯，这是由于一个预言，预言说凡是与她所生的儿子，都会比他的父亲更伟大。首先，在他身上缺乏与神不相称的预知；因为若不是特弥斯向他讲述了未来的事，他自己原本不会知道。但如果他不是神圣的，他确实就不是神；因为神性（divinity）一词从神（god）而来，正如人性（humanity）从人（man）而来。然后，有一种软弱的意识；但那个曾经害怕的人，必然是害怕一个比他更强大的存在。但这样做的人肯定知道他不是最伟大的，因为可能存在更伟大的事物。他还以最庄严的方式向斯提克斯沼泽起誓：\Quote{这是天界诸神唯一敬畏的对象。}这个敬畏的对象是什么？或者由谁设定的？那么，是否有一种强大的力量能够惩罚那些作伪证的神祇？如果众神是不朽的，为什么如此害怕这阴间的沼泽？为什么他们要害怕那只有被死亡必然性束缚的人才会见到的东西？那么，为什么人们要举目望天？为什么他们以天上的神祇起誓，而天上的神祇自己却求助于阴间的神祇，并在他们中间找到可敬畏和崇拜的对象？但那个说法是什么意思：有一些命运女神，所有神祇和朱庇特本人也服从它们？如果\OriginalTerm{帕耳开}{Parcæ}的力量如此巨大，以至于它们比所有天上的神祇及其统治者和主人更强大，那么为什么不说它们才是统治的呢？因为必然性迫使所有神祇服从它们的法律和法令？现在，谁能怀疑一个服从于某事物的人不可能是最伟大的？因为如果他是最伟大的，他就不会接受命运女神，而是会指定它们。现在，我回到另一个我此前遗漏的话题。在一位女神身上，他唯独克制了自己，尽管他深深地爱恋着她；但这并非出于任何美德，而是出于对继承者的恐惧。但这种恐惧显然表明他是一个终有一死且软弱无能、毫无分量的人：因为在他出生的那一刻，他本来可能被处死，就像他的兄长被处死一样；如果他有可能活下去，他绝不会将最高权力让给一个弟弟。但朱庇特本人由于被秘密保存并秘密抚养，被称为宙斯（Zeus）或禅（Zen），并非如他们所想象的那样，是由于天上的火的热力，也不是因为他是生命的赐予者，或者因为他将生命的气息吹入活物——这些能力唯独属于天主；因为他如何能赐予生命的气息，而他自己却从别处接受了它？但他之所以如此称呼，是因为他是萨图尔努斯（Saturn）的男性子女中第一个活下来的。因此，如果萨图尔努斯没有被他的妻子欺骗，人类可能就会有另一位神祇作为他们的统治者。但有人会说，这些是诗人们虚构的。持这种观点的人错了。因为他们谈论的是关于人的事；但为了美化他们习惯以赞美来纪念的那些人，他们说他们是神。因此，他们关于这些人的神性所说的话是虚构的，而关于他们的人性所说的话则不是；这一点从我们将要提出的一个例子中可以明显看出。当他要对达娜厄（Danae）施暴时，他将大量的金币倒在她的怀里。这是他为她的耻辱所付出的代价。但那些将他作为神来谈论的诗人，为了不削弱他所假定的威严的权威，虚构说他本人化身为金雨降临，使用了和他们在描述大量投射物和箭矢时所说的\Quote{铁雨}同样的比喻。据说他通过一只鹰带走了伽倪墨得斯（Ganymede）；这是诗人的描绘。但他要么是通过一个以鹰为军旗的军团带走了他；要么是他所乘坐的船有鹰形的守护神，正如他劫持欧罗巴（Europa）并带她过海时船上有公牛像一样。同样，据说他将伊那科斯（Inachus）之女伊俄（Io）变成了一头小母牛。为了让她逃脱朱诺（Juno）的愤怒，她就那样——全身覆盖着刚毛，以母牛的形态——游过大海，来到了埃及；在那里，她恢复了原来的样子，成为了如今被称为伊西斯（Isis）的女神。那么，有什么论据可以证明欧罗巴没有骑在公牛上，伊俄也没有变成母牛？因为编年史中有一个确定的日期，伊西斯的航行在那天被庆祝；由此我们得知她并没有游过海，而是乘船渡海。因此，那些自以为有智慧的人，因为他们理解天上不可能有活的、尘世的身体，便拒绝了整个关于伽倪墨得斯的故事，认为它是假的，并察觉到这件事发生在地上，因为事情本身及其情欲都是尘世的。因此，诗人们并没有凭空创造这些行为，因为如果他们这样做，他们将是卑鄙的；但他们给这些行为增添了一层色彩。因为他们说这些事情不是为了诋毁，而是出于美化的愿望。因此人们被欺骗了；尤其是因为当他们认为所有这些事情都是诗人们虚构的时候，他们却崇拜他们所不知道的东西。因为他们不知道诗意的许可的界限，不知道在虚构中可以走多远，因为诗人的职责是以某种优雅的方式将实际事件通过间接的变形改变和转移到其他表现形式中。但是完全虚构你所讲述的一切，那是愚蠢和欺骗，而不是做一个诗人。\par

但假使那些被视为神话的事物是诗人所虚构的，那么关于女性神祇和天神婚姻的叙述也是他们虚构的吗？为什么她们这样被塑造、这样被崇拜？除非偶然并非只有诗人，连画家和雕像家也说谎。因为如果这位\Person{朱庇特}是你们所称的神，如果他不生于\Person{萨图尔努斯}和\Person{奥普斯}，那么所有神庙中只应放置他的像。女性的像含义何在？那模糊的性别又是什么？若\Person{朱庇特}以此形像被表现，连石头也将承认他是人。他们说诗人说了谎，却又相信他们；是的，他们用事实本身确证诗人并未说谎；因为他们塑造神的形象，正从性别差异可见诗人所说为真。试问放置在神庙中\Person{朱庇特}脚前、与他同样受崇拜的\Person{伽倪墨得斯}像和鹰像，除了作为不敬的罪过和放荡的记忆永存之外，还能导出什么结论？因此，没有什么是诗人完全虚构的：也许有些被扭曲的塑造转移和遮蔽了，真理被包裹和隐藏其中；如那关于抽签分王国的事。他们说天归\Person{朱庇特}，海归\Person{尼普顿}，冥界归\Person{普路托}。为什么不让大地作第三份，除非这分配发生在地上？所以，他们真的这样分割和分配了世界的治权：东方帝国归\Person{朱庇特}，西方一部分划归\Person{普路托}（他又有别名\Person{阿格西劳斯}）；因为东方地区——光明由此赐予凡俗人——似乎更高，而西方地区较低。这样，他们在虚构下遮蔽了真理，以至于真理本身并不减损公众的信念。关于\Person{尼普顿}的那份是明显的；因为我们说他的王国类似于\Person{马尔克·安东尼}所拥有的无限权力，元老院曾将沿海岸的权力赋予他，以便他惩治海盗并平定整个海洋。这样，所有的海岸连同岛屿都归\Person{尼普顿}抽得。这如何证明？无疑古代的故事可以作证。\Person{欧赫墨鲁斯}，一位古代作家，来自\Place{墨塞涅}城，收集了\Person{朱庇特}及其他被视为神者的事迹，并根据最古老神庙中的题词和神圣铭文，特别是\Person{特里费利亚的朱庇特}圣所中的铭文，编成一部史书。那里有铭文表明\Person{朱庇特}本人曾立下一根金柱，他在柱上记述了自己的功业，以便后世能记住他的作为。这部史书被\Person{恩尼乌斯}翻译并沿用，其言曰：\Quote{朱庇特将海洋的治理交给尼普顿，使他统辖所有海岛和沿海之地。}\par

因此，诗人的记述是真实的，只是披着外表的遮盖和修饰。可能是\Place{奥林波斯山}给了诗人暗示，说\Person{朱庇特}得了天国，因为\Place{奥林波斯}既是山名也是天名。但同一史书告诉我们，\Person{朱庇特}住在\Place{奥林波斯山}上，它说：\Quote{那时朱庇特一生大部分时间度在奥林波斯山上；若有任何争讼，人们便到他那里求裁断。此外，若有人发明了有益于人类生活的新事物，便会去那里向朱庇特展示。}诗人以这种方式转述许多事物，并非为了虚假地抨击他们所崇拜的对象，而是借各种着色的形象为他们的诗歌增添美和雅致。但那些不理解被形象所表现的方式、原因或本性的人，便攻击诗人虚假和亵渎。连哲学家也被这错误欺骗了；因为由于关于\Person{朱庇特}的这些事迹显得与神的特性不符，他们便引入两个\Person{朱庇特}：一个自然的，一个神话的。一方面他们看到了真实：诗人所谈论的那位确实是人；但对于那个自然的\Person{朱庇特}，因受迷信的普遍实践引导，他们犯了错误，将一个人的名字转给了天主——我们已经说过，天主是惟一的，原不需要名字。但不可否认的是，那位由\Person{奥普斯}和\Person{萨图尔努斯}所生的就是\Person{朱庇特}。因此，那些将\Person{朱庇特}的名字赋予至高天主的人，其信念是空虚的。因为有些人惯于以此借口为自己的错误辩护：他们被说服天主的独一性，既然不能否认这一点，便断言他们崇拜祂，但乐意称祂为\Person{朱庇特}。但还有什么比这更荒谬的呢？因为\Person{朱庇特}通常不单独受崇拜，而是伴有妻子和女儿。由此他的真实本性显而易见；也不能将那名字转移到一个既无\Person{密涅瓦}又无\Person{朱诺}的地方。我为何要说这名字的特殊含义不表达神性的权能，而是人的权能？因为\Person{西塞罗}解释\Person{朱庇特}和\Person{朱诺}这两个名字源于帮助；\Person{朱庇特}被称为帮助之父——这个名字不适合天主：因为帮助是人的部分，向陌生人提供某些援助，且受益甚小。没有人恳求天主帮助他，而是恳求祂保护他，赐予他生命和安全，这比帮助重大得多，重要得多。\par

既然我们是在谈论一位父亲，没有人会说父亲在生育或抚养儿子时是在帮助儿子。因为那种说法太微不足道，不足以表达来自父亲的恩惠的宏大。对于天主，我们的圣父，由祂我们存在，完全属于祂，由祂形成、赋予生命、光照，祂赐给我们生命，赐予我们安全，供给我们各种食物，这有多么不适合！一个人若认为他只是得到天主的帮助，就没有领会天主的恩惠。因此，他不仅无知，而且不敬，竟以朱庇特之名贬低至高权能的卓越性。所以，如果我们从他的行为和品格已经证明朱庇特是一个人，曾在世上为王，那么只剩下我们还应探究他的死亡。恩纽斯在他的\WorkTitle{神圣史}中，描述了他一生中所做的一切事后，结尾这样写道：那时，朱庇特五次环游大地，将统治权赐给所有朋友和亲人，为人类设立法律，提供固定的生活方式和谷物，并给予许多其他恩惠，受到不朽的荣耀和纪念，为朋友留下持久的纪念，当他的年岁几乎耗尽时，他在克里特改变了生命，离世到了众神那里。库瑞忒斯，他的儿子们，照顾他并尊敬他；他的坟墓在克里特岛克诺索斯城，据说维斯塔建立了这座城；他的墓上刻有古希腊文字的铭文：\Quote{\GreekText{Ζὰν} \GreekText{Κρόνου}}，即拉丁文的\Quote{Jupiter the son of Saturn}。这无疑不是由诗人传下来的，而是由古代事件的作者传下来的；这些事如此真实，以至于\Person{西比尔}的几行诗句也证实了它们，内容如下：——\par

\Quote{无生命的鬼魔，死者的肖像，\newline{}不祥的克里特以其坟墓自夸。}\par

西塞罗在其著作\WorkTitle{论神性}中，说到神学家列举了三位朱庇特，接着指出第三位是克里特人，\Person{萨图尔努斯}之子，他的坟墓就在那个岛上。因此，一位神怎能在一处活着，在另一处死去；在一处有神庙，在另一处有坟墓？那么让罗马人知道，他们的\Place{卡皮托利山}，即他们公开崇敬对象的首要所在，不过是一座空洞的纪念碑。\par

现在让我们来看在他之前为王的他的父亲，他自身或许拥有更大的权能，因为据说他是由如此巨大的元素结合而生。让我们看看他身上有什么配得上神的地方，尤其是据说他拥有黄金时代，因为在他的统治期间大地上有正义。我在他身上发现了他儿子所没有的东西。因为有什么比正义的统治和虔诚的时代更适合神的品性呢？但当我本着同样的原则思考他是一个儿子时，我不能认为他是至高神；因为我看到有比他更古老的东西——就是天和地。但我寻找的是超越一切存在、万有的根源和本原的天主。那创造天本身、奠定大地根基的必然存在。但如果萨图尔努斯如所认为的那样是从这些产生的，他怎么能是首神？既然他起源于另一位？或者在萨图尔努斯出生之前谁掌管宇宙？但这正如我最近所说，是诗人的虚构。因为无感觉的元素，被如此长的间隔隔开，不可能相遇并生出一个儿子，或者所生的孩子一点也不像父母，却具有父母没有的形状。\par

\Person{米努西乌斯·菲利克斯}在其题为\WorkTitle{奥克塔维乌斯}的论著中提出了这些证据：\Quote{\Person{萨图尔努斯}被其子放逐后，来到\Place{意大利}，被称为\Person{科路斯}（即天）之子，因为我们习惯说，那些我们所钦佩其美德的人，或那些意外到来的人，是从天而降；他被称为地之子，因为我们称那些父母不详之人为地之子。}这些说法确实与真相有几分相似，但并不真实，因为显然他在统治期间就已经被如此尊崇。他本可以这样论证：\Person{萨图尔努斯}是一位极其强大的国王，为了保存其父母的名号，便将他们的名字赋予了天和地，而此前它们另有名称，因此我们知道山岳和河流也被赋予了名字。因为当诗人提到\Person{阿特拉斯}或\Place{伊那科斯河}的后裔时，他们并非绝对地说人可能从无生命的物体中出生，而无疑是指那些从这些人的后裔中出生的人，这些人或在生前或在死后将名字赋予山岳或河流。因为这是古人，尤其是希腊人的常见做法。因此我们听说，一些海得名于那些落入其中的人，如\Place{爱琴海}、\Place{伊卡里亚海}和\Place{赫勒斯滂}。在\Place{拉丁姆}，\Person{阿文提努斯}将其名字赋予他葬身的那座山；\Person{提贝里努斯}（即台伯河）将其名字赋予他溺死的那条河。因此，如果那些生育了最强大国王之人的名字被赋予了天和地，这并不奇怪。因此，似乎\Person{萨图尔努斯}并非从天而生（这不可能），而是从那个名为\Person{乌拉诺斯}的人所生。\Person{特里斯墨吉斯忒斯}证实了这一点；因为当他提到完全有学识的人寥寥无几时，他列举了其亲属中的\Person{乌拉诺斯}、\Person{萨图尔努斯}和\Person{墨丘利}。由于他对此一无所知，他便给出了另一种说法；我已说明他本可以如何论证。现在我将讲述此事是如何、在何时、由谁完成的；因为做这事的不是\Person{萨图尔努斯}，而是\Person{朱庇特}。\Person{恩尼乌斯}在其\WorkTitle{圣史}中这样叙述：\Quote{然后\Person{潘}将他领到一座名为天柱的山上。他登上那里，环视大地，在那山上建造了一座祭坛敬拜\Person{科路斯}；\Person{朱庇特}是第一个在那祭坛上献祭的人。在那里，他仰望天（我们现在这样称呼它），以及那高于世界、被称为穹苍之处，他以他祖父的名字为天命名；\Person{朱庇特}在祈祷中首先将\InnerQuote{天}之名赋予了那被称为穹苍之处，并把他献在那里的牺牲整只焚烧。}不止在此处发现\Person{朱庇特}献过祭。\Person{凯撒}在\Person{阿拉托斯}的作品中也提到，\Person{阿格拉奥斯泰涅斯}说，当他从\Place{纳克索斯}岛出发去对抗提坦，并在岸边献祭时，一只鹰飞到\Person{朱庇特}那里作为预兆，胜利者将此视为吉兆，并将其置于自己的保护之下。但\WorkTitle{圣史}证实，在此之前已有一只鹰落在他头上，预示他将得到王国。那么，\Person{朱庇特}除了献祭给他的祖父\Person{科路斯}，还能献祭给谁呢？据\Person{欧赫墨罗斯}所言，\Person{科路斯}死在\Place{奥克阿尼亚}，葬于\Place{奥拉提亚}镇？\par

% structure: p0080 source=h2 target=subsection reason=relative-heading-rank; base=h2

\subsection{第十二章 斯多葛派将诗人的虚构转化为哲学体系}

既然我们已经揭示了诗人的奥秘，并查明了\Person{萨图尔努斯}的父母，让我们回到他的美德与行为。据说，他在统治时公正。首先，仅从这一事实来看，他现在已不是神，因为他的统治已经终结。其次，他甚至并不公正，不仅对他的儿子（据说是被他吞噬的）不虔敬，而且对他的父亲（据说他将其阉割）也不虔敬。这也许确实发生过。但人们考虑到那被称为天的元素，便拒绝将整个寓言视为极其愚蠢的编造；尽管斯多葛派（按惯例）试图将其转化为一种物理体系，\Person{西塞罗}在其论著\WorkTitle{论诸神的本性}中陈述了他们的观点。他说，他们认为，最高的、以太般的天之本性，即火之本性，它自身产生万物，却没有那包含生殖器官的身体部分。这一理论或许适用于\Person{维斯塔}，倘若她被称为男性的话。正是因此，他们认为\Person{维斯塔}是贞女，因为火是不朽的元素；没有任何东西能从火中出生，因为它会吞噬它所触及的一切。\Person{奥维德}在\WorkTitle{岁时记}中说：\Quote{你们不要认为\Person{维斯塔}不过是活着的火焰；你们看不到有形体从火焰中产生。因此她确实是贞女，因为她不撒播种子，也不接收种子，并且热爱贞洁的侍者。}\par

这也可能归因于\Person{武尔肯}（Vulcan），他确实被认为是火，但诗人并没有使他阉割。也可能归因于太阳，在太阳中有着生殖力的本性和原因。因为没有太阳的火热，任何事物都无法诞生或增长；所以没有哪个元素比热更需要生殖器官，万物都靠热的滋养而受孕、产生和维持。最后，即使情况如他们所愿，为什么我们应认为\Person{乌拉诺斯}被阉割，而不是他生来就没有生殖器官呢？因为他若靠自己产生，显然他不需要生殖器官，既然他生下了\Person{萨图恩}本人；但如果他有生殖器官，并被他儿子阉割，万物的起源和整个自然就毁灭了。我为何要说他们不仅剥夺了\Person{萨图恩}的神性，还剥夺了人性智慧，当他们断言\Person{萨图恩}是包容时空进程与变化者，并且他在希腊语中就有那个名字？因为他被称为\OriginalTerm{克罗诺斯}{Cronos}，这等同于\OriginalTerm{克洛诺斯}{Chronos}，即时间段。但他被称为萨图恩，因为他被岁月充满。这些是\Person{西塞罗}的话，阐述\OriginalTerm{斯托屹派}{Stoics}的意见：\Quote{任何人都会很容易理解这些事的无价值。因为如果萨图恩是乌拉诺斯之子，时间怎能从乌拉诺斯诞生，或乌拉诺斯被时间阉割，或后来时间能被他的儿子\Person{朱皮特}剥夺主权？或者朱皮特如何从时间诞生？永恒能被多少年充满，既然它没有界限？}\par

% structure: p0083 source=h2 target=subsection reason=relative-heading-rank; base=h2

\subsection{第十三章 \OriginalTerm{斯托屹派}{Stoics}的解释是何等虚妄和琐碎，以及其中关于\Person{朱皮特}的起源，关于\Person{萨图恩}和\Person{奥普斯}。}

如果这些哲学家的思辨是琐碎的，那么留下什么呢？除非我们相信这是一个事实：作为一个人，他遭受了另一个人的阉割？除非偶然有人视他为神，而他竟害怕一个共同继承人；然而，如果他具有任何神圣知识，他不应阉割他的父亲，而应阉割他自己，以防止\Person{朱皮特}的诞生，因为\Person{朱皮特}剥夺了他王国的所有权。而且，当他娶了他的姐妹\Person{瑞亚}（在拉丁语中我们称她为\Person{奥普斯}）时，据说他被神谕警告不要抚养他的男孩，因为将会发生他被一个儿子驱逐流放。由于害怕这一点，显然他没有像寓言所报道的那样吞吃他的儿子，而是杀了他们；虽然在神圣历史中记载，\Person{萨图恩}、\Person{奥普斯}和其他人在那时习惯于吃人肉，但\Person{朱皮特}，他给人法律和文明，是第一个通过法令禁止那种食物的人。如果这是真的，那么他身上可能有什么正义？但让我们假设一个虚构的故事，说\Person{萨图恩}吞吃了他的儿子，只是在某种意义上为真；那么我们必须与凡俗人一样认为，他吃掉了他的儿子，就是那些他埋葬了的儿子？但当\Person{奥普斯}生下\Person{朱皮特}时，她偷走了婴儿，秘密地送到\Place{克里特岛}去抚养。再者，我不能不责备他缺乏远见。因为他为何从别人那里接受神谕，而不从他自身\Emph{？}他身居天堂，为何看不见地上发生的事？为何那些\Person{科律邦特斯}和他们的铙钹逃过了他的注意？最后，为什么存在任何更大的力量可能战胜他的能力？无疑，年老的他容易被一个年轻人战胜，并剥夺了主权。因此他被流放和放逐；经过长期漂泊后，乘船来到意大利，正如\Person{奥维德}在他的\WorkTitle{岁时记}中所叙述的：——\par

\Quote{船的原因仍有待解释。\Person{棍刀神}先游历了世界，然后乘船来到\Place{托斯卡纳河}。}\par

\Person{雅努斯}接待了流浪贫苦的他；古钱币证明了这一点，上面有\Person{雅努斯}的双面像，另一面是一艘船；正如同一诗人补充的：——\par

\Quote{但虔诚的后代在钱币上刻了一艘船，见证这位陌生之神的到来。}\par

因此不仅所有诗人，而且古代历史和事件的作家都同意他曾经是一个人，因为他们将他当时在意大利的作为流传下来：希腊作家中有\Person{第奥多罗斯}和\Person{塔勒斯}；拉丁作家中有\Person{内波斯}、\Person{卡西乌斯}和\Person{瓦罗}。因为当时人们在意大利过着质朴的生活，——\par

\Quote{他首先将种族联合，\newline{}先前分散在山巅，\newline{}并给他们法令遵守，\newline{}并愿意他躺卧的土地\newline{}应承载\Place{拉丁姆}之名。}\par

有谁想像他是一个神，曾被驱赶流放、逃走、躲藏？没有如此愚钝之人。因为逃走或躲藏者，必害怕暴力和死亡。\Person{奥费欧斯}，生活在比他更近的时代，公开叙说\Person{萨图恩}曾在地上和人群中为王：——\par

\Quote{首先\Person{克罗诺斯}统治地上众人，\newline{}然后从\Person{克罗诺斯}生出强大的王，\newline{}广布声名的\Person{宙斯}。}\par

我们自己的\Person{马罗}也说：\par

\Quote{金色的萨图恩曾在大地上度过这种生活；}\par

在另一处：——\par

\Quote{那是传说的黄金时代，\newline{}如此和平、宁静地流逝\newline{}在他统治下的岁月。}\par

诗人在前一段中没有说他在天堂过这种生活，在后一段中也没有说他统治天上的诸神。由此可见他是地上的王；他在另一处更清楚地宣告：——\par

\Quote{黄金时代的恢复者，\newline{}在萨图恩古时统治的土地上。}\par

\Person{恩尼乌斯}在他翻译的\Person{欧海默勒斯}中说，\Person{萨图恩}不是第一个统治的，而是他的父亲\Person{乌拉诺斯}。他说，起初，\Person{乌拉诺斯}首先在地上拥有最高权力。他与兄弟们一起建立和准备了那个王国。如果存在疑问，从最权威的方面来说，关于子与父并没有大的争议。但有可能两者都发生了：\Person{乌拉诺斯}首先开始在其他中突出于权力，并拥有首要地位，但并非王国；之后\Person{萨图恩}获得了更大的资源，并取了王的称号。\par

% structure: p0099 source=h2 target=subsection reason=relative-heading-rank; base=h2

\subsection{第十四章  欧赫墨罗斯与恩纽斯的圣史关于诸神所教导的}

如今，由于圣史与我们所述之事在某种程度上有所不同，让我们揭开那些真实记载中的内容，以免我们在指控迷信时，似乎追随并赞同诗人的愚妄。以下是恩纽斯的话：\Quote{此后，\Person{萨图尔努斯}娶了\Person{奥普斯}。比\Person{萨图尔努斯}年长的\Person{提坦}为自己要求王位。于是，他们的母亲\Person{维斯塔}及姐妹\Person{刻瑞斯}和\Person{奥普斯}劝告\Person{萨图尔努斯}不要将王位让给他的兄弟。\Person{提坦}因身材不及\Person{萨图尔努斯}，且见母亲和姐妹都在努力让\Person{萨图尔努斯}为王，便将王位让给了他。因此，他与\Person{萨图尔努斯}约定：若\Person{萨图尔努斯}生下任何男婴，他不得将他们抚养成人。\Person{提坦}这样做，是为了让王位能回到自己的儿子手中。随后，当\Person{萨图尔努斯}的第一个儿子出生时，他们杀死了他。后来，\Person{尤皮特}和\Person{尤诺}这对孪生子出生了。于是，他们将\Person{尤诺}呈现在\Person{萨图尔努斯}眼前，秘密藏匿了\Person{尤皮特}，把他交给\Person{维斯塔}抚养，瞒着\Person{萨图尔努斯}。\Person{奥普斯}也瞒着\Person{萨图尔努斯}生下了\Person{尼普顿}，并秘密藏匿了他。同样，\Person{奥普斯}第三次分娩生下了孪生子\Person{普路托}和\Person{格劳卡}。\Person{普路托}在拉丁语中称为\Person{狄斯帕特尔}；另一些人称他为\Person{奥尔库斯}。于是，他们将女儿\Person{格劳卡}给\Person{萨图尔努斯}看，藏匿并隐藏了儿子\Person{普路托}。随后\Person{格劳卡}在幼年时死去。}这就是\Person{尤皮特}及其兄弟的世系，正如这些所记载的，亲属关系以此方式从圣史传给我们。稍后他又引入了这些：\Quote{然后，\Person{提坦}得知\Person{萨图尔努斯}有了儿子并被秘密抚养，就悄悄带上他的儿子们（他们被称为提坦），擒获了他的兄弟\Person{萨图尔努斯}和\Person{奥普斯}，将他们围困在一道墙内，并派守卫看管他们。}\par

这段历史的真实性由\Person{埃律特莱亚女先知}教导，她几乎说了同样的事，只有一些不影响主题本身的差异。因此，\Person{尤皮特}被免除了最严重恶行的指控——即传说他用锁链捆绑了自己的父亲；因为这是他的叔父\Person{提坦}所为，因为他违背了自己的承诺和誓言，抚养了男婴。历史的其余部分如此编排：据说\Person{尤皮特}长大后，得知父母被守卫包围并被囚禁，便带领一大群克里特人前来，在战斗中击败了\Person{提坦}及其儿子们，救出了父母，将王位还给了父亲，然后返回了克里特。此后，据说有一个神谕给了\Person{萨图尔努斯}，要他提防儿子把他赶出王位；\Person{萨图尔努斯}为了削弱神谕并避免危险，设下埋伏要杀死\Person{尤皮特}；\Person{尤皮特}得知阴谋后，重新为自己索要王位，并放逐了\Person{萨图尔努斯}；而\Person{萨图尔努斯}被\Person{尤皮特}派去抓捕或杀死他的武装人员追捕，在各处漂泊，几乎无法在意大利找到一个藏身之处。\par

% structure: p0102 source=h2 target=subsection reason=relative-heading-rank; base=h2

\subsection{第十五章  凡人是如何获得神之名的}

既然从这些事情中显而易见他们曾是凡人，那么不难看出他们是如何开始被称为神的。因为如果在\Person{萨图尔努斯}或\Person{乌兰诺斯}之前没有国王，由于人数稀少，人们过着没有任何统治者的乡野生活，毫无疑问，那时人们开始以最高的赞誉和新的尊荣来尊崇国王本人及其整个家族，甚至称他们为神；这或许是由于他们奇妙的卓越，当时还粗鄙单纯的人们确实怀有这样的想法，或者如常有的情况，出于对眼前权力的奉承，或者由于他们带来的好处——人们因此被规整并进入文明状态。后来，国王们自己，由于备受他们使之文明的人们爱戴，在死后留下了对他们的怀念。因此，人们塑造了他们的形象，以便从凝视他们的肖像中获得一些安慰；由于对功绩的爱戴更进一步，人们开始敬拜逝者的记忆，以显得对他们的服务心存感激，并吸引继任者怀有善治的愿望。\Person{西塞罗}在他的著作\WorkTitle{论神性}中教导这一点，说：\Quote{然而，人类的生活和共同的交往，导致那些因其恩惠而卓越的人，通过名声和善意被提升到天上。因此，\Person{赫丘利}、\Person{卡斯托耳}与\Person{波鲁克斯}、\Person{埃斯库拉庇乌斯}和\Person{利伯}被列于诸神之中。}在另一处：\Quote{在大多数国家中可以理解，为了激发勇气，或使最勇敢的人更愿意为国家冒险，他们的记忆被献上不朽诸神的荣耀。}无疑正是为此，罗马人献祭了他们的凯撒，摩尔人献祭了他们的国王。就这样，宗教性的尊崇逐渐被献给他们；那些认识他们的人，首先教导自己的儿女和孙辈，然后教导所有后代，遵行这一仪式。然而，这些伟大的国王，由于名声显赫，在各省都受到尊崇。\par

但各邦各族私下以最高的崇敬敬奉其民族或城市的建立者，这些建立者或是因勇武而显赫的男子，或是因贞洁而可钦的女子；如埃及人敬奉\Person{伊西斯}，摩尔人敬奉\Person{朱巴}，马其顿人敬奉\Person{卡比洛斯}，迦太基人敬奉\Person{乌拉诺斯}，拉丁人敬奉\Person{法乌努斯}，萨宾人敬奉\Person{桑库斯}，罗马人敬奉\Person{奎里努斯}。同样，雅典人敬奉\Person{弥涅耳瓦}，萨摩斯人敬奉\Person{朱诺}，帕福斯人敬奉\Person{维纳斯}，利姆诺斯人敬奉\Person{伏尔甘}，纳克索斯人敬奉\Person{利伯尔}，得洛斯人敬奉\Person{阿波罗}。于是在不同民族与国度间，便兴起了各种不同的神圣礼仪，因为人们渴望向自己的君王表示感激，却找不到其他可以向死者奉献的尊荣。此外，继任者的孝心也极大地助长了这一谬误：为显得自己出自神圣的根源，他们向自己的父母奉上神圣的尊荣，并命令他人也如此行。凡读过维吉尔作品中\Person{埃涅阿斯}对友人发令之词的人，岂会怀疑敬奉诸神的礼仪是如何创立的呢？——\par

\Quote{现在满杯奠酒倾献\newline{}给大能\Person{朱庇特}，万民敬拜，\newline{}呼求\Person{安喀塞斯}蒙福的灵魂。}\par

他不仅赋予他不朽，还赋予他掌管风的能力：——\par

\Quote{呼求诸风加速我们的航程，\newline{}并祈求我们如此珍爱的他，\newline{}年年接纳我们的祭品，\newline{}待我们建立那应许之城，\newline{}在颂扬他的神圣殿宇中。}\par

的确，\Person{利伯尔}、\Person{潘}、\Person{墨丘利}和\Person{阿波罗}对待\Person{朱庇特}也是这样，而后来的继承者对待他们也是如此。诗人也推波助澜，借助为取悦人而作的诗歌，将他们升到天上；正如那些用虚假的颂词谄媚君王（即便是邪恶的君王）的人所为。这一弊端起源于希腊人，他们的轻浮因着语言的才能与丰富，竟以难以置信的程度激起谎言的迷雾。因此，出于对他们的仰慕，他们首先创立了神圣的礼仪，并将之传递给万国。因着这虚妄，西彼拉如此斥责他们：——\par

\Quote{希腊啊，你为何信赖属君王的凡人？\newline{}为何向死者奉献空虚的礼物？\newline{}你向偶像献祭；这错误是谁暗示的，\newline{}竟使你离开大能天主的临在，\newline{}而奉献这些祭品？}\par

\Person{马尔库斯·图利乌斯·西塞罗}不仅是一位出色的演说家，也是一位哲学家，因为唯有他模仿了\Person{柏拉图}，在那篇关于女儿去世时安慰自己的文章中，他毫不犹豫地说，那些被公开敬拜的神原本是人。他的这一见证应被视为更有分量，因为他担任\OriginalTerm{占卜官}{augur}这一司祭职务，并证明自己敬礼并尊崇同样的神。因此，在短短几行诗中，他向我们呈现了两个事实。因为当他宣布打算像古人那样将自己女儿的形象祝圣时，他既教导说他们是死人，也揭示了虚妄迷信的起源。\Quote{由于，确实，}他说，\Quote{我们看见许多男女列在众神之中，并敬礼他们的神龛，这些神龛在城市和乡村受到最崇高的尊敬，让我们赞同那些人的智慧，我们凭借他们的才能和发明，生活完全被法律和制度所装饰，并建立在稳固的基础上。如果任何活物配得上被祝圣，那无疑就是她。如果\Person{卡德摩斯}、\Person{安菲特律翁}或\Person{廷达柔斯}的后代配得上被名声颂扬到天上，那么同样的荣誉无疑应当归于她。我确实会这样做；并且在天神的赞同下，我将把你，所有女人中最优秀和最有学问的，安置在他们的聚会中，并将你祝圣给所有人的敬重。}或许有人会说\Person{西塞罗}因过度悲伤而胡言乱语。但实际上，那篇完整的演说，无论在学识、例证还是在表达风格上，都堪称完美，没有显示任何错乱心智的迹象，而是显示出坚定和判断；而这段文字本身也没有表现出悲伤的迹象。因为我认为他不可能写得如此多样、丰富和修饰，除非他的悲伤已经被理性本身、朋友的安慰和时间的流逝所缓解。我何必提及他在\WorkTitle{论共和国}以及关于荣誉的著作中所说的话呢？因为在\WorkTitle{论法律}的著作中，他效仿\Person{柏拉图}，希望阐述他认为一个正义和明智的国家会采用的法律，他这样规定宗教：\Quote{让他们敬畏那些神，既包括那些一直被视为天神的神，也包括那些因其对人类的服务而被安置在天上的神：\Person{赫丘利}、\Person{利柏耳}、\Person{埃斯库拉庇乌斯}、\Person{卡斯托耳}、\Person{波鲁克斯}和\Person{奎里努斯}。}此外，在他的\WorkTitle{图斯库兰论辩}中，当他说天上几乎充满了人类时，他说：\Quote{如果，确实，我应当尝试考察古代记载，并从其中提取希腊作家传下来的东西，即使那些被视为最高等级的神，也会被发现是从我们这里去了天上。查问在希腊哪些人的坟墓被指出：记住，因为你已被领入秘仪，秘仪中传下了什么；然后你终于会明白这种信念流传得多么广泛。}显然，他诉诸\Person{阿提库斯}的良心，以便从秘仪本身就能明白，所有被敬拜的都是人；当他在\Person{赫丘利}、\Person{利柏耳}、\Person{埃斯库拉庇乌斯}、\Person{卡斯托耳}和\Person{波鲁克斯}身上毫不犹豫地承认这一点时，他不敢公开坦率地承认他们的父亲\Person{阿波罗}和\Person{朱庇特}也是如此，同样还有\Person{尼普顿}、\Person{武尔坎}、\Person{玛尔斯}和\Person{墨丘利}，他称这些神为更伟大的神；因此他说这种观点广为流传，以便我们理解关于\Person{朱庇特}和其他更古老的神也是如此：因为如果古人以同样的方式祝圣他们的记忆，正如他说他将祝圣他女儿的形象和名字，那么哀悼者可以原谅，但相信者不能原谅。因为有谁如此愚昧，竟相信天堂会为一群无知的群众的同意和喜悦而向死者开放？或者有人能给予他自己所没有的东西？在罗马人中，\Person{尤利乌斯·凯撒}被奉为神，因为一个罪人\Person{安东尼}喜欢这样；\Person{奎里努斯}被奉为神，因为牧人们认为好，尽管其中一个是其孪生兄弟的谋杀者，另一个是其国家的毁灭者。但如果\Person{安东尼}不曾担任执政官，由于他对国家的贡献，\Person{盖乌斯·凯撒}连死人的荣誉都得不到，而且那是根据他岳父\Person{皮索}和他的亲戚\Person{卢基乌斯·凯撒}的建议，他们反对举行葬礼，以及根据执政官\Person{多拉贝拉}的建议，他推倒了广场上的柱子，即他的纪念碑，并净化了广场。因为\Person{恩尼乌斯}宣称\Person{罗慕路斯}受到人民的哀悼，他让人民因失去国王而悲痛地说：\Quote{哦，\Person{罗慕路斯}，\Person{罗慕路斯}，说你的国家产生了你什么样的守护神？你生了我们到光明区域。哦父亲，哦父，哦种族，由神而生。}由于这种哀悼，他们更轻易地相信了\Person{尤利乌斯·普罗库鲁斯}所说的谎言，他被元老们收买，要向百姓宣布他曾看见国王以一种比人更威严的形象；并且他命令人民为他建造一座神殿，他是神，被称为\Person{奎里努斯}。这一举动立刻使人民相信\Person{罗慕路斯}已升入众神之列，并使元老院免除了杀害国王的嫌疑。\par

% structure: p0111 source=h2 target=subsection reason=relative-heading-rank; base=h2

\subsection{第十六章 以什么论据证明那些以性别区分的不能是神。}

我本可满足于我已述说的那些事，但仍有诸多必要之事，为我所承担的工作。因为虽然通过摧毁迷信的主要部分，我已将其全部清除，但我仍乐意追述剩余部分，更充分地驳斥如此根深蒂固的信念，好让人们最终羞愧并悔改他们的错误。这是一项伟大的事业，配得上一名男子汉。\Quote{我着手将人的心灵从迷信的束缚中释放出来，}正如\Person{卢克莱修}所说；他确实未能做到这一点，因为他没有提出任何真理。这是我们的职责，我们既肯定真天主的存在，又驳斥虚假的神祇。因此，那些认为诗人们编造了关于神祇的寓言，却又相信女神的存在并崇拜她们的人，不知不觉又回到了他们所否认的——即神祇有性交并生育。因为两性的设立不可能不是为了生育。但一旦承认了性别的差异，他们却未察觉到随之而来的必然是受孕。而这对天主而言是不可能的。但姑且按他们所想象的那样；因为他们说\Person{朱庇特}和其他神祇有儿子。因此新的神祇不断出生，而且确实是每天都在出生，因为神祇在生育能力上并不逊于人类。由此推出，万物充满了无数神祇，因为诚然他们没有一个会死亡。因为既然人类的数量难以置信，其数目不可估量——尽管他们出生后必有一死——那么对于那些经历了如此多的世代而出生且保持不朽的神祇，我们该如何设想呢？那么，为何只有如此少的神祇受到崇拜？除非我们以某种方式认为神祇有两性，不是为了生育，而仅仅是为了享乐，并且神祇行那些人类羞于去做和屈从之事。但当有人说某神出生于某神时，随之而来的是他们一直持续出生，如果他们曾在某个时间出生；或者如果他们曾在某个时间停止出生，我们理应知道他们为何以及在何时停止。\Person{塞内卡}在他的道德哲学著作中，不无诙谐地问道：\Quote{为什么被诗人们描绘成最沉迷于情欲的\Person{朱庇特}停止了生育孩子？是因为他成了花甲老人，受帕皮亚法约束吗？还是他获得了三子特权？或者他终于产生了这样的想法：\InnerQuote{你怎样待人，也必怎样被人待；}他害怕有人会像他对待\Person{萨图尔努斯}那样对待他吗？}但让那些主张他们是神祇的人看看，他们能以何种方式回答我将提出的这个论证。如果神祇有两性，那么就有夫妻交合；如果这发生，他们必定有房屋，因为他们并非没有德性和羞耻感，以至于像我们所见到的禽兽那样公开和混杂地行事。如果他们有了房屋，随之而来的是他们也有城市；对此我们有\Person{奥维德}的权威，他说：\Quote{众神占据各自的位置；在此前方，天上强大而显赫的居民安置了他们的住所。}如果他们有了城市，他们也就会有田地。谁能看不到结果——即他们耕种土地？而这样做是为了食物。因此他们是可朽的。这个论证反过来同样有力。因为如果他们没有土地，就没有城市；如果没有城市，也就没有房屋；如果没有房屋，就没有夫妻交合；如果没有夫妻交合，就没有女性性别。但我们看到神祇中也有女性。因此并没有神祇。如果有人能做到，让他推翻这个论证。因为一事如此跟随另一事，以至于不可能不承认这些结论。但连前一个论证也没有人能反驳。两性中一强一弱。因为男性更强壮，女性更虚弱。但神祇不会有虚弱；因此没有女性性别。加上前一个论证的最后结论：既然神祇中也有女性，因此没有神祇。\par

% structure: p0113 source=h2 target=subsection reason=relative-heading-rank; base=h2

\subsection{第十七章 论斯多葛学派的同一观点，以及关于神祇的艰难和可耻行为}

基于这些原因，斯多葛派形成了一种不同的关于神祇的观念；因为他们不明白真理是什么，所以试图将他们与自然事物的体系结合。\Person{西塞罗}追随他们，提出了关于神祇及其宗教的观点。他说道：\Quote{那么，你是否看到，从那些已被发现良好且有用的物理主题中，如何引出了对虚假和虚构神祇的论证？这种情形导致了虚假的意见和动荡的错误，以及几乎老妇般的迷信。因为我们知道神祇的形态、年龄、服饰和装饰；此外还有他们的种族、婚姻、所有亲属关系，以及一切还原到人类软弱相似性的事物。还有什么能说得更明白、更真实呢？}罗马哲学的首领，且担任最可敬的司祭职，驳斥了虚假和虚构的神祇，并证明对他们的崇拜是由几乎老妇般的迷信构成；他抱怨人们陷入了虚假的意见和动荡的错误。因为他关于神祇自然属性的整部第三卷书\WorkTitle{论神性}完全推翻了所有宗教。那么，我们还期望什么更多的事呢？我们能在口才上超越\Person{西塞罗}吗？绝不；但他缺乏自信，因为他不知真理，正如他在同一著作中淳朴地承认。因为他说，他更容易说出什么不是，而不是什么是；也就是说，他意识到所接受的体系是虚假的，但不知真理。因此很显然，那些被认为是神祇的人不过是人，他们的记忆在死后被奉为神圣。也正因为如此，每个神祇被赋予不同的年龄和既定的形态表征，因为他们的雕像被制作成死亡降临时他们各自所穿的服饰和年龄。\par

请容我们考虑这些不幸神祇的苦难。\Person{伊西斯}失去了她的儿子；\Person{色列斯}失去了她的女儿；\Person{拉托娜}被驱逐并驱赶于大地，艰难地找到一个小岛以生产。众神之母既爱慕一位俊美青年，又在发现他与妓女在一起时将其阉割；因此她的祭礼现今由\Person{加吕斯}作为司祭举行。\Person{朱诺}猛烈迫害妓女，因为她无法与自己的兄弟受孕。\Person{瓦罗}写道，\Place{萨摩斯}岛先前名为\Place{帕耳忒尼娅}，因为\Person{朱诺}在那里成长，并在那里与\Person{朱庇特}结婚。因此在\Place{萨摩斯}有她的一座极其古老高贵的神殿，以及一尊新娘装扮的雕像；她的年度祭礼以婚姻的形式举行。所以，如果她成长，如果她起初是少女而后成为妇人，那不明白她曾是凡人的人便承认自己是禽兽。我又何必提\Person{维纳斯}的放荡，她不仅服侍神祇，也服侍凡人的情欲？因为她与\Person{玛尔斯}的臭名昭著的淫乱生了\Person{哈耳摩尼亚}；与\Person{墨丘利}生了\Person{赫耳玛佛洛狄托斯}，兼具两性；与\Person{朱庇特}生了\Person{丘比特}；与\Person{安基塞斯}生了\Person{埃涅阿斯}；与\Person{布忒斯}生了\Person{厄律克斯}；与\Person{阿多尼斯}则未能生子，因为他被野猪击中，在少年时被杀。她首先开创了娼妓之术，正如神圣历史所载；并教导\Place{塞浦路斯}的妇女靠卖淫谋利，她为此命令，好让她自己在众女中不显得唯一不贞和追求男子。她还有任何宗教敬拜的理由吗？她所记载的通奸比生育还多。但即使是那些被颂扬的贞女，也未能保全其童贞不受玷污。因为我们从何源头设想\Person{厄里克托尼俄斯}诞生？难道是从大地，如同诗人们所显示的那样？但事件本身大声宣告。因为当\Person{伏尔甘}为神祇制造了武器，\Person{朱庇特}给他选择任意所愿赏报的权利，并照他的习惯以冥湖起誓不拒绝他任何要求，于是这跛脚工匠要求\Person{密涅瓦}为妻。此时卓越而强大的\Person{朱庇特}被如此重大的誓言所束，无法拒绝；但他却劝告\Person{密涅瓦}抵抗并保卫她的童贞。于是他们说到在争斗中\Person{伏尔甘}将种子洒落地上，由此\Person{厄里克托尼俄斯}诞生；而且这名字来源于\OriginalTerm{争斗}{\GreekText{ἔριδος}}和\OriginalTerm{土地}{\GreekText{χθσνο}}。\par

\SourceRecord{Translated by William Fletcher.；Ante-Nicene Fathers；Vol. 7.；Edited by Alexander Roberts, James Donaldson, and A. Cleveland Coxe.；Buffalo, NY: Christian Literature Publishing Co.,；1886.；<http://www.newadvent.org/fathers/07011.htm>.}

% structure: p0001 source=h1 target=section reason=page-title

\section{神圣训导，第二卷（论错误的起源）}

% structure: p0002 source=h2 target=subsection reason=relative-heading-rank; base=h2

\subsection{第一章：理性的遗忘使人不认识真天主，他们在逆境中敬拜祂，在顺境中藐视祂。}

虽然我已在前一卷中证明众神的宗教仪式是虚假的，因为那些普世人类出于愚昧的信念以不同的仪式所尊崇的对象，乃是凡人，他们完成了生命的期限后，便服从神圣规定的必然性而死去；但为免任何疑窦留存，这卷书将揭开错误的根源，并阐明人类受欺骗的所有原因：起初他们相信那些人是神，后来以根深蒂固的信念坚守他们极错误地奉行的宗教仪式。因为，\Person{君士坦丁}皇帝，既然我已证明了这些事物的空虚，并揭露了人类不虔敬的虚荣，我渴望确立唯一的天主的尊威，承担起更有益、更伟大的职责：将人类从弯曲的道路上召回，使他们与自己重归于好，免得他们像某些哲学家那样过于轻视自己，或认为自己软弱、无用、毫无价值，全然是枉然出生的。因为这种观念驱使许多人走向邪恶的追求。当他们幻想我们没有受到任何天主的眷顾，或死后不再存在时，他们便完全放纵自己的情欲；当他们以为这是被允许的，便急切地投身于享乐，不知不觉中陷入死亡的罗网；因为他们不知道什么是人类的合理行为：如果他们想要明白这一点，首先他们会承认他们的主，并追随德行和正义；他们不会使自己的灵魂受制于尘世编造的虚构，也不会寻求他们情欲的致命诱惑；总之，他们会高度珍视自己，并明白在人类身上有比外表更多的东西；而且，除非人类抛弃邪恶并承担对其真父的敬拜，他们无法保持自己的能力和地位。我确实，如我所应当的，常常思考事物的总纲，习惯惊叹那维持并统御万有的唯一天主的尊威竟被如此遗忘，以至于那唯一应当敬拜的对象反而尤其被忽视；人类竟陷入如此盲目，以至于他们偏爱死者胜过真实永生的天主，偏爱那些属于大地、葬于大地者胜过大地本身的创造者。\par

然而，这种人类的不虔若完全出于对神圣名字的无知，或许还可有些宽容。但我们常看到，其他神的敬拜者自己也承认并认识至高神，那么他们能希望为自己的不虔获得什么宽恕呢？他们不承认那位人类不可完全无知之神的敬拜。因为在起誓、表达愿望和感谢时，他们不提及\OriginalTerm{尤皮特}{Jupiter}或众多神祇，而是提及天主；因此，真理藉着本性的力量，甚至从不情愿的胸怀中自发地涌出。而这，确实，并非人类在顺境中的情况。因为那时，天主最常逃离人类的记忆，当他们正享受祂的恩惠时，却应当尊崇祂的神圣仁慈。但若有任何重大的紧急情况压迫他们，那时他们便记念天主。若有战争的恐怖响起，若有疾病的致命力量笼罩他们，若有长期的干旱剥夺了农作物的滋养，若有猛烈的暴风雨或冰雹袭击他们，他们便投奔天主，向天主祈求援助，恳请天主救助他们。若有人在狂暴的海上颠簸，风浪肆虐，他所呼求的就是这位天主。若有人遭受任何暴力侵扰，他呼求祂的援助。若有人陷入极度贫困乞求食物，他只向天主呼吁，并藉着祂神圣而无比的名号寻求赢得人的同情。如此，他们从未记念天主，除非在患难之中。当恐惧离去，危险消退时，他们便迅速前往众神的庙宇：向他们奠酒，向他们献祭，给他们戴上花冠。但对于那位他们在危难中呼求的天主，他们连口头感谢也不给。因此，从顺境产生奢侈；从奢侈，连同所有其他恶习，产生对天主的不虔。\par

我们能设想这从何而来？除非我们想象有一种邪恶的力量，总是敌视真理，以人类的错误为乐，其唯一任务就是不断散布黑暗，蒙蔽人类的心智，使他们看不见光明——总之，叫他们不仰望天空，不观察自身身体的本性，其起源我们将在适当位置叙述；但现在让我们驳斥谬误。因为既然其他动物身体前屈，俯视地面，因为它们没有获得理性和智慧；而造物主天主赐予我们直立姿势和仰起的脸庞，显而易见，向诸神献祭的仪式不符合人类的理性，因为它们使这天生的存在屈尊敬拜地上之物。因为我们那唯一真正的父，当祂创造人——即有智慧且能运用理性的动物时——将人从地上举起，提升到默观其造物主的高度。正如一位天才诗人所精妙描述的：——\par

\Quote{当其他动物俯身看地时，祂赐予人仰起的脸庞，嘱咐他仰望天空，将面容昂然朝向星辰。}\par

希腊人显然从这一情形得出\OriginalTerm{人}{\GreekText{ἄνθρωπος}}之名，因为他向上看。因此，那些不向上看、反而向下看的人，是自我否定，并弃绝了人之名；除非他们以为我们直立的事实是无缘由地赋予人的。天主愿意我们仰望上天，这无疑不是没有理由的。因为飞鸟和几乎所有的哑巴受造物都能看见天，但以直立之姿仰望天穹，却是特别赐予我们的，为使我们能在那里寻求宗教；既然我们无法用肉眼看见天主，就可用心智默观祂，祂的宝座就在那里。这事必然不是朝拜铜石等属地之物的人所能做到的。然而，最错误的是：属身体的、暂时的本性竟能直立，而属灵魂的、永恒的本性却卑躬屈膝；其实，外形和姿势别无他意，只表明人的心智应当与面容朝向同一方向，灵魂应当如身体般正直，好能效法它所当治理的对象。但人们既忘了自己的名，也忘了自己的本性，将目光从天穹移开，盯在地上，畏惧自己双手的造物，仿佛有什么能大于自己的工匠。\par

% structure: p0008 source=h2 target=subsection reason=relative-heading-rank; base=h2

\subsection{第二章：制作像的首要原因是什么；论天主的真像与对祂的真朝拜。}

那么，要么塑造那些他们自己后来会畏惧的对象，要么畏惧自己所造之物，这是何等疯狂？但他们说：我们畏惧的不是像本身，而是按它们的模样所塑造、并以其名字奉献的那些神明。无疑，你之所以畏惧，是因为你认为他们在天上；因为若他们是神，情况不可能相反。那么，你为何不举目望天，呼求他们的名字，在露天献祭？为何你注目墙壁、木头和石头，而不注目你相信他们所在之处？圣殿和祭台是什么意思？至于像本身，它们不过是逝者或不在场者的纪念物，其意义又是什么？因为制作肖像的意图是由人发明的，为的是可能保留那些因死亡而离去或因分离而不在者的记忆。那么，我们将神明归入哪一类？若归入逝者，谁会愚蠢到去朝拜他们？若归入不在场者，那么他们就不应受朝拜，因为他们既看不见我们的行为，也听不见我们的祈祷。但若神明不能缺席——因为既是神明的，他们就在宇宙的任何部分看见和听见一切——那么像就是多余的，因为神明无处不在，只需以祈祷呼求那些能听见我们的名字即可。而若他们临在，就不能不在自己的像旁。人们想象死者的灵魂游荡于坟墓和他们身体的遗骸周围，完全是这样。但一旦神性开始临近，他的雕像就不再需要了。\par

因为我要问：若有人常注视一个定居异乡者的肖像，以此慰藉自己对那缺席者的思念，那么当那人已归来且在场时，他若仍坚持注视肖像，并宁愿享受肖像而非亲眼见到本人，这岂能算是心智健全？当然不是。因为一个人的肖像在远离时才显得必要，而当他在场时就变成多余。但论到天主，祂的神能弥漫各处，永不能缺席，显然肖像总是多余的。然而他们担心，若看不见任何可朝拜的在场之物，他们的宗教就全然虚空，因此他们立起像；既然这些像是死者的代表，它们也与死者相似，全然无感觉。但永生天主的肖像应当是活而有感觉的。若它从相似得名，怎能设想这些没有感觉和运动的像与天主相似？因此，天主的肖像不是人手用石头、铜或其他材料所塑造的，而是人自己，因为他有感觉和运动，并做出许多伟大的事。那些愚人不明白：若像能有感觉和运动，它们反倒会自己朝拜人——是人为它们装饰打扮，若未经人手塑造，它们不过是粗糙未磨的石头、或原始未成形的木头。\par

因此，人应当被视为这些偶像的根源；因为偶像是藉着人的工具而产生的，并且通过人，它们最初具有了形状、形貌和美丽。因此，造偶像的人比被造的对象更优越。然而，没有人仰望造偶像者本身，也没有人崇敬他；人反而畏惧自己所造的东西，仿佛作品比工匠更有能力。因此，塞内加在他的道德论著中正确地写道：\Quote{他们敬拜神明的偶像，屈膝恳求它们，跪拜它们，终日坐在或站在它们旁边，向它们献供，宰杀祭牲；而他们如此珍视这些偶像，却藐视制作它们的工匠。}还有什么比这更矛盾的呢，藐视雕刻家而崇拜雕像，甚至不让制造你们神明的人进入你们的社交圈？它们能有什么力量或能力呢，既然制造它们的人毫无能力？而且，人甚至不能赋予它们自己所有的那些能力——看、听、说话和行动的能力。难道有人会愚蠢到认为神像中有任何东西，而其中除了仅仅的相似性之外，连人的东西都没有吗？但没有人思考这些事；因为人们被这种信念所浸透，他们的头脑彻底吸收了愚昧的欺骗。于是，有感觉的生物敬拜无感觉的东西，理性的生物敬拜无理性的东西，活着的生物敬拜无生命的东西，出自天上的生物敬拜属地的物体。因此，我乐于仿佛站在一座高耸的瞭望塔上，让所有人都能听到，大声宣告佩尔西乌斯的那句话：——\par

\Quote{哦，俯伏于地、缺乏天上之物的灵魂啊？}\par

不如仰望上天，仰望天主你的创造者将你提升到的那景象。他赐给你高昂的面容；你却将它俯向大地；你将那些崇高的心灵压低到属地的事物上，而它们本应与身体一同提升到它们的父那里，仿佛你后悔自己不是生为四足动物。天上的存在不该使自己与属地的事物平等，并倾向大地。你们为何剥夺自己天上的恩惠，并自愿扑倒在地呢？因为你们悲惨地在地上打滚，当你们在这里寻求本应在上天寻求的东西时。至于那些虚浮易碎的产品，人手的工作，无论由何种材料制成，它们除了出自它们的泥土外，还能是什么呢？那么，你们为何使自己屈从于更低的事物？为何将泥土置于你们的头上？因为当你们俯就泥土并自卑时，你们自愿沉入地狱，并定罪于死亡；因为除了死亡和地狱，没有比泥土更低更卑微的了。如果你们想要逃脱这些，你们就会藐视脚下的泥土，保持你们身体正直的姿势，这是你们所接受的，以便你们能够将你们的眼睛和心灵转向那造天地者。但藐视和践踏泥土，无非就是停止崇拜偶像，因为它们是由泥土制成的；也不贪求财富，并藐视身体的快乐，因为财富，以及我们用作寓所的身体本身，也不过是泥土。敬拜活着的存在，好使你们活着；因为那使自己和他的灵魂屈从于死者的人，必然死亡。\par

% structure: p0014 source=h2 target=subsection reason=relative-heading-rank; base=h2

\subsection{第三章　西塞罗及其他有学问的人未能引导人民脱离错误}

但是，这样对粗俗无知的人说话有什么用呢，当我们看到博学明智之人虽然了解这些仪式的虚妄，却因某种乖谬而仍然崇拜他们所谴责的那些对象时？西塞罗很清楚人们所敬拜的神明是虚假的。因为当他讲了许多倾向于推翻宗教仪式的话之后，他却说这些事不应由平民讨论，以免这种讨论会熄灭公共接受的宗教体系。对于这样的人，你能做什么呢？他明知自己错误，却自愿撞向石头，好让所有人都绊倒；或者挖出自己的眼睛，好让所有人都瞎眼？他既不善待别人——他任凭他们陷于错误，也不善待自己——因为他倾向别人的错误，并毫不运用自己智慧的益处，以便在行动中实现自己头脑的构思，而是明知且有意地将脚伸进陷阱，好让他也与其余的人一同被捉住，而作为一个更谨慎的人，他本应解救他们。不，西塞罗，如果你有任何美德，就努力使人民有智慧吧：那是一个合适的主题，你可以倾注你全部的口才力量。因为在如此善良的事业上，无需害怕你的言辞会匮乏，既然你甚至经常以丰富和活力为恶劣的事业辩护。但你确实害怕苏格拉底的监禁，因此你不敢承担真理的辩护。但作为智者，你本应藐视死亡。而且，为良言而死，远比因辱骂而死更荣耀。你的\WorkTitle{腓力比词}的声望，未必比驱散人类的错误并通过你的辩论将人心召回到健康状态更有利于你。\par

但让我们宽容那在智者身上本不应有的胆怯。那么，为什么你自己也陷于同样的错误？我看见你敬拜人手所造的地上之物：你明白它们是虚妄的，却仍然做那些你所承认为最愚蠢的人所做的事。那么，你看见了真相，却既不捍卫也不跟随，这对你有什么益处呢？如果连那些意识到自己处于错误之中的人还甘愿犯错，那么那些无知的平民，他们喜爱空洞的游行，以孩童般的心智凝视一切，岂不更是如此！他们喜爱琐碎之物，被偶像的形态所迷惑；他们无法在自己的心中权衡每一对象，以致理解不了凡属世人眼睛所见之物都不应受敬拜，因为它必然是必朽的。他们看不见天主，这并不奇怪，因为他们甚至看不见他们所相信看见的人。因为，这落入眼睛注意的东西并不是人，而是人的容器；容器的质地和形状并非从容纳它的器皿的轮廓中看出，而是从行为和品性中看出。因此，那些敬拜偶像的人只是没有灵魂的身体，因为他们将自己委身于物质之物，用身体比用心灵看见的还少；而灵魂的职责正是更清楚地感知那些肉眼所不能见的事物。那位哲学家兼诗人严厉地指责那些卑贱、卑微的人，他们违反自己本性的倾向，俯伏敬拜地上之物；因为他说道：——\par

\Quote{他们因对诸神的恐惧而贬抑自己的灵魂，并将其压伏、压制到地上。}\par

当他这样说时，他的意思其实是不同——即什么都不应敬拜，因为诸神不关心人的事务。\par

在另一处，他终于承认，诸神的礼仪和敬拜是一种徒劳的职事：——\par

\Quote{经常以蒙头转向石头，走近每一个祭台，俯伏在地，在诸神的圣所前摊开双手，用大量兽血洒祭台，并接连不断地献上誓愿，这也不是什么虔诚。}\par

而且，如果这些事是无用的，那么崇高尊贵的灵魂就不应当被召唤并压抑到地上，而应当只思想天上的事。\par

因此，虚假的宗教体系已为更明智的人们所攻击，因为他们察觉了其虚假；但真宗教并未被引入，因为他们不知道它是什么、在哪里。所以，他们视之如不存在，因为他们无法在其真理中找到它。就这样，他们陷入了比那些信奉虚假宗教的人更大的错误。因为那些敬拜易碎偶像的人，无论他们多么愚蠢，既然他们将天上的事物置于地上和可朽坏之物中，却仍保留了些许智慧，并且可以得到宽恕，因为他们持守了人的首要职责——如果不是在实际上，至少是在他们的意图中；因为，人和野兽之间的区别，即使不是唯一的，也肯定是最大的区别，在于宗教。但后一类人，按其优越的智慧来说——因为他们理解了虚假宗教的错误——却使自己显得更加愚蠢，因为他们没有设想有任何宗教是真实的。因此，因为评判他人的事务比评判自己的更容易，当他们看到他人的堕落时，却没有注意到自己脚下的陷阱。双方都可发现最大的愚蠢和某种智慧的痕迹；以至于你可能疑惑，哪一方更应被称为愚蠢——是那些接受虚假宗教的人，还是那些不接受任何宗教的人？但（如我所说）宽恕可以给予那些无知且不自称智慧的人；但不能给予那些自诩智慧却显出愚蠢的人。我当然不是如此不公正，以致认为他们能凭自己占卜而发现真理；因为我承认这是不可能的。但我向他们要求的是他们凭理性本身能够做到的事。因为，如果他们既理解某种宗教是真的，又在攻击虚假宗教时公开宣称人们并未拥有那真宗教，那么他们会行事更加明智。\par

但这个想法或许影响了他们：若存在任何真宗教，它必会施展自身、维护其权威，不容许任何与之对立的事物存在。因为他们完全无法看清，真宗教究竟因何故、被何人、以何种方式受到压制——这真宗教原本分有天上的奥秘和天上的秘密。而人若不蒙教导，绝不可能知道这事。事情的要点如下：无知和愚昧之人视假宗教为真，因为他们既不知道真的，也不理解假的。但较精明的人，由于不知道真的，要么坚持那些他们明知是假的宗教，以便显得拥有什么；要么什么都不崇拜，以免陷入错误，然而这正是最大错误的一部分——在人的形像下模仿牲畜的生活。理解何为假，确实是智慧的一部分，不过是人的智慧。人无法超越这步，因此许多哲学家废除了宗教制度，如我已指出的；但认识真理则是神圣智慧的事。人凭自己无法获得这知识，除非由天主教导。因此，哲学家达到了人的智慧的顶峰，从而理解什么不是；但他们未能达到说出真实存在的能力。西塞罗有一句名言：\Quote{但愿我能像驳斥虚假那样轻易地找到真理。}因为这事超越人的处境能力，这项职务的资格就授予了我们——天主已将真理的知识托付给我们；最后四卷书将致力于解释这真理。现在，让我们暂时按已开始的，将虚假之事揭示出来。\par

% structure: p0024 source=h2 target=subsection reason=relative-heading-rank; base=h2

\subsection{第四章 论雕像、圣殿的装饰，以及连外邦人自身对它们的轻视}

那么，雕像能有什么尊威呢？它们完全处在渺小之人的权下，要么被造成别的东西，要么根本不被制造。因此，\Person{普里阿普斯}在\Person{贺拉斯}的诗中这样说：\par

\Quote{从前我是无花果树的树干，一块无用的木材；木匠不知应做一条长凳还是\Person{普里阿普斯}，最后决定让它成为神。因此我是神，对窃贼和飞鸟是极大的恐怖。}\par

谁会对这样的守护者感到安心呢？窃贼竟愚昧到惧怕\Person{普里阿普斯}的形像；而那些他们认为会被镰刀吓走的飞鸟，却落在制作精巧——亦即完全像人——的雕像上，在那里筑巢并玷污它们。但\Person{弗拉库斯}，作为讽刺诗人，嘲笑了人的愚妄。而那些制作雕像的人却自以为在行正经事。总之，那位在其他方面精明的大诗人，唯独在此事上显出愚昧——不像诗人，倒如老妇一般——甚至在他那些最完美的书中，他命令这样做：——\par

\Quote{愿赫勒斯滂的\Person{普里阿普斯}以柳枝镰刀驱赶窃贼和飞鸟，守护它们。}\par

因此，他们崇拜由可朽者所造的可朽之物。它们可能被打碎、焚烧或毁灭。当房屋因年久倒塌，或被焚毁成灰烬时，它们常易碎裂；在许多情况下，若非凭借其自身的巨大，或受勤勉看守保护，它们会成为窃贼的猎物。那么，惧怕那些本身可能因房屋倒塌、火灾或偷窃而受损的对象，是何等疯狂！指望那些不能保护自身者来保护自己，是何等愚昧！求助于那些受伤害时自身不得报复（除非崇拜者代为复仇）者的守护，是何等乖谬！那么，真理在哪里？在于不能对宗教施加暴力的地方；在于显不出任何可被损伤之物的地方；在于不能犯任何亵渎的地方。\par

但凡隶属于眼目和双手之物，因其可朽，实则与整个不朽之事不相容。因此，人用黄金、象牙和珠宝装饰他们的神，仿佛它们能从中得到任何快乐，这是徒劳的。对无感觉之物，贵重礼物有何用？与对死者岂非一样？因为他们用香料和贵重衣料包裹尸体，埋葬在地里，同样也尊崇那些制造时毫无知觉、受敬拜时毫无所知的神；因为它们在祝圣时并未获得感觉。\Person{佩尔西乌斯}不满于将金器带入圣殿，他认为将那些不是圣洁工具而是贪婪工具的东西算作宗教供品是多余的。因为以下是更适于作为礼物献给那应受正确敬拜之神的东西：——\par

\Quote{成文的法律和良心的神圣律法，以及心灵的隐秘深处，和充满高贵的胸怀。}\par

高贵而明智的意见。但他荒谬地补充了这一点：神庙中有这种金子，如同献给维纳斯的由少女奉献的玩偶；也许他因玩偶的微小而鄙视它们。因为他没有看到，由波利克里托斯、欧弗拉诺尔和菲迪亚斯之手用金和象牙雕成的神像本身，不过是大型玩偶，并非由少女奉献——少女的游戏尚可宽容——而是由长着胡须的男人奉献的。因此塞内卡有理由嘲笑甚至老年人的愚蠢。\Quote{我们并不是（如常言所说）两次成为孩童，而是始终如此。但区别在于，成年时我们有更大的玩耍对象。}因此，人们向这些大型的、如同为舞台装饰的玩偶奉献香料、乳香和芬芳；他们向这些玩偶祭献昂贵而肥美的牺牲，而玩偶虽有口却不适宜进食；他们向玩偶奉上袍子和昂贵衣物，尽管它们不需要衣着；他们向玩偶奉献金银，而接受者与施予者一样匮乏。\par

西西里的僭主狄奥尼修斯并非毫无理由地在一次胜利后成为希腊的主宰时，鄙视、掠夺和嘲笑这样的神祇，因为他以戏谑的言语紧随其亵渎行为。当他从奥林匹亚的朱庇特雕像上取下一件金袍时，他命令给雕像穿上一件羊毛袍，说金袍夏天沉重、冬天寒冷，而羊毛袍则适应每个季节。他还从埃斯库拉庇乌斯雕像上取下金须，说这既不合适也不公正，因为他的父亲阿波罗尚且没有胡须、面容光滑，儿子却应在父亲之前长胡须。他又取走雕像伸出的手中所持的碗、战利品和一些小雕像，并说他不是拿走，而是接受；因为拒绝自愿奉上的好东西——而人们习惯于向这些神祇祈求这些好东西——是非常愚蠢和忘恩负义的。他做这些事却未受惩罚，因为他是一位得胜的国王。此外，他惯常的好运也伴随着他；他甚至活到老年，并将王国传给儿子。因此，在他身上，既然人们无法惩罚他的亵渎行为，神祇们本应亲自复仇。但如果任何卑微之人犯了类似罪行，鞭子、火焰、刑架、十字架以及人们在愤怒和狂怒中所能发明的一切酷刑都等在那里惩罚他。但是，当人们惩罚被当场抓获的亵渎者时，他们自己却怀疑神祇的能力。因为如果他们相信神祇有能力报复，为何不把复仇的机会特别留给神祇呢？此外，他们还想象，亵渎的窃贼被发现和逮捕是神祇的旨意；他们的残忍并非出于愤怒，而是出于恐惧，唯恐自己因未能报复对神祇的伤害而受到惩罚。事实上，他们表现出难以置信的浅薄，竟认为神祇会因为别人的罪过而伤害他们，而神祇自己却无法伤害那些亵渎和掠夺他们的人。然而，事实上，他们自己也常常惩罚亵渎者：那可能是偶然发生的，有时如此，但并不总是。但我稍后将说明那是如何发生的。现在，在此期间，我要问：为什么神祇不惩罚狄奥尼修斯身上那么多、那么大的亵渎行为？他公开侮辱神祇，并非暗中进行。为什么神祇不把这个拥有如此权力的亵渎者驱逐出他们的神庙、礼仪和神像？为什么他即使带走了他们的圣物，仍然航行顺利——正如他本人按惯例以玩笑方式作证的那样？\Quote{你们看，}他对害怕沉船的同伴说，\Quote{不朽的神祇自己是怎样赐予亵渎者一次顺利的航行？}但也许他从柏拉图那里学到，神祇没有权力。\par

卡伊乌斯·维勒斯又如何？他的控告者图利将他与同一位狄奥尼修斯、法拉里斯以及所有暴君相比。难道他没有洗劫整个西西里，运走神像和神庙的装饰吗？逐一列举每个细节是徒劳的：我愿提及一件，其中控告者以全部雄辩之力——简言之，用尽声音和身体的每一努力——哀叹关于卡提纳的色列斯，或亨纳的色列斯：前者具有如此大的神圣性，以致男人不得进入她神庙的隐秘内室；后者具有如此大的古老性，以致所有记载都讲述女神本人在亨纳的土地上首先发现谷物，并且她的童贞女儿从同一地点被带走。最后，在格拉古兄弟时代，当国家因骚乱和异兆而动荡不安时，根据西比拉预言发现应当平息最古老的色列斯，便派遣使者前往亨纳。于是，这位色列斯——要么是最神圣者，男人甚至为了朝拜也不得看见她；要么是最古老者，罗马元老院和人民曾用祭献和礼物平息了她——竟被卡伊乌斯·维勒斯派去的强盗奴隶从她的隐秘古老内室中安然无恙地运走了。那位演说家确实在申明西西里人曾恳求他承担该省的案件时，用了以下话语：\Quote{他们城中如今甚至没有任何神明可以投靠，因为维勒斯已从他们最可敬的神龛中取走了最神圣的像。}仿佛若维勒斯从城中和神龛中取走了它们，便也把它们从天上取走了。由此看来，那些神明除了它们被制成的材料之外，别无他物。西西里人向你——哦，马尔库斯·图利乌斯——即向一个人求助，并非没有理由；因为他们三年来已亲身经历那些神明毫无能力。因为若他们向那些不能为自己对卡伊乌斯·维勒斯发怒者寻求庇护以对抗人的伤害，那将是最愚蠢的。但有人会坚称，维勒斯是因这些行为而被判罪的。因此，他并非受神明惩罚，而是因西塞罗的努力而被惩罚，西塞罗或压垮了他的辩护者，或抵挡了他的影响。何必说在维勒斯本人身上，那与其说是判决，不如说是劳作的喘息？因此，正如不朽之神给了狄奥尼修斯运走神明战利品时的顺风航行，他们似乎也赐予维勒斯安宁的休息，使他能在其中平静地享用其亵渎圣物的果实。因为后来当内战肆虐时，他远离一切危险和恐惧，在判决的伪装下听到了他人灾难性的不幸和悲惨的死亡；而那个在众人皆保其位时似乎已跌倒的人，事实上当众人皆倒时唯独他保住了位子；直到三执政的放逐令——正是那个夺走了神明尊严的复仇者图利的放逐令——将他带走，他同时满足于享用由亵渎圣物所得财富和生命，且因年老而衰朽。此外，他在这一点上也是幸运的：在自己死亡之前，他听到了控告者最残酷的结局；神明无疑安排，这亵渎圣物者和其崇拜的掠夺者，不会在从复仇中得到安慰之前死去。\par

% structure: p0035 source=h2 target=subsection reason=relative-heading-rank; base=h2

\subsection{第五章  唯独创造万有的天主当受敬拜，而非元素或天体；并驳斥斯多葛派认为星辰和行星是神明的观点。}

因此，弃绝虚妄无知之物，将我们的目光转向那真实天主的宝座和居所，是何等更好的事！祂将大地悬于坚固根基之上，用闪亮的星辰点缀天空；为人的事务点起太阳这最明亮无匹的光，以证明祂独一的威严；用海洋环绕大地，命令河流永远奔流不息！\par

\Quote{祂也命令平原延伸，山谷下沉，林木披上绿叶，石山隆起。}\par

这一切确实并非一千七百年前出生的尤庇特所作，而是那同一位\Quote{万物的创造者，更美好世界的根源}，祂被称为天主，其开端无法被人领悟，也不应成为探究的对象。对人而言，足以获得圆满和纯全的智慧，只需明白天主的存在就够了：此理解的力与要旨在于，他应仰望并尊敬人类共同的父及奇妙之物的创造者。由此，有些思想迟钝愚昧的人竟崇拜那些既是受造物又毫无感觉的要素，当他们惊叹天主的工程——即那点缀着各种光体的天空，那平原与山峦的大地，那河流、湖泊与泉源的大海——对这些事物感到惊奇，却忘记了那本身看不见的创造者，便开始崇拜和敬礼祂的工程。他们完全无法理解，那位从无中造出这一切的祂是多么伟大和奇妙。而他们看到这些事物，服从于神圣的法律，以永久的必然性服务于人的用途和利益，却仍然视它们为神，对天主的恩惠忘恩负义，以致宁愿偏爱自己的工程，却不去敬拜那位最慈爱的天主和父亲。但若粗野或无知的凡人犯错，又何足为奇？连斯多葛学派的哲学家也持同样见解，认为所有运行的天体都应算作神。正如斯多葛派的卢奇利乌斯在西塞罗著作中所言：\Quote{因此，众星的这种规律性，在不同轨道上历经永恒的这种巨大时间一致性，若无心智、理性和设计，我是无法理解的；当我们看见这些星座中的事物时，我们不得不仅仅将这些天体本身列入众神之数。}而他在此之前不久又说：\Quote{剩下的就是，}他说，\Quote{众星的运动是自主的；看见这些事物的人，若是否认这一点，就不仅是不学无术，而且是不虔诚了。}我们却是坚决否认这一点，并且证明你们，哦，哲学家们，不仅不学无术和不虔诚，而且是盲目、愚蠢和无知的，你们的浅薄已经超过了未受教育者的无知。因为他们只视太阳和月亮为神，而你们却连星辰也算在内。\par

请将星辰的奥秘启示给我们，好使我们可以为每一颗星辰设立祭台和庙宇；使我们知道以何种礼仪、在何日敬拜每一颗星辰，以何名号、以何祷词呼求它们；除非我们或许应当毫无区别地敬拜如此无数的神明，并整体地敬拜如此微渺的神明。我何必提及，他们据此推断所有天体皆为神明的论证，反而导向相反的结论？因为，如果他们基于此——即星辰有固定的、合乎理性的轨道——而想象它们是神明，那他们就错了。因为由此显然可见它们不是神明，因为它们不被允许偏离既定的轨道。但如果它们是神明，它们就会像地上的生物一样，毫无必然性地到处飘荡，因为它们的意志不受约束，每个都随其倾向所引而移动。因此，星辰的运动不是\OriginalTerm{自愿的}{voluntary}，而是\OriginalTerm{必然性}{necessity}的，因为它们遵守为它们制定的规律。然而，当他在讨论星辰的轨道时，尽管他从事物与时间的和谐中理解到它们并非出于偶然，他却判断它们是自愿的；仿佛除非它们含有一种知晓自身职责的\OriginalTerm{理解力}{understanding}，它们就不能以如此秩序和安排运动。哦，对于无知者而言，真理是多么艰难！对于知晓者而言，真理是多么容易！\Quote{如果，他说，星辰的运动不是出于偶然，那么剩下的只能是它们是自愿的；不，事实上，正如它们显然不是出于偶然，同样清楚它们也不是自愿的。}那么，它们在完成轨道时为何保持其规律呢？毫无疑问，天主，宇宙的创造者，如此安排和设计它们，使它们以神圣而奇妙的秩序在天上运行，完成季节交替的变化。\Person{西西里的阿基米德}难道能在空心铜器中制作出宇宙的相似再现，在其中如此安排太阳和月亮，使它们仿佛每天都产生不等的运动，类似天球的旋转，并且那个球体在旋转时不仅展示太阳的靠近和远离，或月亮的盈亏，还展示星辰（无论是固定星还是游星）的不等轨道吗？那么，当人的技艺能够通过模仿再现它们时，天主难道就不能计划和创造原初之物吗？因此，如果一个\OriginalTerm{斯多亚学派}{Stoic}的人看到了那些星辰在铜器中被描绘和塑造的形象，他会说它们凭自身的\OriginalTerm{设计}{design}运动，而不是凭借工匠的\OriginalTerm{天才}{genius of the artificer}吗？因此，星辰中有一种\OriginalTerm{设计}{design}，以适应完成其轨道；但这是天主的\OriginalTerm{设计}{design}，祂创造并治理万有，而不是那些被如此运动的星辰本身的\OriginalTerm{设计}{design}。因为如果天主愿意太阳保持固定，显然就会有永久的白昼。同样，如果星辰没有运动，谁会怀疑会有永恒的黑夜？但是，为了使昼夜交替，祂愿意星辰运动，并且运动得如此多样，以至于不仅有光明与黑暗的相互交替（由此确立劳作与休息的交替周期），还有寒冷与炎热的交替，使不同季节的力量和影响适应庄稼的生产或成熟。因为哲学家们没有看到天主的德能在设计星辰运动中的这种技巧，就以为它们是活的，仿佛它们凭自己的脚步和主动运动，而不是凭借\OriginalTerm{神圣的智慧}{divine intelligence}。但谁不明白天主为何设计它们呢？无疑是为了避免日光隐退后，过于漆黑的夜晚以其污浊可怖的阴暗变得过于压抑，并对生物造成伤害。因此，祂以奇妙的多样性点缀了天空，并以许多微小的光芒调和了黑暗本身。因此，\Person{纳索}的评判岂不是比那些自以为追求智慧、认为那些光是由天主设置以驱散黑暗之阴霾的人更为明智吗？他以这三行诗句结束了那本简短包含自然现象的著作：——\par

\Quote{这些形象，数量如此众多，形态如此多样，天主将它们安置在天上；又将它们散布在阴沉的黑暗中，吩咐它们为寒冷的夜晚带来明亮的光。}\par

但如果星辰不可能是神明，那么太阳和月亮也不能是神明，因为它们与星辰之光只在大小上有别，而非在\OriginalTerm{设计}{design}上。如果这些都不是神明，那包含它们全部的天空也是如此。\par

% structure: p0042 source=h2 target=subsection reason=relative-heading-rank; base=h2

\subsection{\Emph{第六章 整个宇宙及各元素既非天主，亦非拥有生命。}}

同样，如果我们脚下踩踏、开垦耕作以获取食物的土地不是神，那么平原和山脉也不是神；如果这些都不是神，那么整个大地就不可被视为神。同样，如果用于饮水和沐浴以满足活物需求的水不是神，那么流出水的泉源也不是神。如果泉源不是神，那么由泉源汇聚而成的河流也不是神。如果河流也不是神，那么由河流汇成的海洋也不可被视为神。但如果天、地、海——这些世界的组成部分——都不能是神，那么整个世界就不是神；而斯多葛学派却主张世界既是活的又是有智慧的，因此是神。但他们在这一点上自相矛盾，以致他们所说的每一句话又被自己推翻。因为他们这样论证：那从自身产生有感觉之物者，本身不可能无感觉。但世界产生了有感觉的人，因此世界本身也必须有感觉。他们还论证：凡其一部分有感觉者，本身不能无感觉；因此，因为人有感觉，世界作为人的一部分，也就有感觉。这些命题本身是真的：产生有感觉之物者本身有感觉；其一部分有感觉者本身也有感觉。但他们用以推导结论的前提却是假的：因为世界并不产生人，人也不是世界的一部分。因为创造世界的那同一天主，也从起初创造了人；人并不是世界的一部分，如同肢体是身体的一部分那样；因为世界可以没有人而存在，如同城市或房子可以没有人。正如房子是一个人的居所，城市是一个民族的居所，世界也是全人类的居所；居住者是一回事，被居住者是另一回事。这些人急于证明他们错误假定的东西——即世界有感觉且是神——却没有意识到他们自己论证的后果。因为如果人是世界的一部分，并且因为人有感觉而世界有感觉，那么由此必然得出，因为人是可朽的，世界也必然是可朽的，不仅是可朽的，而且还会遭受各种疾病和痛苦。反之，如果世界是神，那么它的各部分也显然是不朽的：因此人也是神，因为他如你所说，是世界的一部分。如果人是神，那么驮兽、牲畜、其他各类走兽、飞禽和鱼类也都是神，因为它们同样有感觉，且是世界的一部分。这还可以容忍，因为埃及人甚至崇拜这些东西。但结果会发展到连青蛙、蚊子和蚂蚁都成了神，因为这些也有感觉，且是世界的一部分。因此，从虚假源头推出的论证总是导致愚蠢荒谬的结论。何必再提这些哲学家断言世界是为神灵和人类共同的居所而建造的呢？因此，如果世界是被造的，它既不是神，也不是活的：因为活物不是被造的，而是诞生的；如果它被建造，就像房屋或船只被建造一样。所以世界有一位建造者，就是天主；被造的世界与造它的那位不同。他们一面断言天上的火焰和世界的其他元素是神，一面又说世界本身是神，这是多么不一致和荒谬！怎么可能从一大堆神中拼凑出一个神来？如果星辰是神，那么世界就不是神，而是诸神的居所。但如果世界是神，那么其中的一切事物就不是神，而是神的肢体，这些肢体显然不能单独具有神的名号。因为没有人能正确地说一个人的肢体是许多人；然而，活物与世界之间并没有类似的比拟。因为活物有感觉，它的肢体也有感觉；除非它们与身体分离，否则不会失去感觉。但世界与此有何相似之处？实际上，他们自己告诉我们，因为他们并不否认世界是被造的，为要成为神和人共同的居所。因此，如果它被建造为居所，它本身就不是神，构成它的元素也不是神；因为房屋不能统治自己，构成房屋的部分也不能。因此，他们不仅在真理面前，甚至在他们的言论中也被驳倒。因为正如房屋为了居住而被建造，本身没有感觉，并服从于建造或居住它的主人；同样，世界本身没有感觉，服从于创造它的天主，祂为自己的用途创造了它。\par

% structure: p0044 source=h2 target=subsection reason=relative-heading-rank; base=h2

\subsection{第七章 论天主与愚人的宗教仪式；论贪婪与祖先的权威。}

因此，愚昧者在两方面犯错：首先，他们偏重受造物，即天主的化工，而轻视天主本身；其次，他们以人形敬拜这些受造物本身的各种形象。因为他们仿照人的模样，塑造了太阳和月亮的像，还有火、土、海的像，分别称为\OriginalTerm{伏尔甘}{Vulcan}、\OriginalTerm{维斯塔}{Vesta}和\OriginalTerm{尼普顿}{Neptune}。他们并不公然向受造物本身献祭。人们竟如此迷恋偶像，以至于真实的事物反被认为价值较低：他们喜爱的，实是黄金、宝石和象牙。这些东西的美丽和光辉眩惑了他们的眼睛，他们认为，凡没有这些闪耀之处，便没有宗教。于是，在敬拜诸神的幌子下，他们敬拜了贪婪和欲望。因为他们相信，诸神喜爱他们所贪求的一切，正是这一切使人每日发生偷窃、抢劫和凶杀，导致战争毁灭了全世界的民族和城邦。因此，他们将战利品和掠夺物奉献给诸神。这些神必定软弱无力，且缺乏至高之善，若他们竟受欲望支配的话。因为若他们渴望从地上得来什么，我们怎能视他们为天上的？若他们有所匮乏，怎能算为有福？若他们喜爱那些东西——而人追求这些东西时其欲望并非全然无过——他们怎能算为纯洁无瑕？因此，人们接近诸神，与其说是出于宗教（因为在用不正当手段获得且会朽坏的事物中，宗教不可能有位置），不如说是为了观看黄金，欣赏磨光的大理石和象牙的光彩，以不倦的目光审视那些饰以宝石和彩色的衣袍，或是镶满闪闪发光的宝石的杯盏。庙宇越是装饰华丽，神像越是精美，人们就认为它们越有威严：他们的宗教完全局限于人所赞赏的那些欲望之物。\par

这些就是他们的祖先传给他们的宗教仪轨，他们以极大的固执坚持维护并捍卫它们。他们并不考虑这些仪轨的性质如何；仅仅因为古人传下了它们，他们就确信其卓越和真实；古人的权威如此之大，以至于被认为质疑它便是罪过。因此，到处都把它们当作确定无疑的真理来信奉。简而言之，西塞罗书中科塔对卢基里乌斯这样说：\Quote{巴尔布斯，你知道科塔的看法，你也知道大司祭的看法。现在让我听听你的意见：既然你是哲学家，我应当从你那里得到你宗教的理由；但在我祖先的情况下，相信他们是合理的，即使他们没有提出理由。}你若相信，为何还需要理由呢？理由可能使你不再相信。但若你需要理由，并认为这个问题需要探究，那么你并不相信；因为你探究的目的，是查明之后才去遵从。看，理性告诉你，诸神的宗教仪轨并不真实：你将如何？你宁愿遵循传统，还是理性？这一点，其实并非由他人传授给你，而是你自己发现并选择的，因为你完全根除了所有宗教体系。你若偏好理性，就必须放弃祖先的仪轨和权威，因为除了理性所规定的，没有什么是正确的。但若虔敬劝你追随祖先，那么你就得承认，那些遵循违反理性的宗教发明的人是愚蠢的；而你自己，既然崇拜你已证明为虚假的东西，也是无知的。既然祖先的名号如此强烈地被用来反对我们，我请求让我们看看，那些祖先到底是什么人——据说不容背离他们的权威。\par

罗慕路斯在将要建城时，召集了与他一同长大的牧羊人；由于他们的人数似乎不足以建立城市，他便设立了一个避难所。所有最堕落的人，不论身份如何，毫无区别地从邻近各处蜂拥而至。他就这样从这些人中聚集了民众；他挑选最年长者进入元老院，称他们为\Quote{父老}，以便凭他们的建议指导一切。关于这个元老院，哀歌诗人普罗佩提乌斯这样说道：\par

\Quote{号角曾召唤古老的基里特人集会；\newline{} 那百人在田野里常组成元老院。\newline{} 元老院如今高耸，因着盛装的元老而闪耀，\newline{} 却曾接纳身披兽皮、心灵质朴的父老。}\par

这些人就是那些父老，博学明达的人以最大的虔诚服从他们的法令；而所有后代都必须认为，一百个身穿兽皮的老人随心所欲制定的东西，便是真实而不可更改的；然而，正如在第一卷中所提到的，他们被蓬皮利乌斯引诱，相信了他亲自传授的那些神圣仪轨的真实性。有什么理由使他们的权威被后世如此尊崇呢？因为他们生前，无论贵贱，都没有人认为他们值得联姻。\par

% structure: p0050 source=h2 target=subsection reason=relative-heading-rank; base=h2

\subsection{第八章 论理性在宗教中的运用；兼论梦境、占兆、神谕及类似异兆。}

因此，尤其在生命整个计划所系之事上，每个人都应当自信，运用自己的判断和个体能力去探究和权衡真理，而不应因信任他人而被他们的错误所欺骗，好像自己毫无理解力一样。天主已将智慧同样赐予众人，使他们既能探究所未闻之事，也能权衡所已闻之事。时光在前者并不比我们在智慧上更优越；因为如果智慧同样赋予所有人，我们就不可能被先于我们的人所超越。智慧不可缩减，如同太阳的光辉一样；因为太阳是眼睛的光，智慧是人心灵的光。因此，既然智慧——即追求真理——对所有人都是自然的，那些不加任何判断就赞同祖先的发现、像羊一样被他人带领的人，便剥夺了自己的智慧。但他们没有注意到，由于引入了祖先之名，他们认为自己作为后裔不可能有更多的知识，而前者作为祖先也不可能是愚昧的。那么，什么阻止我们以他们为先例呢？既然他们将虚假的发明传给后代，我们也一样，发现真理的我们岂不能将更好的事物传给后代？因此，还有一个伟大的探究主题，其讨论并非来自才能，而是来自知识；这必须更详细地解释，以免留下任何怀疑。因为或许有人会求助于许多无可置疑的权威所传承的事：那些我们已证明不是神的人物，常常通过异兆、梦境、占卜和神谕显示其威严。的确，可以列举许多奇妙之事，尤其是这一件：最精湛的占卜师\Person{阿克修斯·纳维乌斯}，当他警告\Person{老塔奎尼乌斯}未经占卜同意不得着手任何新事时，国王贬低其技艺的可信度，命他去询问飞鸟，然后告知是否能实现他自己心中所设想之事；纳维乌斯断言可能；国王说：\Quote{那么拿着这块磨刀石，用剃刀把它切开。}而纳维乌斯毫不犹豫地拿起并切开了它。\par

其次是\OriginalTerm{卡斯托尔和波吕克斯}{Castor and Pollux}在拉丁战争期间被看见在\Place{尤图尔纳湖}旁洗刷战马的汗水，当时他们毗邻该泉的神庙自行开启。在马其顿战争中，据说同样的神祇，骑着白马，在夜间向前往\Place{罗马}的\Person{普布利乌斯·瓦提埃努斯}显现，宣告珀修斯王已于当日被击败并被俘；几天后从\Person{保卢斯}收到的信件证实了此事。同样奇妙的是，\OriginalTerm{福尔图娜}{Fortune}女神雕像以女性形象据说曾不止一次说话；还有\OriginalTerm{朱诺·莫内塔}{Juno Moneta}的雕像，当\Place{维爱}被攻陷时，一名被派去搬迁它的士兵戏谑地问她是否愿意迁往罗马，她回答说愿意。\Person{克劳迪娅}也被作为奇迹的例子。因为当依据\OriginalTerm{西卜林书}{Sibylline books}，\OriginalTerm{伊代亚母亲}{Idæan mother}被请来，载着她的船在\Place{台伯河}浅滩搁浅，无法用任何力量移动时，传说一直因过度打扮而被视为不贞的克劳迪娅，双膝跪地，恳求女神若她自认贞洁，就跟随她的腰带；于是所有壮汉无法移动的船，竟被一个女子移动了。同样奇妙的是，在瘟疫流行期间，从\Place{埃皮达鲁斯}请来的\OriginalTerm{埃斯库拉庇乌斯}{Æsculapius}据说使罗马城摆脱了长期持续的瘟疫。\par

还可以提到亵圣者，因为神祇被认为通过立即惩罚他们来报复所受的侵害。监察官\Person{阿皮乌斯·克劳狄乌斯}不顾神谕的劝告，将\OriginalTerm{赫丘利}{Hercules}的圣礼转交给公奴，因而双目失明；放弃其特权的波提提乌斯氏族在一年内绝嗣。同样，监察官\Person{富尔维乌斯}，当他从\OriginalTerm{拉基尼亚的朱诺}{Lacinian Juno}神庙取下大理石瓦片，以覆盖他在\Place{罗马}建造的\OriginalTerm{骑士福尔图娜}{equestrian Fortuna}神庙时，失去理智，并因在\Place{伊利里库姆}服役的两个儿子丧生而遭受极大的心灵痛苦。此外，\Person{马克·安东尼}的副将\Person{图鲁利乌斯}，在\Place{科斯岛}砍伐了\OriginalTerm{埃斯库拉庇乌斯}{Æsculapius}的一片圣林并建造了一支舰队后，后来在同一地点被\Person{凯撒}的士兵杀害。这些例子还包括\Person{皮洛士}，他从\OriginalTerm{洛克里斯的普洛塞庇娜}{Locrian Proserpine}宝库中取走金钱，结果船舶失事，被冲到女神神庙附近的海岸上，以至于除了那笔金钱外，无一物完好无损。\OriginalTerm{米利都的刻瑞斯}{Ceres of Miletus}也在人们中获得了极大的敬畏。因为当城市被\Person{亚历山大}攻陷，士兵冲入她的神庙劫掠时，突然喷出的一团火焰使他们都失明了。\par

也发现有梦似乎显示诸神的能力。据说\Person{朱庇特}在睡梦中向平民\Person{提比略·阿提尼乌斯}显现，命令他向执政官和元老院宣布，在最近的赛马竞技中，一位公开舞者使他不悦，因为有位名叫\Person{安东尼·马克西姆斯}的人，在竞技场中央用叉刑架严厉鞭打一个奴隶并带他受罚，因此应重办竞赛。当他忽略了这一命令，据说同日他失去了儿子，自己也被重病缠身；他再次见到同一形象询问他是否因忽略命令而受了足够的惩罚，便被人用轿子抬到执政官那里；他在元老院中解释了一切后，恢复了体力，步行回到了家。那个据说\Person{奥古斯都·凯撒}的性命仰赖的梦也非比寻常。因为当他在与\Person{布鲁图斯}的内战中患重病，并决定不参加战斗时，\Person{弥涅耳瓦}的形象显现于他的医生\Person{阿托里乌斯}，建议他凯撒不应因身体虚弱而留驻营中。于是他被人用轿子抬到军中，同日营地被\Person{布鲁图斯}攻占。许多其他类似性质的例子可以举出；但我担心，如果我在陈述相反主题上拖延太久，我可能显得忘记了自己的目的，或招致多言的指责。\par

% structure: p0055 source=h2 target=subsection reason=relative-heading-rank; base=h2

\subsection{第九章 论魔鬼、世界、天主、眷顾、人及其智慧。}

因此我将阐述这一切的方法，使困难和隐晦的主题更易理解；我将揭露所有假神的欺骗，世人受其引导已远离真理之道。但我要从源头追溯此事；若有任何不谙真理、无知之人阅读此书，他可得教导，明白这些灾祸的\Quote{根源与起源}究竟是什么；并藉着光明，认出自己及全人类的错误。\par

既然天主在谋划上拥有最大的远见，在行动执行上拥有最大的技能，在祂开始这世界的工程之前——因为祂内存在并永远存在圆满、最完美良善的泉源——为使良善如溪流从祂而出，流向远方，祂产生了一位与自己相似的神体，充满天主圣父的完美。至于祂如何愿意如此，我将在第四卷中努力说明。然后祂造了另一位，神圣起源的气质并未保存在他内。因此他如中毒般被自己的嫉妒感染，从善转为恶；并凭自己的意志（天主赐予他的自由意志）他为自己得了一个相反的名称。由此可知，一切邪恶的根源是嫉妒。因为他嫉妒祂的前辈，后者因坚毅而蒙天主圣父悦纳亲爱。这个由善变恶的存在，希腊人称之为\OriginalTerm{魔鬼}{diabolus}；我们称他为控告者，因为他向天主报告那些他引诱我们去犯的过失。因此，当天主开始建造世界时，祂将那位首要而最伟大的圣子置于整个工程之上，同时用祂作为谋士和工匠，在规划、安排和完成时，因为祂在知识、判断和能力上都是圆满的；关于祂我现在谈得较少，因为在别处祂的卓越、祂的名字和祂的本性都必须由我们阐述。不要有人问天主用什么材料造了这些伟大而奇妙的工程：因为祂从无中造了一切。\par

不应听信诗人们说起初是一团混沌，即物质和元素的混乱；但天主后来分开了那整团，将每个物体从混乱堆积中分开，并按秩序排列，建造并装饰了世界。如今很容易回答这些不了解天主能力的人：因为他们相信祂除了从已存在和已准备好的材料之外，不能产生任何东西；哲学家们也陷入了这种错误。因为\Person{西塞罗}在讨论诸神的本性时这样说：\Quote{首先，因此，说万物所由来的质料是由天主眷顾所造，这是不合理的；而是它具有并一直具有自身的能力和本性。因此，正如建造者准备盖房时，自己不制造材料，而是使用已备好的，雕刻家也使用蜡一样；同样，天主眷顾本应有现成的材料，非自己所造，而是已备好供使用的。但如果质料不是天主所造的，那么地、水、气、火也不是天主所造的。}喔，这十行里有多少错误！首先，那个在他几乎所有的其他论辩和著作中都维护天主眷顾、并用非常犀利的论证攻击那些否认眷顾存在的人，如今他自己却像叛徒或逃兵一样试图取消眷顾；对于这样的人，你若想驳斥他，既无需思考也无需费力：只需提醒他自己的话。因为任何人都不可能比\Person{西塞罗}自己更有力地反驳\Person{西塞罗}了。但让我们对\OriginalTerm{学院派}{Academics}的习惯和实践作出让步，即允许人们非常自由地说话，并持有他们愿意的任何观点。让我们考察这些观点本身。他说，质料由天主所造，这是不合理的。你用什么论据证明这一点？因为你对它不合理没有给出任何理由。因此，相反，我觉得它非常合理；而且我这样觉得也并非毫无理由，当我想到在天主内有某种更多的东西，而你确实将祂降低到人的软弱，除了一般的施工外，你不允许祂有任何其他东西。那么，如果天主也像人一样需要别人的帮助，那神圣的能力与人有何区别？但祂确实需要帮助，如果祂除非由别人提供材料，否则就无法建造任何东西。但如果是这样，显然祂的能力不完全，而准备材料的人必定更有能力。那么，用什么名字来称呼那在能力上超越天主的人呢？——因为制造自己的东西比安排别人的东西更伟大。但如果不可能有任何东西比天主更有能力——祂必须具有完美的能力、力量和智慧——那就可以推出，那制造由质料组成之物者也制造了质料。因为任何东西若没有天主的权能运作或违背祂的意愿而存在，既不可能也不相宜。但他又说，质料具有并一直具有自身的能力和本性，这是合理的。若没有人赋予它，它哪来的能力？若没有人产生它，它哪来的本性？如果它有能力，那能力是来自某人。但它能从谁那里得到能力，除非来自天主？此外，如果它有本性——这显然因被产生而得名——它必然是被产生的。但它能从谁那里获得存在，除非从天主？因为你说万物有它们的起源的本性，如果它没有理智，就不能造出任何东西。但如果它有产生和制造的能力，那么它就有理智，必定是天主。因为那既有预见去计划、又有技巧和力量去实现的能力，不能称为别的名字。因此，所有\OriginalTerm{斯多葛派}{Stoics}中最聪明的\Person{塞内加}说得更好，他看出\Quote{自然无非就是天主。}因此他说：\Quote{我们岂不应当赞美那拥有自然美善的天主？}因为祂没有向任何人学习。是的，真的，我们要赞美祂；因为虽然这对祂来说是自然的，但祂将这自然赐给了自己，因为天主本身就是自然。因此，当你将万物的起源归于自然，而从天主那里拿走时，你就处于同样的困境：——\par

\Quote{你通过借钱来还债，\Person{格塔}。}\par

因为仅仅改变名称，你就清楚地承认它是由你所否认的那一位造的。\par

接着是一个最无意义的比较。\Quote{正如建筑者，}他说道，\Quote{当他将要建造任何建筑物时，并不自己制造材料，而是使用已经准备好的材料，同样雕塑家也使用蜡；因此，天主的\OriginalTerm{眷顾}{providence}应该手边有材料，不是自己生产的，而是已经准备好供使用的。}绝非如此；因为如果天主从已经提供的材料中制造，那是人的作为，祂的能力就会更小。如果没有木头，建筑者什么也建不起来，因为他自己不能制造木头；不能这样做正是人类软弱的特征。但天主为自己制造材料，因为祂有能力。因为有能力是天主的属性；如果祂不能，祂就不是天主。人从已经存在的东西中产生他的作品，因为由于他的可朽性他是软弱的，并且由于他的软弱他的能力是有限且适度的；但天主从不存在的东西中产生祂的作品，因为由于祂的永恒祂是强壮的，并且由于祂的力量祂的能力是无限的，没有终点或极限，如同造物主自己的生命一样。那么，如果天主在将要创造世界时，首先准备了用来创造世界的材料，并从不存在中准备它，这有什么可惊奇的呢？因为天主不可能从其他来源借用任何东西，既然万有都在祂内并出自于祂。因为如果有任何东西在祂之前，并且有任何东西被造却不是由祂造的，那么祂将失去天主的能力和名号。但也许有人说，\OriginalTerm{质料}{matter}从未被造，如同天主一样，天主用质料创造了这个世界。那样的话，就会建立两个永恒的根源，并且彼此对立，这不可能不发生冲突与毁灭。因为那些具有相反力量和方式的必然要碰撞。这样，如果它们彼此对立，两者都不可能是永恒的，因为一个必须压制另一个。因此，永恒之物的本性不可能是别的，只能是一，以致万有从那根源如同从泉源流出来。所以，要么天主出自质料，要么质料出自天主。这两者中哪一个更真实，容易理解。因为在这两者中，一个具有感觉，另一个是\OriginalTerm{无感觉的}{insensible}。制造任何东西的能力，只能存在于具有感觉、理智、思想和运动能力的东西中。任何东西除非在它存在之前就由理性预见了如何制造，并且在它被制造之后如何持续，否则就不能开始、制造或完成。简而言之，只有那愿意制造并有手来完成他所愿意的人才能制造东西。但那无感觉的东西总是静止不动，麻木；在没有任何自愿运动的地方，什么也不能起源。因为如果每个动物都有理性，那么可以肯定它不能从缺乏理性的东西中产生，也不能从其他来源接受那在原初源泉中不存在的东西。然而，不要让任何人感到困扰，有些动物似乎从土里生出。因为土地本身并不生出这些，而是靠天主的圣神，没有祂什么也不能产生。所以天主不是从质料中产生的，因为一个\OriginalTerm{有感觉的}{sensible}存在绝不能从一个无感觉的存在中产生，一个有智慧的存在从一个无理性的存在中产生，一个不能受苦的存在从一个能受苦的存在中产生，一个无形体的存在从一个有形体的存在中产生；而是相反，质料出自天主。因为任何由坚固且可触摸的物体组成的东西都容许外力。那容许外力的东西能够分解；那分解的东西就毁灭；那毁灭的东西必然有一个起源；那有起源的东西有一个它由之起源的源头，即某个制造者，他是有理智的、有预见力的、并精于制造。有一个确定无疑的，不是别的，就是天主。既然祂具有感觉、理智、眷顾、能力和活力，祂能够创造和制造有生命的和无生命的事物，因为祂有制造一切的方法。但质料不能一直存在，因为如果它一直存在，它就是不可改变的。因为那一直存在的东西不会停止一直存在；那没有开端的东西必然没有终结。而且，那有开端的东西没有终结，比那没有开端的东西有终结更容易。所以，如果质料不是被造的，那么什么也不能从它被造。但如果什么也不能从它被造，那么质料本身就不能存在。因为质料就是那用来制造某种东西的东西。但一切用来制造任何东西的东西，既然它接受了工匠的手，就被摧毁，并开始成为别的东西。因此，既然在质料被用来制造世界时它有了终结，它也就有了开端。因为那被毁灭的东西先前被建造；那被松解的东西先前被结合；那被终结的东西是被开始的。所以，如果从它的变化和终结中推断出质料有开端，那么那个开端能来自谁，除了来自天主？因此，只有天主是未被造的存在；所以祂能毁灭其他东西，但祂自己不能被毁灭。那在祂内的一直会持久，因为祂不是从任何其他源头产生或出现；祂的出生也不依赖于任何其他对象，这对象的改变可能导致祂的分解。祂来自祂自己，如我们在第一卷所说；所以祂就是祂愿意祂是的那样，不能受苦、不变、不朽、有福、永恒。\par

但是现在，\Person{图理乌斯}用来结束这一思想的结论，要荒谬得多。他说：\Quote{但是，如果物质不是由天主造的，那么地、水、气、火，也就不是由天主造的。}他多么巧妙地避开了危险！因为他说前一点时，仿佛它不需要证明，实际上它比那条引发此说法的论据更不确定。他说：如果物质不是由天主造的，那么世界就不是由天主造的。他宁愿从虚假的前提引出虚假的推论，而不从真实的前提引出真实的推论。虽然不确定的事情应当由确定的事情来证明，他却从一个不确定性中提取证据，来推翻一个确定的事实。\newline{}因为世界是由天主眷顾所造——且不提宣告此事的\Person{特里斯墨吉斯图斯}，也不提宣告同一事的\Person{西彼拉}的诗句，也不提那些异口同声、和谐一致地见证世界被造且是天主之工的众先知——就连哲学家们也几乎都一致赞同；因为这是每一学派的主要代表——毕达哥拉斯学派、斯多亚学派和逍遥学派的意见。简而言之，从最初的七位贤人直到\Person{苏格拉底}和\Person{柏拉图}，这一点都被视为公认且确定无疑的事实；直到许多世代之后，疯狂的\Person{伊壁鸠鲁}出现了，他独自一人敢于否认这一最明显的事实，无疑是为了追求新奇，好以他自己的名字创立一个学派。因为他找不到任何新东西，为了显得与他人不同，他想要推翻旧有的观点。但是，所有在他周围狂吠的哲学家都驳斥了他。\newline{}因此，世界由天主眷顾安排，比物质由天主眷顾收集，更为确定。所以他本不应因为世界的物质不是由天主眷顾造的，就假定世界不是由天主眷顾造的；而应当因为世界是由天主眷顾造的，就推论出物质也是由天主造的。因为相信物质是由天主造的——因为他是全能的——比相信世界不是由天主造的——因为没有任何东西能在没有心智、理智和设计的情况下被造——更合理。但这并非\Person{西塞罗}的过错，而是他所属学派的过错。因为当他承担了一场辩论，意欲取消哲学家们胡言乱语的神明本性时，由于他对真理无知，他想象必须完全取消神性。因此他能够取消那些神明，因为它们并不存在。但是当他试图推翻那唯一天主的天主眷顾时，因为他开始与真理作对，他的论证失败了，他必然陷入了这个他无法自拔的陷阱。\newline{}那么，在这里，我牢牢地抓住了他；我把他钉在了原地，因为争辩另一方的\Person{卢西利乌斯}沉默了。那么，这里就是转折点；一切都取决于此。让\Person{科塔}自己从这一困境中解脱出来吧，如果他能够的话；让他提出一些论据，证明物质是永恒存在的，没有任何天主眷顾创造了它。让他展示任何沉重或轻飘的东西如何能够没有作者而存在，或如何改变，以及那始终存在的东西如何停止存在，使那从未存在的东西开始存在。如果他能证明这些，那么我才会承认世界本身并非由天主眷顾建立；然而即使这样承认，我也会用另一个圈套抓住他。因为他会再次转回到他不想回到的同一个点上，即他会说世界所由构成的物质和由物质构成的世界，都是由本性存在的；尽管我主张本性本身就是天主。因为除了拥有理智、远见和能力的那一位，没有人能做出奇妙的事，即具有最大秩序的事物。因此最终会看到，天主造了万物，没有任何东西能够存在而不从天主获得其根源。\par

但是同一个人，每当他追随伊壁鸠鲁派，不承认世界是由天主所造时，就习惯于追问：天主用什么样的手、用什么样的机器、用什么样的杠杆、用什么样的装置，造出了如此宏伟的工程？假如他能活在天主造世界的时候，他或许能看到。但是，为了使人不能窥探天主的工程，天主不愿在他完成一切之前把他带到这个世界里来。然而，人当时是不可能被带进来的：因为当上天正在建造、大地的基础正在奠定之时，当潮湿的东西可能因极度的僵冷而凝结成冰，或因烈火煮沸而变硬时，人如何能存在呢？或者，当太阳尚未建立、谷物和动物尚未产生时，人如何能生活呢？因此，必须最后造人，当时对世界和万物的最后润色已经完成。最后，\WorkTitle{圣经}教导说，人是天主的最后工程，他被带进这个世界，如同进入一座预备妥当、布置好的房屋；因为万物都是为他造的。诗人们也承认这一点。\Person{奥维德}在描述了世界的完成和其他动物的形成之后，补充道：\par

\Quote{还有一种比这些更神圣、更能容纳崇高心灵的动物尚付阙如，它能够统治其余的一切。人产生了。}\par

所以我们应认为，探查天主希望保密的事，是极为不敬的！但他（指异端者）的探询并非出于渴望听闻或学习，而是为了驳斥；因为他确信无人能断言这一点。仿佛真的可以认为，这些事物不是由天主创造的，只是因为我们无法清楚看见它们是如何被造的！如果你从小在一栋建造精良、装饰华美的房屋中长大，却从未见过工场，你会因为不知道那房屋是如何建造的，就认为它不是人所建造的吗？关于房屋，你必然会问同样的问题，正如你现在问关于世界的问题——用什么样的手、什么样的工具，人才能完成如此伟大的工程？尤其是如果你看到巨大的石块、宏伟的巨石、庞大的柱廊、整个建筑高耸巍峨，这些难道不会在你眼中远超人类力量的程度吗？因为你不知道，这些事物的完成更多依靠技巧和才智，而非蛮力。\par

但若人（他里面没有什么是完美的）尚能凭借技巧完成超出其脆弱力量所能及的事，那么当有人声称世界是由天主（祂是完美的，智慧无边，力量无限）所造时，这为何在你看来竟不可信呢？祂的工程为眼目所见；但祂如何造成它们，甚至心智也不得见，因为——正如\Person{Hermes}所说——有死者不能接近（即不能更近地趋近并以理解力跟随）不朽者，暂时的不能接近永恒者，可朽的不能接近不可朽者。正因如此，属地的生物暂时还不能感知天上的事，因为它被身体所禁锢和束缚，仿佛被关在监牢里，无法以自由和不受限制的知觉辨别一切。因此，让他明白，那些探究不可描述之事的举动是何等愚妄。因为这就是超越自身本性的界限，不明白人允许接近到什么程度。简言之，当天主向人启示真理时，祂只愿意我们知道那些为了获得生命而关乎人的事；但对于那些属于渎神和焦躁好奇的事，祂保持沉默，使它们成为秘密。那么，你为何探究那些你不能知道的事呢？即使你知道，你也不会更幸福。在人身上，完美的智慧在于知道只有一位天主，并且万物都是由祂所造。\par

% structure: p0067 source=h2 target=subsection reason=relative-heading-rank; base=h2

\subsection{\Emph{第10章 论世界及其部分、元素与季节}}

如今，既然已经驳斥了那些对世界及创造世界的天主持有错误见解的人，让我们回到关于世界神圣构造的讨论，这是我们神圣宗教的圣经所教导我们的。因此，首先，天主创造了天，并将它高悬在上，作为祂自己——创造者——的居所。然后祂奠定了地，将它安置在天下，作为人与其他动物族类的住所。祂愿意地由水环绕并维系。祂用明亮的光体装饰并充满了自己的居所；用太阳、月亮的灿烂圆盘，以及闪烁星群的闪耀标记来装饰天；但在世上，祂安置了与这些光体相反的黑暗。因为地本身不含光明，除非从天那里接收；在祂的居所中，祂安置了永恒的光、众神和永生的生命；相反，祂在地上安置了黑暗、下界的居民和死亡。因为这些事物与前者相隔多远，恶事与好事、恶习与美德就相隔多远。祂又设立了地本身两个相反且性质不同的部分——即东和西；其中东归天主，因为祂自己是光明之泉，万物的启明者，并且祂使我们复活升入永生。而西归给那扰攘败坏的心智，因为它隐藏光明，总是带来黑暗，并使人在罪恶中死去和灭亡。因为光明属于东，生命的全部历程依赖于光；所以黑暗属于西；而死亡与毁灭包含在黑暗中。然后祂同样度量了其他部分——即南和北，这两个部分与前两者密切相连。因为那更被太阳温暖照耀的部分，最临近并与东方紧密相连；而那因寒冷和永久的冰而麻木的部分，则属于与极西相同的区域。因为黑暗与光明相对，寒冷与炎热相对。因此，正如炎热最接近光明，南也最接近东；正如寒冷最接近黑暗，北也最接近西。祂又为这些部分各自指定了自己的时间——即春给东方，夏给南方，秋给西方，冬给北方。在这两个部分（南和北）中，也包含了生命和死亡的图像，因为生命在于热，死亡在于冷。热由火产生，冷由水产生。按照这些部分的划分，祂又创造了昼夜，使它们通过交替更迭完成岁月的时间历程和永恒周转，我们称之为年。由最初东方所提供的白昼，必须归于天主，正如所有性质更优的事物都归于祂一样。但由极西所带来的黑夜，则确实归于我们所说过的天主的敌手。\par

天主在造这些事物时，也顾念了将来；祂如此造它们，是为了借此显示真宗教与假迷信的样式。因为正如每日升起的太阳——虽然它只有一个（\Person{西塞罗}因此认为它被称为太阳，因为众星被遮蔽，唯它独现）——然而，由于它是真光，完美圆满，热力极强，以最明亮的光辉照耀万物；同样，天主虽是唯一，却拥有完美的尊威、能力和光辉。但黑夜——我们说是归于天主的那位败亡对手的——由相似性表明属于他的众多而不同的迷信。因为虽然无数的星辰似乎闪烁发光，但由于它们不是圆满稳固的光，也不发出热力，其数量也不能战胜黑暗，因此这两样东西被视为最重要，它们具有彼此不同且对立的能力——热与湿，天主奇妙地安排它们以维持和产生万物。因为天主的德能在于热与火，若祂不借掺和冷湿之物来缓和其炽烈和力量，则无物能生或存，凡开始存在者必立即被焚烧而毁。因此有些哲学家和诗人说世界是由不协调的和谐构成的；但他们并未彻底领悟此事。\Person{赫拉克利特}说万物由火产生；\Person{米利都的泰勒斯}说由水产生。两人都看到了真理的某些方面，但也都错了：因为若只有一种元素，则火不能从水产生，水也不能从火产生；但更真实的是万物由二者混合而生。火确实不能与水混合，因为它们彼此对立；若它们碰撞，则较强的一方必毁灭另一方。但它们的实质可以混合。火的实质是热，水的实质是湿。所以\Person{奥维德}说得对：\par

\Quote{因为湿与热一旦混合，便受孕，万物由此二者而生。虽然火与水相争，湿气却产生万物，不协调的和谐适于生成。}\par

因为一种元素可说是阳性的，另一种可说是阴性的：一种主动，一种被动。为此古人规定，婚姻契约须以火与水的仪式批准，因为动物的幼体由热和湿获得身体，并由此获得生命。\par

因为既然每个动物由灵魂和身体组成，身体的质料在于湿，灵魂的质料在于热：这可以从鸟类的后代得知；尽管它们充满浓湿，但若非被创造性的热所温暖，湿不能成为身体，身体也不能获得生命。流放者通常也被禁止使用火和水：因为当时似乎对任何有罪者处以死刑都是不合法的，因为他们终究是人。因此，当禁止使用构成人类生命的事物时，这被视为等同于对如此被判的人实际执行死刑。这两种元素被视为如此重要，以至于人们认为它们对于人的产生和维持生命是必不可少的。其中之一我们与其他动物共有，另一种则唯独赋予人。因为我们是一个属天且不死的族类，我们使用火——这赐给我们作为不死的证明，因为火来自天上；其性质是活动且上升的，含有生命的原理。但其他动物，因为它们完全是可死的，只使用水——这是一个有形和属地的元素。水的性质，因为是活动的且有向下趋势，显示死亡的形状。因此牲畜不仰望上天，也不怀有宗教情感，因为火的使用被剥夺了。但天主从何源头或以何种方式点燃或使这两个主要元素——火和水——流动，唯有造它们的主知道。\par

% structure: p0073 source=h2 target=subsection reason=relative-heading-rank; base=h2

\subsection{第十一章 论生物、人；\Person{普罗米修斯}、\Person{丢卡利翁}、\Person{帕尔开}。}

因此，在完成世界后，祂命令创造各种不同种类和不同形状的动物，有大的也有小的。它们被造成成对的，即每种一雌一雄；由它们的后代充满天空、大地和海洋。天主又按它们的种类从地上赐给所有这些动物食物，使它们能服务于人：有些作为食物，有些作为衣物；而那些力量强大的，祂赐予它们帮助耕种土地，因此被称为负重的牲畜。这样，当万物以奇妙的秩序安排好后，祂决定为自己预备一个永恒的国度，并创造无数的灵魂，将不死赐予它们。于是祂为自己造了一个具有知觉和理智的形象——即按祂自己肖像的样式——没有什么比这更完美的了：祂用地上的尘土造了人，由此人称\Quote{人}，因为他是从土造的。最后，\Person{柏拉图}说人的形体是神圣的；正如\Person{西比尔}所说：\par

\Quote{你是我的肖像，人啊，拥有正确的理智。}\par

关于人的形成，诗人们也没有给出不同的叙述，尽管他们可能歪曲了它；因为他们说，人是由\OriginalTerm{普罗米修斯}{Prometheus}用泥土造的。他们在这事本身没有弄错，但在造物者的名称上弄错了。因为他们从未接触过真理的线索；那些由先知的神谕流传下来并载于天主圣书中的事；那些从寓言和模糊意见中收集并歪曲的事（正如真理往往在众人以各种交谈传播时被败坏，每个人都对他所听到的加添一些东西）——他们将这些事纳入他们的诗篇；而在这件事上，他们确实行为愚蠢，竟将如此奇妙而神圣的工程归功于人。因为，如果人可以通过与普罗米修斯自身从\OriginalTerm{伊阿珀托斯}{Iapetus}所生相同的方式产生，那又何必用泥土造人呢？如果他是一个人，他能够生一个人，但不能造一个人。但他在\OriginalTerm{高加索山}{Mount Caucasus}上的刑罚表明他不是诸神中的一位。但没有人将他的父亲伊阿珀托斯或叔父\OriginalTerm{提坦}{Titan}算作神，因为王国的崇高尊严唯\OriginalTerm{萨图尔努斯}{Saturn}所拥有，他因而连同他所有的后裔获得了神性的尊荣。诗人们的这一杜撰可通过许多论证加以驳斥。所有人都同意，洪水发生是为了毁灭邪恶并把它从地上清除。现在，哲学家、诗人以及古代史的作家都断言相同的事，在此他们尤其与先知的语言一致。因此，如果洪水发生是为了毁灭因人数过多而增长的邪恶，那么普罗米修斯怎能是人的创造者？因为同一作家们说，他的儿子\OriginalTerm{丢卡利翁}{Deucalion}因他的义德是唯一被保存下来的人。一个单独的后裔和一代人怎么可能如此迅速地充满世人？但显然，他们与之前的叙述一样也歪曲了这一点；因为他们既不知道洪水何时发生在地上，也不知道当人类灭亡时谁因他的义德配得被救，以及如何并与谁一起被救——凡此种种都是由受启示的著作教导的。因此，显然，他们关于普罗米修斯作为的叙述是虚假的。\par

但是因为我曾说过，诗人不习惯于说全然不真实的话，而是以隐喻包裹并由此模糊他们的叙述，我并非说他们在这一点上说了假话，而是说，首先，普罗米修斯用丰润柔软的泥土造了一个人的形象，并且他首创了制造雕像和形象的技艺；因为他生活在\OriginalTerm{朱庇特}{Jupiter}的时代，那时开始建造庙宇，并引入了敬拜诸神的新方式。这样，真理就被虚假败坏了；而本应说成是由天主造成的东西，也开始归功于仿效神圣工作的人。但是，用泥土造成真正而活生生的人，这是天主的工作。这一点\OriginalTerm{赫耳墨斯}{Hermes}也提到过，他不仅说人是天主按照天主的形象造成的，甚至试图解释祂如何巧妙地形成了人体的每个肢体，因为没有哪个肢体不是既在使用需要上也在美观上同样有用的。就连\OriginalTerm{斯多葛学派}{Stoics}，当他们讨论天主眷顾的主题时，也试图这样做；\OriginalTerm{图利乌斯}{Tully}在不少地方追随了他们。但是，他简略地处理了一个如此丰富而富有成果的主题，我现在之所以略过它，是因为我最近专门写了一本书论述这个主题，给我\OriginalTerm{德梅特里亚努斯}{Demetrianus}弟子。但我不能在这里忽略某些错误的哲学家所说的话，即人和其它动物是从地里产生的，没有任何创造者；由此产生了\OriginalTerm{维吉尔}{Virgil}的这句话：\par

\Quote{从坚硬的土地上抬起了头的大地所生的人类种族。}\par

这种意见尤其被那些否认天主眷顾存在的人所持有。因为斯多葛学派将动物的形成归于神圣的技巧。但\OriginalTerm{亚里士多德}{Aristotle}通过说世界一直存在，从而摆脱了辛劳和麻烦；因此，人类和其中的其它事物没有开端，而是一直存在，且将永远存在。但是，当我们看到每一个动物个体，以前并不存在，却开始存在并停止存在时，那么整个类族必然在某个时间开始存在，并且必然在某个时间停止存在，因为它曾有一个开端。\par

一切事物必然包含在三个时间段中——过去、现在和未来。开始属于过去，存在属于现在，解体属于未来。这一切在个人身上都能看到：我们出生时开始，活着时存在，死亡时终止。因此他们设想有三位\OriginalTerm{命运三女神}{Parcae}：一位为人类缠绕生命之线，第二位编织它，第三位切断并完成它。但在全人类中，因为只看到现在，却从它推断出过去（即开始）和未来（即解体）。因为它存在，显然它曾在某个时间开始存在，因为没有开始就没有事物能存在；因为它有开始，显然它将在某个时间结束。因为由必朽者构成的整体不能是不朽的。既然我们各自都会死，有可能由于某种灾祸，所有人都同时灭亡：或因土地不产（这有时局部发生），或因瘟疫蔓延（常蹂躏个别城邦和国家），或因世界火灾（如据说发生在\Person{法厄同}身上的事），或因洪水（如\Person{丢卡利翁}时期所传，当时除一人外全族被灭）。如果这场洪水是偶然发生的，那么很可能那唯一的幸存者也会灭亡。但如果他是出于天主眷顾的旨意被保留（这是不可否认的）以复兴人类，那么显然人类生命和毁灭都在天主的能力中。如果它可能完全死亡，因为它部分死亡，那么它显然在某时有起源；正如易于朽坏表明了开端，它也证明了终结。如果这些是真的，\Person{亚里士多德}将无法坚持世界本身也没有开端。但如果\Person{柏拉图}和\Person{伊壁鸠鲁}从\Person{亚里士多德}那里强求这一点，那么\Person{柏拉图}与\Person{亚里士多德}（他们认为世界是永恒的）尽管雄辩，也将被\Person{伊壁鸠鲁}剥夺这一点，因为随之而来的是，既然它有开端，也必然有终结。但我们将在最后一卷更详细地谈论这些事。现在让我们回归到人类的起源。\par

% structure: p0081 source=h2 target=subsection reason=relative-heading-rank; base=h2

\subsection{第12章 动物并非自发产生，而是出于天主的安排，若知道它对我们有益，天主本会让我们知晓。}

他们说，在天空的某些变化和星辰的运动中，存在一种产生动物的成熟时机；这样，新生的土地保留着生产性的种子，自身产生出某些类似子宫的容器，\Person{卢克莱修}论及此事说——\par

\Quote{子宫由根须附着在大地上生长；}\par

那么，它们如何能够忍受或避免寒冷或炎热的力量，或甚至出生呢？因为太阳会灼烧它们，寒冷会收缩它们？但是，他们说，在世界之初，没有冬天也没有夏天，而是气温均衡的永恒春天。那么，为什么我们现在看不到这些事情中的任何一件发生呢？因为，他们说，这必须发生一次，以便动物可以出生；但在它们开始存在、并且繁殖能力被赋予它们之后，大地就停止生产了，时间的状况也改变了。哦，反驳谎言是多么容易啊！首先，在这个世界中，没有任何事物能够不保持它开始时那样的永久存在。因为太阳、月亮和星辰并非那时尚未受造；也不是在被造之后没有它们的运动；也没有那管理和统治它们运行的天主治理在它们开始时就未能运行。其次，如果正如他们所说，必然存在一种天主的眷顾，而他们恰好落入他们特别回避的那种情况。因为当动物尚未出生时，显然有人提供了它们应该被出生，以使世界不致因荒芜和凄凉而显得阴暗。但是，为了使它们无需父母的作用而从大地产生，必须有极大的智慧来提供；其次，从大地凝结的湿气形成各种身体形状；并且，从覆盖它们的容器中获得生命和感觉的能力后，它们被倾倒出来，仿佛从母亲的子宫中，这是一种奇妙而不可描述的提供。但是，让我们假定这也是偶然发生的；随后的事情肯定不能是偶然的——大地立即流出奶水，大气温度均匀。如果这些事情明显地发生了，以便新生的动物可以有营养，或免受危险，那么必定有人通过某种天主的旨意提供了这些事情。\par

但是，除了天主，谁能提供这种照顾呢？然而，让我们看看他们主张的那件事本身是否可能发生——人从大地中诞生。如果有人考虑一个婴儿需要多么长时间和以何种方式被抚养，他必定会明白，那些土生土长的孩子不可能在没有人抚养的情况下被养大。因为他们必须被丢弃躺卧许多个月，直到筋腱强壮到能够自己移动和改变位置，而这在一年的时间里几乎不可能发生。现在请想一想，一个婴儿能否以同样的方式和同样的地点躺卧许多个月，即它被丢弃的地方，而不死亡——被它为了营养而供应的泥土湿气以及它自身排泄物的混合体所淹没和腐烂？因此，它不可能不由某个人抚养；除非所有动物出生时都不是柔弱的，而是已经长大了；而他们从未想到要说这些。因此，那整个方法是不可能的和虚妄的；如果那可以被称为方法的话，而它却试图使没有任何方法。因为凡是说万物都是自行产生、不归因于任何天主眷顾的人，他肯定不是在主张，而是在推翻方法。但如果没有设计，任何事情都不能做成或产生，那么就明显有天主眷顾，那被称为设计的特别属于它。因此，天主，万物的创造者，造了人。甚至西塞罗，虽然对圣经无知，也看到了这一点，他在其《论法律》第一卷中，传下了和先知们相同的东西；我引用他的话：\Quote{这种有预见、睿智、多样、敏锐、富有记忆、充满方法和设计的动物，我们称之为人的，是由至高天主在非凡的环境下产生的；因为在如此多种类和本性的动物中，惟独它享有判断和反思，而所有其他动物都缺乏这些。} 你看到了吗？那人虽远离真理的知识，但因为他持有智慧的形象，就懂得人只能由天主产生？但无论如何，需要神圣的见证，以免人的见证不够。西彼拉见证人是天主的工作：——\par

\Quote{祂是唯一的天主，无敌的创造者，祂亲自固定了人形的样式，祂亲自混合了属于生命产生的一切本性。}\par

圣经的记载也含有同样的论述。因此，天主履行了真正父亲的职责。祂亲自形成了身体；祂亲自注入了我们用以呼吸的灵魂，这灵魂来自祂自己圣神生命之源。无论我们是什么，完全是祂的工程。祂本会教导我们祂是如何成就这事的，如果我们应当知道的话；正如祂教导我们其他事物，这些事物将古代错误和真光的认识都传达给了我们。\par

% structure: p0088 source=h2 target=subsection reason=relative-heading-rank; base=h2

\subsection{为何人有两性；何为人的第一次死亡、第二次死亡，以及我们原祖的过犯与惩罚。}

因此，当祂首先按自己的肖像造了男人之后，接着又按男人的形象造了女人，好让二人通过结合能繁衍后代，以众多人口充满全地。但在造人本身时，祂完成并终结了两种材料——我们说过它们是相互对立的——即火与水的本性。因为造了身体之后，祂将来自自己圣神生命之源的灵魂吹入其中，这灵魂是永久的，好使它能承载世界本身的相似，而世界是由对立元素构成的。因为人由灵魂和身体组成，即如同由天和地组成：因为我们赖以生活的灵魂，其根源仿佛是从天来自天主，身体则来自地，我们说过是由尘土造成的。恩培多克勒——你难以断定该把他归于诗人还是哲学家，因为他像罗马的卢克莱修和瓦罗那样用诗写关于万物的本性——确定了四种元素，即火、气、水、土；或许他追随赫耳墨斯·特里斯墨吉斯忒斯，后者说我们的身体是由天主用这四种元素组成的，因为他说它们自身包含某种火、某种气、某种水、某种土，然而它们既不是火，也不是气，也不是水，也不是土。这些事确非虚假；因为土的本性包含在肉身中，湿气的本性在血中，气的本性在气息中，火的本性在生命热中。但血不能与身体分离，正如湿气不能与土分离；生命热也不能与气息分离，如火不能与气分离：这样，万有中只发现两种元素，其全部本性都包含在我们身体的形成中。因此，人是由不同且对立的本体造成的，如同世界本身由光明与黑暗、生命与死亡构成；祂告诫我们，这两样东西在人内彼此争斗：所以如果灵魂——它来自天主——取得上风，它便是不朽的，活在永光之中；反之，如果身体压倒灵魂，使之服从其统治，便处于永久的黑暗与死亡之中。其效力并非完全消灭不义者的灵魂，而是使它们受永罚。\par

我们称那刑罚为第二次死亡，它本身也是永恒的，正如永生一样。我们如此界定第一次死亡：死亡是有生命之物的本性的分解；或如此：死亡是灵魂与肉身的分离。但我们如此界定第二次死亡：死亡是永恒痛苦的承受；或如此：死亡是灵魂因其罪罚而受永罚的定罪。这并不延及无言的牲畜，它们的精神并非由天主所造，而是出自普通空气，因而死亡使之消散。因此，在这天与地的结合中——其形象在人身上发展出来——属于天主的那些事物占据较高的部分，即灵魂，它掌管身体；但属于魔鬼的那些事物占据较低的部分，显然就是身体：因为身体是属土的，理应服从灵魂，正如大地服从于天。因为身体仿佛是一个器皿，这属天的精神可将其用作暂时的居所。两者的职责是——后者来自天上和天主，应发号施令；前者来自大地和魔鬼，应服从。这一点，确实也未能逃过一个放荡之人\Person{撒路斯提乌斯}的注意，他说：\Quote{但我们全部的力量在于灵魂和身体；我们用灵魂来命令，用身体来服从。}倘若他言行一致就好了；因他竟是极卑劣享乐的奴隶，并以生活的败坏摧毁了自己言论的效力。但若灵魂是火，如我们已表明的那样，它应如火焰般向天上升，免得熄灭；也就是说，它应当升至那在天上的永生。正如火若不被某种富含燃料之物滋养以为其提供维持，便不能燃烧和保持存活，同样，灵魂的燃料和食物唯独是义德，它籍此被滋养以得生命。在这些事之后，天主按我所指出的方式造了人，将他安置在乐园里——即一个极其丰饶悦目的园子，祂在东方地区栽种了各种树木，使他能以其各样的果实为生；并免除一切劳苦，使他能完全献身于事奉他的天主圣父。\par

然后祂赐给他明确的诫命，藉遵守这些诫命，他可继续不朽；若违反之，则受死亡之罚。他受命不可吃园中唯一一棵树上的果子，那树位于园中央，天主将善恶知识置于其中。随后，那控告者嫉妒天主的工程，便使出一切诡计和伎俩来诱惑那人，以剥夺他的不朽。他首先用诡计诱骗女人取那禁果，并藉她为工具，也说服了那人自己违犯天主的法律。因此，那人在获得善恶知识后，便开始因自己的赤身露体而羞耻，并躲避天主的面——这是他素来不惯于做的。于是天主将那男人逐出园子，对罪人宣判了刑罚，使他必须劳苦才能维持生计。祂又用火环绕那园子，以防那人靠近，直到祂在地上执行最后的审判，并除去死亡，将义人——祂的崇拜者——召回同一地方；正如圣作者们所教导，以及\Person{厄立特里亚西比拉}所说：\Quote{而那些尊敬真天主的人，承受永生，他们自己一同住在乐园里，那美丽的园子，直到永远。}但这些事属于末后的事，我们将在本作品最后部分谈论。现在让我们解释那些首先的事。因此，死亡跟随人而来，正如天主的判决，连\Person{西比拉}也在她的诗句中教导说：\Quote{人由天主亲手所造，蛇却诡诈地引诱他，使他临到死亡的命运，并领受善恶知识。}于是人的生命变得有限；然而，仍然长久，因为它被延长至一千年。\Person{瓦罗}并非不知此事，这事在圣经中传下，且因众人的知识而流传，他仍试图为古人被认为活了一千年提供理由。他说在埃及人中，以月为年，因此太阳穿过黄道十二宫的循环并不构成一年，而是月亮——它在三十天内穿越那承载星座的圆环——才构成一年；这论据明显是虚假的。因为那时无人超过千年。但如今那些达到百年的人（这常常发生）无疑活了一千二百个月。权威人士报告说，人通常可活到一百二十岁。但因\Person{瓦罗}不知道人的寿命为何及何时被缩短，他便自己缩短了它，因为他知道人有可能活到一千四百个月。\par

% structure: p0092 source=h2 target=subsection reason=relative-heading-rank; base=h2

\subsection{第十四章  论\Person{诺厄}——酒的发明者，他首先认识星辰，以及虚假宗教的起源。}

但后来，天主看见大地充满邪恶与罪行，便决意用洪水毁灭人类；然而，为重新繁衍人类，祂拣选了一人，当众人皆败坏时，此人卓然屹立，成为义德的非凡典范。他在六百岁时，遵照天主的命令建造了一只方舟，当洪水淹没所有最高的山峰时，他本人与妻子、三个儿子及儿媳得救。大地干涸后，天主憎恶前世的邪恶，为使长寿不再成为图谋恶事的原因，便逐代逐渐缩短人的寿命，并将上限定为一百二十年，不得超越。但那人从方舟出来后，据圣经所载，殷勤耕作大地，亲手栽种了一个葡萄园。由此，那些视巴克斯为酒神的人被驳斥。因为他不但在巴克斯之前，而且在萨图尔努斯和乌拉诺斯之前许多世代。当他初次从葡萄园取果，畅饮至醉，赤身躺卧。他的一个儿子名叫含，看见后，没有遮盖父亲的裸体，反而出去告诉兄弟们。但兄弟们拿着衣服，背脸进去，遮盖了父亲。\BibleRef{创 9:23} 父亲得知所发生的事后，遗弃并逐出了他的儿子。那儿子流亡，定居在现今称为阿剌伯的那片土地上；那地从他得名客纳罕，他的后裔称为客纳罕人。这是第一个不认识天主的民族，因为其首领和始祖被父亲诅咒，没有从父亲那里领受对天主的崇拜；这样，他留给后代对天主本性的无知。\par

从这个民族，随着人数增多，所有邻近的百姓都涌流而出。但他父亲的子孙称为希伯来人，真正天主的宗教建立在他们中间。后来，当希伯来人数目大增，因住地狭小无法容纳，青年人或受父母派遣，或自愿因贫穷所迫，离开故土寻找新定居地，便向各方散居，充满了所有岛屿和整个大地；这样，他们从神圣根源的干茎上被扯离，凭自己的判断建立了新的风俗和制度。首先占据埃及的人最先开始仰望并崇拜天体。由于该地气候特点，他们不住在房屋里，且天空没有云遮蔽，他们观察星辰的运行与隐没，在频繁的朝拜中更仔细、更自由地注视它们。后来，受某些异象诱导，他们发明了怪异的动物形象来崇拜；其始创者我们稍后将揭示。其他人散居大地，惊叹世界的元素，开始崇拜天、太阳、地、海，没有偶像和庙宇，在露天向它们献祭，直到后来为最强大的君王建造庙宇和雕像，并开始以牺牲和馨香尊崇他们；这样，他们偏离了对天主的认识，开始成为外教人。因此，那些主张对诸神的崇拜从世界起初就有，且外教信仰先于天主宗教的人错了：他们认为后者是后来发现的，因为他们不知道真理的源头和起源。现在让我们回到世界的开端。\par

% structure: p0095 source=h2 target=subsection reason=relative-heading-rank; base=h2

\subsection{第十五章 论天使的败坏与两种魔鬼}

因此，当人类人数开始增多时，天主出于预见，恐怕起初曾将大地权柄赐予他的魔鬼，以其诡计败坏或毁灭人类，如同起初所做的那样，便派遣天使来保护和完善人类；祂既赐给了他们自由意志，便特别命令他们不可沾染地上的污秽，以免失去天上本性的尊严。祂明确禁止他们做那祂知道他们将要做的事，好使他们不存有获得宽恕的希望。因此，当他们住在人间时，那极其诡诈的大地统治者，借着与他交往，逐渐引诱他们陷入恶习，并用与女人交合玷污了他们。于是，因他们深陷的罪恶，他们不被接入天上，便坠落于地。这样，天使变为魔鬼的随从和仆役。从这些所生者，既非天使亦非人，而带有一种混合本性，不被允许进入地狱，正如他们的父辈不被接入天上。于是便有两种魔鬼：一种属天，一种属地。后者是邪恶的灵，一切恶行的制造者，那魔鬼便是他们的首领。因此特里斯墨吉斯忒斯称他为魔鬼的统治者。但文法学家说，他们被称为\OriginalTerm{恶魔}{daemones}，犹如\OriginalTerm{精明的}{dœmones}，即精通和熟悉事物：因为他们认为这些是神。他们确实知道许多未来的事，但并非全部，因为不准他们完全知晓天主的计划；因此他们惯常使自己的回答顺应模棱两可的结局。诗人们也知道他们是魔鬼，并如此描述。赫西俄德这样说：—\par

\Quote{这些是按宙斯旨意的魔鬼，善良的，居住在大地上，守护凡人的。}\par

这是为此而说的，因为天主曾派他们作为人类族类的守护者；但他们自己虽然是人类的毁灭者，却仍希望显得是人类的守护者，好使他们自己被崇拜，而不崇拜天主。哲学家们也讨论这些存在。因为\Person{柏拉图}甚至在他的\WorkTitle{会饮篇}中试图解释他们的本性；\Person{苏格拉底}说有一个魔鬼常在他身边，从幼年起就与他联结，他的生活就是由这魔鬼的意愿和引导来支配的。\OriginalTerm{贤士}{Magi}的技艺和能力也完全在于这些存在的影响；被他们呼召而来的这些存在，以欺骗性的幻象蒙骗人们的眼睛，使人看不到存在的事物，反而以为看到了不存在的事物。这些污秽而堕落的灵魂，正如我所说，遍行大地，以毁灭人类来寻求自己沦亡的慰藉。因此，他们到处设下陷阱、诡计、欺诈和错谬；因为他们附着于个人，挨家挨户占据整座房屋，并自称为\OriginalTerm{守护神}{genii}；因为在拉丁语中，他们就是用这个词来翻译\Quote{魔鬼}的。他们把守护神供奉在家中，每日向他们奠酒，把这些狡诈的魔鬼当作地上之神和避祸之神（这些祸患本是由他们自己带来并施加的）来崇拜。而这些灵，由于没有实体且不可捉摸，就潜入人的身体；并在人体内部暗地运作，败坏健康，加速疾病，以噩梦惊吓人的灵魂，以疯狂折磨人的心智，好通过这些恶事强迫人们求助于他们。\par

% structure: p0099 source=h2 target=subsection reason=relative-heading-rank; base=h2

\subsection{第十六章　魔鬼对稳固于信德的人毫无能力。}

所有这一切欺骗的本性，对于没有真理的人来说是隐晦的。因为他们认为那些魔鬼在停止伤害时对他们有益，而魔鬼除了伤害之外毫无能力。或许有人会说，因此应当崇拜魔鬼，以免他们伤害人，因为他们有能力伤害。他们确实伤害人，但只伤害那些惧怕他们的人，不被天主强有力而崇高之手保护的人，未受真理之奥秘启迪的人。但义人，即天主的崇拜者，他们惧怕；因义人的名字被呼求，魔鬼便从附魔者的身上离去。因为，被义人的言语如同鞭子抽打，他们不仅承认自己是魔鬼，甚至说出自己的名字——那些在神庙中被朝拜的名字——这些事通常就发生在他们自己的崇拜者面前。显然，这并非宗教的耻辱，而是他们自身荣誉的耻辱，因为他们不能对天主说谎，因天主和义人的声音使他们受折磨。因此，他们常常发出极大的哀嚎，呼喊自己受鞭打，被焚烧，即将出来：认识天主和正义竟有如此大的力量！那么，他们能伤害谁呢？除非那些在他们自己权力之下的人。总之，\Person{赫耳墨斯}断言，认识天主的人不仅不受魔鬼的攻击，甚至不受\OriginalTerm{命运}{fate}的束缚。\Quote{唯一的保护，}他说，\Quote{是虔诚，因为对于虔诚之人，无论是恶魔鬼还是命运，都毫无能力；因为天主拯救虔诚之人脱离一切邪恶；因为在人中唯一且至善的就是虔诚。}他也在一旁证明虔诚是什么，说：\Quote{因为虔诚就是对天主的认识。}他的弟子\Person{阿斯克勒庇俄斯}在他写给国王的那篇完美论说中更充分地表达了同样的观点。确实，他们两人都断言魔鬼是人类的敌人和折磨者，\Person{特里斯墨吉斯忒斯}因此称他们为邪恶的天使；他并非不知道他们从天上存在变为败坏，开始成为地上的。\par

% structure: p0101 source=h2 target=subsection reason=relative-heading-rank; base=h2

\subsection{第十七章　占星术、占卜术及类似技艺是魔鬼的发明。}

这些就是占星术、占卜术、预测术、所谓的神谕、招魂术、魔术以及人在公开或暗中所行的任何邪恶行径的发明者。这一切本身就是虚假的，正如\Person{厄立特里亚女先知}所见证的：——\par

\Quote{所有这些都是错谬的，\newline{}愚昧之人日日追寻。}\par

但这些权威人士以他们的支持使人相信它们是真实的。因此，他们以谎言的占卜欺骗人的轻信，因为向他们揭露真相并不有利。是这些人教导人制作偶像和雕像；为了将人的思想从崇拜真天主转向他们，他们使死去的君王的容貌被塑造并饰以精美，竖立并祝圣，并自取他们的名字，好像扮演某些角色。但术士和被民众真正称为巫术者的人，在施行他们可憎的技艺时，以这些魔灵的真名呼求他们，就是在圣经中读到的那些天上的名字。而且，这些不洁而游荡的灵，为了混乱一切、以错误覆盖人的思想，将虚假与真实交织混杂。因为他们自己捏造说有许多天上的存在，还有一个万王之王，\Person{朱庇特}；因为天上有许多天使的灵，却有一位万有的父和主，即天主。但他们以虚假的名字隐藏了真理，使之远离视线。\par

因为，如我在起初所表明的，天主不需要名字，因为祂是唯一者；天使们，由于他们是不死不灭的，既不愿也不容许被称为神：因为他们唯一且仅有的职责是服从天主的旨意，除了奉祂的命令外，绝对不做任何事。因为我们说，世界由天主治理，如同一个行省由它的总督治理；没有人会说他的随从是他治理行省的共享者，尽管事务是通过他们的服务来进行的。然而，这些随从有可能因总督的不知情而做出违背他命令的事——这是人类状况的结果。但那世界的守护者、宇宙的统治者，祂知道一切，祂神圣的眼目无物可隐藏，唯有祂与祂的圣子拥有万有的权柄；天使身上别无所有，唯有服从的必要。因此，他们不愿受人尊敬，因为他们所有的尊荣都在天主内。但那些背叛天主服务的天使，因为他们是真理的仇敌、天主的背叛者，就企图为自己夺取神的名号和崇拜；不是因为他们渴望任何尊荣（因为对于丧亡者有何尊荣？），也不是为了伤害不能受伤害的天主，而是为了伤害人——他们竭力使人偏离对真正威严的敬拜和认识，好使人不能获得永生，而他们自己因邪恶已失掉了永生。因此，他们散布黑暗，以晦暗遮掩真理，使人不认识自己的主和父。为了轻易诱惑人，他们藏在圣殿中，在一切祭献时近在身旁；他们屡行奇迹，使人大为惊异，从而将神性权能与影响归于雕像。因此，鸟卜者能用剃刀切割石头；维爱的朱诺回答说她愿意迁往罗马；女贵族的幸运女神预告了迫近的危险；船跟随了克劳狄亚的手；被劫掠的朱诺、洛克里的珀耳塞福涅、米利都的刻瑞斯惩罚了亵渎者；赫剌克勒斯向阿庇乌斯复仇，朱庇特向阿提尼乌斯复仇，弥涅耳瓦向凯撒复仇。因此，从厄庇道鲁斯请来的巨蛇使罗马城摆脱了瘟疫。因为魔鬼的首领自身以本形被运到那里，毫无伪装；确实，为此目的派遣的使节带回了庞大的巨蛇。\par

但他们尤其在神谕中欺骗人，其骗术连非信徒也无法与真理区分；因此，人们想象命令、胜利、财富和事务的成功结果是由他们赐予的——简言之，国家常因他们的干预而摆脱迫近的危险；这些危险他们既已预告，又因祭献平息而消除。但这一切都是欺骗。因为他们既然预知天主的安排——由于他们曾是祂的仆役——就在这些事中介入，使得凡由天主完成或正在完成之事，他们自己似乎尤其在做或已经做了；每当按照天主的旨意，有某种益处将要降临某个民族或城市时，他们就通过奇事、梦或神谕应许必将成就，只要有人献给他们圣殿、尊荣和祭献。而一旦这些奉献完成，当必然的结果到来时，他们便为自己赢得了最大的崇敬。因此，人们许愿建造圣殿，祝圣新的雕像；宰杀成群的牺牲；然而，当这一切做完后，行这些事的人的生命和安全却不见得更有保障。但每当危险逼近时，他们就声称因某些琐细微不足道的原因而发怒；例如，朱诺对瓦罗发怒，因为他将一名俊美的男孩放在朱庇特车上守护袍服，因此罗马之名在坎尼几乎覆灭。但如果朱诺害怕第二个伽倪墨得斯，为何罗马青年要受惩罚？如果众神只关心首领，忽略其余大众，为何如此做的瓦罗独自逃脱，而无辜的保卢斯被杀？确实，当汉尼拔以诡计和勇武消灭了罗马人的两支军队时，罗马人所遭遇的并非\Quote{敌对朱诺的命运，}。因为朱诺既不敢保卫迦太基——那里有她的武器和战车——也不敢加害罗马人；因为\par

\Quote{她曾听说，特洛伊的子孙\newline{}生来就是要摧毁她的迦太基。}\par

但这些是那些人的欺骗，他们隐藏于死者名下，设计陷害活着的人。因此，无论即将来临的危险能否避免，他们都希望使人相信是被他们平息后消除的；或者若无法避免，他们就设法使人相信是因不敬他们而发生的。这样，他们从不认识他们的人那里获得了权威和恐惧。凭借这种诡计和这些伎俩，他们使对唯一真天主的认知在万民中失传了。因为他们已为自己的恶习所毁灭，就狂怒施暴，图谋毁灭他人。因此，这些人类之敌甚至发明了人祭，以便尽可能多地吞噬生命。\par

% structure: p0109 source=h2 target=subsection reason=relative-heading-rank; base=h2

\subsection{第十八章  论天主的忍耐与惩罚、魔鬼的敬拜及虚假宗教。}

有人会说：那么，天主为何允许这些事发生，而不纠正如此灾难性的错误呢？为使邪恶与良善对立，使恶习与美德抗衡；使祂有可惩罚的人，也有可尊荣的人。因为祂已决定在最后的时候审判生者死者，关于这审判我将在最后一卷中讲述。因此，祂延迟直到时日的终结来临，那时祂将以天上的权能与大力倾泻祂的义怒，正如\par

\Quote{虔诚先知的预言\newline{}在惊慌的耳中回响着恐惧。}\par

但现在祂容许人犯错，甚至对祂自己不虔敬，祂是正义、温和且忍耐的。因为具有完美美善者，不可能不具完美的忍耐。有些人因此猜想，天主完全不受愤怒所动，因为祂不受情感影响——情感是心灵的扰动；因为任何易受情感和情绪影响的生物都是软弱的。但这种信念完全抹杀了真理和宗教。然而，关于天主义怒的讨论暂且搁置，因为这题目十分丰富，将在专论此事的著作中更广泛地处理。凡敬拜且追随这些最邪恶的恶灵者，既不享有天堂，也不享有属于天主的光明，反将沦入我们在万物分配中所说归于恶者之首的那些事物——即黑暗、地狱和永罚。\par

我已证明诸神的宗教仪式在三方面是虚妄的：第一，因为那些被敬拜的偶像是死人的形象；这是错误且矛盾的事，即人的形象被神的形象所敬拜，因为敬拜者低于且弱于被敬拜者；再者，离弃活人而去服侍死人的纪念物，是无可补救的罪行，因为这些纪念物不能给任何人生命或光明，它们自身也缺乏这些。并且除了一位天主外别无他神，每个灵魂都受祂的审判和权能所支配。第二，那些神圣的偶像本身——最无知的人向其服役——完全缺乏知觉，因为它们只是泥土。但谁能不明白，一个直立的生物屈身去敬拜泥土是不合法的？泥土被置于我们脚下，正是为了被践踏，而不是被我们敬拜；我们本从泥土中升起，并获得高于其他生物的崇高地位，为使我们不再向下转身，也不将这属天的面容抛向地面，而是将目光转向其本性的状况所指引之处，并且只敬拜和崇拜我们唯一创造者和圣父单独的神性，祂使人直立，好叫我们知道我们被召唤向崇高和天上的事物。第三，因为那些主持宗教仪式本身的恶灵，被天主定罪并抛弃，在地上打滚，它们不仅不能给敬拜者带来任何益处——因为万物的权能都在独一者手中——反而以致命的诱惑和错误毁灭他们；因为他们日常的事务就是将人卷入黑暗，使人不去寻求真天主。因此它们不应被敬拜，因为它们处于天主的判决之下。因为若你跟随正义，你就能在权能上胜过它们，并借呼号神圣之名驱逐和驱散它们，却将自己献给它们的权能，这是极大的罪行。但如果这些宗教仪式如我所证实在多方面是虚妄的，那么显然，那些向死人祈祷、或敬拜大地、或将灵魂交给不洁的灵的人，行为不合乎人之为人，且将因他们的不虔敬和罪恶受罚；这些人背叛了人类之父天主，行了无可补救的仪式，并违反了所有神圣的法律。\par

% structure: p0114 source=h2 target=subsection reason=relative-heading-rank; base=h2

\subsection{论崇拜偶像与地上之物}

因此，凡渴望遵守人应负的义务并持守其本性尊严的人，当从地面升起，以高举的心灵，将目光转向上天；不要在天主脚下寻找祂，也不要从自己的足迹中挖掘出敬拜的对象，因为凡在人之下的，必然低于人；而要在高处寻找祂，在至高之处寻找祂；因为没有比人更大的，除非那在人之上者。但天主大于人；因此祂在高处，不在低处；也不应在最底之处寻找，而应在最崇高之处。因此，毫无疑问，凡有偶像之处就没有宗教。因为如果宗教包含神圣事物，而除了天上的事物外没有神圣之物，那么偶像便是没有宗教的，因为从泥土所造之物中不可能有天上之物。这一点，有智慧的人从其名称本身就可以明白。因为凡是摹仿，必然虚假；凡以欺骗和摹仿伪造真理的，不能冠以真物之名。但如果一切摹仿并非特别严肃之事，而如同游戏和玩笑，那么偶像中就没有宗教，只有宗教的戏仿。因此，真实的应优先于一切虚假的；应践踏地上之物，好获得天上之物。因为情况是这样：凡将自己的灵魂（其本源来自上天）俯伏于阴府和最低之物的人，必坠入他自投之处。因此，他当牢记自己的本性和境况，常奋力追求天上的事物。凡如此行的人，将被断为完全智慧、正义、真正的人；简而言之，他将被视为配得上天堂，他的圣父将认出他并非卑微、如野兽般伏地，而是如祂所造般直立向上。\par

% structure: p0116 source=h2 target=subsection reason=relative-heading-rank; base=h2

\subsection{论哲学与真理}

我所承担的这一工作，若我没有受骗，已经完成了一个伟大而艰难的部分；天主的威严供给我言说的能力，我们已经驱散了根深蒂固的谬误。但现在，一场与哲学家们更伟大、更艰巨的论争摆在我们面前，他们的学识和口才的高峰，犹如一座庞大的建筑，阻挡着我。因为正如前一种情况中我们受到众多民族几乎一致的压迫，在本题中，我们受到在各类赞美中都出类拔萃之人的权威的压迫。然而，谁能不知道，较少的有学识之人比较多的无知之人更有份量呢？但我们不可绝望，在天主和真理的引导下，这些人也可能放弃他们的见解；我也不认为他们会固执到否认他们以健全睁开的眼睛注视光彩照人的太阳。但愿他们自己常声称的——他们充满探索的渴望——是真实的，那么我必定会成功使他们相信那长久寻求的真理终于被找到，并承认它不可能由人的才能找到。\par

\SourceRecord{Translated by William Fletcher.；Ante-Nicene Fathers；Vol. 7.；Edited by Alexander Roberts, James Donaldson, and A. Cleveland Coxe.；Buffalo, NY: Christian Literature Publishing Co.,；1886.；<http://www.newadvent.org/fathers/07012.htm>.}

% structure: p0001 source=h1 target=section reason=page-title

\section{\WorkTitle{神圣训诫，第三卷（论哲学家的虚假智慧）}}

% structure: p0002 source=h2 target=subsection reason=relative-heading-rank; base=h2

\subsection{第一章 真理与雄辩的比较：哲学家为何未能达到真理。论圣经的简朴风格。}

由于人们认为真理仍隐藏于晦暗之中——或因庶民的错误与无知，他们受制于各种愚昧迷信；或因哲学家的谬误，他们心灵的乖谬使真理更加混乱而非明朗——我但愿自己能拥有雄辩之力，虽不及\Person{马尔库斯·图利乌斯}那般超凡卓绝，但至少稍近之；这样，藉着才能的力量与自身之重的支持，真理终能显现，驱散并驳斥公众的错误以及被视为智慧者的错误，为人类带来灿烂光明。我愿如此，原因有二：要么，人们因着修饰而更易相信真理，因为他们甚至被言辞的装饰和话语的诱惑所俘而相信谬误；要么，至少哲学家们会被我们击败，尤其是被他们自己惯常引以为傲并信赖的武器所击败。\par

但既然天主命定事情的本质如此：纯朴无伪的真理因其自有充分之美而更加清晰，因此若加以外在装饰反而被败坏；而谬误却因其自身败坏而消散融化，除非借助外在的修饰来装点，才得以取悦于人；我便安然接受赋予我的中等才华。但我着手这项工程并非倚赖雄辩，而是倚赖真理——这工程或许庞大到非我之力所能承担；然而，即便我失败，真理本身亦将在天主的助佑——这是祂的职分——下完成。因我知道，最伟大的演说家也常被能力中等的辩护者击败，因为真理的力量如此强大，连在小事上也能凭自身的清晰来捍卫自己：我为何要想象，在如此重大的事业中，真理会被那些我承认是机巧善辩、却言说虚假之事的人所压倒？为何不认为真理会显现得光辉灿烂，即使不是藉着我们非常微弱、从细小泉源流淌的言辞，而是藉着它自身的光明？此外，即使曾有因其文学学识而值得钦佩的哲学家，我也不应将真理的知识与学问让给他们，因为没有人能藉沉思或辩论获得真理。我现在也无意贬低那些希望认识真理之人的追求，因为天主使人性最渴望达至真理；但我反对并坚持的是：他们的效果并未随其正直而正确的意愿而来，因为他们既不知道真理本身是什么，也不知道如何、在何处、以何种心灵寻求它。因此，当他们想要纠正人们的错误时，却陷入了陷阱和最大的错误。我之所以从事驳斥哲学的任务，正是由我所承担的主题本身的顺序所引导。\par

因为一切谬误皆来自虚假宗教或虚假智慧，在驳斥谬误时，必须将两者都推翻。因为圣经已传授给我们：哲学家的思念是愚拙的；而这一点正须以事实和论证来证明，以免有人被智慧的美名所诱，或被空洞雄辩的光彩所惑，宁愿相信人的事而不信神圣的事。这些事情确实以简明质朴的方式叙述。因为当天主对人说话时，祂不应以论证来确认自己的话，仿佛不如此便不能获得信任；但正如所当行的，祂作为万有的全能审判者说话，审判者的职责不是辩论，而是宣判。祂本身，作为天主，就是真理。但我们，既然拥有一切的神圣见证，必定能显示真理能以何等更确凿的论证被捍卫，即便是虚假之事也常被如此辩护而显得真实。因此，我们毫无理由给予哲学家如此尊荣，以至于畏惧他们的口才。因为他们作为有学问的人可能说得动听，却不能说得真实，因为他们并未从掌握真理的那一位那里学到真理。诚然，我们指证他们的无知——他们自己也屡次承认——也并非什么大事。既然他们仅仅在应当被相信的那一点上不被相信，我将努力证明，他们从未像在表达对自己无知的看法时那样说得真实。\par

% structure: p0006 source=h2 target=subsection reason=relative-heading-rank; base=h2

\subsection{第二章 论哲学，及其在阐述真理上的徒劳。}

现在，在之前的两卷书中，我们已经揭示了迷信的虚假，并阐述了整个错误的起源；因此，本卷的任务是也要展示哲学的空虚和虚假，以便消除一切错误，使真理得以显现并变得明朗。因此，让我们从哲学的普遍名称开始，这样一旦头部被摧毁，我们就能更容易地着手摧毁整个身体；如果确实可以称之为身体的话，其各部分和成员彼此不一致，没有任何联结，反而好像四散分离，颤动着而非活着。哲学（正如其名称所示，也是他们自己定义的那样）是对智慧的爱。那么，我有什么论据能证明哲学不是智慧，而不是从名称本身的意义来证明呢？因为专心致志于智慧的人显然还不是智慧者，而是专心致志于那件事物以便成为智慧者。在其他技艺中，可以看出这种专心致志的效果及其趋向：因为当一个人通过学习掌握了这些技艺后，他就不再被称为该行业的虔诚追随者，而是工匠。但据说，他们出于谦逊而自称为智慧的追求者，而非智慧者。不，实际上，\Person{毕达哥拉斯}是第一个发明这个名称的人，他比古代那些自认为智慧的人多了一点智慧，他明白靠任何人类研究都不可能获得智慧，因此不应该把一个完美的名称应用于一个无法理解和不完美的对象。所以，当有人问他从事什么职业时，他回答说他是哲学家，即智慧的探索者。因此，如果哲学探索智慧，它本身就不是智慧，因为探索者与所探索的事物必然是两种不同的东西；而且探索本身也不正确，因为它什么也找不到。\par

但我甚至不准备承认哲学家们致力于追求智慧，因为通过那种追求无法获得智慧。因为如果发现真理的能力与这种追求相关，并且这种追求是通往智慧的一条道路，那么它最终就会被找到。但是，既然如此多的时间和才能浪费在寻找智慧上，却仍未获得，那么很明显那里没有智慧。因此，从事哲学的人并不是在致力于追求智慧；而只是他们自己想象自己这样做，因为他们不知道他们所寻找的东西在哪里，或具有什么性质。因此，无论他们是否致力于追求智慧，他们都不是智慧者，因为一个要么以不当方式寻求、要么根本不寻求的东西，是永远无法被发现的。让我们来看这件事本身：通过这类追求是否可能发现任何东西，还是什么都发现不了。\par

% structure: p0009 source=h2 target=subsection reason=relative-heading-rank; base=h2

\subsection{第三章 哲学由哪些主题组成，以及谁是学园派的首席创始人。}

哲学似乎由两个主题组成：知识和猜测，仅此而已。知识不能来自理解，也不能通过思想把握；因为将知识作为自身特有属性来拥有，并不属于人，而属于天主。但凡人的本性不能接受知识，除了来自外部的知识。为此，神圣的理智打开了身体的眼睛、耳朵和其他感官，以便通过这些入口，知识能够流入心灵。因为探究或想要知道自然事物的原因——太阳是否像它看起来那么大，或者比整个地球大很多倍；月亮是球形的还是凹的；恒星是固定在天空上，还是自由地在空中运行；天空本身有多大，由什么物质构成；它是静止不动的，还是以难以置信的速度旋转；地球的厚度有多大，它依靠什么基础保持平衡和悬浮——我说，通过争论和猜测来理解这些事物，就好像我们想要讨论我们在某个非常遥远的国家从未见过、只听说过名字的城市可能具有什么特征。如果我们声称自己知道这种不可知的事情，我们不就是发疯了吗，竟敢断言可能被反驳的事情？那些自以为知道人类不可知的自然事物的人，更应该被视为疯狂和愚昧！因此，\Person{苏格拉底}和追随他的学园派正确地取消了知识，知识不属于争论者，而属于占卜者。剩下的只有哲学中的猜测；因为缺乏知识的地方，完全被猜测占据。每个人都在猜测他所不知道的东西。但讨论自然事物的人，猜测事物就像他们所讨论的那样。因此，他们不知道真理，因为知识关乎确定的事，而猜测关乎不确定的事。\par

让我们回到前面提到的例子。来吧，让我们来推测那个我们除了名字以外一无所知的城市的状况和特征。它很可能位于平原上，城墙是石头砌的，建筑高耸，街道众多，庙宇宏伟而装饰华丽。请允许我们描述一下市民的风俗和举止。但当我们描述完这些后，另一个人会提出相反的说法；当他结束时，第三个人会出现，接着还有其他人；他们会做出与我们完全不同的推测。那么所有中哪一个更真实？也许一个也没有。但已经提到了所有情况所允许的一切，所以其中某一个必然是真实的。但谁说了真话却不得而知。可能所有人都在某种程度上描述有误，也可能所有人都在某种程度上触及了真理。因此，如果我们通过辩论来寻求这个，我们是愚蠢的；因为可能会有人出现，嘲笑我们的推测，并认为我们疯了，因为我们想推测我们不知道的东西。但没必要去寻求遥远的案例，也许没有人会来反驳我们。来吧，让我们推测一下现在广场里发生了什么，元老院里发生了什么。那也太远了。让我们说说在只有一墙之隔的地方发生了什么事；除了听到或看到的人，没有人能知道。因此没有人敢说这个，因为他会立即被事实本身的存在所驳倒，而不是被言语。而这正是哲学家们所做的，他们讨论天上正在发生的事情，却认为他们可以免于惩罚，因为没有人来反驳他们的错误。但如果他们想到有人会降下来证明他们是疯狂和虚假的，他们就永远不会讨论那些他们不可能知道的事情。不过，他们因为没有被反驳而认为自己的无耻和鲁莽更成功，这并不是事实；因为天主反驳他们，只有天主知道真理，尽管他可能看似纵容他们的行为，并且他把这样的人的智慧视为最大的愚蠢。\par

% structure: p0012 source=h2 target=subsection reason=relative-heading-rank; base=h2

\subsection{第四章 知识被\Person{苏格拉底}取消，推测被\Person{芝诺}取消。}

那么，\Person{芝诺}和斯多葛学派拒绝推测是对的。因为推测你知道你不知道的事，不是智者的行为，而是鲁莽和愚蠢之人的行为。因此，如果什么都不能知道，正如\Person{苏格拉底}所教导的，或者不该推测，正如\Person{芝诺}所教导的，那么哲学就被完全消除了。何须我多说呢？不仅这两位哲学的领袖推翻了它，而且所有学派都推翻了它，以致如今看来哲学早已被自己的武器摧毁了。哲学被分成许多派别，它们各持己见。真理在哪里？当然不可能在所有派别中。让我们找出一个，那么其他所有派别都将没有智慧。让我们逐一审查它们；同样，我们给予一个派别的，就会从其他派别夺走。因为每个特定的派别都推翻所有其他派别，以确认自身及其学说；它不允许任何其他派别拥有智慧，以免承认自己是愚蠢的；但正如它夺走其他派别，它也被所有其他派别夺走。因为它们毕竟是哲学家，指责它愚蠢。无论你称赞并声称哪个派别是真的，它都会被其他哲学家斥为假的。那么，我们是相信那个称赞自身及其学说的派别，还是相信那些相互指责无知的大多数呢？多数人所持的必然比单独一人所持的更好。因为没有人能正确判断自己，正如著名诗人所证实的；因为人的天性如此安排，他们看见并辨别他人的事务胜过自己的。既然一切都不确定，我们必须要么相信所有人，要么谁都不信：如果我们谁都不信，那么智者就不存在，因为他们分别肯定不同的事物却自以为智慧；如果所有人我们都信，同样也没有智者，因为所有人都各自否认他人的智慧。因此，所有人在这种方式下都被摧毁了；正如诗人们传说中的\OriginalTerm{播种者}{sparti}一样，这些人互相残杀，以致没有一个人留下；这发生的原因是，他们有剑，却没有盾。因此，如果各派别因许多派别的判断而被判为愚蠢，那么所有派别都是虚妄和空洞的；这样，哲学就消耗并毁灭了自己。由于\Person{阿尔克西拉乌}，学园的创始人，理解了这一点，他收集了所有派别的相互谴责和著名哲学家对无知的承认，并武装自己对抗所有派别。这样，他建立了一种新的不搞哲学的哲学。因此，从这位创始人开始，便有了两种哲学：一种是旧的，自诩拥有知识；另一种是新的，与前者对立并贬低它。在这两种哲学之间，我看到了分歧，仿佛是内战。智慧将站在哪一边？智慧不能被撕成碎片。如果事物的本性可以被认识，这支新兵队伍将灭亡；如果不能，老兵将被消灭；如果它们势均力敌，那么作为一切向导的哲学仍将灭亡，因为它被分裂了；因为没有什么能反对自身而不自我毁灭。但是，正如我已经表明的，由于人类状况的脆弱，人不可能有内在而特有的知识，那么\Person{阿尔克西拉乌}一派就占了上风。但即使这一派也不能稳固，因为不可能什么都不知道。\par

% structure: p0014 source=h2 target=subsection reason=relative-heading-rank; base=h2

\subsection{第五章 许多事物的知识是必要的。}

因为有许多事，是自然本身、频繁的使用以及生活的需要迫使我们知道的。因此，除非你知道哪些事对生活有益而应当追求，哪些事危险而应当回避和避免，否则你必然灭亡。此外，还有许多事物是经验所发现的。因为太阳和月亮的各种运行、星辰的运动、时间的计算已被发现，医学研究者发现了身体的本性和草药的效力，土地耕种者收集了土壤的本性以及未来雨和风暴的征兆。简言之，没有一门技艺是不依赖于知识的。因此，阿尔克西拉乌（\Person{阿尔克西拉乌}）若有点智慧，就应当区分哪些是可以知道的，哪些是不可知道的。但他若这样做，就会把自己降低到普通人的水平。因为普通人有时更有智慧，因为他们只明智到必要的程度。你若问他们是否知道任何事或一无所知，他们会说他们知道他们所知道的事，并会承认他们不知道他们所不知道的事。因此，他取消别人的体系是对的，但他建立自己的体系却不对。因为对一切事物的无知不能成为智慧，智慧的特性是知识。这样，当他战胜了哲学家，教导他们一无所知时，他自己也失去了哲学家的名号，因为他的体系就是什么都不知道。因为责备别人无知的人，自己应当有知识；但当他什么都不知道时，因着同一件事而自称哲学家——他正是凭这件事取消别人的哲学——这是何等乖张或何等无礼！因为他们有权这样回答：若你因为我们什么都不知道而断定我们无知，从而不智慧，那末，因为你承认自己什么都不知道，难道你不也不智慧吗？因此，阿尔克西拉乌取得了什么进步呢？除了消灭了所有哲学家，他也用同一把剑刺伤了自己。\par

% structure: p0016 source=h2 target=subsection reason=relative-heading-rank; base=h2

\subsection{第六章  论智慧、学院派与自然哲学}

那么，智慧无处可寻吗？是的，它确实在他们中间，但无人看见。有些人认为一切皆可知：他们显然不智慧。另一些人认为无物可知：他们也不智慧。前者对人所赋太多，后者对人赋太少。两者都缺了限度。那么，智慧在哪里？在于既不认为你知道一切——这属于天主；也不认为你一无所知——这属于野兽。它属于一种中间特性，即知识与无知联合且结合。我们身上的知识来自灵魂，灵魂源自天上；无知来自肉体，肉体来自地上：因此我们与天主有共同之处，也与动物有共同之处。这样，既然我们由这两元素组成，其一拥有光明，另一拥有黑暗，那么一部分知识给了我们，一部分无知。可以说，我们可在这座桥上安全走过而不致坠落；因为所有倾向于任何一侧——无论是左还是右——的人都坠落了。但我将说明每一方如何犯错。学院派依据晦涩的事物反对自然哲学家，说无知识；他们满足于少数不可理解之事的例子，便拥抱无知，仿佛取消了全部知识，因为他们取消了部分知识。但自然哲学家另一方面依据明显的事物立论，推断一切皆可知；他们满足于明显之事，便保留了知识，仿佛完全捍卫了知识，因为他们捍卫了部分知识。因此，一方看不见明显之事，另一方看不见晦涩之事；但两派在极力争辩要么保留要么取消知识时，都没有看见中间存在着一条引向智慧的道路。\par

但阿尔克西拉乌（\Person{阿尔克西拉乌}）教导无知识，他贬损斯多葛学派的首领芝诺（\Person{芝诺}），为要借苏格拉底的权威彻底推翻哲学，便采纳了这意见，断言无物可知。这样，他就否定了哲学家的判断——那些哲学家曾以为真理凭其才智被引出并发现——理由是因为那智慧是必朽的，且在几个世纪前建立，现已达到最大增长，如今必然衰老而消亡；\OriginalTerm{学园}{Academy}突然兴起，仿佛是哲学的老龄，可将其枯萎的哲学终结。阿尔克西拉乌正确地看出，那些想象凭猜测能达到真理知识的人是狂妄，或更确切说，是愚蠢。但无人能驳斥说谎者，除非他先前知道什么是真的；然而阿尔克西拉乌在不知真理的情况下试图这样做，便引入了一种可称为不稳定或无常的哲学。因为为了无物可知，必须知道一些事物。因为若你一无所知，那\Quote{无物可知}的知识本身也会被取消。因此，那发言说\Quote{无物可知}的人，仿佛是宣布了一个已经达到和已知的结论：所以，有些事物是可知的。\par

与这一推理类似的，是通常在学园中被用作所谓\OriginalTerm{自相矛盾的推理}{asystaton}谬误的例子：某人曾梦见自己不应该相信梦。因为如果相信了梦，那么他就应当不相信梦；如果不相信梦，那么他就应当相信梦。这样，如果什么都不能知道，那么必须知道这一点：什么都不知道。但如果知道什么都不能知道，那么\Quote{什么都不能知道}这个陈述就必然是假的。这样就引入了一个自相矛盾、自我否定的论点。但这位诡辩者想要剥夺其他哲学家的学识，以便将其藏在自己家中。因为他的确不是要剥夺自己的学识——他断言某事是为了从他人那里剥夺它——但他并未成功；因为学识显露出来，暴露了它的掠夺者。如果他做出一个例外，说只有天上的或自然界的原因和系统（因为它们是隐藏的）是不能知道的（因为无人教导），并且不应探究（因为无法通过探究找到），那他会做得多么明智和正确啊！如果他提出这个例外，他既能告诫自然哲学家不要探究超出人类反思界限的事物，又能使自己免于诽谤而来的恶意，并且肯定会给我们留下一些可以遵循的东西。但现在，既然他引导我们不再跟随他人，以免我们想知道超出我们能力所知的东西，他也同样引导我们不再跟随他自己。因为谁愿意辛苦劳作以避免知道任何事呢？或者愿意从事这种学习，以至于甚至失去普通的知识？因为如果这种学习存在，它必然由知识构成；如果它不存在，谁会愚蠢到认为那种东西值得学习，在其中要么什么也学不到，要么甚至还会失去一些东西？因此，如果不能知道一切（如自然哲学家所认为的），也不能什么都不知道（如学园派所教导的），那么哲学就完全被消灭了。\par

% structure: p0020 source=h2 target=subsection reason=relative-heading-rank; base=h2

\subsection{第七章 论道德哲学与至善}

让我们现在转向哲学的另一个部分，他们称之为道德哲学，其中包含了整体哲学的方法，因为在自然哲学中只有愉悦，而在道德哲学中还有效用。既然在安排生活条件和形成品格时犯错更为危险，我们必须更加勤勉，以便知道应当如何生活。因为在前一个主题中，可以给予一些宽容：无论他们说什么，都不带来益处；或者如果愚妄地胡言，也不造成损害。但在这个主题中，没有意见分歧的余地，也没有错误的余地。所有人都必须有相同的见解，哲学本身必须如同用同一张口来指导；因为如果犯任何错误，生命就被彻底颠覆。在之前的部分中，危险较小，困难却较大；因为主题的晦涩迫使我们抱有不同和多样的观点。但在这里，危险更大，困难却较小；因为主题本身的运用和日常经验能够教导什么是更真实和更好的。因此，让我们看看他们是否一致，或者他们为更好地指导生活提供了什么帮助。没有必要扩展开来论述每一点；让我们选择一个，尤其是那个最首要和主要的点，所有智慧都围绕它并依赖它。\Person{伊壁鸠鲁}认为至善在于心灵的愉悦，\Person{亚里斯提卜}认为在于肉体的愉悦。\Person{卡利丰}和\Person{狄诺玛库}将愉悦与德行结合，\Person{狄奥多罗斯}将愉悦与痛苦的缺失结合，\Person{希罗尼穆斯}将至善置于痛苦的缺失中；逍遥学派则置于心灵、身体和运气的善中。\Person{赫里鲁斯}的至善是知识；\Person{芝诺}的是合乎自然地生活；某些斯多葛学派的是追随德行。\Person{亚里士多德}将至善置于完整性和德行中。这些几乎是所有人的观点。在如此多样的意见中，我们跟随谁？我们相信谁？所有都有同等权威。如果我们能够选择更好的，那就意味着哲学对我们不是必需的，因为我们已经智慧，足以判断智者们的意见。但既然我们来是为了学习智慧，我们怎能判断，我们尚未开始成为智慧的人呢？特别是当学园派就在近旁，揪住我们的衣角，阻止我们相信任何人，而不提出我们可以跟随的东西。\par

% structure: p0022 source=h2 target=subsection reason=relative-heading-rank; base=h2

\subsection{第八章 论至善、灵魂与肉体的愉悦，以及德行}

那么，还有什么剩下的呢？无非是离开那些狂吠不休、顽固不化的争辩者，来到那位审判者面前——祂实在是淳朴而宁静智慧的真赐予者。这智慧不仅能塑造我们、引领我们走上正道，还能对那些人的争论作出评判。它教导我们，什么是人真正的、最高的善；但在开始谈论这个主题之前，必须先驳倒所有那些意见，从而表明那些哲学家中没有一人是明智的。既然探讨的是人的本分，那么作为最高受造物的人，其至善应当置于他无法与其他动物共有的东西之中。正如牙齿是野兽的特性，角是牛的特性，翅膀是鸟的特性，同样，人必须拥有某种专属于他自己的东西，否则他就会失去其状态的固定秩序。因为凡是为生命或繁衍而赐予一切受造物的东西，固然是一种自然的善，但还不是至善，除非它是该类所特有的。因此，那种相信心灵的快乐是至善的人并非智者，因为无论那是无忧无虑还是喜悦，都普遍存在。我认为\Person{亚里斯提卜}甚至不值得回答；因为他总是纵情于肉体的快乐，只是感官享乐的奴隶，没有人能把他当作一个人：他的生活方式与禽兽无异，唯一的区别是他有言语能力。但是，如果言语能力赐给了驴、狗或猪，而你询问它们为什么如此疯狂地追逐雌性，以至于几乎无法分离，甚至不顾饮食；为什么它们要么赶走其他雄性，要么即使战败也不放弃追逐，反而常常在被更强壮的动物打伤后更加执著；为什么它们既不怕雨也不怕寒；为什么它们吃苦耐劳、不畏危险——它们除了回答\Quote{至善是肉体的快乐}之外，还能给出什么答案呢？它们热切地追求快乐，以便体验到最愉悦的感觉；这些感觉如此重要，以至于为了获得它们，它们认为任何的劳苦、创伤甚至死亡都不应拒绝。那么，我们要从这些只有非理性受造物之感受的人那里寻求生活准则吗？\par

昔兰尼学派说，德行本身之所以值得称赞，正是因为它能产生快乐。\Quote{确实如此，}那条污秽的狗或泥沼中打滚的猪说。\Quote{正因此，我竭尽全力与对手搏斗，我的勇敢可以为我带来快乐；如果我被打败，就必然会失去它。}那么，我们要从这些在感觉上而非言语上与牲畜和禽兽无异的人那里学习智慧吗？把无痛苦视为至善，这与其说是逍遥派和斯多亚派的观点，不如说是临床派哲学家的观点。有谁会不认为这场讨论是由那些生病且受某种痛苦折磨的人进行的呢？还有什么比把医生能给的东西当作至善更可笑的事呢？因此，我们必须先感受痛苦，才能享受善；而且痛苦要剧烈且频繁，这样随后的无痛苦才能带来更大的快乐。所以，从未感受过痛苦的人是最悲惨的，因为他缺乏那种善；而我们原本却认为他是最幸福的，因为他没有恶。那个说完全无痛苦是至善的人，离这种愚蠢并不远。因为，除了每个动物都避免痛苦这一事实，谁能把这种善赐给自己呢？获得它，我们只能渴望而已。但至善如果不能随时为人所拥有，就不能使人幸福；而赋予人这种善的，不是德行、学问或劳作，而是自然本身将其赐予一切生灵。那些把快乐与德行原则结合起来的人，本想避免这种一切皆然的混杂，却制造了一种自相矛盾的善；因为一个沉溺于快乐的人必然缺乏德行原则，而一个追求原则的人必然缺乏快乐。\par

逍遥派的至善可能显得过分、多样，而且——除了心灵之善（它们究竟是什么仍是重大争议）——与野兽共通。因为身体之善——即安全、无痛苦、健康——对哑巴受造物与对人都同样必要；而且我不知是否对它们更为必要，因为人可以靠治疗和照料得到缓解，哑巴动物却不能。对于他们所谓的幸运之善也是如此；因为人需要生活资源，它们也需要猎物和牧场。因此，通过引入一种不在人能力范围内的善，他们使人完全受制于他人。我们也要听听\Person{芝诺}，因为他有时梦想德行。他说，至善就是顺应自然而生活。那么，我们必须像禽兽那样生活。因为在它们身上，可以发现所有应当远离人的东西：它们渴求快乐，恐惧，欺骗，埋伏，杀害；而尤其切题的是，它们不认识天主。那么，为什么他要教导我按自然生活呢？自然本身倾向于更坏的道路，并在某种更诱人的甜言蜜语之下，使人一头扎进恶习。如果他说禽兽的本性不同于人的本性，因为人生来是为了德行，那他说的倒有些道理；然而，这仍不是至善的定义，因为没有一种动物不按照自己的本性生活。\par

凡以知识为至善者，乃赋予人特有之物；然人求知乃为别有所用，非为知识本身。盖有谁满足于单纯求知，而不欲从其知识中获取某种益处？技艺之习得，旨在付诸实践；而实践之目的，或为维持生计，或为享乐，或为求荣。是故，若非为自身而求者，不得谓之至善。如此，无论吾人视知识为至善，抑或视知识所生之物——即生计、荣耀、享乐——为至善，又有何区别？况且此类事物非人所独有，故非至善；盖享乐与饮食之欲，不独存于人，亦存于禽兽。至于求荣之欲，岂非亦见于马？盖马因获胜而欢跃，因败北而忧伤。\Quote{彼等爱赞颂之情如此之深，求胜之心如此之切。}该卓越诗人言之有理，谓须试探\Quote{彼等败时何其忧伤，胜时何其欢欣。}然若知识所生之物与人兽所共有，则知识非至善明矣。且此定义之另一大弊，在于提出纯知识。若然，凡通晓任何技艺者——纵使知晓有害之事者——皆将视为有福；于是习得配毒药者，与习得施药救者，其福乐相等。吾故问：知识当指向何种对象？若指向自然事物之原因，则吾即便知晓尼罗河之源，或自然哲学家关于天象之虚妄梦幻，又有何福乐可言？况吾为何提及此类事物并无知识，惟有猜测，且因人而异？余下惟有可能：至善乃知善恶之事。然则，彼何以称知识为至善，而不称智慧？二者字义相同。然从未有人谓智慧为至善，虽此言更合宜。盖知识单独不足以担当趋善避恶之事，尚需以德行辅之。诸多哲学家纵论善恶之性，然因本性所迫，生活方式与其言论相悖，盖因欠缺德行。德行与知识结合，方为智慧。\par

余下须驳斥另一些哲人，彼等视德行本身为至善，\Person{Marcus Tullius}亦持此见；此辈甚为轻率。盖德行本身并非至善，乃至善之缔造者与母亲；因无德行则不能达至善。二者皆易理解。吾问：彼等以为那卓越之善易得，抑或须历经艰难劳苦？请彼等发挥才智，维护谬误。若易得而无须劳苦，则非至善。盖吾人何苦自寻烦恼？何苦日夜操劳？所追求者近在咫尺，任何人稍加心意即可得之？然若不经劳苦，连普通中善亦不可得——因善物本性艰难险阻，恶物则趋向堕落——则得至善须最大劳苦。若此最为真确，则需另有德行，方能达至彼称为至善之德行；然此不合情理且荒谬，盖德行不能借自身而达于自身。若无劳苦即无善可达，则显然借德行方能达至善，因德行之力量与职责在于承担并完成劳苦。故至善不能为必须借以达于另一善者。然彼等不明德行之功效与趋向，又寻不到更尊贵之物，遂止步于德行之名，谓当追求之，纵使从中无利可图；如是彼等为自己树立了一种本身仍需善之善。\Person{Aristotle}与之相去不远，彼以为至善乃德行与荣誉并存；彷佛有德行者能不荣誉？或倘若有些许羞耻，德行便不复存在？然彼见可能因错误判断而对德行有恶评，故认为当顾及人之评判中何为正确与良善之背离，因德行按其应得被尊重，非吾等所能左右。盖荣誉若非藉众人好评而恒久赋予一人，又为何物？然若因人之错误与乖戾而致恶名，又将如何？岂因愚者以为卑鄙可耻而弃德行乎？德行虽可受压迫与骚扰，然欲其为自身独特持久之善，则当不需外在援助，而能自立于自身力量，恒久不变。如是，不应从人望其善，亦不应拒绝任何恶。\par

% structure: p0028 source=h2 target=subsection reason=relative-heading-rank; base=h2

\subsection{论至善、真天主的敬拜，并驳斥\Person{阿那克萨戈拉}}

我现在来论真智慧的至善，其性质须以如下方式确定：首先，它必须是人所独有，而不属于任何其他动物；其次，它必须只属于灵魂，而不与身体共享；最后，没有知识和德行，任何人都不能获得它。这个限定排除并取消了上述所有人所持的意见，因为他们的言论中毫不包含这类内容。我现在要说出它是什么，以便如我所计划的那样表明，所有哲学家都是盲目而愚蠢的，他们既不能看见、也不能理解、更从未揣测到什么被确定为人的至善。\Person{阿那克萨戈拉}被问及他为何而生时，他回答说，他可以看到天和太阳。这句话受到所有人的赞赏，并被视作哲学家的名言。但我认为，他由于对答案毫无准备，为了不沉默而随口说出了这句话。但若他有智慧，他本该思考并反省自身；因为如果一个人不知道自己的状况，他甚至不能算是一个人。但我们设想这句话并非一时冲动说出的。让我们看看他在三个词中犯了多少、多大的错误。首先，他错误地把人的全部职责只归于眼睛，什么都不归于心智，而一切归于身体。但倘若他瞎了，他便会失去人的职责吗？那除非灵魂毁灭，否则不会发生。身体的其它部分呢？它们会各有其职责而一无所用吗？我为什么说耳朵比眼睛更重要呢？因为只有靠耳朵才能获得学问和智慧，而仅靠眼睛却不能。你是为了看天和太阳而生的吗？谁带你来看这景象？或者你的视觉对天和万物的本性有何贡献？无疑是为了让你赞美这浩大而奇妙的工程。因此要承认天主是万物的造物主，祂带你进入这个世界，作为祂伟大工程的见证者和赞美者。你相信观看天和太阳是件大事：那么你为何不感谢那施恩的作者呢？你为何不用心智去衡量你所赞美之工的卓越、眷顾和大能的那位？因为那创造值得赞美之物的，祂自己更当受赞美。若有人邀请你赴宴，你享用了美餐，你若看重享受本身胜过施恩的主人，你显得明智吗？哲学家们就这样完全把一切归于身体，丝毫不归于心智，并且看不到眼前之物以外的东西。但放下身体的一切职份，人的事业应只在于心智。因此我们不是为了看见受造之物而生的，而是为了默观，即用心智观看万物的造物主本身。因此，若有人问一个真正智慧的人他为何而生，他会毫不畏惧、毫不迟疑地回答：他生来就是为了敬拜天主，是祂为了自己的缘故而使我们存在，好让我们事奉祂。而事奉天主无非是通过善行持守和保存正义。但他，作为对神事无知的人，把至关重要的事缩减到最小，只选择了他说要观看的两样东西。但若他说他生来是为了观看世界，尽管这包含了万物，而且使用了更响亮的措辞，他仍未完成人的职责；因为灵魂超越身体多少，天主就超越世界多少，因为天主创造并管理世界。因此，不是世界应当被眼睛默观，因为二者都是身体；而是天主应当被灵魂默观：因为天主本身不死不灭，祂也愿意灵魂永存。而默观天主就是对人类共同之父的敬畏和敬拜。而如果哲学家们缺少这个，并因对神事无知而俯伏于地，我们必须假定\Person{阿那克萨戈拉}既没有看见天也没有看见太阳，尽管他说他生来就是为了看见它们。所以，为人所提出的目标是浅显而容易的，只要他有智慧；而人性尤其属于它。因为人性本身若不是正义，又是什么？正义若不是虔敬，又是什么？而虔敬无非是认天主为父。\par

% structure: p0030 source=h2 target=subsection reason=relative-heading-rank; base=h2

\subsection{第十章 认识并敬拜天主是人的特有属性}

因此，人的至善只在于宗教；因为其他事物，即使是那些被认为专属于人的，在其他动物身上也可以找到。当它们用彼此特有的标志辨别并区分自己的声音时，它们似乎也在交谈；它们也似乎有一种微笑，当它们竖起耳朵、收缩嘴巴、眼睛放松而嬉戏地向人、或向自己的伴侣和幼崽讨好时。它们所表达的问候难道不像是一种相互的爱与宽容吗？同样，那些预见未来并为自己储备食物的动物，显然具有远见。在许多动物中也可以找到理性的迹象。因为既然它们渴望对自己有益的事物，防范邪恶，躲避危险，为自己准备带有不同出口的隐蔽处，它们无疑具有一定的理解力。谁能否认它们具有理性，因为它们常常欺骗人自己？那些负有产蜜职责的动物，当它们居住在分配的地方时，构筑营地，用难以言表的技巧建造居所，并服从它们的王——我不知道它们是否不具有完美的审慎。因此，那些赐予人的事物是否与其他生物所共有，是不确定的：它们肯定没有宗教。我确实这样判断：理性赐予所有动物，但对于哑巴动物，只是为了保护生命，而对于人，则是为了延续生命。而且因为理性本身在人内是完美的，它被称为智慧，这使人在这一点上卓越，即唯独人得以领悟神圣事物。关于这一点，\Person{西塞罗}的意见是真实的：\Quote{在如此众多的动物种类中，}他说，\Quote{除了人，没有一种具有对天主的任何认识；而在人自身中，没有一个民族如此未开化或野蛮，以至于即使它不知道对天主的应有概念，也知道应当持有对祂的某种概念。}由此导致，他承认天主，好像他回忆起自己由之而来的根源。因此，那些希望让心灵摆脱一切恐惧的哲学家，甚至取消宗教，从而剥夺了人独特而卓越的善，这种善与正直的生活和一切与人相关的事物不同，因为天主使一切生物服从于人，也使人人服从于祂自己。他们还有什么理由主张，心灵应当转向面容所朝向的同一方向？因为如果我们必须仰望苍穹，那无疑只为了宗教的缘故；如果宗教被取消，我们与苍穹便毫无关系。因此，我们或者朝那个方向看，或者俯向大地。即使我们愿意，我们也不能俯向大地，因为我们的姿态是直立的。我们必须仰望苍穹，这是身体的本性所呼唤的。如果承认这是必须做的，那么这样做要么是为了投身于宗教，要么是为了认识天体的性质。但我们绝不可能认识天体的性质，因为此类事物无法通过推理发现，如我先前所示。因此，我们必须投身于宗教，而不承担此事的人便俯伏于地，模仿野兽的生活，放弃人的职责。因此，无知者更明智；因为尽管他们在选择宗教时犯错，但他们仍记得自己的本性和处境。\par

% structure: p0032 source=h2 target=subsection reason=relative-heading-rank; base=h2

\subsection{第十一章　论宗教、智慧与至善}

因此，全人类普遍同意，应当投身于宗教；但我们须解释在这一主题上犯下何种错误。天主愿意人的本性是这样的：他应当渴望并热衷两件事：宗教与智慧。但人们在此错误：要么投身宗教而不关注智慧，要么只致力于智慧而不关注宗教，尽管二者不可分离地相连。结果是，他们陷入多种宗教，但都是虚假的，因为他们舍弃了智慧，而智慧本可教导他们不可能有许多神；或者他们致力于智慧，但却是虚假的智慧，因为他们没有关注至高天主的宗教，而天主本可指示他们认识真理。因此，从事这两种道路的人走上歧途，充满极大的错误，因为人的职责与一切真理都包含在这两件不可分离的事中。因此，我惊奇，在哲学家中竟无一人发现至善的居所与所在地。因为他们本可以这样寻求：无论至善是什么，它必须是向所有人提出的一个目标。有快乐，这是所有人都渴望的；但这与野兽相同，没有高尚力量，带来饱和感，过度则有害，随年龄增长而减少，且不是众多人的份：因为缺乏资源的人（占大多数）也必须没有快乐。因此，快乐不是至善；它甚至不是一种善。我们该说什么关于财富呢？这更适用于它们。因为它们落在更少的人身上，而且通常是偶然的；它们常常落在懒惰者身上，有时因罪过而得，并且已被拥有者所渴望。我们该说什么关于统治权本身呢？它并不构成至善：因为并非所有人能统治，但所有人都必须能够达到至善。\par

因此，让我们寻求那向众人展示的东西。是德行吗？不可否认，德行是一种善，而且无疑是众人的善。但若德行不能幸福，因为它的力量和本性在于忍受邪恶，那它就确实不是至善。让我们寻求别的。但没有什么比德行更美，更配得上智者。因为若因丑陋而避免恶习，那么也当因美丽而渴慕德行。那么，如何？那公认为善而可敬的事，竟得不到任何回报，如此贫瘠，以致不能从自身带来益处吗？那伴随着这充满邪恶的生命中的巨大辛劳、困难和争斗，必定要产生某种至善。但我们会称它为什么？快乐吗？但从可敬的事中不能生出卑贱的事。我们会称它为财富吗？或权位吗？但这些都是脆弱而不确定的。是光荣吗？或尊荣？或长久的名声？但所有这些都不包含在德行本身之内，而是依赖于他人的意见和判断。因为德行常被憎恨并遭受恶遇。但由它产生的善应当与它如此紧密相连，以至于无法与它分离或割裂；而除非这善专属于德行，且是如此，以至于不能增删什么，它就无法显得是至善。为什么我要说德行的职责在于轻视这一切？因为不渴望、不追求、不爱慕快乐、财富、统治、尊荣，以及那些被视作善的事物，不像他人被欲望所压倒，那确实是德行。因此，它成就了某种更崇高、更卓越的事；也没有任何东西与这些现世的善争斗，除非那渴慕更伟大、更真实之事的。我们不要对能够找到它感到绝望，只要我们向各方向转动我们的思想；因为所寻求的不是轻微或琐碎的酬报。\par

% structure: p0035 source=h2 target=subsection reason=relative-heading-rank; base=h2

\subsection{第十二章 论身灵的双重战斗；并为永生而渴慕德行。}

但我们所探究的是我们为何而生：这样我们就能追溯出德行的效果。人由两部分组成：灵魂和身体。有许多属于灵魂的，有许多属于身体的，有许多两者共有的，正如德行本身；而每当德行指向身体时，为区别起见，它被称为勇毅。因此，既然勇毅与两者相连，就向每一方提出竞赛，并从竞赛中向每一方提出胜利：身体，因为它是坚实的、可把握的，必须与坚实而可把握的对象争斗；但灵魂，另一方面，因为它轻灵、精妙且不可见，与那些看不见、摸不着的仇敌争斗。但灵魂的仇敌是什么，不就是私欲、恶习和罪过吗？若德行克服并驱散了这些，灵魂就会纯净无染。那么，我们如何能收集到灵魂勇毅的效果呢？无疑，从与它紧密相连并相似的东西，即身体的勇毅；因为当身体面临任何遭遇和竞赛时，它除了从胜利中寻求生命外，还寻求什么？因为无论你与人还是与兽争斗，竞赛都是为了安全。因此，正如身体借胜利从毁灭中获得保全，同样灵魂获得其存在的延续；正如身体被仇敌征服时遭受死亡，同样灵魂被恶习压倒时也必须死亡。那么，灵魂所进行的竞赛与身体所进行的竞赛之间将有什么差别呢？除了身体寻求暂时的生命，而灵魂寻求永恒的生命外？因此，若德行本身不幸福，因为（如我所说）它的全部力量在于忍受邪恶；若它轻视所有被欲求为善的事物；若在其最高境界中它暴露于死亡，因为它常拒绝他人所渴慕的生命，并勇敢地承受他人所恐惧的死亡；若它必然从自身产生某种至善，因为直至死亡的劳苦、忍受和克服不能不获得报酬；若在世上找不到配得上它的报酬，因为它轻视所有脆弱易逝的事物，那么还有什么别的可能，除非它能成就某种天上的报酬，因为它鄙视所有尘世的事物，并能追求更高的事物，因为它轻视卑微之物？而这报酬只能是永生。\par

因此，那位创建了\OriginalTerm{麦加拉学派}{Megareans}体系、与众不同的著名哲学家\Person{欧几里得}{Euclid}很有道理地说，那不变且始终如一者便是\OriginalTerm{至善}{chief good}。他确实理解了至善的本性，尽管他没有说明它由什么构成；但它在于\OriginalTerm{不朽}{immortality}，而非任何其他东西，因为唯有它不能被减少、增加或改变。\Person{塞内加}{Seneca}也不经意地承认，\OriginalTerm{德性}{virtue}除了不朽之外没有别的赏报。因为在他所写的关于过早死亡的文章中赞扬德性时，他说：\Quote{德性是唯一能赐予我们不朽、使我们与神明相等的事物。}但同样地，他所追随的\OriginalTerm{斯多亚派}{Stoics}也说，没有德性，无人能变得幸福。因此，德性的赏报是\OriginalTerm{幸福生活}{happy life}，如果德性正如所言造就幸福生活的话。因此，德性不应如他们所说为自身而追求，而应为幸福生活——它必然跟随德性——而追求。而这个论证本可教导他们至善在于什么。但眼下这个属肉体的生命不能是幸福的，因为它通过身体而遭受邪恶。\Person{伊壁鸠鲁}{Epicurus}称天主是幸福且不朽坏的，因为祂是永恒的。因为幸福状态应当是完美的，以致没有任何事物可以困扰、减少或改变它。除非是不朽坏的，否则任何事物都不能在其他方面被判断为幸福。但除了那不朽的以外，没有事物是不朽坏的。因此，不朽独自是幸福的，因为它既不能朽坏，也不能被毁灭。但如果德性属于人的能力（无人能否认），那么幸福也属于他。因为一个有德性的人不可能是不幸的。如果幸福属于他的能力，那么拥有幸福属性的不朽也属于他。\par

因此，\OriginalTerm{至善}{chief good}唯独在于不朽，这不属于任何其他动物或身体；没有知识的\OriginalTerm{德性}{virtue}——即没有对天主和正义的认识——也不能临到任何人。追求此事的真确与合理，恰由对今生的渴望所表明：因为虽然它只是暂时的，且充满劳苦，却仍为众所寻求与渴望；无论老人幼童、君王与最底层者，总之，智者与愚者都渴望它。正如\Person{阿那克萨戈拉}{Anaxagoras}所认为，对天与光本身的默观如此珍贵，以至于人们甘愿为此承受一切苦难。因此，这短暂而劳苦的生命，不仅为人类，也为其他动物所普遍同意视为一大善事，显而易见，如果它没有终结且免于一切邪恶，它也会成为一种极其伟大而完美的善事。简而言之，若非出于对更长生命的希望，绝不会有人藐视这生命（无论多么短暂）或承受死亡。因为那些为了同胞的安全而自愿赴死的人——如\Place{底比斯}{Thebes}的\Person{墨诺叩斯}{Menœceus}、\Place{雅典}{Athens}的\Person{科德鲁斯}{Codrus}、\Place{罗马}{Rome}的\Person{库尔提乌斯}{Curtius}和两位\Person{穆尔}{Mures}——若非以为将通过同胞的称誉而达到不朽，绝不会选择死亡而放弃生命的益处；尽管他们对不朽的生命无知，但事实本身并未逃过他们的注意。因为如果德性藐视富贵，因为它们是脆弱的，也藐视享乐，因为它们是短暂的，那么它当然会藐视脆弱而短暂的生命，以获取那坚实而持久的生命。因此，反思本身，依照正当次序推进并衡量一切，引导我们到达那卓越而超越的善，我们正为此而生。如果哲学家们曾如此行，如果他们不曾固执地坚持他们一度领会的观点，他们无疑会到达这真理，如我最近所示。如果这不是那些将属天灵魂连同身体一同熄灭之人的做法，那么那些讨论灵魂不朽的人应当明白，德性之所以摆在我们面前，是为了在情欲被制伏、对尘世之物的渴望被克服之后，我们纯洁而胜利的灵魂得以回归天主，即回到它们最初的源头。正因如此，我们作为唯一被高举以仰望天的受造物，使我们相信我们的至善在至高之处。因此，唯独我们领受宗教，为从此知道人的灵魂是不朽的，因为它渴望并承认那不朽的天主。\par

因此，在所有哲学家当中，那些将知识或德行视为至善的人，确实走上了真理的道路，却未能达致圆满。因为这两者共同构成了所追求的目标。知识使我们知道必须通过什么方式、为了什么目的而达成；德行则使我们实际达成。二者缺一不可；因为从知识产生德行，从德行产生至善。因此，哲学家们始终寻求、至今仍在寻求的幸福生活，既不存在于敬拜诸神之中，也不存在于哲学之中；正因如此，他们无法找到它，因为他们不在至高之处寻求至善，反而在至低之处寻求。因为什么是至高？岂非天，以及灵魂由之而出的天主？什么是至低？岂非大地，以及身体由之而受造之物？因此，尽管有些哲学家将至善归于灵魂而非身体，但既然他们将至善关联于这个随身体而终结的生命，他们便又回到了身体——这整个在大地上度过的时光都与身体相关。所以，他们未能获得至善并非没有理由；因为凡仅指向身体、且缺乏不朽者，必定是至低之物。因此，幸福并非如哲学家所想的那样降临到人的境况中；而是这样降临于人：并非当人活在身体中时——这身体为了解体必定要被朽坏——而是当灵魂脱离与身体的交往、仅活在精神之中时。唯独在这一件事上，我们此生可能幸福：如果我们显得不幸福；如果我们避开享乐的诱惑，献身于德行的服役，活在一切劳苦与磨难之中——这些正是操练并坚固德行之手段；简而言之，如果我们坚持那条为我们开辟的通往幸福、崎岖艰难的道路。因此，使人幸福的至善不可能存在，除非在那种与不朽希望相连的宗教与教义之中。\par

% structure: p0040 source=h2 target=subsection reason=relative-heading-rank; base=h2

\subsection{第十三章　论灵魂不朽、智慧、哲学与雄辩。}

此处，主题似乎要求我们——既然已经教导不朽是至善——同时证明灵魂是不朽的。关于这一点，哲学家们有过重大争论；而那些对灵魂持正确见解的人，也未能解释或证明任何事：因为他们缺乏神圣知识，既未能提出足以制胜的真实论证，也未能提供足以说服的证据。但我们将在最后一卷更恰当地讨论这个问题，届时我们将论述幸福生活的主题。剩下的那第三部分哲学，他们称之为逻辑，其中包含整个辩证法的主题和整个言说方法。神圣学问并不需要它，因为智慧的所在不是舌头，而是心怀；你使用何种语言并不重要，因为问题不在于言辞，而在于事实。我们并非争论文法学家或修辞学家——他们的知识关乎恰当的表达方式——而是争论智者，其学识关乎正确的生活方式。但如果前面所说的自然哲学不是必要的，这逻辑学也不是必要的（因为它们不能使人幸福），那么哲学的全部力量就只包含在伦理部分之中（据说\Person{苏格拉底}舍弃其他部分而专心于此）。既然我已经证明，哲学家们在这部分也犯了错误——他们没有把握住我们为之而生的至善——那么哲学就显得完全虚假而空洞，因为它既未预备我们履行正义的职责，也未坚固人生义务的稳定进程。因此，让他们明白：那些以为哲学就是智慧的人是错误的；不要被任何人的权威所牵引，而要倾向真理，接近它。此处容不得鲁莽：如果我们被一个空洞的品性或虚假的意见所欺骗，我们将永永远远承受愚昧的惩罚。然而，人——若他信赖自己，即信赖人——不仅（且不说愚昧，因他看不见自己的错误）无疑是傲慢的，竟敢为自己要求人的境况所不容许的东西。\par

我们也能从这位最伟大的罗马语言作者的这段感言中看出他受骗多深；因为当他在\WorkTitle{论义务}一书中说哲学无非是对智慧的渴望，而智慧本身是对神圣事物与人间事物的知识时，他接着又说：\Quote{若有人指责这种渴望，我确实不明白他认为还有什么值得称赞。因为若寻求心灵的享受和免于忧虑，还有什么享受能与那些总是探究与美好幸福生活相关并有助于促进它的人的研究相比呢？或者若考虑到一贯性和德行，要么是这种研究使我们得以获得它们，要么根本没有任何研究。说最重要之事没有体系，而最小之事却都有体系，这是轻率发言之人的做法，在最重要之事上犯了错误。但如果德行有任何学问，当你离开那种学问时，又该去哪里寻找呢？}就我而言，虽然我因渴望教导他人而在某种程度上努力获取学问，但我从未擅长辞令，因为我甚至从未公开演说；但事业之善本身必能使我善辩，而神圣的知识和真理本身足以清晰而充分地辩护。因此，我希望西塞罗暂时从死中复活，好让这位如此完美的辩者由一个无辞令之才的卑微之人教导，首先，什么是指责那称为哲学的研究之人认为值得称赞的；其次，德行和正义并非由那种研究习得，也不是像他所想的任何其他研究；最后，既然德行有一门学问，当他放下那种学问时，他可能被教导该去哪里寻找它，而他并非为了听闻和学习而寻求它。因为他能从谁那里听到呢？当时无人知道它。但正如他通常的做法，在诉讼辩护中，他希望通过诘问来迫使对手，从而引导他承认，仿佛他确信无人能回答以表明哲学不是德行的教授。在\WorkTitle{图斯库卢姆论辩}中，他公开宣称了这一点，将言谈转向哲学，仿佛通过雄辩的演讲风格来炫耀自己。\Quote{哦，哲学，生命的向导，}他说，\Quote{哦，德行的探索者，邪恶的驱逐者；没有你，不仅我们，连人类生活又能做些什么呢？你是法律的发明者，你是道德和纪律的教师；}——仿佛她真能自己感知什么，而他不更应该称赞那位赐予她的人。同样，他本可感谢饮食，因为没有它们生命无法存在；然而这些虽服务于感官，却不带来益处。但这些是身体的滋养，智慧却是灵魂的滋养。\par

% structure: p0043 source=h2 target=subsection reason=relative-heading-rank; base=h2

\subsection{第十四章　论\Person{卢克莱修}等人及\Person{西塞罗}本人在确定智慧起源上的错误}

因此，\Person{卢克莱修}的做法更正确，他在赞美那第一个发现智慧的人；但他愚蠢之处在于，他假设智慧是由人发现的——仿佛他赞美的那个人像诗人们的传说中那样在泉边找到了像笛子一样的东西。但若他赞美智慧的发现者如神——因为他如此说：——\par

\Quote{我认为，没有一个是血肉之躯的人。因为如果我们必须说，正如主题本身的公认庄严所要求的，他曾是神，他曾是神，最高贵的\Person{梅米乌斯}，}——\par

然而天主不应因此被赞美，因祂发现了智慧，而是因祂创造了能接受智慧的人。因为只赞美整体的一部分，就贬低了赞美。但他赞美祂如同一个人；而祂正因发现了智慧而本应被视为天主。因为他如此说：——\par

\Quote{难道不该把这人与诸神并列吗？}\par

由此看来，他要么想赞美\Person{毕达哥拉斯}（如我所说，他是第一个自称哲学家的人），要么想赞美\Person{米利都的泰勒斯}（据称他是第一个讨论自然本性的人）。这样，当他试图提高时，反而贬低了那事物本身。因为若它能被人类发现，便不伟大。但作为诗人，他可被宽恕。然而那位同样完美的演说家、同样完美的哲学家，也指责希腊人，他常指责他们的轻浮，却又效仿他们。智慧本身，他时而称之为诸神的恩赐，时而称之为诸神的发明，他照诗人的方式塑造它，因它的美丽而赞美它。他还沉痛抱怨有人贬低它。\Quote{难道有人，}他说，\Quote{敢指责生命之母，并以这种弑亲之罪玷污自己，如此不敬地忘恩负义吗？}\par

那么，\Person{马尔库斯·图利乌斯}，我们就是弑亲者，在你看来应该被缝进皮袋里，因为我们否认哲学是生命之母吗？或者你，对天主如此不虔敬地忘恩负义（不是这位坐在\Place{卡皮托利}山上、你崇拜其形象的那位神，而是那位创造世界、造人并将智慧也作为天上恩赐赐予的天主），你称她为德行之师或生命之母，向她学习的人，难道不会比从前更加不确定吗？因为她教的是什么德行？因为哲学家们至今未能解释德行在哪里。她又是哪条生命之母？因为教师们自己在决定合适的生活道路之前，就已经被衰老和死亡耗尽。你能宣称她是哪条真理的探索者？既然你经常证明，在如此众多的哲学家中，迄今没有一个真正的智者。那么，那位生命的女大师教了你什么？是用责骂攻击最有权力的执政官，用你的毒舌使他成为国家的敌人吗？但我们姑且放过那些可以用命运之名来辩解的事。实际上，你确实投身于哲学研究，而且没有人比你更勤奋；因为你熟悉所有哲学体系，正如你自己常夸耀的那样，并且用拉丁文阐明了这门学问本身，显示自己是\Person{柏拉图}的模仿者。因此告诉我们，你学到了什么，或者你在哪个学派中发现了真理。无疑是你所遵循并赞同的学园派。但这什么也不教，除了让你知道自己的无知。因此，你自己的书驳斥了你，并显示出从哲学中可能获得的学问对生命的虚无。这是你的话：\Quote{但在我看来，我们不仅对智慧是盲目的，而且对那些似乎在一定程度上可以被辨认的事物也是迟钝和愚钝的。}因此，如果哲学是生命之师，为什么你觉得自己是盲目、迟钝和愚钝的？而你本应在她教导下，既能看见、又能明智，并处于最清晰的光明中。但我们从写给你儿子的信中得知，你是如何承认了哲学的真理，你在信中劝告他，哲学的原则应该知道，但我们必须作为社会成员生活。\par

有什么话能如此矛盾？如果哲学的原则应该知道，那么知道它们就是为了使我们生活得美好和明智。或者，如果我们必须作为社会成员生活，那么哲学就不是智慧，如果按照社会而不是按照哲学生活更好的话。因为如果被称为哲学的是智慧，那么不按哲学生活的人肯定是愚蠢地生活。但如果按社会生活的人不愚蠢，那么按哲学生活的人就愚蠢。因此，根据你自己的判断，哲学被谴责为愚蠢和空虚。而且，在你的\WorkTitle{安慰论}中，即不是一部轻率和欢快的作品，你引入了关于哲学的这种情感：\Quote{但我不知道什么错误占据了我们的心，或对真理的可悲无知。}那么，哲学的指引在哪里？或者，那生命之母教了你什么，如果你对真理是可悲地无知？但如果这种对错误和无知的认可是几乎违背你的意志从你内心深处被迫挤出的，为什么你最终不向自己承认这个真理：哲学，尽管它什么也不教，你却对它大加赞扬直到天上，它不可能是德行之师？\par

% structure: p0051 source=h2 target=subsection reason=relative-heading-rank; base=h2

\subsection{第十五章 \Person{塞内卡}在哲学中的错误，以及哲学家的言语如何与他们的生活不一致。}

受同一错误的影响（因为当\Person{西塞罗}犯错时，谁能保持正确的方向呢？），\Person{塞内加}说：\Quote{哲学无非是正确的生活方法，或体面生活的科学，或度过美好生活的艺术。我们若说哲学是美好而体面生活的法律，也绝不会错。那称它为生活准则的人，给了它应有的地位。}他显然不是指哲学的普通名称；因为，既然哲学散布到许多派别和体系之中，没有任何确定的东西——简言之，没有众口一词、同心同意的事——那么，还有什么比把哲学称为生活的准则更虚假的呢？既然其训条的多样性阻碍正道、引起混乱？或是美好生活的法律，而它的服从者却广泛不一？或是度过生活的科学，在其中反复的矛盾除了普遍的不确定性之外一无所成？因为我要问，他是否认为\OriginalTerm{学园派}{Academy}是哲学？我想他不会否认。若是如此，那么这些说法中没有一个与哲学相符；哲学使一切不确定，废除法律，视艺术为无物，颠覆方法，扭曲规则，完全取消知识。因此所有那些都是虚假的，因为它们与一个始终不确定、至今解释不了任何东西的体系不符。因此，除了这唯一真实的、属天的智慧——哲学家们对此并不知晓——之外，没有任何体系、科学或美好生活的法律被建立起来。因为那属地的智慧，由于它是虚假的，变得多样多变，且完全自相矛盾。正如世界的创造者与管理者天主只有一个，真理是一；同样，智慧也必须是一且纯一，因为任何真实良善的事物，除非是同类中唯一的，否则不能完美。但如果哲学能够塑造生活，那么除了哲学家之外就没有好人，所有未学过哲学的人将始终是恶的。然而，有并且总有无数的没有学问却曾是或现在是善良的人，而哲学家中却很少有人在生活中做过任何可称赞的事；请问，谁能看不出那些人不是美德的教师——他们自己尚且缺乏美德？因为若有人仔细查考他们的品性，便会发现他们易怒、贪婪、好色、傲慢、放荡，并以智慧的外表掩盖自己的恶习，在家里做那些他们在学校中谴责的事。\par

也许我说谎是为了指控。难道\Person{图留斯}不也承认并抱怨同样的事吗？他说：\Quote{多么少的哲学家被发现具有这样的品格，如此在灵魂和生活中构成，正如理性所要求的！多么少的人认为真正的教导不是知识的展示，而是生活的法律！多么少的人顺从自己，服从自己的决断！我们可以看到有些人如此轻浮和炫耀，以至于他们最好根本没有学习过；另一些急切渴望金钱，另一些渴望荣耀；许多人是情欲的奴隶，以至于他们的言语与生活惊人地不一致。}\Person{科尔内利乌斯·内波斯}也给同一位\Person{西塞罗}写信说：\Quote{我远不认为哲学是生活的导师和幸福的完成者，反而认为许多从事这门学科讨论的人比任何人都更需要生活的导师。因为我看到，那些在学校中给出最精细教训——关于谦逊和自制——的人，同时却在所有情欲的无节制欲望中生活。}\Person{塞内加}也在他的\WorkTitle{劝勉集}中说：\Quote{许多哲学家正是这种描述，雄辩地谴责自己：因为如果你听他们反对贪婪、反对情欲和野心，你会认为他们是在公开揭露自己的品格，他们在公开场合的谴责如此完全地流回到自己身上；因此，除了把他们视为医生——其广告包含药物，但药箱里却是毒药——之外，没有别的看法是正确的。有些人不以自己的恶习为耻；反而为自己的卑鄙编造辩护，以便他们甚至可以体面地犯罪。}\Person{塞内加}还说：\Quote{智者甚至会做他不赞同的事，以便找到通往完成更大事物的途径；他也不会放弃良好的道德，而是使之适应场合；别人为荣耀或快乐所使用的东西，他将为了行动而使用。}然后他不久后说：\Quote{奢侈者和无知者所做的一切，智者也会做，但不是以同样的方式和同样的目的。但你行动时的意图并不重要，因为行为本身是邪恶的；因为行为可见，意图不可见。}\par

昔居勒尼学派之师\Person{亚里斯提卜}与名妓赖伊斯有淫乱之私。彼哲学家以严师自居，为此过辩护曰：彼与赖伊斯其他情人之别，在于彼拥有赖伊斯，而他人则为赖伊斯所拥有。呜呼！此等智慧，何其光辉，堪为正人君子所效法乎？汝果愿以子女托付此人受教，俾其学习占有娼妓乎？彼自称与放荡之徒有别：彼等耗其家产，而已则无需费财而纵情逸乐。然此娼妓实更明智，盖彼视哲学家为己之弄臣，使所有少年，因师之榜样与权威而腐化，毫无羞愧地群集于彼。是故，哲学家以何动机投身于彼声名狼藉之娼妓，有何区别？当众民及敌手见其比一切堕落之人都更卑劣时，已无可辩。彼不惟如此生活，更进而传授情欲，将其习气从妓院移入学堂，力主肉体快乐为至善。此等有害而可耻之学说，非源于哲学之心，乃出于娼妓之怀。\par

何须言及犬儒学派？彼等当众纵欲，岂足为奇？彼等既效犬之生活，则其名号取自犬类，又何怪乎？故此派无德行教诲，即令倡导较高尚之事者，或自身不行其所劝；或偶行之（此事甚罕），亦非其体系导之向正，乃自然常引未受教者归于赞许。\par

% structure: p0056 source=h2 target=subsection reason=relative-heading-rank; base=h2

\subsection{第十六章　哲学家常训善而行恶，西塞罗可证；故吾人不应专务哲学，而当追求智慧。}

然彼等长年耽于怠惰，不行任何德操之操练，毕生以言说为事，则吾人可视之为何如人？岂非虚浮之辈？盖智慧若不施行于其可用力之行动，则为空洞而虚假。\Person{图利乌斯}（西塞罗）确然以从事国务、治理邦国、创建新城或以正义维护已建之城、以良法或有益之议或庄重裁判保全公民之安全与自由者，置于哲学家之师之上。盖使人成善，胜于向隐居斗室者训以职责，况且训者自身亦不遵行；彼等既远离真实行动，则显然其哲学体系之发明，仅为练习口舌或辩护之故。但仅教而不行者，自减其训诫之分量；盖发令者自身教人不服从，谁肯服从？且给予正确荣耀之训诫，诚为善事，但若汝不行之，则属欺骗；口善而心不善，乃矛盾而虚浮。\par

故彼等求于哲学者，非效用，乃享乐。此亦为\Person{西塞罗}所证。彼云：\Quote{诚然，}彼谓\Quote{其全部辩论虽含极丰之德性与知识之源，然较之其行为与成就，吾恐似未予人事之事务如许利益，反予闲暇之娱乐。}彼不当惧，因其言真；然彼似恐因泄露玄秘而遭哲学家控告，故未敢坦然宣告真理：即彼等并非为教而辩，乃为闲暇自娱；且彼等既为行动之谏者而自身毫无行动，当视为空谈者。然确因彼等未给生活带来益处，既未遵其自家之法，亦无人历经如此多世纪而按其律生活。故哲学必须全然摒弃，因吾人不应专求智慧——此无界限与中庸——而必须立时成为智者。盖第二生命未赐予我们，以致今生求智，来世方智；每项结果当在此生达成。智慧当速寻，以便速取，以免生命任何部分流逝——其终结不确定。\Person{西塞罗}《\WorkTitle{霍尔滕修斯}》中反驳哲学者，为巧妙论证所迫：彼谓人不当哲学，然彼似仍在哲学，盖讨论人生所当为与不当为者，乃哲学家之事。吾人则免此责难，因吾人废除哲学——人为思想之发明——而维护智慧——神圣传统，且证其为众人所当取。彼废除哲学而未引入更佳者，故被认为废除智慧；因此其意见更易被驳倒，因公认人非生于愚昧，乃生于智慧。\par

且同一\Person{霍尔滕修斯}所用之论证，于哲学亦甚有力：即由此可知哲学非智慧，因其起源与开端昭然。彼问：哲学家何时始有？首位者，我想是\Person{泰勒斯}，而其时代甚近。然则古代人中，求真理之爱隐藏何处？\Person{卢克莱修}亦云：\par

\Quote{此外，此事物之本性及体系，新近始被发现，而我乃首先唯一能将其转为母语者。}\par

塞内加说：\Quote{智慧的开端至今尚未满千年。}因此人类在很多世代中没有体系。佩尔西乌斯嘲讽此事说：\newline{}\par

\Quote{当智慧来到城中，\newline{}连同胡椒和棕榈枝；}\par

仿佛智慧是随着美味商品一同被引入城中的。因为如果智慧符合人的本性，它必然与人类同时起源；但如果它不符合本性，人性就无法接纳它。然而，既然人性接纳了它，那么智慧就是自古就有的：因此，哲学既然不是自古就有的，就不是同一真正的智慧。但事实上，希腊人因为未能达到真理的圣书，不知道智慧是如何被败坏的。因此，他们认为人类生活缺乏智慧，便发明了哲学；也就是说，他们想通过讨论来发掘那隐藏在他们面前未知的真理；由于对真理的无知，他们便视这种努力为智慧。\par

% structure: p0064 source=h2 target=subsection reason=relative-heading-rank; base=h2

\subsection{第十七章：从哲学转向哲学家，始于伊壁鸠鲁；以及他如何将留基伯和德谟克利特视为错误的始作俑者。}

我已经尽可能简洁地论述了哲学本身；现在让我们来谈哲学家，不是要与这些站不住脚的人争辩，而是要追击那些从我们战场上逃跑和溃败的人。伊壁鸠鲁的体系比其它体系更为普遍地被人追随；这并不是因为它提出了任何真理，而是因为快乐的诱人名称吸引了众人。因为每个人都自然地倾向于恶习。此外，为了吸引众人归向他，他针对每个不同的性格说出特别合适的话。他禁止懒惰者投身学问；他使贪财者免于向民众广施财物；他禁止怠惰者承担国事，阻止迟钝者锻炼身体，阻止怯懦者服兵役。不虔敬的人被告知诸神不关心人的行为；无情和自私的人被命令不要给任何人任何东西，因为智者一切都为自己。对避世的人，孤独受到赞扬；对过于俭省的人，他得知清水和面食可以维持生命；如果一个人憎恨妻子，便向他列举独身的福气；对有不肖子女的人，宣扬无子女者的幸福；对违背自然的父母，他说没有自然的纽带。对娇弱和不能忍耐的人，他说痛苦是万恶之首；对刚毅的人，他说智者即便在酷刑下也是幸福的。对追求权势和荣耀的人，他命令巴结国王；对不能忍受烦扰的人，他命令避开王宫。因此这个狡猾的人从各种各样、彼此不同的人中聚敛了一群听众；当他力求取悦所有人时，他自身比他们所有人彼此之间更加矛盾。但我们必须说明这整个体系源自何处，有何起源。\par

伊壁鸠鲁看到善人总是遭遇逆境、贫困、劳苦、流放、失去挚友。相反，他看到恶人享福，权势显赫，荣誉加身；他看到无辜者得不到保护，罪行不受惩罚；他看到死亡肆虐，毫无品性之分，也没有任何年龄的安排或区别，有些人活到老年，而另一些在幼年就被夺走，有些人在强壮有力时死去，其他人则在青春初期夭折；他看到在战争中，更优秀的人尤其被征服和杀戮。但特别触动他的是，他看到敬神的人尤其遭受更重的灾祸，而那些完全忽视诸神或用不敬方式崇拜诸神的人，却遭受较轻或毫无灾祸；甚至神庙本身也常常被闪电击中起火。卢克莱修对此抱怨，论及神明时说：\par

\Quote{于是他可能投掷闪电，常常摧毁自己的神庙，退入荒郊，在那里发泄他的霹雳，这霹雳常常放过有罪者，却击毙无辜和无过之人。}\par

但如果他能收集到丝毫真理，他就绝不会说神明摧毁自己的神庙，因为他摧毁它们正是由于它们不属于他。卡皮托利山，作为罗马城和宗教的主要所在地，曾不止一次被闪电击中并起火。但聪明人对这事的看法可以从西塞罗的话中看出，他说那火焰从天而降，不是为了摧毁朱庇特这地上居所，而是为了求得更崇高壮丽的住所。关于此事，在他关于自己执政官任期的著作中，他说了与卢克莱修类似的话：\par

\Quote{因为高踞奥林帕斯的天父在至高天打雷，亲自袭击自己的堡垒和著名的神庙，并将烈火投向他在卡皮托利的住所。}\par

因此，在他们愚顽的固执中，他们不仅不理解真天主的权能和威严，甚至还增加他们错误的邪妄，试图违背一切神圣律法，去修复一个屡次被上天的审判定罪的神庙。\par

因此，当\Person{伊壁鸠鲁}思量这些事时，仿佛被这些事的不公所诱导（因他在无知于原因和主题时，事情在他看来如此），他认为没有天主眷顾。他使自己确信这一点后，还着手为它辩护，从而将自己纠缠于无法解决的错误中。因为若没有天主眷顾，世界怎能被造得如此有序且有条理？他说：没有条理，因为许多事物被造成的方式不同于它们应当被造成的方式。这位神圣之人找到了可指责之处。如今，若我有闲暇逐条驳斥这些，我本可轻易表明此人既不明智也不神志健全。再者，若没有天主眷顾，动物的身体如何被如此远见地安排，使得各个肢体以奇妙的方式配置，各自履行其职责？他说，天主眷顾的系统在动物的产生中并未筹划任何事；因为眼睛并非为看而生，耳朵并非为听而生，舌头并非为说话而生，脚并非为行走而生；既然这些在可能说话、听见、看见和行走之前就已产生。因此它们并非为使用而产生，而是使用从它们中产生。若没有天主眷顾，为何降雨、果实生长、树木发芽？他说，这些事并非总是为生物的缘故而做，因它们对天主眷顾无益；但一切事物必自行产生。那么，它们从何而来，或一切进行之事如何发生？他说，无需假设天主眷顾；因为有种子飘浮在虚空中，从这些无序聚集的种子中，一切事物产生并成形。那么，我们为何不察觉或区分它们？因为他说，它们既无颜色，也无热，也无气味；也无滋味和湿度；它们如此微小，以至于不能被切割和分割。\par

这样，因为他一开始就接受了一个错误的原则，后续问题的必然性将他引向了荒谬。因为这些原子在何处或从何而来？为何除了\Person{留基伯}一人之外，没有人想到它们？\Person{德谟克利特}从他那里接受了教导，将他的愚昧遗产留给了\Person{伊壁鸠鲁}。若这些是微小的物体，且如他们所说确实是坚固的，它们当然能够落入眼睛的注意中。若万物的本性相同，它们如何组成各种对象？他说，它们以多样的次序和位置相遇，正如字母，虽数量有限，却通过多样的排列组成无数的词。但有人反驳说，字母有各种形状。他说，这些第一原理也是如此；因为它们粗糙，有钩子，光滑。因此，若其中有任何突出的部分，它们能被切割和分割。但若它们光滑且无钩，就不能粘合。因此它们应当有钩，以便相互连接。但由于它们据说微小到任何武器的锋刃都不能将其切断，它们又如何能有钩子或棱角？因为既然它们突出，就必须能够被扯开。其次，它们以何种相互契约、何种辨识力相遇，以便从中构造出任何事物？若它们没有理智，就不能以如此次序和条理相遇；因为除理性外，没有东西能按理性完成任何事情。多少论证可以驳斥这琐碎之言！但我必须继续我的主题。这就是他。\par

\Quote{他在才智上超越了人类，并熄灭了所有人的光明，如同升起的以太太阳熄灭了星辰。}\par

这些诗句我从未能不笑而读。因为这不是关于被尊为哲学王之\Person{苏格拉底}或\Person{柏拉图}说的，而是关于一个虽然心智健全、身体强健，却比任何患病的人更无意义地狂言的人。因此，那最虚浮的诗人，我不说是装饰，而是以对狮子的赞美压倒压碎了老鼠。但同一个人也使我们摆脱对死亡的恐惧，关于此他的原话如下：\par

\Quote{当我们存在时，死亡不存在；当死亡存在时，我们不存在：因此死亡与我们无关。}\par

他多么巧妙地欺骗了我们！仿佛现在完成的死亡——感觉已被剥夺——才是恐惧的对象，而实际上真正的恐惧是正在死去的行动——感觉正被从我们身上剥夺。因为有一个时间，我们自己甚至还存在，死亡尚未存在；而那个时间似乎是悲惨的，因为死亡开始存在，而我们正停止存在。\par

说死亡并不可悲，并非毫无理由。可悲的是死亡的临近：即因疾病而衰弱，忍受刺伤，身体被兵器击中，被火烧灼，被野兽的牙齿撕裂。这些才是人所畏惧的，并非因为它们带来死亡，而是因为它们带来巨大的痛苦。然而，\Person{伊壁鸠鲁}倒宁愿说痛苦并非邪恶。\Person{伊壁鸠鲁}却说痛苦是万恶之中最大的。既然如此，如果那伴随或导致死亡的事物是邪恶的，我怎能不畏惧呢？我为何要说这论证是假的，既然灵魂并不消亡？但\Person{伊壁鸠鲁}说，灵魂是会消亡的；因为与身体一同诞生的，必与身体一同消亡。我已经说过，我宁愿推迟讨论这个问题，把它留到我著作的最后部分，以便用论证和神圣的证言来驳斥\Person{伊壁鸠鲁}的这一信念（无论是来自\Person{德谟克利特}还是\Person{迪凯亚尔库斯}）。但也许他自以为可以在放纵恶习中逍遥法外；因为他提倡最可耻的快乐，并说人生来就是为了享受快乐。听到这种说法的人，谁会戒绝恶习和邪恶的行为呢？因为如果灵魂注定要消亡，那就让我们热切追求财富，以便能够享受各种放纵；如果我们没有财富，那就让我们偷窃、用计谋或暴力从拥有者那里夺取，尤其如果没有神察看人的行为的话：只要不受惩罚的希望对我们有利，我们就掠夺和杀人。因为聪明人做恶事，只要对他有利且安全，就是聪明的；既然天上若有神，他不向任何人发怒。同样，愚昧人行善也是如此；因为他既不被愤怒激动，也不被恩惠影响。因此，让我们以各种可能的方式享受快乐；因为过不了多久，我们就根本不存在了。因此，不要让任何一天、甚至一刻时光从我们身边溜走而没有快乐；免得我们自身注定要消亡，连我们已经度过的生命也一同消亡。\par

虽然他并没有在言语上这样说，但他实际上是在这样教导。因为他主张聪明人做一切事都是为了自己，他把所做的一切都归于自己的利益。因此，听到这些可耻之事的人，就不会认为应当行任何善事，因为施恩是为了他人的利益；也不会认为应当戒除罪行，因为作恶会带来利益。如果某个海盗头目或强盗首领要鼓动他的手下使用暴力，他还能用什么别的语言呢？他只能说出与\Person{伊壁鸠鲁}相同的话：诸神并不留意；他们既不被愤怒影响，也不被慈爱影响；不必惧怕未来的惩罚，因为死后灵魂就死了，根本没有未来的惩罚；快乐是至善；人与人之间没有社会；每个人都只为自己打算；没有人爱别人，除非是为了自己；勇敢的人不应畏惧死亡或任何痛苦；因为即使他被折磨或被焚烧，他也应当说他不介意。显然，有充分理由让人把这视为一个聪明人的说法，因为它最适合应用到强盗身上！\par

% structure: p0079 source=h2 target=subsection reason=relative-heading-rank; base=h2

\subsection{第十八章　毕达哥拉斯学派和斯多葛学派虽持灵魂不朽之说，却愚蠢地劝人自愿死亡}

另有一些人则持有相反的观点，即灵魂死后仍存；这些人主要是毕达哥拉斯学派和斯多葛学派。虽然因为他们领悟了真理而可以宽容，但我不得不责备他们，因为他们并非通过自己的见解，而是偶然地撞上了真理。因此，即使在正确领悟的那件事上，他们也在某种程度上犯了错误。因为，他们害怕那个推断出灵魂必然随身体死亡的论证，即灵魂与身体一同诞生，因而宣称灵魂并非与身体一同诞生，而是被引入身体，并从一具身体迁移到另一具身体。他们没有想到，灵魂在身体之后仍能存续，除非灵魂似乎先于身体存在。因此，双方都犯了同样的、几乎相似的错误。但一方在过去上受了欺骗，另一方在未来上受了欺骗。因为没有人看到最真实的一点：灵魂既被创造又不会死，因为他们不知道这事为何发生，也不知道人的本性是什么。因此，他们中有许多人，由于猜想灵魂是不朽的，便对自己施以暴力，好像他们要升天一样。例如\Person{克莱安忒斯}、\Person{克吕西波}、\Person{芝诺}，以及\Person{恩培多克勒}（他在深夜投身于燃烧的\Place{埃特纳}火山口，以便他突然消失时可以被相信他已归向众神）；罗马人中还有\Person{加图}，他一生都在模仿\Person{苏格拉底}的做作。但\Person{德谟克利特}持另一种信念。然而，\par

他自愿把自己的头颅献给了死亡；\par

再没有比这更邪恶的事了。因为如果杀人犯因为毁灭人而有罪，那么自杀的人也同样有罪，因为他杀死了一个人。不仅如此，这个罪行可以算是更大的，因为惩罚只属于天主。因为我们不是自愿进入此生；同样，我们只能按照那把我们放在这身体里、让我们居住直到他命令我们离开的那位的命令，才能离开这被指定给我们看守的身体居所；如果有人向我们施暴，我们必须以平静之心忍受，因为无辜者的死亡不会不受报复，而且我们有一位伟大的审判官，他独自永远掌握着报复的大权。\par

因此，所有这些哲学家都是杀人犯；连罗马智慧之首加图也不例外，据说他在自杀前通读了柏拉图论灵魂不朽的著作，并因这位哲学家的权威而犯下这一大罪；不过，他似乎因仇恨奴役而有某种死亡的理由。我又何必提起那位安布拉基亚人呢？他读了同样的著作，就投海自尽，没有别的原因，只因为他相信了柏拉图——这种教义全然可憎，应当避免，若它驱使人离开生命的话。但倘若柏拉图知道并教导过，由谁、以何种方式、赐给何人、因何种行为、在何时赐予不死，他就不会驱使克莱欧布洛图斯或加图自愿赴死，而会训练他们正义地生活。因为在我看来，加图寻求死亡的理由，与其说是为了逃避凯撒，不如说是为了服从他所追随的斯多葛派的信条，并借某一壮举使其声名显赫；我看不出他若活着会遭遇什么不幸。因为盖乌斯·凯撒，尽管身处内战最激烈之时，其仁慈竟如此之大，其目的无非是借保全两位杰出公民——西塞罗和加图——以显得对国家有功。但让我们回到那些赞美死亡为好处的人吧。你抱怨生命，仿佛你曾活过，或曾为自己定下为何出生的目的。难道一切共同的真正圣父不能正当地指责泰伦提乌斯的这句名言：——\par

\Quote{首先，学习生命是由什么构成的；然后，你若对生命不满，就可求助于死亡。}\par

你愤恨自己遭受灾祸，仿佛你配得上任何好处，你竟不认识你的圣父、主和君王；你虽然眼睛看见明亮的光，却心灵盲目，躺在无知黑暗的深处。而这种无知使一些人竟不羞愧地说：我们为此而生，就是要承受我们罪过的惩罚；但我看不出还有什么比这更愚蠢的了。因为我们根本不存在时，在哪里或犯了什么罪过呢？除非我们碰巧相信那位糊涂的老头，他谎称自己以前活过，并在前世曾是欧福耳布斯。我相信，他因出身卑微，便从荷马的诗歌中为自己选了一个家族。啊，毕达哥拉斯多么奇妙卓越的记忆！啊，我们大家多么可悲的遗忘，因为不知道在前世我们是谁！但也许是因某种错误或偏爱，唯独他没有触及勒忒河的深渊，或品尝遗忘之水；无疑，这位轻浮的老头（正如无事可做的老妇常有的那样）仿佛为轻信的幼童编造了寓言。但如果他对他说话的那些人评价甚好，如果他认为他们是人，他就绝不会妄自容许自己说出如此悖谬的谎言。但这最轻浮之人的愚行是值得嘲笑的。对西塞罗我们又该怎样呢？他在其\WorkTitle{安慰}的开头说，人生来就是为了赎罪，后来他又重复这一论断，仿佛在责备那人不认为生命是一种惩罚？因此，他先前说自己被错误和可悲的无知所困，是对的。\par

% structure: p0086 source=h2 target=subsection reason=relative-heading-rank; base=h2

\subsection{第十九章 西塞罗和其他最聪明的人教导灵魂不死，但以一种不信的方式；并且，善死或恶死必须从先前的生活来衡量。}

但那些主张死亡有益的人，因为他们不认识真理，如此推论：如果死后一无所存，死亡并非灾祸；因为它消除了对灾祸的感知。但如果灵魂存续，死亡甚至是有益的；因为随之而来的是不朽。\Person{西塞罗}在\WorkTitle{论法律}中如此表述这一见解：\Quote{我们可以自贺，因为死亡将带来的状态要么比生命中的更好，至少不会更糟。因为如果灵魂在没有身体的情况下充满活力，那便是神圣的生命；如果它没有感知，那就肯定没有灾祸。} 在他看来，这论证很巧妙，仿佛不可能有其他状态。但每个结论都是错误的。因为圣经教导我们，灵魂并非消灭；而是根据其义德获得赏报，或根据其罪行受到永罚。因为让过邪恶生活而享顺遂的人逃脱应得的惩罚是不对的；让因义德而受磨难的人被剥夺赏报也是不对的。这一点是如此真实，以致\Person{西塞罗}也在他的\WorkTitle{安慰}中宣称，义人与恶人不居住在同一处所。他说，那些智者并不认为所有人进入天堂的途径相同；因为他们教导说，那些被恶习和罪行玷污的人被推入黑暗，躺在泥淖中；而另一方面，贞洁、纯洁、正直、无玷的灵魂，又因修习和践行德行而得以精炼，便轻快而容易地飞向诸神，亦即飞向与他们相似的本性。但这一见解与前一个论点是矛盾的。因为前者基于每人出生时就被赋予不朽的假设。那么，美德与罪愆之间将有何区别，如果一个人是\Person{阿里斯提德}还是\Person{法拉里斯}，是\Person{加图}还是\Person{喀提林}都毫无区别？但除非人拥有真理，否则他看不到这种见解与行为之间的矛盾。因此，如果有人问我死亡是善是恶，我将回答，它的性质取决于生命的过程。因为正如生命本身若是善行度过则为善，若是邪恶度过则为恶，同样，死亡也应根据生命的过往行为来权衡。因此，如果生命是在事奉天主中度过的，死亡就不是恶，因为它是迁移至不朽。但若不是，死亡就必然是恶，因为如我所说，它将人转移到永罚。\par

那么，我们该说什么？岂不是说他们错了，要么把死亡当作善而渴望，要么把生命当作恶而逃避？除非他们极不公正，不以较少的恶与较多的福比较。因为当他们一生享受各种精选的满足时，如果任何苦楚偶然接踵而来，他们就想死；而且他们如此看待，以致似乎从未顺利过，只要他们偶尔遭遇不顺。因此他们谴责整个生命，认为它除了充满邪恶之外别无他物。由此产生了那个愚蠢的见解：我们认为生命的状态其实是死亡，我们畏惧的死亡其实是生命；以至于首要的善是不出生，其次便是早死。为使这一见解更有分量，它被归给\Person{西勒诺斯}。\Person{西塞罗}在他的\WorkTitle{安慰}中说：\Quote{不出生是最好的事，不要落入这些生命的礁石上。其次的事情是，如果你已经出生，就尽快死去，逃离命运的暴力如同逃离一场大火。} 从这一点可以看出他相信这种最愚蠢的说法：他自己还添加了一些东西加以修饰。因此我问，他认为对谁来说不出生是最好的？当根本没有任何有感知的人时；因为正是感知使事物成为善或恶。其次，他为何把整个生命看作不过是礁石和大火？仿佛不出生是我们能力所及，或者生命是由命运赐予我们，而不是由天主，或者生命的历程似乎与大火有任何相似之处？\par

\Person{柏拉图}的说法不无相似：他感谢\OriginalTerm{自然}{nature}，首先因为他生而为人而非哑畜；其次因为他为男而非女；因为他为希腊人而非蛮族人；最后因为他为\Place{雅典}人，且生于\Person{苏格拉底}的时代。无法言说真理的无知会产生多么大的盲目和错误，有人会完全辩称，在人类事务中从未有过比这更愚蠢的话。仿佛如果他生为蛮族人、或女人、或甚至驴子，他仍是同一个\Person{柏拉图}，而不是那被产生的存在本身。但他显然相信了\Person{毕达哥拉斯}，后者为阻止人食用动物，说灵魂从人的身体转入其他动物的身体；这既愚蠢又不可能。愚蠢，因为当同一\OriginalTerm{造物主}{Artificer}既已一次性造了最初的灵魂，总能造新的，就没有必要将长期存在的灵魂引入新身体；不可能，因为具有\OriginalTerm{正当理性}{right reason}的灵魂不能改变其状态的本性，如同火不能向下冲，或像河流那样斜倾其火焰。因此这位智者想象，可能发生这样的事：当时在\Person{柏拉图}内的灵魂会被关在某个其他动物内，并被赋予人的感受力，从而理解并悲伤于它被不相称的身体所累。如果他曾说，他感谢是因为他生来具有好的禀赋且能接受教导，并且拥有那些使他能接受自由教育的资源，那他会多么理性得多啊！生为\Place{雅典}人有什么好处？难道不是有许多才华和学问卓越的人住在其他城市，他们个人比所有\Place{雅典}人都强？我们必须相信有成千上万的人，尽管生在\Place{雅典}和\Person{苏格拉底}的时代，却仍是无知和愚蠢的。因为并非城墙或出生之地能赋予人智慧。为自己生于\Person{苏格拉底}的时代而庆幸有何用？\Person{苏格拉底}能给学习者提供才能吗？\Person{柏拉图}没有想到，\Person{阿尔西比亚德}和\Person{克里底亚}也是同一位\Person{苏格拉底}的常听者，其中一人是他国家最活跃的敌人，另一人是所有暴君中最残忍的。\par

% structure: p0090 source=h2 target=subsection reason=relative-heading-rank; base=h2

\subsection{第二十章 \Person{苏格拉底}在哲学上比其他人更有知识，尽管在许多事上他行事愚蠢。}

现在让我们来看，\Person{苏格拉底}自己身上有什么如此伟大，以致一位智人理应感谢自己生于他的时代。我不否认他比那些以为事物本性可由心灵领悟的人略为精明。而我认为，他们不仅愚蠢，而且不敬；因为他们想将好奇的眼睛投向那\OriginalTerm{上智安排}{heavenly providence}的秘密。我们知道在\Place{罗马}和许多城市，有些神圣事物被认为人观看是不敬的。因此不被允许亵渎那些物件的人避免观看；若因失误或意外有人看见了，他的罪过先由惩罚、然后由重复祭献来赎除。对那些想窥探不允许之事的人，你能做什么？实在，那些以不敬的争辩试图亵渎世界和这天上的圣殿的奥秘的人，比进入\OriginalTerm{灶神}{Vesta}、\OriginalTerm{善神}{the Good Goddess}或\OriginalTerm{谷神}{Ceres}神庙的人更邪恶得多。而这些圣所虽不容男人接近，却仍由人所建。但这些人不仅逃脱不敬的指控，而且更不体面的是，他们获得了口才的声名和才华的荣耀。即便他们能探究到什么呢？因为他们既愚蠢于断言，又邪恶于查究；既然他们既不能发现任何事，即使发现了也不能捍卫它。因为若偶然他们看见了真理——这是常发生的事——他们行事如此，以致它被他人当作虚假而驳斥。因为没有人从天上降下来对个人的意见作出判决；因此无人能怀疑，那些追求这些事的人是愚蠢、无感觉且疯狂的。\par

因此\Person{苏格拉底}有某种人的智慧，当他明白这些事不可能被确定时，便远离这类问题；但我担心他只在此事上如此行。因为他许多行为不仅不值得赞扬，而且最应受责备，在这些事上他最像他的同类。从其中我选一件让所有人判断。\Person{苏格拉底}用了这句众所周知的谚语：\Quote{我们之上的事与我们无关。}因此让我们趴在地上，把那赋予我们用于产生优秀作品的手当作脚来用。天与我们无关，我们本被提升去默观它；总之，光本身与我们无干；无疑我们赖以生存的原因来自天上。但如果他认识到我们不应讨论天界事物的本性，他也无法理解他脚下那些事物的本性。那么如何？他言语有误吗？不太可能；但他无疑有意表达他所说的，即我们不应投身于宗教；但如果他公开这样说，无人会容忍。\par

因为谁能看不见，这个以如此奇妙的方式完成的世界，是由某个天主眷顾所管理的？因为没有任何东西可以在没有指导者的情况下存在。同样，一所被居民遗弃的房子就会衰败；一艘没有舵手的船就会沉没；一个被灵魂抛弃的身体就会腐朽。我们更不能认为，如此巨大的建筑，要么没有工匠就能建造，要么没有统治者就能存在这么久。但如果他想要推翻那些公开的迷信，我并不反对；相反，如果他找到了什么更好的东西来代替它们，我反而会称赞他。但同一个人却以狗和鹅起誓。哦，小丑（正如伊壁鸠鲁派的\Person{齐诺}所说），一个没有头脑、堕落、绝望的人，如果他想要嘲笑宗教的话；一个疯子，如果他认真这样做，以至于把最卑贱的动物当作天主！因为谁敢指责埃及人的迷信，当\Person{苏格拉底}在\Place{雅典}以自己的权威证实了它们时？但是，他在死前请朋友们为他献上一只他向\Person{阿斯克勒庇俄斯}许愿的公鸡，这难道不是极度虚浮的标志吗？他显然害怕，因为他许的愿，会被\Person{阿斯克勒庇俄斯}在判官\Person{剌达曼堤斯}面前告发。如果他是因病而死，我会认为他极其疯狂。但既然他在神志健全的情况下这样做，那么认为他是智者的人，他自己就是神志不健全的。看啊，智者为自己出生在这样的时代而庆幸的人！\par

% structure: p0094 source=h2 target=subsection reason=relative-heading-rank; base=h2

\subsection{第二十一章 论\Person{柏拉图}的体系，这会导致国家的倾覆。}

然而，让我们看看他从\Person{苏格拉底}那里学到了什么，苏格拉底完全摒弃了自然哲学，转而致力于美德与责任的研究。因此，我不怀疑他教导听众关于正义的规诫。所以，在\Person{苏格拉底}的教导下，\Person{柏拉图}并没有忽视这一点：正义的力量在于平等，因为所有人都是在平等的条件下诞生的。因此（他说），他们必须没有私有或自己的东西；但为了平等，正如正义的方法所要求的，他们必须拥有一切公有。这还是可以忍受的，只要看起来是指金钱而言。但这有多么不可能和不义，我可以用许多东西来证明。不过，让我们承认其可能性。因为假设所有都是智者，并且轻视金钱。那么，这种公有将他引向何处？他说，婚姻也应该是公有的；以至于许多男人可以像狗一样涌向同一个女人，而那个力量更强的男人可以成功地得到她；或者如果他们有哲学家的耐心，他们可以像在妓院里一样排队等候。哦，\Person{柏拉图}奇妙的平等！那么，贞洁的美德在哪里？夫妻的忠诚在哪里？如果你夺走了这些，一切正义都被夺走了。但他还说，如果一个国家由哲学家作王，或者国王是哲学家，这个国家就会繁荣。但是，如果你把这个主权交给这样一个正义和公平的人，他剥夺了一些人的财产，又把别人的财产给了另一些人，他还会践踏妇女的羞耻心；这不仅是国王从未做过的事，甚至暴君也未曾做过。\par

但他提出这个最可耻的建议的动机是什么？如果所有人都成为所有人的丈夫、父亲、妻子和子女，国家就会和谐，并以相互之爱的纽带联系在一起。这是怎样的人类的混乱？当没有确定的东西可以爱时，感情怎么可能被保持？什么男人会爱一个女人，或者什么女人会爱一个男人，除非他们一直生活在一起——除非心灵的奉献和彼此保持的忠诚使他们的爱不可分割？但这一美德在那种滥交的快乐中没有位置。再者，如果所有人都是所有人的子女，谁能像爱自己的子女一样爱孩子，当他不确定或者怀疑是不是自己的时候？谁会尊崇某人为父亲，当他不认识生他的人时？结果，他不仅把陌生人当作父亲，也把父亲当作陌生人。为什么我要说可能有一个共同的妻子，但不可能有一个共同的儿子？儿子只能由一个人所生。因此，自然本身大声反对，公有对他是失落的。剩下的是，他只为和睦而主张妻子公有。但没有什么比许多男人对一个女人的渴望引起更强烈的纷争了。对此，\Person{柏拉图}本可以被提醒，如果不是因为理性，那么也肯定通过例子：既通过哑巴动物，它们为此最激烈地争斗，也通过人，他们总是为了此事互相发动最残酷的战争。\par

% structure: p0097 source=h2 target=subsection reason=relative-heading-rank; base=h2

\subsection{第二十二章 论\Person{柏拉图}的规诫及对其的谴责。}

由此可见，我们所谈的公有制度除了通奸和淫欲之外，别无他物，而德行尤其为彻底消灭这些所必需。\OriginalTerm{和谐}{concord}，他并未找到他所寻求的和谐，因为他没有看到和谐从何而来。\OriginalTerm{正义}{justice}并不在于外在事物，甚至不在于身体，而是完全作用于人的心灵。因此，凡希望使人平等的人，不应剥夺婚姻和财富，而应剥夺傲慢、骄傲和高傲，使那些有权势、高高在上的人知道，他们与最贫困的人也是平等的。因为，若是从富人身上除去侮慢和不义，那么有人富有人贫便无关紧要，因为他们将在精神上平等，而只有对天主的敬畏才能产生这种结果。\Person{柏拉图}认为自己找到了正义，实际上他完全消除了正义，因为公有不应是易逝之物的，而应是心灵的。因为若正义是诸德之母，那么当这些德行逐一被剥夺时，正义本身也就被推翻了。但柏拉图首先剥夺了\OriginalTerm{节约}{frugality}，当人没有自己的财产可以占有时，节约便不存在；他剥夺了\OriginalTerm{节制}{abstinence}，因为不再有属于他人的东西可以让人节制；他剥夺了两性中最伟大的德行——\OriginalTerm{节德}{temperance}和\OriginalTerm{贞洁}{chastity}；他剥夺了\OriginalTerm{自重}{self-respect}、\OriginalTerm{羞耻}{shame}和\OriginalTerm{谦逊}{modesty}，倘若那些通常被视为卑劣可耻之事开始被视为光荣和合法。因此，当他希望赋予所有人德行时，他却从所有人那里夺走了德行。因为财产的持有既包含恶习的材料，也包含德行的材料；而财产的公有除了恶习的放纵外，别无他物。因为拥有众多情妇的男人只能被称为奢侈和挥霍；同样，被众多男人占有的女人必然不是奸妇——因她们没有固定的婚姻——而是娼妓和妓女。因此，他使人类生活——我不说——与哑巴动物相似，而是与牲畜和野兽相似。因为几乎所有的鸟类都结成婚姻，成对结合，以和谐的心灵守护它们的巢穴——如同婚床——并哺育自己的幼雏，因为它们熟悉自己的后代；你若将别的幼雏放入，它们会将其赶走。但这贤人违背人类习俗，违背自然，选择了更愚蠢的模仿对象；他看到在其他动物中，雄性和雌性的职分并无分别，便认为女人也应当从军，参与公共议政，担任官职，执掌统帅之位。因此他给女人配了马匹和武器；由此推论，他应为男人配羊毛、织布机和怀抱婴儿。他也没有看到自己所言的不可能性，因为世上从未存在过如此愚蠢或虚妄的民族，竟会这样生活。\par

% structure: p0099 source=h2 target=subsection reason=relative-heading-rank; base=h2

\subsection{第二十三章　论某些哲学家之谬误，兼论太阳与月亮}

既然哲学家中的头面人物本身已被发现如此空洞，那么对于那些次要人物，我们又当如何看呢？他们往往在夸耀自己蔑视金钱时才显得如此智慧。勇敢的精神！但我等着看他们的行为，以及那蔑视的结果。他们视祖先传下的财产为恶而避开并放弃；惟恐在风暴中遇难，便在平静中自行坠入深渊，其坚决并非出于德行，而是出于乖谬的恐惧——如同那些因惧怕被敌人杀死而自杀的人，以求一死免死。这样，他们既无荣誉也无影响，丢弃了本可藉以获得慷慨之誉的资财。\Person{德谟克利特}因放弃田地、任其成为公共牧场而受称赞。若是他将田地赠与他人，我当赞成。但若所作之事无益且人人行之则为恶，便不算明智。但这疏忽尚可容忍。对于那位将财产换成金钱并投入海里之人，我该说什么？我怀疑他是否清醒，还是癫狂。他说：\Quote{去吧，你们这些邪恶的欲望，沉入深海！我要将你们抛弃，免得我被你们抛弃。}你若如此蔑视金钱，不如用于仁慈博爱之事：施舍给穷人。你将要抛弃的，本可成为许多人的援助，使其免于饥渴赤身。至少效法\Person{图迪塔努斯}的疯癫和暴怒：散尽你的财物，让民众抢夺。你既可摆脱金钱，又能善用其利；因为凡有益于众人的，便是妥善的置用。\par

但谁赞同\Person{芝诺}提出的\OriginalTerm{过失平等}{equality of faults}？不过我们还是略过那向来遭众人嘲笑之事。足以证明这位狂人之谬误的是：他将怜悯归入恶习和疾病之列。他剥夺了我们的一种情感，而此情感几乎涉及整个人生。因为人的本性较其他动物更为软弱，天主眷顾以天然保护方式武装了其他动物，使之能忍耐严寒酷暑或抵御来自身体的侵袭；既然人未获得任何此类禀赋，便代之以怜悯之情——这真可称为\OriginalTerm{人道}{humanity}——藉此我们得以彼此保护。因为若一人见他人而变得凶残——正如我们在某些独居性动物身上所见——则人间便无社会，无筑城之关切与制度；如此，生命甚至不得安全，因为人的软弱既易受其他动物攻击，人又会如野兽般互相残杀。他在其他事情上的疯狂也不遑多让。\par

论到那声称雪是黑色的人，还有什么可说呢？他自然也会宣称沥青是白色的！\Person{色诺芬尼}就说，他生来就是为了观看天空和太阳，然而当太阳照耀时，他却在地上什么也看不见。\Person{色诺芬尼}最愚蠢地相信了数学家的说法，即月球的球体比地球大十八倍；与此愚行相称，他说在月球的凹面上有另一个地球，那里有另一族人类，其生活方式与我们在这地球上的相似。因此，这些月球居民有另一个月球，为他们夜间照明，如同这个月球为我们照明一样。或许我们这个球体对于下方另一个地球而言，也是一个月球。\Person{塞内卡}说，斯多葛派中曾有一个人，他思忖是否也该为太阳指派其居民；他对此犹豫不决，实属愚昧。因为若他真指派了居民，又会造成什么损害呢？但我认为，是酷热使他却步，以免如此众多的人口因酷热而灭亡，这样大的灾祸就会被说成是因他的过失而发生了。\par

% structure: p0103 source=h2 target=subsection reason=relative-heading-rank; base=h2

\subsection{第24章 论对跖点、天与星辰}

那些想象在我们脚下存在对跖点的人，情况又如何呢？他们说的有什么道理吗？或者有谁会如此愚蠢，竟相信有人脚步高于头顶？或者在我们这里平放的东西，在他们那里却倒悬着？庄稼和树木向下生长？雨水、雪和冰雹向地上方降落？当哲学家们制造出悬挂的田野、海洋、城市和山脉时，还有人惊异于空中花园被列为世界七大奇迹之一吗？我们也必须阐明这一错误的根源。因为他们总是以同样的方式受骗。当他们一开始研究时假设了某些错误的东西，受真理相似性的引导，便必然陷入那些随之而来的事物。于是他们陷入许多可笑的事情；因为与错误的东西一致的事物，其本身也必然是错误的。但由于他们信任最初的假设，便不考虑后续事物的性质，而以各种方式为之辩护；然而他们本应从后续事物来判断最初假设是真是假。\par

那么，是怎样的论证思路将他们引向对跖点的概念呢？他们看见星辰的轨迹向西运行；他们看见太阳和月亮总是在同一方向落下，从同一方向升起。但由于他们不明白是什么机制调节其运行，也不明白它们如何从西方返回东方，于是便假定天穹本身向各个方向倾斜——因其广阔无垠，必然呈现这种外观——他们便认为世界像一个球体那样圆，并想象天穹随天体的运动而旋转；因此，星辰和太阳在落下时，因世界运动的极快速度而被带回东方。于是，他们制造了铜制球体，仿佛按照世界的形状，并在其上刻下某些奇形怪状的图像，称之为星座。根据天穹的球形，必然得出地球被包在其曲面的中央。但若如此，地球本身也必然像一个球体；因为被球形所包围的东西不可能是别的形状，只能是圆的。但若地球也是球形的，就必然会出现这样的情形：它对天穹的各个部分呈现相同的外观；也就是说，它应高山耸立，平原伸展，海洋平坦。若如此，则最后一个结果也必然出现：地球上没有部分无人居住或其他动物。因此，地球的球形还导致了那些悬空的对跖点的发明。\par

但若你询问那些为这些奇异虚构辩护的人，为什么所有事物不向天穹较低的那部分坠落，他们回答说：事物的本性是，重的物体被引向中央，并且都向中央聚合，如同我们在车轮的辐条上所见的；而轻的物体，如雾、烟和火，则被从中央驱离，以寻求天穹。对于那些一旦犯错便始终执着于其愚行，并以此一虚妄之事为彼一虚妄之事辩护的人，我不知道该说什么；只是有时我想，他们或是为了开玩笑而讨论哲学，或是故意且有意识地承担起为虚假之事辩护的任务，仿佛要在虚假的论题上锻炼或展示其才华。但我本可以用许多论据证明天穹不可能低于地球，若不是此书必须在此结束，而且仍有一些对当前著作更为必要的事项有待处理。既然逐一检视每个人的错误非一卷书所能完成，那么列举少许便足矣，从中可理解其余错误的性质。\par

% structure: p0107 source=h2 target=subsection reason=relative-heading-rank; base=h2

\subsection{第25章 论研习哲学及其所需之重要资质}

现在我们必须就哲学说几句概括性的话，以加强我们的论点，然后作结。在我们作者中，柏拉图最伟大的模仿者认为哲学不属于大众，因为只有有学问的人才能达到它。西塞罗说：\Quote{哲学满足于少数评判者，它自愿刻意避开群众。}因此，如果它避开人群，它就不是智慧；因为，如果智慧赐予人，它就是毫无区别地赐予所有人，以至于根本没有不能获得它的人。但那些如此拥抱德行（德行乃是赐予人类的）的人，却唯独他们自己似乎想要享用那公共的善物；他们嫉妒得好像他们想要捆住或挖出别人的眼睛，使他们看不见太阳。因为，拒绝将智慧赐予人，难道不是剥夺他们心灵中那真实而神圣的光明吗？但如果人的本性能够接受智慧，那么工人、农夫、妇女，总之所有具有人形的人都应当被教导成为智慧者；并且，应当从每一种语言、身份、性别和年龄中召集民众。因此，一个非常有力的论据是：哲学既不是导向智慧，其本身也不是智慧，因为它的奥秘只通过哲学家的胡须和斗篷才得以揭示。此外，斯多葛派察觉到了这一点，他们主张奴隶和妇女也应学习哲学；伊壁鸠鲁也是，他邀请完全不识字的人学习哲学；还有柏拉图，他想建立一个由智慧者组成的城邦。\par

他们确实试图做真理所要求的事，但他们无法超越言语。首先，因为许多学问的训练对于致力于哲学是必要的。为了阅读的练习，必须获得普通学识，因为在如此多样的主题中，不可能所有事情都通过听来学习，或保留在记忆中。也必须给予文法学家不少关注，以便知道正确的说话方法。这需要占据许多年。修辞学也不可忽视，以便能够说出和表达所学到的东西。几何学、音乐和天文学也是必要的，因为这些学问与哲学有某种联系；所有这些科目，妇女无法学习，她们必须在成熟之年学习日后为家庭所用的事务；仆人也不能，他们在尤其能够学习的那些年岁里必须生活在服侍中；穷人、劳动者或农夫也不能，他们必须靠劳动维持每日的生计。正因如此，图利乌斯说哲学厌恶群众。但是伊壁鸠鲁却愿意接受无知者。那么，他们将如何理解那些关于万物本原的论述呢？这些论述的复杂和纠结，即使是受过良好教育的人也几乎难以达到。\par

因此，在那些晦涩难懂、因多种智力而混乱、并由雄辩者精心语言修饰的论题中，无知和未受教育者有何位置？最后，在整个有记载的人类记忆中，除了特弥斯忒一人外，他们从未教导任何妇女学习哲学；除了斐多一人外，他们从未教导奴隶学习哲学。据说斐多在遭受压迫的奴役中生活时，被刻柏斯赎身并教导。他们还列举了柏拉图和狄奥根尼：然而，这些人并非奴隶，尽管他们曾沦为奴役，因为他们是战俘。据说一个叫阿尼刻里斯的人用八千塞斯特斯赎回了柏拉图。对此，塞内加严厉斥责了赎身者本人，因为他把柏拉图估价得如此之低。他（塞内加）在我看来是个疯子，竟然对一个没有扔掉大笔钱财的人生气；毫无疑问，他本该像赎回赫克托耳的尸体那样称量黄金，或者坚持支付比卖主所要求的更多的钱。此外，他们没有教导任何蛮族人，只有斯基提亚的阿那卡西斯一人例外，而他若不曾先学会语言和文学，也绝不会梦想哲学。\par

% structure: p0111 source=h2 target=subsection reason=relative-heading-rank; base=h2

\subsection{第二十六章　唯有天主的教导赐予智慧，天主的法律有何效力}

因此，他们察觉到本性要求公正地要求的事，但他们自己却无力实现，且看到哲学家们也无法成就的事，唯有藉着天主的教导才得以完成；因为那才是智慧。毫无疑问，他们能够说服任何连自己都未曾说服的人吗？或者，当他们自己屈服于恶习，并承认被本性压倒时，他们能抑制任何人的欲望、克制愤怒、约束情欲吗？然而，天主的诫命因其纯朴和真理而对人的灵魂有何种影响，日常的证据已经显明。给我一个暴躁、粗鲁、放纵的人：用天主的极少数言语，\par

\Quote{我必使他驯良如绵羊。}\par

请给我一个贪婪、悭吝、固执的人，我立刻就能将他变得慷慨大方，双手满满地施舍钱财。请给我一个畏惧痛苦和死亡的人，他立刻就会蔑视十字架、火刑和法拉里斯的铜牛。请给我一个好色、通奸、贪吃的人，你立刻就会看到他变得清醒、贞洁、节制。请给我一个残忍嗜血的人：那暴怒立刻会转变为真正的宽仁。请给我一个不义、愚蠢、作恶的人：他立刻就会变得正义、智慧、无辜；因为藉着一次洗涤，他的一切邪恶都被除去。天主的智慧何其伟大，一旦注入人的胸怀，只须一次冲动，便一劳永逸地驱散愚昧——那过错的根源；为此，无需代价、书籍或夜间苦读。这些成果是白白获得的，轻松而迅速，只要耳朵敞开，心灵渴慕智慧。无人应当害怕：我们不卖水，也不以报酬出售太阳。天主那最丰沛、最充盈的泉源向众人敞开；这天上之光为一切有眼的人升起。有哪位哲学家曾成就这些事，或者，如果他愿意，能成就吗？尽管他们终生研究哲学，却既不能改善他人，也不能改善自己（如果本性设置了任何障碍）。因此，他们的智慧竭尽全力，并非根除恶习，而是隐藏恶习。但天主的几条训诫如此彻底地改变整个人，脱去旧人，使他成为新人，以致你认不出他是同一个人。\par

% structure: p0115 source=h2 target=subsection reason=relative-heading-rank; base=h2

\subsection{第二十七章 哲学家的训诫对真智慧贡献何其微小——这智慧只有在宗教中才能找到。}

那么，怎么？难道他们没有吩咐类似的事吗？是的，确实有许多；而且他们时常接近真理。但这些训诫没有分量，因为它们是属人的，缺乏那更大的、即天主的权威。因此无人相信它们，因为听者认为发令者同自己一样是人。此外，他们毫无确凿之处，没有出于知识的事。既然一切都是凭推测，并且提出许多互相歧异、各不相同的主张，最愚蠢的人才愿意服从他们的训诫，因为存疑它们是否真确；因此无人服从，因为无人愿为不确定之事劳苦。斯多葛派说，唯有德行能产生幸福生活。再没有比这更真实的说法了。但假如他受折磨，或因痛苦而受苦？在刽子手手中，人还能幸福吗？但真正说来，身体所受的痛苦是德行的材料；因此即使受酷刑，他也不悲惨。伊壁鸠鲁说得更强硬。他说，智者总是幸福的；即使被关在法拉里斯的铜牛里，他也会说出这话：\Quote{这是愉快的，我毫不在乎。}谁不嘲笑他？尤其因为一个耽于享乐的人承担了刚毅之士的角色，而且到了过分的程度；因为不可能有人将身体的酷刑当作快乐，因为尽德行之职只需忍受和承受它们就够了。斯多葛派啊，你们怎么说？伊壁鸠鲁啊，你怎么说？智者甚至在受酷刑时也是幸福的。如果是因忍受的荣耀，他享受不到，因为或许他在酷刑中死去。如果是因对事迹的回忆，要么灵魂消灭，他感觉不到；要么他感觉到，也一无所获。\par

那么，德行还有什么别的益处？人生的幸福是什么？难道是为了人能平静地死去？你向我呈现的不过是某一个时辰、或许片刻的益处，为了它，不值得在整个人生中被苦难和辛劳耗尽。但死亡占据多少时间？它到来时，你曾否平静地承受它，已无所谓。于是，从德行中只求得了荣耀。但这荣耀要么是多余的、短暂的，要么因人有偏见的判断而不会随之而来。因此，当德行受制于死亡和腐朽时，德行便没有果实。所以说这些话的人看见了德行的某种影子，却没有看见德行本身。因为他们将眼睛盯在地上，没有举目仰望，好能看见她——\par

\Quote{她从天的四方显现。}\par

这就是为什么无人遵守他们的训诫；因为他们要么维护享乐而把人引向恶习，要么维护德行，却仅以耻辱作为对罪的惩罚，仅以荣誉和称赞作为对德行的赏报，因为他们说，德行当为自身之故而被追求，而非为任何其他对象。因此，智者即使在酷刑下也是幸福的；但当他因他的信德、因正义或因天主而遭受酷刑时，那痛苦的坚忍将使他最幸福。因为唯独天主能尊崇德行，德行的赏报唯独是不朽。凡不追求此不朽，也不拥有与永生相连的宗教的人，确实不知道德行的力量，他们对德行的赏报一无所知；也不仰望天堂，尽管他们自己想象自己是在仰望，当他们探究那些不容探究的问题时，因为仰望天堂的唯一原因，要么是信奉宗教，要么是相信自己的灵魂是不朽的。因为若有人明白天主当受敬拜，或怀有置于他面前的不朽之望德，他的心就在天堂；虽然他也许不能以肉眼看见，却以灵魂之眼看见。但那些不接纳宗教的人属于地，因为宗教来自天堂；而那些认为灵魂与肉身一同消亡的人，同样向下看地：因为在属于地的肉身之外，他们看不见任何属于不朽之物。因此，人这样被造，以挺直的身体仰望天堂，是没有益处的，除非他以高扬的心智辨识天主，并且他的思绪完全沉浸于永生的望德。\par

% structure: p0120 source=h2 target=subsection reason=relative-heading-rank; base=h2

\subsection{第二十八章 论真宗教与自然。命运是否是一位女神，并论哲学。}

因此，生命中没有别的能使我们的计划和状况倚靠，除了认识创造我们的天主，以及虔诚而恭敬地敬拜祂；既然哲学家们偏离了这一点，显然他们并不智慧。他们确实寻求智慧；但因未以正确的方式寻求，便堕入更远的距离，陷入如此大的错误，以至于他们甚至连普通的智慧都不具备。因为他们不仅不愿持守宗教，甚至还将其废除；同时，在虚假德行的外表引导下，他们竭力使心灵摆脱一切恐惧：而这种对宗教的推翻竟获得了\Quote{自然}之名。因为他们或是不知道世界是由谁造的，或是想要说服人相信没有什么事物是藉上智完成的，便说自然是万有之母，仿佛是说万有都是自发的：藉此一词，他们完全承认了自己的无知。因为自然，离开天主的眷顾和权能，完全就是虚无。但若他们称天主为自然，那么用\Quote{自然}而非\Quote{天主}之名，又是何等悖逆！但若自然是计划、必然性或出生的条件，它本身并不具有感觉；而必然应有一个神圣心智，以其预见为万有提供存在的开端。或者，如果自然是天和地，以及一切受造之物，那么自然不是天主，而是天主的化工。\par

他们因类似的错误相信命运存在，视之为一位以各种偶然事件戏弄世人的女神，因为他们不知道善事恶事从何而来。他们认为自己是被召来与她战斗的；他们也不指明是谁、为了什么而将他们如此匹配；他们只夸耀自己每时每刻都与命运进行生死搏斗。如今，凡因朋友去世和离别而安慰过他人者，都以最严厉的指控斥责命运之名；在他们关于德行的任何讨论中，命运无不受鞭挞。\Person{马尔库斯·图利乌斯}在其\WorkTitle{Consolation}中说，他始终与命运抗争，且每当他英勇击退敌人的攻击时，命运总是被他战胜；甚至当他被逐出家园、失去祖国时，他也没有被命运征服；但当他失去最亲爱的女儿时，他羞耻地承认自己被命运战胜。他说：\Quote{我屈服，并举起我的手。}这个如此俯伏在地的人，还有什么比他更可悲的呢？他说：他行事愚蠢；但此人却自称是智者。那么，这个称号的采用意味着什么呢？那种以华丽言辞声称的对事物的蔑视又意味着什么？那与众不同的服饰呢？或者，既然至今未找到任何智者，你为何还要传授智慧之训？当我们否认哲学家是智者时，还有人怨恨我们吗？他们自己也承认既没有知识也没有智慧。因为若有时他们失败得甚至无法伪装任何事物——正如他们在其他情况下的做法——那么，确实，他们就被提醒了自己的无知；并且，仿佛在疯狂中，他们跳起来喊道，他们是瞎眼的、愚蠢的。\Person{阿那克萨戈拉}宣称万物都被黑暗笼罩。\Person{恩培多克勒}抱怨感官的道路狭窄，好像为了他的思考需要四马战车似的。\Person{德谟克利特}说真理沉在井底，深不见底；的确，正如他说的其他事情一样，是愚蠢的。因为真理并非如同沉在井里，允许他下去甚至跌进去，而是如同放在高山的最高峰，或在天上，这是最真实的。因为有什么理由使他说真理沉在下面而不是高高在上呢？除非他碰巧宁愿把心智也放在脚上，或脚后跟的底部，而不是放在胸中或头部。\par

他们与真理本身相距如此遥远，甚至连他们自己的身体姿势都未能告诫他们，真理必须在最高之处被他们寻求。由此绝望产生了\Person{苏格拉底}的那句承认，他说他一无所知，只知道这一件事，即他一无所知。由此产生了\Place{学园}的体系，如果那种既学习又教导无知的体系可以被称为体系的话。但即使是那些自认为拥有知识的人，也无法始终连贯地捍卫他们认为自己知道的东西。因为他们彼此不一致，由于对神圣事物的无知，他们如此前后矛盾和不确定，常常断言相互矛盾的事情，以至于你无法确定和判断他们的意思。因此，你为什么要与那些自取灭亡的人战斗？你为什么要劳神反驳那些被自己的言语反驳和压迫的人？\Person{西塞罗}说，\Person{亚里士多德}在指责古代哲学家时，宣称他们要么是最愚蠢的，要么是最虚荣的，因为他们认为哲学已因他们的才能而完善；但他看到，由于在短短几年内有了巨大进步，哲学将在短期内完成。那么，那个时间是什么？哲学以何种方式、何时、由谁完成？因为他说他们因认为哲学由他们的才能完善而最为愚蠢，这是真的；但他自己说话也不够谨慎，他认为哲学要么由古人开始，要么由较近的人发展，要么将由后来的人在短期内完成。因为那些不用正确方法寻求的事物永远无法被探究。\par

% structure: p0124 source=h2 target=subsection reason=relative-heading-rank; base=h2

\subsection{第二十九章 再论命运与德行}

但让我们回到我们先前搁置的主题。因此，命运本身什么都不是；我们也不应将它视为具有任何知觉，因为命运是意外事件的突然和出乎意料的发生。但哲学家们，为了有时不犯错，希望在愚蠢的事情上变聪明；他们说命运不是一个女神，如同人们通常相信的，而是一个男神。然而，有时他们将这个神称为自然，有时称为命运，\Quote{因为他带来，}同一\Person{西塞罗}说，\Quote{许多我们意想不到的事情，由于我们缺乏智慧和不知原因。}因此，既然他们不知道事物发生的原因，他们也必定不知道行这事的那一位。同一作者在一部非常严肃的著作中，在其中他给予他的儿子从哲学中吸取的生活准则，说：\Quote{谁能不知道命运的力量在双方都是巨大的？因为当我们遇到她顺风时，我们得到我们渴望的结果，而当她向我们吹逆风时，我们被撞在岩石上。}首先，那个说无物可知的人，说这话时似乎他自己和所有人都有知识。其次，那个力图使即使明显的事情也变得可疑的人，认为这件事是明显的，而对他而言这件事本应特别可疑；因为对一个智者来说，这完全是假的。\Quote{谁不知道呢？他说，我确实不知道。让他教我，如果他能的话，那力量是什么，那微风是什么，以及那逆风是什么。}\par

因此，对于一个有才华的人来说，说出如果无人否认他就无法证明的话，是可耻的。最后，那个说必须保留赞同，因为鲁莽地赞同未知之事是愚人的行为的人，他完全相信了粗俗和未受教育者的意见，他们认为命运给予人好与坏的事物。因为他们将她的形象描绘为拿着丰饶之角和舵，仿佛她既赐予财富又掌管人类事务。\Person{维吉尔}赞同这种意见，他称命运为全能者；还有那位历史学家，他说：\Quote{但确实，命运主宰一切。}那么，其他神的位置在哪里？为何不可以说她独自统治，如果她比其他人更有力量？或者为何她不是独自被崇拜，如果她在一切事物中拥有权力？或者，如果她只施加邪恶，让他们提出一些理由，为什么如果她是一位女神，她会嫉妒人类，并希望他们毁灭，尽管他们虔诚地崇拜她；为什么她对恶人更有利，对善人更不利；为什么她策划、折磨、欺骗、消灭；谁任命她为人类种族永久的折磨者；简言之，为什么她获得了如此恶毒的力量，以至于她按照自己的任性而非按照真理使一切事物显赫或暗淡。我说，哲学家们本应探究这些事情，而不是鲁莽地指责无辜的命运：因为虽然她存在，但他们无法提出任何理由说明为什么她应该像人们认为的那样对人类怀有敌意。因此，所有那些他们谴责命运不公、并反对命运傲慢夸耀自己德行的言论，无非是轻率胡言。\par

因此，那些天主已将真理启示给他们的人，不要嫉妒我们。我们知道命运算不得什么，也同样知道有一个邪恶诡诈的灵，与善为敌，是正义的仇敌，与天主对立；这敌意的原因我们已在第二卷中解释了。他向众人设下陷阱；但对于那些不认识天主的人，他用错误阻挡他们，用愚昧淹没他们，用黑暗笼罩他们，使无人能够达到对天主之名的认识，惟独在这名中包含了智慧和永生。另一方面，对于认识天主的人，他用诡计和奸诈袭击他们，好以欲望和贪恋诱捕他们，当他们被罪恶的谄媚腐蚀时，便驱使他们走向死亡；或者，若他阴谋不成，就试图用暴力和强权将他们打倒。因为天主起初在原始过犯时没有立即将他推入惩罚，正是为了使他以恶意锻炼人的德行；因为德行若不经常搅动，若不通过不断的考验而坚强，就不能成为完善的，因为德行就是在忍受邪恶时无畏而不屈的忍耐。由此得出，若没有仇敌，就没有德行。因此，当这些人察觉到这乖戾势力与德行为敌的力量，却不知其名时，便为自己发明了命运这个毫无意义的名称；\Person{尤维纳利斯}在以下诗句中说明了这远离智慧的程度：——\par

\Quote{若有审慎，神能无不临在；但我们却使你成为女神，哦命运，并将你置于天上。}\par

因此，是愚昧、错误、盲目，以及如\Person{西塞罗}所说的，对事实和原因的无知，引入了自然和命运这些名称。但正如他们不认得自己的仇敌，同样他们也不真正认识德行——对德行的认识是从对仇敌的观念而来。若德行与智慧相连，或如他们所说，德行本身就是智慧，那么他们必定不知道德行存在于哪些对象中。因为一个人若不知道必须用哪件武器武装自己，就不可能拥有真正的武器；若他在战斗中攻击的不是真正的敌人，而是影子，就不能战胜仇敌。因为他将会被击倒，当他将注意力投向另一对象时，未能预先预见或防范对准他要害的攻击。\par

% structure: p0130 source=h2 target=subsection reason=relative-heading-rank; base=h2

\subsection{第三十章 前面所论之事的结论；以及我们必须以何种方式从哲学家的虚妄过渡到真智慧和对真天主的认识，在这认识中唯独有德行和幸福。}

我已按照我那微薄才能所允许的，教导了哲学家们远离真理的道路。然而，我看出我遗漏了许多事情，因为我本不该与哲学家进行辩论。但我有必要借此题外话表明，如此众多伟大的才智都白白耗费在虚伪的主题上，以免有人被腐化的迷信所隔绝，想要投奔他们，以为会找到某种确定性。因此，人所唯一盼望、唯一安全，就存于我们所护卫的这教义中。人的全部智慧唯独在于认识并敬拜天主：这是我们的主张，我们的见解。因此，我要用我全部声音的力量作证、宣告、声明：这里，这里就是所有哲学家毕生所求的；然而，他们既不能探究、领悟，也不能达到它，因为他们要么固守腐化的宗教，要么完全弃绝了宗教。所以，让那些不教导人生活、反而扰乱人群的人统统离开吧。因为他们教导什么呢？或他们指导谁呢？他们连自己都未指导过。病人能医治谁？瞎子能领谁？所以，我们一切关注智慧的人，都当投身于这事业。难道我们要等到\Person{苏格拉底}有所知？或等到\Person{阿那克萨哥拉}在黑暗中找到光明？或等到\Person{德谟克利特}从井中汲出真理？或等到\Person{恩培多克勒}拓展他灵魂的道路？或等到\Person{阿塞西劳斯}和\Person{卡尔内亚德}看见、感觉、知觉吗？\par

看，天上有声音教导真理，并为我们展示比太阳本身更亮的光明。我们为何对自己不公，迟迟不愿拿起智慧呢？有学识的人虽穷尽一生追求，却从未能发现它。愿凡渴望智慧与幸福的人，听从天主的声音，学习正义，领悟自己出生的奥秘，蔑视人间之事，拥抱神圣之事，好获得那为他所生的至善。推翻了所有虚假宗教，并驳斥了为它们辩护而惯常或可能提出的一切论据之后，又证明了哲学体系的虚妄，我们现在必须转向真宗教和真智慧，因为正如我要教导的，二者是相连的；好使我们能以论证、例子或可靠的见证来维护它，并表明那些异教崇拜者不断以愚昧指责我们的事，并不在我们这里，而全在他们那里。尽管在前几卷中，当我与虚假宗教争辩时，以及在这一卷中，当我推翻虚假智慧时，我已指出了真理何在，但下一卷将更清楚地指明什么是真宗教、什么是真智慧。\par

\SourceRecord{Translated by William Fletcher.；Ante-Nicene Fathers；Vol. 7.；Edited by Alexander Roberts, James Donaldson, and A. Cleveland Coxe.；Buffalo, NY: Christian Literature Publishing Co.,；1886.；<http://www.newadvent.org/fathers/07013.htm>.}

% structure: p0001 source=h1 target=section reason=page-title

\section{\WorkTitle{天主的训示，卷四（论真智慧与真宗教）}}

% structure: p0002 source=h2 target=subsection reason=relative-heading-rank; base=h2

\subsection{第一章　论人类先前的宗教，以及谬误如何蔓延至每个时代，并论希腊七贤}

君士坦丁皇帝啊，当我反思并时常在心中思量人类最初的状态时，我常常感到既惊奇又不配：因着一个时代的愚妄，人们接纳了各种迷信，相信存在众多神明，他们突然陷入如此的自昧，以至真理从眼前被夺走，对真天主的宗教不被遵守，人性的状况也不被认识，因为人不在天上，而是在地上寻求至善。正因如此，古代的幸福无疑被改变了。因为，离弃了万物的父母与创造者天主，人们开始崇拜自己双手所造的无灵之物。这种败坏产生了什么后果，或引入了何种邪恶，题目本身已足以说明。因为，人们背离了至善——这至善是有福的、永生的，正因它不能被看见、触摸或理解——以及与此善相合的、同样不朽的美德，堕落到那些腐朽脆弱的神祇那里，并投身于那些只装饰、滋养和取悦肉体的东西，从而为自己连同他们的神祇和肉身的财物一同寻求永死，因为所有肉身都隶属于死亡。因此，这种迷信必然带来不义与不敬。因为人不再举目向天；反而心思向下，紧贴地上的财物，犹如紧贴土生的迷信。随之而来的是人类的不和、欺诈及一切邪恶；因为，轻视了永恒不朽的、人唯一当追求的财物，他们反而选择了短暂易逝的东西，并且人更多地信任邪恶，因为他们宁可选择恶习而非美德，因为恶习显得更近在眼前。\par

因此，先前时代曾沐浴于最明亮之光的人类生活，如今被阴霾和黑暗笼罩；与此败坏相应，当智慧被夺走后，人们终于开始为自己妄称智慧之名。因为在众人皆有智慧之时，无人被称为此名。但愿这曾为全体人类共有的名称，虽缩减为少数人，仍保持其效力！因为那少数人或许能凭才干、权威或不断的劝诫，使人民摆脱恶习与谬误。然而智慧竟完全消亡，以至从这名称的狂妄本身便可看出，那些被称为智慧者中，无一人真正有智慧。然而，在所谓的哲学发现之前，据说曾有七人，因敢于探究并讨论自然事物，而被认为配得智慧和智慧人之名。\par

啊，可怜而可悲的时代！在全世界中，仅有七人被称为人——因为除非有智慧，无人能正当地被称为人！但若所有其他人都愚昧，连他们自己也不是有智慧的，因为无人能在愚昧人的判断中真正有智慧。他们离智慧如此遥远，以至于后来学问增长、众多伟大的心智始终专注于此事时，真理仍不能被察觉和确定。因为在七贤的声誉之后，整个希腊被探究真理的巨大热忱所点燃，真是难以置信。起初，他们认为智慧之名本身过于傲慢，故不自称为智慧人，而称为\OriginalTerm{爱智者}{desirous of wisdom}。由此，他们既谴责了那些冒然妄称智慧人之名者的错误与愚昧，也谴责了他们自己的无知——他们确实没有否认这一点。因为每当事物的本性仿佛将手按在他们的心智上，使他们无法给出任何解释时，他们惯于坦承自己一无所知、毫无分辨。因此，那些在某种程度上看见自己的人，远比那些相信自己有智慧的人更为智慧。\par

% structure: p0006 source=h2 target=subsection reason=relative-heading-rank; base=h2

\subsection{第二章　智慧应在何处寻找；为何毕达哥拉斯与柏拉图未接近犹太人}

因此，如果那些被称为智慧的人并不智慧，后来的那些也不犹豫承认自己缺乏智慧，那么除非智慧是在别处寻找，因为它并未在被寻找的地方找到，否则还剩下什么呢？但我们能猜想，为何它虽被如此多的心智、在如此漫长的岁月中以最大的热忱和辛劳寻求，却仍未被找到？除非是因为哲人们在自己的界限之外寻找。既然他们遍历并探索了所有地方，却未在任何地方找到智慧，而智慧又必然存在于某处，那么显然，它尤其应当在那看似愚昧的地方寻找——天主在其遮掩下隐藏了智慧与真理的宝藏，以免祂神圣工程的奥秘暴露于世。因此，我常感到惊奇：当\Person{毕达哥拉斯}，以及在他之后的\Person{柏拉图}，因热爱探究真理而深入到埃及人、\OriginalTerm{玛吉}{Magi}和波斯人那里，以了解他们的宗教礼制时（因为他们怀疑智慧与宗教有关），他们却没有接近犹太人——那时唯有犹太人拥有智慧，而且他们本可以更容易地去到那里。但我认为，是天主的眷顾使他们远离了犹太人，以免他们得知真理，因为那时还不允许真天主的宗教和正义为外邦人所知。因为天主决定，在末期临近时，将从天上派遣一位伟大的领袖，向外邦民族揭示那从背信弃义、忘恩负义的百姓那里夺走的事。我将在本书中努力讨论这主题，如果我先能证明智慧与宗教如此紧密相连，以致两者不可分离的话。\par

% structure: p0008 source=h2 target=subsection reason=relative-heading-rank; base=h2

\subsection{第三章 智慧与宗教不可分离：自然之主必然是每个人的父}

正如我在前书中所教导的，敬拜诸神并不蕴含智慧；这不仅因为它使人——一个神圣的受造物——屈从于世俗和脆弱的事物，更因为它其中并无任何能有助于陶冶品性、塑造生活的事物；也不包含对真理的任何探究，而只有敬拜的仪式，这仪式不在于心灵的侍奉，而在于身体的行为。因此，那不能以正义和德行的任何规诫来教导和改善人的宗教，不应被视为真宗教。同样，哲学，因其不拥有真宗教——即最高的虔敬——也不是真智慧。因为如果主宰此世界的天主以难以置信的恩惠扶助人类，并以父的慈爱眷顾人类，祂确实希望得到感恩和尊荣，那么人若对天恩忘恩负义，就不能保持他的虔敬；这当然不是智慧人的行为。因此，正如我所说，哲学和诸神的宗教体系是分离的，并且彼此相距甚远；既然有些人是智慧的教师，通过他们显然没有通往诸神的途径，而另一些人是宗教的司铎，通过他们并不学习智慧；那么显而易见，前者不是真智慧，后者也不是真宗教。因此，哲学未能领会真理，诸神的宗教体系也无法说明自身，因为它缺乏智慧。但哪里智慧与宗教以不可分割的联系结合在一起，哪里两者必然都是真实的；因为在我们的敬拜中，我们应当有智慧，即知道敬拜的合宜对象和方式，并在我们的智慧中进行敬拜，即通过行动和实行来完成我们的知识。\par

那么，智慧在哪里与宗教结合呢？确实，就在那唯一的天主被敬拜的地方，在那里生命和每一个行动都归于一个源头和一个至高的权威：简言之，智慧的教师也是天主的司铎。然而，不要让人因以下情况而困扰：即常常发生，也可能发生，某些哲学家担任诸神的司祭；当这发生时，哲学并未与宗教结合；哲学在神圣仪式中将无所作为，而宗教在讨论哲学时也将无所作为。因为那种宗教仪式体系是哑默的，不仅因为它涉及哑默的神祇，还因为它的遵守是通过手和手指，而非通过心灵和舌头，就像我们的仪式那样，它是真实的。因此，宗教包含在智慧之中，智慧也包含在宗教之中。两者不可分离；因为智慧无非就是以正义和虔敬的敬拜来朝拜真天主。但多神崇拜不符合本性，这一点甚至可以通过以下论证来推断和领悟：每个人所敬拜的每一位神祇，在庄严的仪式和祈祷中，都必须被呼求为父，不仅为表尊敬，也是出于理性；因为他比人更古老，并且他赐予生命、安全、和滋养，如同父亲一样。因此，朱庇特被祈祷者称为父，正如萨图尔努斯、雅努斯、利贝尔以及其余众神依次被称为父；路奇里乌斯在诸神的会议中嘲笑这一点：\Quote{所以我们中没有一个不被称作众神卓越的父亲；所以父亲尼普顿、利贝尔、父亲萨图尔努斯、玛尔斯、雅努斯、父亲奎里努斯，都以一个名字被称呼。}但如果本性不允许一个人有许多父亲（因为他只由一人所生），那么多神崇拜就是违背本性、违背虔敬的。\par

因此，只有一位值得敬拜，祂才能真正被称为父。祂也必须必然是主，因为正如祂有赐予的能力，祂也有约束的能力。祂被称为父，是因为祂将许多伟大的事物赐予我们；被称为主，是因为祂有最大的惩罚和责戒的能力。而那为父的也是主，甚至可以从民法中看出。因为谁能抚养儿子，除非他对他们拥有如同主人般的权力？他被称为家主并非没有理由，虽然他只有儿子：因为显然，\Quote{父亲}这个名称也包括奴隶，因为\Quote{家}随之而来；而\Quote{家}的名称也包括儿子，因为\Quote{父亲}的名称在先：由此可知，同一个人既是奴隶的父，也是儿子的主。最后，儿子被释放如同奴隶；被释放的奴隶从其庇护主那里获得名字，如同儿子。但如果一个人被称为家主，表明他拥有双重权力，因为作为父亲他应当宽待，作为主人他应当约束，那么儿子也是奴隶，父亲也是主。因此，正如本性的必然性所规定，不能有多于一位父亲，同样也只能有一位主。因为如果许多主人发出相互矛盾的命令，奴隶该怎么办？因此，多神崇拜既违背理性也违背本性，因为不可能有许多父或主；但必须将诸神视为父和主。\par

因此，当同一个人服从于许多父和主，当心灵被引向多个对象而四处游荡时，真理便无法持守。当宗教没有固定和安定的居所时，它也不能有任何稳固性。因此，多神崇拜不可能是真实的敬拜；正如一个女人有多个丈夫不能被称作婚姻，而应被称为娼妓或奸妇。因为当一个女人缺乏端庄、贞洁和忠信时，她必然没有美德。同样，诸神的宗教体系也是不贞洁和不圣洁的，因为它缺乏忠信，因为那种不稳定和不确定的尊荣没有来源或根源。\par

% structure: p0013 source=h2 target=subsection reason=relative-heading-rank; base=h2

\subsection{第四章 论智慧与宗教，以及父与主的权柄}

由这些事可以看出\OriginalTerm{智慧}{wisdom}与\OriginalTerm{宗教}{religion}是如何紧密相连。智慧关涉儿子，这种关系需要爱；宗教关涉仆人，这种关系需要敬畏。因为正如前者应当爱戴并尊敬他们的父亲，后者也应当尊重并敬拜他们的主人。但就唯一的天主而言，既然祂同时具有圣父与主的双重身份，我们作为儿子应当爱祂，作为仆人应当敬畏祂。因此，宗教不能与智慧分离，智慧也不能与宗教分离；因为同一天主应当被理解——这是智慧的部分，也应当被敬拜——这是宗教的部分。但智慧在先，宗教在后；因为认识天主先于敬拜，敬拜是认识的结果。因此，这两个名称虽有不同，却只有一个意义。一个关乎理解，另一个关乎行动。然而，它们如同从同一泉源流出的两条溪流。但智慧与宗教的泉源是天主；若这两条溪流偏离祂，必然干涸；因为不认识祂的人不可能有智慧或宗教。\par

因此，哲学家和那些崇拜多神的人，要么像被剥夺继承权的儿子，要么像逃跑的奴隶，因为前者不寻求他们的父亲，后者不寻求他们的主人。正如被剥夺继承权的人不能获得父亲的遗产，逃跑的奴隶不能免于惩罚，同样，哲学家也不能领受永生，即天国的继承，也就是他们特别寻求的\OriginalTerm{至善}{chief good}；众神的崇拜者也不能逃脱永死的刑罚，那是真实的主对背叛祂威严和名号之人的惩罚。但天主是圣父也是主，这一点对两者——众神的崇拜者以及智慧的教授者——都是未知的：因为他们要么认为根本不该敬拜什么；要么认可虚假的宗教；要么，虽然他们理解至尊天主的权能和力量（如\Person{柏拉图}，他说有一位天主，是世界的创造者；以及\Person{马尔库斯·图利乌斯}，他承认人类被至尊天主以优越的状态创造），但他们却没有向祂如同向至尊圣父应尽的敬拜，这是他们相称且必要的职责。但众神不能成为父亲或主人，不仅因他们的众多，如我以上所示，也因理性：因为众神创造人这一点并无记载，也没有发现众神本身先于人类起源，因为在地球上有人的时候，似乎已有伏尔甘、利伯尔、阿波罗以及朱庇特本人出生。但人的创造通常不归功于萨图尔努斯，也不归功于他的父亲刻路斯。\par

但若被崇拜的众神中，没有一位据说当初形成并创造了人，那么结论就是：没有一位可被称为人的父亲，因此没有一位可以是神。所以，人不应当崇拜那些未曾产生人的神，因为人不可能由众神产生。因此，独一的天主应当被崇拜，祂存在于朱庇特、萨图尔努斯、刻路斯和大地之前。因为祂必定创造了人，祂在造人之前已经完成了天地。唯有那创造我们的才应被称为圣父；唯有那统治的主才应被视为主，祂拥有真实且永恒的生与死的权能。那不爱慕祂的人，是愚昧的仆人，逃避或不认识他的主人；是不孝之子，憎恨或不知道他真正的圣父。\par

% structure: p0017 source=h2 target=subsection reason=relative-heading-rank; base=h2

\subsection{第五章 必须查考先知的预言；论他们的时期，以及民长与君王的时期}

既然我已表明智慧与宗教不可分离，那么接下来我们当谈论宗教本身和智慧。我确实知道讨论天上的主题何其困难，但仍必须冒险尝试，为使真理得以显明并光照人，使许多因真理被愚昧遮掩而藐视拒绝它的人，得以脱离错误与死亡。但在开始论及天主及其作为之前，我必须先就先知说几句话，因为从现在起我要使用他们的见证，这在以前各卷中我未曾这样做。首先，渴望理解真理的人，不仅应当专心领会先知的言论，而且应当最勤奋地查考每位先知生活的时期，以便知道他们预言了哪些未来之事，以及他们的预言在多少年后得以应验。进行这些计算并无困难；因为他们见证了每位先知在哪位君王统治下领受圣神的默感。并且有许多人撰写并出版了关于各时期的书籍，从先知\Person{梅瑟}开始，他大约生活在特洛伊战争前七百年。\Person{梅瑟}治理人民四十年后，由\Person{若苏厄}继任，他执掌最高职位二十七年。\par

此后，他们在民长治下三百七十年。然后情况改变，他们开始有君王；当他们统治了四百五十年，直到\Person{漆德克雅}王朝，犹太人被巴比伦王围困，被掳去，忍受长期的奴役，直到七十年后，被掳的犹太人才被波斯人\Person{居鲁士}（他获得至高的权力）恢复自己的土地和居所，当时在罗马是\Person{塔奎尼乌斯·苏佩布}在位。因此，既然全部时间序列可以从犹太史以及希腊、罗马史中收集，各先知的时代也可以收集；最后一位是\Person{匝加利亚}，公认他在\Person{达理阿}王时代，即他执政第二年第八个月说预言。可见先知比希腊作家古老得多。我提出这一切，好让他们明白自己的错误——他们企图驳斥圣经，好像它是新近写成的，不知道我们神圣宗教的源流出自何泉。但如果有人汇总并考察这些时代，恰当地奠定学问的基础，并充分确定真理，他也会在获得真理知识后放弃错误。\par

% structure: p0020 source=h2 target=subsection reason=relative-heading-rank; base=h2

\subsection{全能天主生了祂的圣子；以及西比尔和特里斯美吉斯托斯关于祂的见证}

因此，万物的设计者与建立者天主，如我们在第二卷所说，在开始这美妙的世界工程之前，生了一个纯粹、不朽的灵，祂称之为圣子。虽然祂后来由自己创造了无数其他存有，我们称之为天使，但这首生者却是唯一被祂认为配得神圣名号的一位，因为祂在祂父亲的卓越与威严中大有能力。然而，至高大天主有一位拥有最大能力的圣子，这不仅由先知们异口同声的言论所表明，也由\Person{特里斯美吉斯托斯}的宣告和西比尔的预言所证明。\Person{赫尔墨斯}在其名为\WorkTitle{完美之言}的书中用了这些话：\Quote{万物的主和创造者，我们认为应当称之为神，因为祂使第二位神可见并可感。但我使用\InnerQuote{可感}一词，并非因为祂自己感知（问题不是祂是否自己感知），而是因为祂引导至感知和理智。因此，祂首先创造了祂，唯一且独一，祂在祂眼中是美丽的，充满一切美善；祂祝圣了祂，并全然爱祂如自己的圣子。}埃律特莱亚西比尔在她以至高神开头的诗篇开头，宣告天主子为万有的领袖和统帅，诗云：——\par

\Quote{万物的滋养者和创造者，祂将甜美的气息置于万物之中，并使神成为万有的领袖。}\par

又在同一首诗篇的结尾：——\par

\Quote{但神赐予忠信之人以尊崇的那位。}\par

另一位西比尔吩咐说应当认识祂：——\par

\Quote{认识祂为你们的神，祂是神的圣子。}\par

无疑他就是天主子，由那最智慧的君王\Person{撒罗满}（此人充满神圣灵感）说出了我们附加的这些话：\BibleQuote{天主在祂道路的开始就奠基了我，在祂万古之前的工作中。在起初，在祂造地之前，在祂奠定深渊之前，在水泉涌出之前，祂设立了我：主在众山之前生了我；祂造了区域，和天下无人居住的边界。当祂预备诸天时，我就在祂旁边；当祂分开自己的宝座，当祂在风之上造了坚固的云，当祂坚固山岳，将它们安置在天之下；当祂奠定大地的坚固根基时，我就在祂身旁安排一切。我就是祂所喜悦的那位：我天天喜悦，当世界完成，祂欢欣。}但\Person{特里斯美吉斯托斯}因此称祂为\Quote{神的工匠}，西比尔称祂为\Quote{参议}，因为天主圣父赐予祂如此的智慧和力量，以至天主在创造世界时同时使用了祂的智慧和双手。\par

% structure: p0028 source=h2 target=subsection reason=relative-heading-rank; base=h2

\subsection{第七章：关于圣子的名号，以及祂为何被称为耶稣和基督}

有人或许会问，这位如此有能力、如此蒙天主所爱的是谁，祂有什么名字，祂不仅在起初、在世界之前受生，还以其智慧安排、以其大能建造了世界。首先，我们应当知道，祂的名字连住在天上的天使也不知道，只有祂自己和天主圣父知道；在圣经中所记的，这天主的旨意成就之前，那名字不会公布。其次，我们必须知道，这名字不能用人的口说出，正如\Person{赫尔墨斯}教训说：\Quote{这原因的原因乃是神圣美善的旨意，它产生了神，神的名字不能用人的口说出。}不久后他对圣子说：\Quote{圣子啊，有一句智慧的密语，神圣地关于万有的唯一主和首先被心智感知的神，论及祂超出了人的能力。}但是，虽然至高圣父从起初赐给祂的名字除了祂自己无人知晓，然而祂在天使中有一个名字，在人间有另一个名字，因为在人间祂被称为耶稣；因为基督不是专名，而是权力与统治的称号；犹太人惯用此称他们的君王。但这个名字的含义必须说明，以免无知者因拼写改变而称祂为\OriginalTerm{克雷斯图斯}{Chrestus}。犹太人先前受命制备圣油，以傅油祝圣被召为司祭或君王的人。正如现在朱袍是罗马人承受王权的标记，在他们那里，以圣油傅油赋予了君王的称号和权能。但既然古希腊人使用\OriginalTerm{}{\GreekText{χρι}}……\par

\SourceRecord{Translated by William Fletcher.；Ante-Nicene Fathers；Vol. 7.；Edited by Alexander Roberts, James Donaldson, and A. Cleveland Coxe.；Buffalo, NY: Christian Literature Publishing Co.,；1886.；<http://www.newadvent.org/fathers/07014.htm>.}

% structure: p0001 source=h1 target=section reason=page-title

\section{\WorkTitle{神圣训导录，第五卷（论正义）}}

% structure: p0002 source=h2 target=subsection reason=relative-heading-rank; base=h2

\subsection{第一章：未经审问不得定罪；哲学家轻蔑圣经的原因；基督教信仰的首批辩护者。}

\Emph{我毫不怀疑}，啊，伟大的\Person{君士坦丁}皇帝——因为他们因过度的迷信而不耐烦——若那些愚蠢地虔诚的人中有人拿起我们这部作品，其中宣告了万有的无与伦比的创造者和这无限世界的统治者，他甚至会以辱骂之言攻击它，或许，还未读完开头，就会把它扔在地上，抛掷出去，咒骂它，并认为自己被污染了，若他耐心阅读或听到这些，就犯了无法赎偿的罪，被不洁所束缚。然而，我们根据人性之权利，若可能的话，要求这个人，他不应在了解整件事之前就定罪。因为如果辩护权被给予亵渎者、叛徒和行巫术者，并且不允许任何人在未经审理的情况下被预先定罪，那么我们的请求似乎并非不义：若有人偶然遇到这个主题，若他阅读，就应通篇阅读；若他听到，就应将判断推迟至终了。但我知道人的固执；我们永远无法获得这点。因为他们害怕被我们战胜，最终在真理本身呼喊时被迫屈服。因此，他们打断并制造障碍，以免听见；并闭上眼睛，以免看见我们呈现给他们的光明。因此他们自己清楚地显示出对自己已弃绝之体系的怀疑，因为他们既不敢研究，也不敢与我们较量，因为他们知道自己容易被击败。因此，讨论被排除后，\par

\Quote{智慧从他们中间被驱逐，他们诉诸暴力，}\par

正如\Person{恩尼乌斯}所说。并且因为他们热切地努力将那些他们明明知道是无辜的人定罪，他们不愿就无辜本身达成一致；仿佛，事实上，在证明无辜之后定其罪，比未经听审定其罪更为不义。但是，如我所说，他们害怕，若是他们听了，就无法定罪。\par

因此他们折磨、处死、流放至高天主的敬拜者，即义人；并且那些如此强烈憎恨的人自己也无法说出他们仇恨的原因。因为他们自己处于错误之中，他们对那些追随真理之路的人生气；当他们能够改正自己时，却通过残酷的行为大大增加他们的错误，他们被无辜者的血玷污，并用暴力将被奉献于天主的灵魂从撕裂的身体中扯出。这就是我们如今试图与之较量和辩论的人：这些人，我们想把他们从愚蠢的信仰引向真理，他们更愿意喝血，而不愿吸收义人的话语。那么如何？我们的劳动会白费吗？绝不。因为如果我们不能将这些正以极大速度奔向死亡的人从死亡中解救出来；如果我们不能将他们从那条弯曲之路召回生命和光明，因为他们自己反对自己的得救；然而我们将坚固那些属于我们的人，他们的意见尚未坚定，尚未以坚固的根基建立和稳固。因为他们中许多人在摇摆，尤其是那些对文学有所了解的人。因为在这方面，哲学家、演说家和诗人是有害的，因为他们能够轻易地用他们甜美的言辞和悦耳节奏的诗歌诱惑不小心的灵魂。这些是藏有毒药的甜蜜。出于这个原因，我希望将智慧与宗教联系起来，以便那虚妄的体系一点也不会伤害好学的人；如此，文学知识不仅不会对宗教和正义造成伤害，反而会带来最大的益处，如果学习它的人应在美德上更加受教，在真理上更为智慧。\par

此外，即使它可能对别人无益，对我们却肯定有益：良心将自悦，心灵将欢喜，因为它沉浸在真理之光中，这光是灵魂的食物，被一种不可思议的愉悦所覆盖。但我们不可绝望。或许\par

\Quote{我们不向聋者歌唱。}\par

因为事务尚未恶劣到没有健全心智的地步，以致真理不能令人喜悦，并且他们也能看见并追随指给他们的正途。只需将智慧的天上蜂蜜涂抹在杯沿，使得苦涩的药液在不知不觉中被他们饮下，毫无厌烦，而最初的甘甜以其诱人之处，在愉悦的遮盖下隐藏了辛辣滋味的苦涩。这正是智慧者、学者及此世权贵不信任圣经的主要原因，因为先知们以通俗简单的语言说话，仿佛是对民众讲话。因此，那些只愿听闻或阅读精炼雄辩之作的人轻视它们；除了以更悦耳之声愉悦他们耳朵的，没有什么能在他们心中留存。而那些看似粗鄙的，则被视为老妇之谈、愚昧和庸俗。他们如此地认为只有悦耳者才是真的，只有能激发快感者才是可信的；没有人依据真理判断事物，而是依据其修饰。因此，他们不信圣经，因为圣经毫无矫饰；他们甚至不信那些解释圣经的人，因为那些人要么全然无知，要么至少学识浅薄。因为很少有完全雄辩的；原因显而易见。雄辩服务于世界，它渴望向民众展示自己，并在邪恶之事上取悦人；因为它常常竭力压倒真理，以显示其力量；它追求财富，渴望尊荣；简言之，它要求至高的尊严。因此，它鄙视这些低微的题材；它避免隐秘之事，因其与自身相悖，因为它以公开为乐，渴求群众与名声。由此导致智慧与真理缺少合适的传报者。若偶有学者投身其中，他们也不足以捍卫之。\par

据我所知，\Person{米努丘·菲利克斯}在律师中并非等闲之辈。他的著作《\WorkTitle{屋大维}》显示，若他完全致力于那项追求，他本该是多么合适的真理维护者。\Person{塞普蒂米乌·德尔图良}也精通各类文学；但在雄辩方面，他欠缺敏捷，不够润饰，且十分晦涩。因此，他也没有获得足够的声誉。\Person{西彼廉}则格外卓越和著名，因为他从演说术的职业中为自己谋求了巨大荣耀，并写了许多在其类别中值得赞叹的作品。他思维敏捷、丰富、令人愉悦，并且（这是风格的最大优点）平易明晰；以至于你无法判断他是更擅长言辞修饰，还是更擅长解释说明，或是更善于说服人。然而，对于那些不了解奥秘的人，他仅凭言语无法取悦他们；因为他所说的是神秘之事，并为此目的而准备，以便只让信徒听闻。简言之，同时代的学者们若偶然读到他的作品，常常嘲笑他。我听说一位颇有技巧的人，通过改动一个字母称他为\Person{科普利亚努斯}，仿佛他将一颗文雅且适合更高事物的心思用在了老妇的寓言上。但如果这发生在他身上（他的雄辩并非令人不悦），那么对于那些言辞贫乏、不悦耳、既无说服力、也无论证的巧妙、更无任何驳斥力量的人，我们该如何设想呢？\par

% structure: p0011 source=h2 target=subsection reason=relative-heading-rank; base=h2

\subsection{第二章  基督徒真理受到何等鲁莽之人的攻击}

因此，由于我们当中缺乏合适且有技巧的教师，能有力地、尖锐地驳斥公共谬误，并能以优雅和丰富来捍卫全部真理，正是这种缺乏促使一些人冒险攻击他们未知的真理。我略过那些先前徒然攻击它的人。当我在\Place{比提尼亚}教授修辞学（我应召去那里），且恰好天主的殿堂被推翻时，同地住着两个人，他们以傲慢或苛刻（我不知道何者更甚）侮辱那被推翻倒地的真理：其中一人自称为哲学的大司祭；但他如此沉溺于恶习，以致身为克己的导师，他既贪财又纵欲；生活如此奢侈，以致虽在学校中他是美德的维护者、节俭和贫穷的赞美者，他在王宫中用餐还不如在自己家中丰盛。然而，他用他的长发和斗篷，以及最大的掩饰——财富，来掩护他的恶习；为了增加财富，他以非凡的努力渗透到法官的友谊中；他通过一个虚假名声的权威突然吸引他们归向自己，不仅是为了交易他们的判决，而且是为了通过这种影响阻止他的邻居（他正将他们赶出家园和土地）收回他们的财产。这个人，实际上，用自己的品格推翻了自己的论据，或者用自己的论据谴责了自己的品格，一个严厉的批评者和最尖锐的控告者，就在那同一时刻（当一群正义之民被不虔敬地攻击时），他吐出了三本书来反对基督徒的宗教和名声；他特别声称，哲学家的职责是纠正人的错误，并将他们引回真道，即崇拜诸神（正如他所说，世界靠神的能力和威严统御），并且不容许无经验的人被任何人的诡计所引诱，以免他们的单纯成为狡诈者的猎物和养分。\par

因此他说，他承担了这一配得上哲学的职责，为要向那些看不见智慧之光的人指明道路，不仅是使他们通过敬拜诸神而恢复健康的心智，而且还要他们放下顽固的执拗，避免身体的酷刑，不愿徒劳地忍受残忍的肢体摧残。但为了表明他为何辛劳地完成那项任务，他大肆赞扬诸位君王，他们的虔敬和远见——正如他本人所言——在其他事务中，尤其是在维护诸神的宗教礼仪方面，显得卓著；总之，他关心人类的利益，为要制止不虔敬和愚蠢的迷信，使所有人都有闲暇从事合法的神圣礼仪，并体验到诸神对他们的眷顾。然而，当他试图削弱他所反对的宗教的根基时，他显得愚蠢、虚妄和可笑；因为那位为他人利益着想的重量级顾问，不仅不知道要反对什么，甚至不知道要说什么。因为若有我们宗教中人在场，虽然因时势而沉默，却仍会在心中嘲笑他；因为他们看到一个人宣称要启发他人，而自己却瞎眼；要引他人出离错误，而自己却不知将脚放在何处；要教导他人真理，而自己却从未在任何时候见过真理的一丝火花；这位自称智慧之师的人，竟试图推翻智慧。不过，所有人都指责他，尤其因为他偏在这一可憎的残暴肆虐之时承担这项工作。啊，哲学家，谄媚者，趋炎附势者！但这个人正如其虚妄所应得的，被鄙弃；因为他未能获得所期望的声望，而他急切追求的光荣也变成了指责和责备。\par

另一个人以更尖刻的笔调写了同样的话题，他当时属于审判官之列，并且是建议发动迫害的主要人物；他不满足于这一罪行，还用文字迫害那些他所迫害过的人。因为他写了两卷书，不是反对基督徒（以免显得以敌意攻击他们），而是写给基督徒（好让人认为他是出于人性和善意为他们着想）。在这些著作中，他竭力证明\WorkTitle{圣经}是虚假的，仿佛它全然自相矛盾；因为他阐释了一些看似彼此矛盾的章节，列举了如此多且如此隐秘的事情，以至于他有时显得像是同一教派中的人。但若果真如此，哪个\Person{德摩斯梯尼}能为他辩护，免于不虔敬之罪？他背叛了他所赞同的宗教，背叛了他所信奉的信德之名，背叛了他所领受的奥迹——除非那些\WorkTitle{圣经}偶然落在他手中。那么，胆敢摧毁无人向他解释的东西，是何等鲁莽！幸好他要么什么也没学到，要么什么也没理解。因为矛盾离圣经之遥远，正如他离信德和真理之遥远。不过，他主要攻击了\Person{保禄}和\Person{伯多禄}，以及其他门徒，视他们为欺骗的散布者，同时他又证实他们是无技巧、无学问的人。因为他说其中一些人靠渔夫的手艺谋利，仿佛他愤慨于某个\Person{阿里斯托芬}或\Person{阿里斯塔尔科斯}没有构思出那个主题。\par

% structure: p0015 source=h2 target=subsection reason=relative-heading-rank; base=h2

\subsection{第三章 论\Person{基督}教义之真理及其反对者之虚妄，并论\Person{基督}并非术士}

因此，这些人是没有技巧的，故缺乏发明的欲望和狡诈。或者，一个无学问的人怎能发明出彼此适应、连贯一致的东西呢？而最有学问的哲学家\Person{柏拉图}、\Person{亚里士多德}、\Person{伊壁鸠鲁}和\Person{芝诺}，他们所说的却彼此矛盾、相互对立？因为虚假的本性就是无法连贯一致。但他们的教导，因为是真理，处处一致，且完全自洽；正因如此，它能产生说服力，因为它基于一以贯之的计划。因此，他们并非为了利益和好处而杜撰那宗教，因为他们既在训诫中，也在实际上遵循了那种没有享乐的生活，并藐视一切被认为是善的东西；而且他们不仅为信德忍受死亡，还知道并预言自己将要死去，此后所有追随他们体系的人也将遭受残酷和不虔敬的事。但那人断言，\Person{基督}本人被犹太人驱逐，并集合了一支九百人的队伍，从事抢劫。谁敢反对如此伟大的权威？我们当然必须相信这一点，因为或许某个\Person{阿波罗}在他梦中向他宣告了此事。如此多的强盗，从古至今都灭亡了，而且每天都在灭亡，你本人无疑也定罪了许多人：他们之中谁在钉十字架后被称为——我不说神——但被称为人呢？但你或许相信了这一点，因为你将杀人犯\Person{玛尔斯}尊为神，虽然如果\Person{亚略巴古人}把他钉了十字架，你就不会这样做。\par

此人试图推翻\Person{基督}的奇妙作为，却并未否认这些作为，他想要证明\Person{阿波罗尼乌斯}行了同样甚至更伟大的事。奇怪的是，他竟未提及\Person{阿普雷乌斯}，关于此人常有许多奇事流传。因此，无知的啊，为何无人以\Person{阿波罗尼乌斯}代替天主来敬拜？除非偶然只有你一人这样做，你确实配得上那位神，真天主将用他来永远惩罚你。如果\Person{基督}因行了奇妙之事而被当作术士，那么\Person{阿波罗尼乌斯}——按你的描述，当\Person{多米提安}想要惩罚他时，他在受审中突然消失——显然比那位被捕并被钉十字架的更狡猾。但或许他正想借此证明\Person{基督}的狂妄，因他自称是神，好使另一位显得更为谦虚，他虽行了更伟大的事（如这位所想），却并未将此归于自己。我暂不比较这些作为本身，因为在第二卷及前卷中我已论及巫术的欺骗与伎俩。我要说，无人不希望死后能获得那连最伟大的君王也渴望的事。因为人为何为自己建造宏伟的坟墓、雕像和肖像？为何以显赫的功绩，甚至为国捐躯，来赢得人的好评？总之，你为何以这可憎的愚昧，如同以泥巴，建造你才智的纪念碑，若不是你期望从你名声的纪念中获得永生？因此，认为\Person{阿波罗尼乌斯}不渴望那显然是他若能得到就必定渴望的事，这是愚蠢的；因为无人拒绝永生，尤其当你说他被人当作神来敬拜，他的像以避凶者\Person{赫拉克勒斯}之名被竖立，甚至如今仍受厄弗所人敬奉。\par

因此，他死后不能被认为是神，因为显然他既是一个凡人，又是一个术士；为此他以他人的名号装作神，因为用他自己的名号他无法达到，也不敢尝试。但我们所说的这位，却能被当作神，因为他不是术士，并且他之所以被相信是神，是因为他确实是神。他说：我并非说\Person{阿波罗尼乌斯}不被当作神是因为他不愿意，而是为了表明，我们——并未因奇事而立即相信他的神性——比你们更明智，你们因些许奇事就相信他是神。你们远离天主的智慧，对所读之事一无所知，这并不奇怪，因为犹太人从起初就常读先知书，且天主的奥秘已托付给他们，他们却仍不明白所读的。因此，若你们尚有理智，要知道：我们相信\Person{基督}是天主，并非因他行了奇事，而是因我们看见在他身上所成就的一切，都是藉先知的预言向我们宣告的。他行了奇事：我们本可把他当作术士，正如你们如今以为的，以及那时犹太人所想的，若非所有先知一致宣告\Person{基督}将行那些事。因此，我们相信他是天主，与其说因他的奇事和作为，不如说因那你们如狗舔舐的十字架，因为这十字架也同时被预言了。所以，他并非凭自己的见证（因为人谈论自己时谁能被相信？），而是凭先知的见证——他们早已预言了他所行和所受的一切——他获得了神性的信仰，这既不可能发生在\Person{阿波罗尼乌斯}、\Person{阿普雷乌斯}或任何术士身上，也永不可能发生。因此，当他倾吐了如此无知而荒谬的胡言乱语，当他急切地要彻底毁灭真理时，他竟敢将他那邪恶、与天主为敌的著作冠以\Quote{爱真理}之名。盲目的心智！比所谓的西米里黑暗更昏暗的头脑！他或许曾是\Person{阿那克萨戈拉斯}的门徒，对他而言雪黑如墨。但这是同样的盲目：称真理为谎言，称谎言为真理。这狡猾的人无疑想披着羊皮隐藏狼，以便用欺骗性的标题诱捕读者。就算这是真的，就算你这样做是由于无知而非恶意：然而，你给我们带来了什么真理呢？除了作为诸神的辩护者，你最终竟出卖了那些神本身？因为，你陈述了至高天主的赞美——你承认他是君王、至大能者、万有的创造者、尊荣的泉源、万物的父亲、一切活物的创造者和保存者——你从你自己的\Person{朱庇特}那里夺走了王国；你将他从至高权位驱赶，贬低到仆役的行列。这样，你自己的结论就定了你的愚蠢、虚妄和错误。因为你肯定诸神存在，却又将他们置于你试图推翻其宗教的那位天主之下，使他们成为奴仆。\par

% structure: p0019 source=h2 target=subsection reason=relative-heading-rank; base=h2

\subsection{第四章　为何出版这部著作，再论\Person{戴尔都良}与\Person{西彼廉}。}

因此，既然我上面所提及的那些人当着我的面、令我痛心地宣扬了他们亵圣的著作，我就被这些人的傲慢不敬以及真理本身的意识（且如我所想，也被天主）所激励，承担起这一职责，竭尽我心神之力驳斥正义的控告者；这并不是要专门反驳这些人——用几句话就可以击溃他们——而是为了通过一次攻击，一举推翻所有在各处从事或已从事同一工作的人。因为我不怀疑，有极多其他人在许多地方，不仅在希腊著作中，也在拉丁著作中，竖立了他们自己邪恶的纪念碑。既然我无法逐一对这些人作答，我认为应当如此辩护这一案件，以致能连同他们所有的著作推翻前期的作者，并且切断后来作者所有写作与答辩的能力。只让他们注意，我必将达成：凡认识这些事的人，要么接受他先前所谴责的，要么（与此相近）至少停止再嘲笑它。虽然戴尔都良在其名为\WorkTitle{护教篇}的论著中充分阐述了同一案件，但是，答复控告者——这仅在于辩护或否认——与教训人（我们所作的事，其中必须包含整个体系的实质）是两回事，我没有回避这一劳苦，为要完成西彼廉在他那篇意图驳斥德默特利亚努斯（正如他自己所说）诋毁并叫嚣反对真理的演说中所未能充分完成的内容。他本应恰当地处理这一主题；但他不应用圣经的见证来反驳他——因为德默特利亚努斯显然视之为虚空、虚构和虚假——而应用论证和理性。因为既然他在与一个不认识真理的人争辩，他就应当暂时放下神圣的经文，从头开始塑造一个完全无知的人，并将光明的开端逐步展示给他，以免将全部亮度呈现给他时使他眼花缭乱。\par

就像婴儿因其胃的娇嫩，无法接受固体和强壮食物的营养，而只能靠流质和柔软的乳汁维持，直到力量增强，可以吃较强壮的食物；同样，适合这个人的做法是：因他尚不能接受神圣事物，就应以人的见证——即哲学家和历史学家的见证——呈现在他面前，好让他尤其被自己的权威所驳倒。既然他没有这样做，却被他对圣经的卓越知识所带走，以致只满足于那些信仰所在之事，我就蒙天主恩宠承担此事，同时为他人效仿铺平道路。如果因我的劝勉，博学善辩的人开始投身于这一主题，并选择在真理的这块田地上展示他们的才华和口才，那么无人可以怀疑虚假宗教将迅速消失，哲学将完全垮掉，只要所有人都被说服：唯有这才是宗教和唯一的真智慧。但我偏离主题已超出我的愿望。\par

% structure: p0022 source=h2 target=subsection reason=relative-heading-rank; base=h2

\subsection{第五章  撒杜尔努斯时代曾有真正义，但被尤庇特驱逐。}

现在必须进行所承诺的关于正义的讨论；正义本身要么是最伟大的美德，要么是美德的源泉，不仅哲学家寻求它，而且更早的诗人也寻求它——他们在哲学之名产生之前就被视为有智慧。这些人清楚地明白，这种正义已离开人间事务；他们虚构说，正义被人性的恶行所冒犯，离开大地，升天而去；并且为了教导什么是按正义生活（因为他们习惯用周回婉转的方式给予训诲），他们从撒杜尔努斯时代——他们称之为黄金时代——重复正义的实例，并描述当正义仍留在地上时人类生活的状况。这不应被视为诗的虚构，而是真理。因为当撒杜尔努斯统治时，尚未建立对诸神的宗教敬拜，也没有任何种族被划归为神性信仰，天主被清晰地崇拜。因此，那时没有纷争、没有仇敌、没有战争。\par

\Quote{狂怒尚未抽出疯狂的剑，}\par

正如盖尔玛尼库斯·凯撒在其翻译自阿拉托斯的诗中所说：\par

\Quote{亲戚之间也未曾知道纷争。}\par

不，甚至陌生人之间也没有；但根本没有剑可抽出。因为当正义在场并活跃时，谁会想到自己的防护呢？既然无人谋害他；谁会想到毁灭他人呢？既然无人贪求任何东西？\par

\Quote{他们宁愿满足于简朴的生活方式，}\par

正如西塞禄在其诗中所叙述的；而这正是我们宗教的特点。\Quote{甚至不允许划界或分割平原：人们共同寻求一切；}因为天主将大地共同赐给所有人，让他们共同生活，而非让疯狂和贪婪的贪欲为自己索取一切，使为众人所出产的不缺乏任何人。对诗人的这句话应当如此理解，并非暗示那时个人没有私有财产，而应被视为诗的修辞：我们当理解为人们是如此慷慨，不把为他们所出产的果实封闭起来，也不独自守藏所储之物，而是允许穷人分享他们劳动的果实：——\par

\Quote{如今奶流如注，如今琼浆涌流。}\par

这并不奇怪，因为善人的仓库慷慨地向所有人开放。贪欲没有拦截天主的恩赐，从而造成共同的饥渴；但所有人同样富足，因为拥有财物的人慷慨地分给那些没有的人。但在撒杜尔努斯被逐出天界、到达拉丁姆之后，——\par

\Quote{被朱庇特，他更强有力的继承人，\newline{}从王位放逐，}——\par

因为人民或因惧怕新王，或出于自愿，已败坏而停止敬拜天主，并开始把君王当作天主来尊崇，因他本人几乎是个弑亲者，成为他人损害虔敬的榜样——\par

\Quote{最公义的处女急忙离开了大地；}\par

但并非如西塞罗所言，\par

\Quote{并在朱庇特的王国和天堂的一部分定居下来。}\par

她怎能定居或停留在那位将父亲逐出王国、以战争折磨他、将他驱逐到世界各地的君王的王国呢？\par

\Quote{他赐给黑蛇有毒的毒液，\newline{}并命令恶狼四处游荡；}\par

也就是说，他在人类中引入了仇恨、嫉妒和阴谋，使他们如蛇一样有毒，如狼一样贪婪。那些迫害正直且忠于天主的人，并赋予法官以暴力对待无辜者的权力的人，确实如此行。朱庇特也许曾做了这类推翻和消除正义的事；因此据记载，他使蛇变得凶猛，并磨砺了狼的性情。\par

\Quote{于是战争那不可驯服的暴怒，\newline{}以及贪婪的获取欲，}\par

并非没有理由。因为对天主的敬拜被除去，人类丧失了善恶的知识。于是，人与人之间的共同交往灭亡了，人类社会的纽带被摧毁。他们开始互相争斗、密谋，并从流人血中为自己获取荣耀。\par

% structure: p0042 source=h2 target=subsection reason=relative-heading-rank; base=h2

\subsection{第6章 正义被放逐后，情欲、不义的法律、胆大妄为、贪婪、野心、骄傲、不敬神以及其他恶习猖獗。}

所有这些邪恶的根源是情欲；它确实是从对真正威严的蔑视中爆发出来的。因为不仅那些有余的人没有分给其他人，他们甚至还夺取别人的财物，将一切都据为己有；以前连个人也努力为人类公用而获取的东西，如今被运到少数人的家中。为了用奴役来制服他人，他们开始特别从生活中撤出并收集必需品，严密地锁起来，以使上天的恩赐成为他们自己的；并非出于仁慈（他们心中毫无这种情感），而是为了搜集情欲和贪婪的一切工具。他们也以正义的名义，通过了最不平等和不公的法律，以维护他们的掠夺和贪婪，对抗大众的力量。因此，他们既凭权威，也凭力量、资源或恶意得势。既然在他们身上没有一点正义的痕迹（正义的职责是仁爱、公平、怜悯），他们便开始以骄傲而膨胀的不平等为乐，并通过随从、刀剑和华丽的衣裳使自己高于他人。为此，他们为自己发明了荣誉、紫袍和束棒，以便借由斧头和刀剑造成的恐怖，像主人一样以权力统治那些战战兢兢、惊恐不安的人。人类生活就处于这样的状况中，由于那位战胜并驱逐了父亲的君王，他没有夺取王国，而是以暴力和武装的人建立了不敬神的专制，并夺走了那正义的黄金时代，迫使人们变得邪恶和不敬神，甚至由此将人从敬拜天主转向敬拜他自己；他过分权力的恐怖迫使人如此行。\par

因为谁会不惧怕那位全副武装、被不常见的钢铁和刀剑闪光围绕的人呢？他连自己的父亲都不饶恕，又怎会饶恕陌生人呢？实际上，他战胜并屠杀了一族强大且勇武的泰坦，谁不应惧怕他呢？难怪全体大众被不寻常的恐惧所压迫，竟对一个人阿谀奉承。他们尊崇他，给予他最大的荣誉。既然模仿君王的习俗和恶习被视为一种谄媚，所有人都放下了虔敬，以免如果他们活得虔敬，似乎会责怪君王的邪恶。因此，被持续的模仿所败坏，他们放弃了神圣的权利，邪恶生活的习惯逐渐成为习性。现在，前一代虔敬而优越的状况已荡然无存；正义被放逐，并带着真理一同离去，留给人类谬误、无知和盲目。因此，那些歌唱正义逃往天堂、到朱庇特王国去的诗人们是无知的。因为如果正义在他们称为\Quote{黄金}的时代存在于地上，那么很明显，她是被改变了黄金时代的朱庇特赶走的。但时代的变化和正义的驱逐，正如我所说，无非是舍弃了对天主的宗教敬拜；唯有这种敬拜使人视他人为可贵的，并知道自己因同一的父——天主而与他有兄弟之谊，从而与那些没有的人分享共同的天主和父的恩赐；不伤害任何人，不压迫任何人，不对陌生人关闭家门，不对恳求者掩耳，而是慷慨、仁慈、宽厚。图利乌斯认为这些是适合君王的赞美。这确实是正义，这就是黄金时代。这个时代最初在朱庇特统治时被败坏，不久之后，当他本人和他的所有后代都被奉为神明，许多神祇的敬拜被采纳时，就被完全除去了。\par

% structure: p0045 source=h2 target=subsection reason=relative-heading-rank; base=h2

\subsection{第7章 论耶稣的来临及其果实；以及那时代的德行与恶习。}

然而，天主作为最宽容的慈父，当最后时刻临近时，派遣了一位使者，使那旧时代和已被驱散的正义复归，免得人类被极重大且永久的谬误所搅动。因此，那黄金时代的景象复返，正义重临大地，但只赐予少数人；而这正义无非就是对唯一的天主虔诚与宗教性的敬拜。但或许有人会问：如果这就是正义，为何不赐予全人类，为何全体民众不都同意它？这是一个大争论的问题，为何天主在将正义赐予大地时保留了差异；这件事我已在别处表明，每遇合适机会都将予以解释。现在只需简明扼要地指出：德行若不与恶习对立，便无从辨别；若不通过逆境操练，便不能完美。因为天主设计了善恶之间的这种区分，好让我们从恶中认识善的性质，也从善中认识恶的性质；若一方被除去，另一方的性质便不能理解。因此，天主并未排除恶，好让德行的性质得以显明。因为，若没有必须承受之事，坚忍如何能保持其意义和名称？若没有谁想使我们背弃天主，忠于天主的信德又如何能得称赞？为此，祂容许不义者更为强大，好使他们能强迫人行恶；也容许他们为数更多，好使德行因稀有而珍贵。\Person{昆体良}在\Quote{蒙着头的人}中极精彩地简要指出了这一点。他说：\Quote{因为若纯洁的稀有性未给它带来赞美，纯洁中还有什么德行可言？但因本性使然，仇恨、欲望和愤怒将人盲目地驱向他们所趋向的对象，那么免于过犯似乎超乎人的能力。否则，若本性赋予所有人同等的感情，虔诚便算不得什么。}\par

这有多真实，事实本身的必然性就教导我们。因为，若德行就是刚毅地抵抗邪恶与恶习，那么显然，没有邪恶与恶习，就没有完美的德行；而天主为使此德行完全且完美，保留了与它对立者，与之争战。因为，德行被折磨它的邪恶所激动，就获得稳固；它被推动得越频繁，就越坚强。显然，这就是原因，它导致虽然正义被赐予人，但还不能说有一个黄金时代存在；因为天主未曾除去恶，以保留那唯独保存了神圣宗教奥迹的差异性。\par

% structure: p0048 source=h2 target=subsection reason=relative-heading-rank; base=h2

\subsection{第八章　人人皆知正义，却未接纳；论天主的真正殿宇及祂的崇拜，使诸恶习被征服。}

因此，那些认为无人是正义的人，他们眼前有正义，却不愿辨识它。因为，如果他们愿意，为善对他们而言很容易，那他们有什么理由在诗歌或一切言谈中描述正义，抱怨它的缺席呢？你们为何将正义描绘为无价值的，并希望她仿佛以某种形象从天上降下？看，她就在你们眼前；若你们有能力，就接纳她，将她安置在你们心灵的居所；不要以为这是困难的，或不适合这时代。要正义且良善，你们所寻求的正义会自愿跟随你们。从你们心中除去一切恶念，那黄金时代就会立刻回到你们中间，除开始敬拜真天主外，别无他法可达到。但你们渴望大地上的正义，而同时敬拜假神继续存在，这不可能实现。即使在那你们臆想的时代，也不可能，因为那时你们不虔敬地崇拜那些神祇尚未出现，而唯一的天主——我是指那憎恨邪恶、要求良善的天主——的敬拜本应遍及全地；祂的殿宇不是石头或泥土，而是人自身，即那带着天主肖像的人。这殿宇不是用可朽坏的金银宝石装饰，而是用德行的持久事工装饰。所以，如果你们还有任何理智，要知道：人之所以邪恶不义，是因为敬拜假神；一切祸患之所以天天增加在人间事务上，是因为创造并治理这世界的大天主被忽视了；因为违背正直，接受了不虔的迷信；最后，因为你们连让少数人来敬拜天主也不容许。\par

但是，如果唯独天主受到敬拜，就不会有纷争和战争，因为人们将知道他们是同一天主的子女；因此，在被神圣而不可侵犯的神圣亲缘纽带联系起来的人之间，就不会有阴谋，因为他们知道天主为灵魂的毁灭者预备了何种惩罚，天主看穿秘密的罪行，甚至内心的思想本身。如果他们通过天主的教导学会满足于自己所有的、虽然微薄的东西，从而宁愿选择坚固永恒的事物而非脆弱易逝的，就不会有欺诈和掠夺。如果所有人都知道，凡超出生育欲望所求的事都被天主谴责，就不会有奸淫、放荡和卖淫。女人也不会因生计所迫而玷污自己的贞洁，寻求最可耻的谋生方式，因为男人也会克制自己的私欲，而富人虔诚仁爱的施舍会救济穷人。因此，正如我所说，如果人们共同遵守天主的法律，如果所有人都做我们的人民独自做的事，那么地上就不会有这些邪恶。如果温良、虔诚、和平、无辜、公正、节制和信德遍布世界，人间的事务将多么幸福、多么黄金般啊！总之，将不需要如此众多且各式各样的法律来治理人，因为仅凭天主的法律就足以达成完美的无辜；也不需要有监狱、统治者的刀剑或惩罚的恐怖，因为神圣诫命的健康性注入人心，自然教导他们行正义之事。但现在，人们因无知于正直与良善而邪恶。西塞罗确实看到了这一点；他在论述法律时说：\Quote{正如世界以其各部分彼此一致，基于同一本性而凝聚和依存，同样，所有人，虽然原本彼此交融，却因败坏而分歧；他们不明白自己有血缘关系，且都受同一保护权的管辖：如果牢记这一点，人类必定会过神明的生命。}因此，对诸神不义不敬的敬拜带来了所有邪恶，人类因此互相毁灭。因为他们既已如同浪子与悖逆的子女，弃绝了共同父——天主的权威，就不能保持自己的虔诚。\par

% structure: p0051 source=h2 target=subsection reason=relative-heading-rank; base=h2

\subsection{第九章　论恶人的罪行与加于基督徒的酷刑。}

然而，有时他们意识到自己是邪恶的，赞美先前的时代状况，并猜想正义因他们的品性和功过而缺席；因为，虽然正义出现在他们眼前，他们不仅不接受或不认识她，反而猛烈地憎恨、迫害、并试图驱逐她。让我们暂时假设我们所追随的并非正义：倘若她所认为的真正义来了，他们会如何接待她？因为他们折磨并杀害那些他们自己承认是义人的效仿者，只因这些人行善行义；而如果他们只杀死恶人，他们本应配得上不被正义临到，因为正义离开人间的唯一理由就是流血。更何况他们杀害义人，视正义的追随者本身为仇敌，甚至比仇敌更甚？这些人虽然用刀剑与烈火急切地夺取义人的生命、财产和子女，但被征服时却得到宽恕；甚至在战争中也有仁慈之地；或者如果他们决定将残忍进行到底，对他们所做的无非是处死或掳为奴隶！但对待那些无罪之人，这却是不可言说的，没有人比那些在所有人类中最无辜的人被视为更有罪。因此，最邪恶的人竟胆敢提及正义，这些人比野兽更凶残，蹂躏天主最温驯的羊群——\par

\Quote{如同骨瘦如柴的狼从巢穴冲出，\newline{}无法无天的饥饿发出阴沉嗥叫，\newline{}驱赶它们到黑夜中四处游荡。}\par

但这些使他们疯狂的并非饥饿的狂怒，而是内心的狂怒；他们不是在黑雾中潜行，而是公开掠夺；罪行的意识从未阻止他们用以亵渎正义的神圣与圣名的那张嘴，如同野兽的颚，沾满无辜者的鲜血。我们该说这种过度而持久的仇恨的原因是什么？\par

\Quote{真理产生仇恨，}\par

正如诗人所说，仿佛受到神圣的默感，或者他们羞于在正直善良之人面前作恶？还是两者兼有？因为真理总是可恨的，原因在于：犯罪者希望有犯罪的自由，并认为唯有当没有人对他的过失不满时，才能更安全地享受邪恶行为的快感。因此，他们竭力彻底消灭并除掉这些作为他们罪行和邪恶见证的人，视他们为负担，仿佛自身的生活受到谴责。因为为何有人要在公共道德败坏时不合时宜地行善，以善良的生活谴责他人？为何不所有人都同样邪恶、贪婪、不贞、奸淫、发假誓、贪婪与欺诈？为何不干脆除掉那些人——在他们面前，恶人羞于过邪恶的生活？那些人虽不用言语（因他们沉默），却以其与恶人截然不同的生活本身攻击并击打罪人的额头。因为凡与他们意见不合者，似乎在责备他们。\par

若这些事是向人做的，倒也不足为奇，因为为了同一个理由，那处于希望中且并非不认识天主的百姓，竟起来反抗天主自己；同样的必然追随着义人，曾攻击过正义的创立者本身。因此，他们以精心设计的酷刑折磨义人，并认为杀死他们憎恨的人还不算完，除非残忍也嘲弄他们的身体。但若有人因惧怕痛苦或死亡，或由于自身的背信弃义，放弃了天上的誓言，并同意参加致命的祭献，这些人他们就赞美，并以荣誉加身，好以他们的榜样引诱他人。但对于那些高度重视自己的信德、且没有否认自己是天主崇拜者的人，他们便使出全部屠戮之力扑向他们，仿佛他们渴求鲜血；他们称这些人为绝望者，因为他们绝不怜惜自己的身体；仿佛还有什么比折磨并撕裂一个你明知无辜的人更为绝望。因此，那些已失去一切善情的人中，毫无羞耻感，他们将以公正之人应受的辱骂反加到他们身上。因为他们称他们为不虔敬，而他们自己自然是虔敬的，且不愿流人血；然而，如果他们思量自己的行为以及那些他们斥为不虔敬者的行为，他们就会明白自己是多么虚假，多么应受一切他们向善人所说或所做之事。因为他们不是我们中的一员，而是属于那些武装劫道、海上行盗的人；或者，若他们无力公开袭击，便暗中下毒；他们杀害妻子以获取嫁妆，或杀害丈夫以便与奸夫结婚；他们或勒死亲生的儿子，或若他们过于虔敬，便将之遗弃；他们对女儿、姐妹、母亲、女司祭均不克制乱伦的激情；他们阴谋反对自己的同胞和祖国；他们不怕被装入袋中；总之，他们犯下亵渎之罪，劫掠他们所崇拜的神明的庙宇；且说些轻微而常行的事：他们争抢遗产，伪造遗嘱，或驱逐或排斥合法的继承人；他们卖淫以供人纵欲；总之，他们忘记了自己的出身，与女子在被动行为上争胜；他们违反一切体统，玷污和侮辱自己身体最神圣的部位；他们阉割自己，而更不虔敬的是，为了成为宗教的司祭；他们甚至不惜自己的生命，在公开场合卖命求死；他们若身为法官，因受贿而腐化，或毁灭无辜，或无惩罚地释放有罪之人；他们以巫术觊觎上天本身，仿佛大地容不下他们的恶行。这些罪行，我说，且不止这些，显然是由那些崇拜诸神的人所犯的。\par

在如此众多且重大的罪行中，公义还有什么位置？我仅从众多中选取了少数，并非为了指责，而是为了显示其本性。凡欲知全貌者，可取阅\Person{塞内加}的书卷，他同时是最真实的记述者和最猛烈的公共道德与恶习的指控者。但\Person{卢基里乌斯}也以简练的诗句描绘了那黑暗的生活：\Quote{但如今从早到晚，无论节日或平日，全体人民和元老们一致聚集在广场，从不离去。他们都投身于同一个追求和技艺：谨慎地欺骗、诡诈地争斗、阿谀地竞争，各人假装自己是好人，相互埋伏，仿佛所有人都是所有人的敌人。}但这些事中，哪一件能归咎于我们的人民？对他们的整个宗教在于无罪无玷地生活。既然他们看到自己和他们的人民都做我们所说的事，而我们的人只行正义与良善之事，他们若有任何理解，便能由此明白：行善的人是虔敬的，而行恶的人是渎神的。因为那在一生所有行为上犯错的人，不可能在最重要的事，即宗教上犯错，宗教是万有之首。因为不虔敬若在最重大的事上被接受，就会贯穿其余。因此，一生全然犯错的人，不可能不在宗教上也被欺骗；因为虔敬若在首要之事上持守其准则，就会在其他事上保持其进程。于是发生这样的事：每一方，其首要主题的特性可从所行行为的状况得知。\par

% structure: p0059 source=h2 target=subsection reason=relative-heading-rank; base=h2

\subsection{第十章　论虚假的虔敬，以及论虚假与真实的宗教。}

值得考察他们的虔敬，好从他们仁慈和虔敬的行为中，理解那些违反虔敬法律的事具有何种特性。为免我显得严厉攻击某人，我将从诗人中选取一个角色，一个最伟大的虔敬榜样。在\Person{玛罗}笔下，那位君王\par

\Quote{比他更勇敢、\newline{}更虔敬、\newline{}更真实的人，从未呼吸过生命的气息。}\par

他向我们提出了什么公义的证明呢？\par

\Quote{那里有受害的俘虏，\newline{}双手反绑于背后，\newline{}注定要献祭于火焰之上。}\par

有什么比这虔敬更仁慈呢？有什么比向死人献上人祭，以人血如油般喂养火焰更仁慈呢？但也许这并不是英雄本身的过错，而是诗人的过错，他以显著的恶行玷污了\Quote{一位以虔敬著称的人}。诗人啊，那么你屡屡赞美的虔敬在哪里呢？看那虔敬的\Person{埃涅阿斯}：——\par

\Quote{他活捉了四个来自\Place{苏尔莫}种的不幸青年，\newline{}和四个以\Person{乌芬斯}为父的青年，\newline{}判他们流血以祭\Person{帕拉斯}的亡魂，\newline{}在\Person{帕拉斯}的火葬堆上。}\par

因此，为什么就在他送那些人捆绑去屠杀时，他说：\par

\Quote{我但愿赐予生者和平。}\par

当他下令将他权力范围内还活着的人像牲畜一样处死时，难道这就是 \OriginalTerm{虔敬}{piety} 吗？但正如我所说，这并非他的过错——因为他或许未受过良好教育——而是你们的过错；因为你们虽有学识，却不知虔敬的本性，竟认为他那邪恶可憎的行为是虔敬的合宜表现。他之所以被称为虔敬，仅仅是因为他爱他的父亲。我为何要说\par

\Quote{善良的\Person{埃涅阿斯}应允了他们的恳求，}\par

却仍将他们杀死？因为，虽以同一父亲的名义起誓，且以\par

\Quote{年轻\Person{尤卢斯}初升的朝阳，}\par

他仍不饶恕他们，\par

\Quote{狂怒燃遍每根血管}\par

怎么！有人能想象他有什么美德吗？他被疯狂如枯草般点燃，忘记父亲的亡灵向他恳求，无法抑制愤怒。因此他绝不虔敬，他不仅杀戮无力抵抗者，甚至杀戮求饶者。此时有人会说：那么 \OriginalTerm{虔敬}{piety} 是什么？在哪里？有何特征？诚然，它存在于那些不知战争、与众人和睦、甚至对敌人友善、视众人如兄弟、懂得克制愤怒、以平静的治理安抚一切情感冲动的人中。因此，一片多么深厚的迷雾，多么巨大的黑暗与错误的乌云笼罩了人的心胸！他们自以为特别虔敬时，却变得尤其不虔敬。因为他们越是虔敬地敬拜那些尘世的造像，就越是邪恶地对待真天主的名号。因此他们常遭受更大的灾祸，作为其不虔敬的报应；又因不知这些灾祸的缘由，便全然归咎于命运，给\Person{伊壁鸠鲁}的哲学留下余地——他认为诸神不受任何事物影响，既不为恩惠所动，也不为愤怒所激，因为他们常见藐视者幸福而敬拜者遭难。发生此事的原因在于：当他们看似虔敬而本性良善时，便认为自己不应遭受所常忍受的那种苦难。然而他们以指责命运来安慰自己，却未察觉若命运存在，她绝不会伤害自己的敬拜者。因此，这种虔敬理应受到惩罚；天主因人们敬拜歪曲宗教而犯下的邪恶而恼怒，便以沉重的灾祸惩罚他们。这些人虽以最大的 \OriginalTerm{信德}{faith} 和纯洁度着圣洁的生活，却因敬拜那些亵渎而渎神的仪式为真天主所憎恶的神祇，而与 \OriginalTerm{正义}{justice} 及真正虔敬之名隔绝。不难指出为何拜偶像者不能成为善良正义的人。他们敬拜嗜血的神祇\Person{玛尔斯}和\Person{贝娄娜}，又如何能戒除流血？他们敬拜驱逐父亲的\Person{朱庇特}，又如何能孝敬父母？他们敬拜\Person{萨图尔努斯}，又如何能怜惜自己的婴孩？他们敬拜裸体且奸淫、仿佛在诸神中卖淫的女神，又如何能持守贞洁？他们熟知\Person{墨丘利}的偷窃，他教导欺骗并非欺诈而是机敏，又如何能远离掠夺与欺诈？他们敬拜\Person{朱庇特}、\Person{赫拉克勒斯}、\Person{利伯}、\Person{阿波罗}及其他诸神，其与男女的奸淫放荡不仅为学者所知，甚至在戏院中上演并编成歌曲，以致人人皆知，他们又如何能克制情欲？在这些事中，人还能正义吗？即使他们本性良善，也会被这些神祇本身训练成不义。因为要取悦你所敬拜的神，就需要那些你知道他喜悦和满足的事物。于是，神按照自己意愿的特征塑造敬拜者的生活，因为最虔敬的敬拜就是效仿。\par

% structure: p0075 source=h2 target=subsection reason=relative-heading-rank; base=h2

\subsection{第十一章　论异教徒对基督徒的残忍}

因此，因为 \OriginalTerm{正义}{justice} 是沉重而不愉快的对于那些性情与他们的神祇相符的人，他们以在其他事上表现出的同样不虔敬，暴力地对待义人。先知们称他们为野兽，并非无因。因此，\Person{马库斯·图利乌斯}说得极好：\Quote{因为若无人宁愿死也不愿变成野兽的形状，尽管他将拥有人的心智，那么人具有人的形状却拥有兽性的心智，岂非更可悲？在我看来，这比前者更糟，正如心智比身体更卓越。}因此他们蔑视野兽的身体，而自己却比野兽更残忍；他们以此自豪：生而为人，却除了容貌和崇高的形体之外，毫无人性可言。因为\Place{高加索}、\Place{印度}、\Place{希尔卡尼亚}曾养育过如此野蛮和嗜血的野兽吗？所有野兽的狂怒只持续到食欲满足为止；饥饿一平息，立即平静下来。那真正是野兽，仅凭其命令\par

\Quote{杀戮的溪流汩汩，\newline{}严阵的战场血腥弥漫。}\newline{} \Quote{可怕的痛苦，狂暴的恐惧蜂拥，\newline{}死亡以多种形态狰狞怒视。}\par

无人能恰当地描述这野兽的残忍：它卧在一处，却以铁的利齿肆虐全世界，不仅撕裂人的四肢，还粉碎他们的骨头，甚至在他们骨灰上肆虐，使葬身之地无处可寻，仿佛那些宣认天主的人所求的乃是坟墓受人瞻仰，而非自己得以到达天主的面前。\par

何等野蛮！何等愤怒！何等疯狂——拒绝给活人光明，给死人土地？\newline{}我因此断言，没有什么比那些被需要所利用或造就为他人暴怒的执行者、不敬命令的帮凶的人更加悲惨。因为那并非荣誉或地位的提升，而是将人判处酷刑，并承受天主的永罚。\newline{}然而，无法描述他们个人在全世界所作所为。有多少卷帙才能容纳如此无限、如此多样的残忍？因为一旦获得权力，每个人都按其本性恣意妄为。有些人出于极度胆怯，所作所为远超命令；另一些人则出于对义人个人的仇恨；有些人出于天生的心灵残忍；有些人出于取悦上司的欲望，并借此服务为更高的官职铺路。有些人嗜杀成性，例如在\Place{弗里吉亚}有一人将整群集会者连同他们的聚会场所一同焚烧。但他越残忍，反倒显得越仁慈。\newline{}然而，最糟糕的迫害者是那些被虚假的仁慈外表所迷惑的人；那决定不处死任何人的法官，才是更严厉、更残忍的施刑者。因此，无法言说这类法官为了达成目的而发明了多少巨大而痛苦折磨的方式。\newline{}但这样做不仅是为了吹嘘自己未曾杀害无辜——因为我亲耳听到有人吹嘘其治理在这方面是\Quote{不流血的}——而且出于嫉妒，免得自己被打败，或他人获得其德行应得的荣耀。如此，在发明刑罚方式时，他们除了胜利之外别无所念。因为他们知道这是一场竞赛和战斗。我在\Place{比提尼亚}见到总督欣喜若狂，仿佛征服了某个蛮族，因为一个曾以极大精神抵抗两年的人终于屈服了。\newline{}因此，他们争斗以求胜利，并对他们的身体施加精妙的痛苦，唯一避免的就是受害者死于酷刑之下；仿佛唯有死亡能使人幸福，仿佛越严酷的酷刑反而会带来更大的德行荣耀？但他们以固执的愚昧命令对受刑者予以细心照料，使其肢体得以恢复以承受其他酷刑，并供给新鲜的血液以供惩罚。有什么能如此虔诚、如此仁慈、如此人道？他们对所爱之人也不会付出如此精心照料。这是诸神的训导；他们以此训练其崇拜者；这是他们所要求的神圣仪式。\newline{}此外，最邪恶的杀人犯针对虔诚者制定了不敬的法律。因为渎神的法令和不公的法学家论辩也被诵读。\Person{多米提乌斯}在其第七卷\WorkTitle{论行省总督的职责}中收集了君主们邪恶的谕旨，以表明那些承认自己是天主的崇拜者的人应受何种惩罚。\par

% structure: p0080 source=h2 target=subsection reason=relative-heading-rank; base=h2

\subsection{第十二章 真正的德行；好公民与坏公民的评价}

对于那些把古代暴君折磨无辜者称为正义，并且虽是邪恶与残忍的教师却想显得公正明智的人，你们将如何对待？你们盲目、迟钝、不谙世事与真理。正义对你们如此可恨，喔堕落的心灵，以至于你们视之与最重大的罪行等同？天真在你们眼中如此全然失落，以至于你们不仅认为它不配死亡，而且视不犯罪、怀抱纯洁无染之胸为超越一切罪行？\newline{}既然我们普遍与那些崇拜诸神的人说话，请准许我们与你们一同行善；因为这是我们的律法，我们的事业，我们的宗教。若我们在你们眼中显得智慧，就效法我们；若显得愚蠢，就轻视我们，甚或嘲笑我们，随你们意；因为我们的愚蠢对我们有益。你们为何撕裂我们，为何加苦于我们？我们不嫉妒你们的智慧。我们偏爱这我们的愚蠢——我们拥抱它。我们相信这对我们有益——爱你们，并施予一切于你们这些恨我们的人。\par

在\Person{西塞罗}的著作中，有一段与真理并不矛盾的文字，出自\Person{弗里乌斯}反对正义的辩论：\Quote{我问，}他说，\Quote{如果有两个人，其中一个是卓越之人，具有最高的正直、最大的正义和非凡的信实，而另一个则以犯罪和胆大妄为著称；如果国家犯下如此错误，竟视那善人为邪恶、败坏、可憎，却认为那最恶之人有最高的正直和信实；并且，根据全体公民的这种看法，那善人遭受骚扰、被拖走、被砍去双手、挖去双眼、被定罪、被捆绑、被烙印、被流放、陷于匮乏，最后，在所有人看来，他理应是最悲惨的；但另一方面，如果那恶人被所有人称赞、尊敬、爱戴——如果所有荣誉、所有权力、所有财富、所有丰裕都赐予他——简言之，如果他在众人眼中被判定为卓越之人，最配得一切好运——那么，我请问，谁会疯狂到怀疑自己宁愿选择两人中的哪一个？}可以肯定，他提出这个例子，仿佛预见到了为了公义我们将遭遇何等灾难，以及以何种方式；因为我们的人民因那些谬误之人的乖谬而遭受这一切。看哪，国家，甚或整个世界本身，都陷于如此错误，竟迫害、折磨、定罪、处死善良公义之人，如同他们是邪恶不敬之人。至于他所说，没有人如此糊涂以至于怀疑自己宁愿选择哪一个，他，作为反对正义的一方，认为智者宁愿在有好名声的情况下作恶，也不愿在有坏名声的情况下为善。\par

但愿这种愚蠢远离我们，让我们不至于宁愿选择虚假而非真实？或者说，我们的善人之品质，更取决于民众的错误，而非我们自己的良心和天主的审判？或者有任何繁荣会引诱我们，使我们宁愿选择带着一切邪恶的真善，而不愿选择带着一切繁荣的假善？让国王保留他们的王国，富人保留他们的财富，正如\Person{普劳图斯}所说，智者保留他们的智慧；让他们把我们的愚拙留给我们，而我们的愚拙显然被证明是智慧，正因为他们嫉妒我们拥有它：谁会嫉妒一个愚拙的人呢？除非他自己是极其愚拙的？但他们并不愚拙到嫉妒愚拙人；但因为他们如此关切和焦虑地追随我们，他们承认我们不是愚拙的。因为他们为何如此残忍地发怒，除非是害怕，随着正义日复一日增强，他们会连同他们衰败的众神一起被抛弃？因此，如果崇拜众神的人是明智的，而我们是愚拙的，他们为何害怕明智的人会被愚拙的人引诱呢？\par

% structure: p0084 source=h2 target=subsection reason=relative-heading-rank; base=h2

\subsection{第十三章 论基督徒的增长与惩罚}

既然因着外教人的皈依，我们的人数不断增加，而从未减少——即使在迫害中也是如此——因为人可能犯罪，被祭献所玷污，却不能离弃天主；因为真理以其自身的力量得胜——请问，有谁如此愚蠢和盲目，看不出智慧在哪一边呢？但他们被恶意和愤怒所蒙蔽，以致看不见；他们以为那些能够避免惩罚却宁愿受折磨和被杀的人是愚蠢的；然而，他们本可以从这件事看出，全世界成千上万的人同心同意，这并非愚蠢。因为如果妇女因性别的软弱而陷入错误（这些人有时称这为妇孺的迷信），那么男子无疑是有智慧的。如果男孩、青年因年龄而轻率，那么成年人和老年人无疑有固定的判断。如果一座城市无知，那么其他无数的城市显然不能是愚蠢的。如果一个省或一个民族缺乏谨慎，那么其余的必定理解正确的事。但是，既然神圣的法律从日出之地到日落之地被接受，每一个性别、每一个年龄、每一个民族和国家都同心同意服从天主——到处都是同样的忍耐，同样的轻看死亡——他们本应明白其中必有道理，它并非毫无原因地被维护至死，它有着基础和坚固性，不仅使那宗教免于伤害和骚扰，而且总是增长并变得更强大。因为在这方面，那些人的恶意也暴露出来，他们以为如果玷污了人，就彻底推翻了天主的宗教，而人还可以向天主赎罪；并且没有一个如此邪恶的天主崇拜者，当有机会时，不会回来平息天主的愤怒，并且更加虔诚。因为罪的意识和惩罚的恐惧使人更虔诚，而通过忏悔恢复的信德总是更坚强。因此，如果他们自己，当他们以为众神对他们发怒时，仍然相信众神可以被礼物、祭献和香料平息，那么他们有什么理由认为我们的天主如此不仁慈和不可和解，以致于一个被迫违反自己意愿向他们的神奠酒的人，似乎不可能成为基督徒呢？除非他们偶然以为那些一度被污染的人将要改变心意，以致现在开始自愿做他们曾在酷刑下所做的事。谁会愿意自愿承担那始于伤害的职责呢？谁看到自己肋旁的伤疤，不会更恨那些神，为他们的缘故他带着持久惩罚的痕迹和刻在肉上的印记呢？因此，当来自天上的和平赐予时，那些疏远我们的人回来，而由于所显示的美德的奇妙性质，又有一批新的人加入。因为当人们看到人被各种酷刑撕裂，并在行刑者疲倦时仍保持不屈的忍耐，他们就会想，正如实际情况那样，这么多人的同意和垂死者的坚定并非无意义，而忍耐本身若无天主的帮助，不可能克服如此巨大的酷刑。强盗和强壮的男人无法忍受这样的撕裂：他们呼叫、呻吟；因为他们被痛苦征服，因为他们缺乏灌注在他们里面的忍耐。但在我们当中（且不说男人），男孩和柔弱的妇女沉默地制服了行刑者，甚至火也不能从他们那里逼出一声呻吟。让罗马人去夸耀他们的\Person{穆提乌斯}或\Person{雷古鲁斯}吧——前者将自己交给敌人杀死，因为他羞于作为俘虏活着；后者被敌人俘虏，当他看到无法逃脱死亡时，将手放在燃烧的炉灶上，以向他想杀死的敌人赎罪，并因那惩罚得到了他不应得的宽恕。看，软弱的性别和脆弱的年龄忍受全身被撕裂和被焚烧：并非出于必要，因为如果他们愿意，可以逃脱；而是出于自愿，因为他们信赖天主。\par

% structure: p0086 source=h2 target=subsection reason=relative-heading-rank; base=h2

\subsection{第十四章 论基督徒的刚毅}

但这就是真正的德性，那些夸口的哲学家们也以它自夸，不是通过行动，而是空谈，说什么没有什么比能不被任何刑讯者赶离自己的见解和目的更适合智者的庄重和坚定，而是值得忍受酷刑和死亡，而不出卖信任或离弃职守，或不因怕死或剧痛而犯任何不义。除非偶然，\Person{弗拉库斯}在他们的眼中是胡言乱语，当他说：\par

\Quote{不为那命令邪恶事物的百万群众的愤怒\newline{}不为那暴君眉宇间近在咫尺的厄运\newline{}动摇那正直而坚定的人\newline{}在他灵魂的坚固完全中。}\par

再没有比这更真实的了，如果这是指那些拒绝任何酷刑、任何死亡方式，以免背离信德和正义的人；他们不惧怕暴君的命令，也不畏惧统治者的刀剑，以便以心灵的恒毅保持真实而坚固的自由——智慧仅在于此。因为谁如此傲慢，谁如此高傲，竟然禁止我举目向天？谁能强迫我，要么崇拜我所不愿崇拜的，要么放弃对我所愿崇拜的敬礼？如果连这必须出于自己意愿的事竟被他人任性强夺，那我们还剩下什么？只要我们有勇气藐视死亡和痛苦，没有人能得逞。但如果我们拥有这种恒毅，为什么当我们做哲学家所称赞的事时，却被视为愚妄？塞内加在指责人们言行不一时，正确地说：最高德性在他们看来在于精神的伟大；然而同一些人又把藐视死亡的人视为疯子，这显然是极大的乖谬。但这些追随虚妄宗教的人，以同样的愚昧行事，也不认识真天主；厄立特里亚女巫称他们为\Quote{聋聩的}，因为他们既听不见也不领会神性之事，却惧怕并崇拜自己手指塑造的泥像。\par

% structure: p0090 source=h2 target=subsection reason=relative-heading-rank; base=h2

\subsection{第十五章  论愚妄、智慧、孝爱、公平与正义。}

但他们视智者为愚妄的理由有其充分的依据（因为他们并非无端受骗）。我们必须尽力阐明这一点，使他们或许（若可能）认识到自己的错误。正义按其本性具有某种愚妄的外表，我能既以神圣的见证又以人的见证证实这一点。但我们或许无法说服他们，除非我们从他们自己的权威中教导他们：没有人能是正义的——这事与真智慧相联——除非他也显得愚妄。卡尔内亚德是学院派哲学家；不知道他在论辩、口才和洞察力方面有多大能力的人，也会从西塞罗或卢基利乌斯的赞词中了解此人的品格；在卢基利乌斯的著作中，尼普顿讨论一个极其困难的问题时，表明即使奥耳库斯使卡尔内亚德本人复活，也无法解释。这位卡尔内亚德曾受雅典人派遣作为使节前往罗马，在加尔巴和曾任监察官的卡托面前——两人当时是最伟大的演说家——详尽地谈论了正义的主题。但次日，此人又以相反的论辩推翻了自己的论证，夺去了他前日所称赞的正义；这并非以哲学家的庄重（其审慎应坚定，意见应确定），而是以或然式的、正反两方论证的演说练习。他惯于这样做，以便能反驳那些断言某事的人。卢基乌斯·福利乌斯在西塞罗著作中提及那场推翻正义的讨论。我认为，由于他讨论的是国家事务，他这样做是为了引入对正义的辩护和赞美；没有正义，他认为国家无法治理。但卡尔内亚德为了反驳主张正义的亚里士多德和柏拉图，在那第一次辩论中收集了所有为正义辩护的论点，以便能够推翻它们——正如他所做的。因为动摇没有根基的正义很容易，因为当时世上并无正义，以使哲学家能认识其性质或特性。我但愿如此众多且如此杰出的人们，若其口才和精神也足以完成对这最伟大德性的辩护——这德性源于宗教，其原则在于公平！但那些不认识第一部分的人，也无法拥有第二部分。但我希望首先简要地说明正义是什么，以便理解哲学家们对正义无知，且无法辩护他们所不认识的。虽然正义包含所有德性，但其中有两个最主要且无法与正义分离：孝爱与公平。因为忠实、节制、正直、无邪、廉洁以及其他此类，或因天性或因父母教养，可能存在于那些不认识正义的人身上，正如它们一直存在；因为古罗马人惯于以正义为荣，显然是以那些（如我所说）可能源于正义并与源头本身分离的德性为荣。但孝爱与公平好比正义的血脉：因为正义的全部包含在这两个泉源中；其源头和起源在于前者，其全部力量和方式在于后者。孝爱无非是对天主的观念，正如特里斯墨吉斯忒斯最真实地定义的，我们在别处已说过。因此，若孝爱是认识天主，而此认识的总和是崇拜祂，那么显然，不认识天主的人就不认识正义。因为不认识其源头的人，如何能认识正义本身？柏拉图确实谈论了许多关于独一的天主，他说天主创造了世界；但他没有谈论宗教：因为他梦见了天主，却不认识他。但如果他本人或任何其他人想要完成对正义的辩护，他首先应当推翻诸神的宗教，因为它们与孝爱对立。而因为苏格拉底确实尝试这样做，他被投入监狱；即使在那时，人们也能看到，那些开始捍卫真正正义并事奉独一天主的人将会遭遇什么。\par

义的另一部分因而是公平；显然，我所说的不是审判公正的公平（尽管这在义人中也值得称赞），而是使人自己与他人平等，\Person{西塞罗}称之为均等。因为天主产生和赐予人气息，祂愿意所有人都平等，即相匹配。祂给所有人规定了相同的生活条件；祂使所有人走向智慧；祂应许永生给所有人；没有人被排除在祂天上的恩惠之外。因为正如祂同样地分布祂同一的光，赐予所有人祂的泉水，供给食物，并给予最愉快的休息睡眠；同样，祂赐予所有人公平和德行。在祂眼中，没有奴隶，没有主人；因为如果所有人都有同一父亲，我们就以同等的权利都是子女。在天主眼中，没有人是贫穷的，除非他没有正义；没有人是富有的，除非他充满德行；简而言之，没有人是卓越的，除非他善良无辜；没有人是最著名的，除非他大量地行慈悲的工作；没有人是最完美的，除非他填满了德行的所有阶梯。因此，罗马人或希腊人不能拥有正义，因为他们有在诸多程度上彼此不同的人，从穷人到富人，从卑微者到有权势者；简而言之，从平民到君王的最高的权威。因为凡不是所有都相匹配的地方，就没有公平；而不平等本身就排除正义，正义的全部力量在于此：它使那些以平等命运到达今生状况的人平等。\par

% structure: p0093 source=h2 target=subsection reason=relative-heading-rank; base=h2

\subsection{第十六章 论义人的职责与基督徒的公平}

因此，既然正义的那两个泉源被改变了，一切德行和一切真理都被撤去了，正义本身返回了天上。所以真正的善没有被哲学家们发现，因为他们既不知道它的起源也不知道它的效果：这除了向我们的人民，没有向其他人启示过。有人会说：在你们中间不是有人贫穷、有人富有；有人是仆人、有人是主人吗？个人之间难道没有一些区别吗？没有；也没有任何其他原因使我们彼此互称兄弟，除了我们相信自己是平等的。因为既然我们不以身体，而以精神来衡量一切人类事物，虽然身体的条件不同，但我们在精神上、在宗教上把他们视为同仆，没有仆人。财富也不能使人显赫，除非它们能使人因善行更加显著。因为人富有，不是因为他们拥有财富，而是因为他们将财富用于正义的工作；而那些看似贫穷的人，正因为此是富有的，因为他们不缺乏什么，也不渴望什么。\par

因此，虽然在心灵的谦卑中我们是平等的，自由人与奴隶，富人与穷人，然而在天主眼中我们因德行而被区分。每个人都越正义就越被高举。因为如果正义是使人将自己置于甚至与更低等级的人平等的地位，尽管他在此本身就卓越，即他使自己与下级平等；然而，如果他不仅作为平等者，甚至作为下级行事，那么他在天主的审判中显然会获得更高得多的尊严等级。因为的确，既然此世生命中的一切事物都是脆弱和易于腐朽的，人们就偏爱自己胜过他人，并为尊严而争辩；没有比这更污秽、更傲慢、更远离智者的行为了：因为这些属地的事物完全与属天的事物相对立。因为正如人的智慧在天主面前是最大的愚昧，而愚昧（正如我所示）是最大的智慧；同样，在世上显赫和高升的人，在天主面前是卑微和低贱的。因为，且不提这些当前属地的、被赋予极大荣誉的善与德行相反，并削弱心灵的力量，我请问，什么高贵能如此稳固，什么资源，什么权力，既然天主能使君王自己变得甚至比最低贱的人还要低贱？因此天主为我们的利益而在神圣的规诫中特别放置了这条：\BibleQuote{因为凡高举自己的，必被贬抑；凡谦抑自己的，必被高举。} 这条规诫的健全性教导我们，凡单纯地把自己放在与他人平等地位并谦卑行事的人，在天主眼中被视为卓越和显赫。因为在\Person{欧里庇得斯}那里提出的这句名言并非虚假：—\par

\Quote{在世间被视为邪恶的事，在天上却被认为是善的。}\par

% structure: p0097 source=h2 target=subsection reason=relative-heading-rank; base=h2

\subsection{第十七章 论基督徒的公平、智慧与愚昧}

我已经解释了哲学家们既不能发现义德也不能捍卫义德的原因。现在我要回到我原先计划的内容。因此，\Person{卡内亚德}既然看到哲学家的论证薄弱，便承担了反驳它们的勇敢任务，因为他明白这些论证是可以被反驳的。他的辩论要旨如下：\Quote{人们为自己制定法律，着眼于自身利益，确实因各自品性而异，且在同样的人身上常因时代而改变；但并没有自然法；所有人以及其他动物，都受自然的引导而趋向自身利益；因此没有义德，即便有，也是最大的愚蠢，因为它通过促进他人利益而损害自身。}他又提出这些论证：\Quote{所有拥有统治权的国家，甚至统治全世界的罗马人自身，如果他们希望成为正义的，即归还他人的财产，就必须回到茅舍，在贫困和苦难中安身。}然后，他离开一般性话题，转向具体事例。\Quote{如果有一个好人，}他说，\Quote{有一个逃跑的奴隶，或者一所不卫生且有害健康的房屋，唯独他知道这些缺陷，并因此将其出售，他会说明这个奴隶是逃跑的，以及他所出售的房屋是有害的，还是会向买主隐瞒？如果他说明，他确实是好的，因为他不会欺骗；但他仍会被视为愚蠢，因为他要么低价出售，要么根本卖不出去。如果他隐瞒，他确实是明智的，因为他会维护自身利益；但他也是邪恶的，因为他会欺骗。再如，如果他遇到某人以为是在卖黄铜矿实则是金子，或以为是铅实则是银子，他会保持沉默以便低价买入，还是会告知实情以便高价买入？显然，宁愿高价购买是愚蠢的。}由此他希望人们明白：既正义又好的人是愚蠢的，而明智的人是邪恶的；然而，人们或许可以在不遭毁灭的情况下满足于贫困。因此，他转向更大的事情，在其中没有人能在不危及生命的情况下保持正义。因为他说：\Quote{当然，正义就是不杀人、不夺取他人的财产。那么，如果一个正义的人遭遇海难，而一个比他更弱的人抢到了一块木板，他会不会把那个人推下木板，以便自己爬上去，借此逃生，尤其是在大海中央没有见证者的情况下？如果他是明智的，他就会这样做；因为除非他如此行动，否则他自己必死。但如果他宁愿死也不愿对他人施加暴力，在这种情况下，他是正义的，但也是愚蠢的，因为他在挽救他人生命的同时却不珍惜自己的生命。同样，如果他自己军队被打败，敌人开始追击，那个正义的人遇到一个骑在马上受了伤的人，他会饶恕他而让自己被杀，还是会把他从马上推下来，以便自己逃脱敌人？如果他这样做，他是明智的，但也是邪恶的；如果他不这样做，他是正义的，但也必然是愚蠢的。}因此，当他这样把正义分为两部分，说一部分是公民的，另一部分是自然的时，他颠覆了二者：因为公民部分是智慧，而非义德；而自然部分是义德，而非智慧。这些论证非常精巧而敏锐，是\Person{马尔库斯·图利乌斯}无法反驳的。因为当他让\Person{莱利乌斯}回应\Person{弗里乌斯}，并为义德辩护时，他避开了这些论证，犹如避开陷阱而不加反驳；以至于同一位\Person{莱利乌斯}似乎没有为自然义德辩护——它已陷入愚蠢的指控——而是为\Person{弗里乌斯}承认是智慧但非正义的公民义德辩护。\par

% structure: p0099 source=h2 target=subsection reason=relative-heading-rank; base=h2

\subsection{第十八章 论义德、智慧与愚昧}

关于我们当前的讨论，我已经表明了正义如何带有愚昧的外表，以便显明那些认为我们宗教的人在做他(\Person{卡内亚德})所提出的事情时显得愚昧，并非毫无理由。现在我感到需要一项更重大的任务，即首先答复\Person{弗里乌斯}，因为\Person{莱利乌斯}并未充分答复他；在答复之后，我将说明为什么天主愿意将正义隐藏在愚昧的外表下，并从人的眼前除去。\Person{莱利乌斯}虽被称为智者，却无法成为真正正义的代言人，因为他没有掌握正义的根源和源头。但这个辩护对我们来说更容易，因为我们藉天主的恩赐，对正义既熟悉又了解，不仅知道其名，更认识其实。因为\Person{柏拉图}与\Person{亚里斯多德}曾诚心诚意地想要捍卫正义，如果他们善良的努力、他们的口才和智力活力也曾得到对神圣事物的认识之助，他们必会有所成就。然而，他们的工作徒劳无益，被人忽视；他们也未能说服任何人按照他们的规条生活，因为那体系没有来自天上的根基。但我们的工作必须更加确定，因为我们是由天主教导的。因为他们只是用言语描绘正义，当正义不在眼前时为其画像；他们也无法用眼前的实例来证实他们的论断。因为听众可能会回答说，他们不可能像他们在论辩中所规定的那样生活；所以至今没有一个人遵循那种生活方式。但我们不仅用言语，还以来自真理的实例来证明我们言论的真实性。因此，\Person{卡内亚德}明白了正义的本性，只是他没有充分认识到那不是愚昧；虽然我似乎理解他这样做的意图。因为他并非真的认为正义的人是愚昧的；但他知道事实并非如此，却不明白为何看起来如此的原因，他想要表明真理是隐藏的，以便维持他自己学派的教义，其首要观点是\Quote{没有什么能被完全理解}。\par

所以，让我们看看正义是否与愚昧有任何一致之处。他说，正义的人，如果他不从受伤者手中夺走马匹，不从遇难者手中夺走木板以保全自己的性命，他就是愚昧的。首先，我否认一个真正正义的人会处于这种情形；因为正义的人不与任何人为敌，也不渴望任何属于他人的东西。因为他为何要航海，或从异地寻求什么，当他自己的土地足以供给他的时候？或者他为何要发动战争，将自己卷入他人的激情中，当他的心灵与人类保持永久和平时？无疑，他会对外国商品或人类之血感到喜悦吗？他不懂得追求利益，满足于自己的生活方式，并且认为不仅自己杀人是不合法的，就连与杀人者在一起或观看杀人也是不合法的！但我且不谈这些，因为人可能被迫甚至违背自己的意愿去经历这些事。那么，\Person{弗里乌斯}——或者不如说\Person{卡内亚德}，因为这番话都是他的——你认为正义如此无用、多余，且被天主如此鄙弃，以至于它自身毫无力量和能力来维持自己的生存吗？但显然，那些不明白人的奥秘、因此将一切事物都归于今生的人，无法知道正义的力量有多么大。因为当他们讨论德性时，虽然他们明白德性充满劳苦和苦难，但他们说德性应当为其本身而追求；因为他们绝没有看到其酬报，那是永恒不朽的。因此，把一切都归于今生，他们便把德性完全变成了愚昧，因为它枉然且无谓地承受今世如此巨大的劳苦。关于这一点，另有机会再谈。\par

与此同时，让我们按开始的方式谈论正义；正义的力量如此之大，以至于当它举目向天时，它应从天主那里得到一切。因此\Person{弗拉库斯}正确地说过，无辜的力量如此之大，无论它走到哪里，都不需要武器或力量来保护自己：——\par

\Quote{生活纯洁无瑕，不存诡诈的人，无需借用摩尔的弓或箭，亲爱的\Person{富斯库斯}！他自卫不靠装满毒箭的箭囊。即使他的道路经过炎热的非洲锡尔特斯，或高加索峡谷，那里无客可栖身，或那神话般的河流希达斯普斯懒洋洋地流过的河岸。\newline{}\newline{}\newline{}\newline{}\newline{}\newline{}\newline{}}\par

所以，在风暴和战争的危险中，正义的人不可能没有天主的护佑；即使他与弑亲者和有罪的人一同航海，恶人也不得幸免，以便这位唯一公正无辜的灵魂得免危险，或至少唯独他得蒙保存，而其余的人丧亡。但让我们假定这位哲学家提出的情形是可能的：那么，正义的人会做什么，如果他遇到一个受伤者骑在马上，或一个遇难者抱着木板？我不愿意承认他会宁愿死，而不是置他人于死地。正义是人的至善，它不会因此被称为愚昧。因为什么比纯洁无辜对人更好、更珍贵？而纯洁无辜越是被推至极点，就越完美，宁愿死也不愿减损纯洁无辜的名声。他说，在涉及自身生命毁灭的情况下保全他人生命是愚昧。那么，你认为为友谊而死也是愚昧吗？\par

为什么，那些毕达哥拉斯学派的朋友们被你们赞美？其中一人把自己交给暴君作为另一人生命的担保，另一人在约定的时刻，当他的担保人正被带去处决时，他出现并通过自己的干预解救了他？如果他们是愚蠢的，那么他们的德性就不会享有如此光荣——当一人愿意为朋友而死，另一人甚至为他所承诺的诺言而死。最后，正因为这种德性，暴君通过保全两人而奖赏了他们，就这样，一个最残忍之人的性情被改变了。而且，甚至据说他恳求他们接纳他作为他们友谊的第三者，由此可知，他并不认为他们是愚蠢的，而是善良和明智的。因此，我不明白，既然为友谊和诺言而死被认为是最高光荣，那么一个人为自己的清白而死为什么就不光荣呢？所以，那些指责我们愿意为天主而死是罪恶的人是最愚蠢的，因为他们自己以最高的赞美把那些愿意为一个人而死的人捧上天。简而言之，为了结束这个讨论，理性本身教导我们，一个人不可能既正义又愚蠢，既明智又不义。因为愚蠢之人不认识正义和良善，因此总是犯错。因为他仿佛被自己的恶习所俘虏；他无法以任何方式抵抗它们，因为他缺乏他所不认识的德性。但正义之人远离一切过错，因为他不能做别的，尽管他有对与错的知识。\par

但是，除了明智之人，谁能分辨善恶呢？因此，这就导致，愚蠢之人永不可能正义，不义之人永不可能明智。如果这是最真实的，那么很明显，那个没有从遇船难者手中拿走木板、或从受伤者手中拿走马匹的人，并不是愚蠢的；因为做这些事是罪，而明智之人远离罪。然而，我自己也承认，这看起来如此，是由于人的错误，他们不知道一切事物的特性。因此，整个这个问题与其说是通过论证，不如说是通过定义而被驳斥。因此，愚蠢是因不知善恶而在言行上犯错。因此，那个为了不伤害他人（这是一种恶）甚至不宽恕自己的人，并不是愚蠢的。而这正是理性和真理本身所教导的。因为我们看到，在所有动物中，由于它们缺乏智慧，自然为自身提供所需。因此，它们伤害他人以使自己获益，因为它们不明白行伤害是恶。但人，拥有善恶的知识，甚至在有损于自己的情况下也不伤害他人，这是无理性的动物无法做到的；因此，无辜被视为人的首要德性之一。由此看来，最明智的人是宁愿毁灭也不愿行伤害的人，以保存那使他与无言的受造物区别开来的义务感。因为那些为了自己的利益或好处，不指出出售黄金者的错误以便低价购买，或不承认自己在出售一个逃跑的奴隶或有病害的房子的人，并不是明智的，如\Person{卡尔内亚德}希望其显现的那样，而是狡猾和诡诈的。而狡猾和诡诈也存在于无言的动物中：无论是它们潜伏等待其他动物，用欺骗捕获它们以便吞食，还是它们以各种方式避开其他动物的陷阱。但智慧只属于人。因为智慧是理解，其目的是行善和正义，或避免不当的言行。明智之人从不致力于追求利益，因为他轻视这些属世的优点；他也不允许任何人被欺骗，因为纠正人的错误并引导他们归回正道是善人的义务；因为人的本性是社会性和慈善的，唯有在这方面，他与天主相似。\par

% structure: p0107 source=h2 target=subsection reason=relative-heading-rank; base=h2

\subsection{第十九章 论德性与基督徒的折磨，以及父亲和主人的权利}

但毫无疑问，这就是为什么那个宁愿贫困或死亡也不愿伤害他人或夺取他人财产的人显得愚蠢的原因——因为他们认为人因死亡而毁灭。从这一信念产生了平民和哲学家的一切错误。因为如果我们死后不存在，那么最愚蠢的人就是不促进现世生活的利益，使其长久并充满各种益处。但若有人如此行，他必然要偏离正义的规则。但如果人还有更长更好的生命——我们既从伟大哲学家的论证，也从先见者的回答和先知的神圣话语中得知这一点——那么明智之人就会轻视这现世生活及其利益，因为它的完全丧失由永生补偿。正义的同一辩护者\Person{莱利乌斯}在\Person{西塞罗}中说：\Quote{德性全然渴望荣誉；德性没有其他奖赏。}确实有另一种奖赏，且最配得上德性，哦，\Person{莱利乌斯}，你从未能想到过；因为你不认识圣经。这奖赏它轻易获得，并不苛刻地要求。你大大错了，如果你认为人能给德性付上奖赏，因为你在另一处极其真实地说过：\Quote{你给这个人提供什么财富？什么命令？什么王国？他把这些视为属世之物，判断自己的利益是神圣的。}因此，哦，\Person{莱利乌斯}，谁能认为你是一个明智之人，当你自相矛盾，并在片刻之后从德性那里夺走你曾给予她的？但很明显，对真理的无知使你的意见不确定和摇摆不定。\par

其次，你加上什么？\Quote{但如果所有忘恩负义的人，或许多心怀嫉妒的人，或有权势的敌人，剥夺了德行应得的赏报。}哦，你表现得德行是多么脆弱，多么无价值，如果它能被剥夺赏报！因为如果它判断自己的善是神圣的，如你所说，那么怎么可能有任何如此忘恩负义、如此嫉妒、如此有权势的人，能够剥夺那些由诸神赐予德行的善呢？\Quote{它确实以许多安慰自娱，}他说，\Quote{并特别以自身的美善自持。}以什么安慰？以什么美善？既然那美善常被当作过失责备，并转为惩罚。因为如果，像\Person{弗留斯}所说，一个人被拖走、折磨、放逐、穷困、被砍去双手、被挖去双眼、被定罪、被锁链捆绑、被焚烧、甚至被残酷折磨，德行会失去它的赏报，还是它自身会消灭？决不会。相反，它将从天主审判者那里领受赏报，它将活着，永远兴盛。如果你除去这些，人生中没有什么比德行更显得无用、更愚蠢的了，德行的自然美善和尊贵可以教导我们灵魂不是可朽的，并且天主为它预备了神圣的赏报。但为此，天主愿意德行本身隐藏在愚钝的外表下，以使真理和祂宗教的奥迹保持隐秘；祂要显示这些迷信以及那种世俗智慧的虚妄和错误，这世俗智慧自视甚高，极为自满，从而使其困难最终被阐明，那最狭窄的道路通向不朽的崇高赏报。\par

我想我已经表明了，为什么我们的人被愚昧者视为愚昧。因为选择被折磨和杀害，而不愿用三根手指拿起乳香投入炉火，这在生命危险的情况下，似乎与更关心他人生命而非自己的生命一样愚昧。因为他们不知道，崇拜天主以外的任何对象是多么严重的亵渎，祂创造了天地，塑造了人类，向他们嘘入生命的气息，并赐予他们光明。但是，如果一个奴隶被认为是极无价值的，他逃跑并离开主人，被认为最该受鞭打、锁链、监禁、十字架和一切恶事；同样，如果一个儿子被认为是背弃和不虔敬的，他离弃父亲，不服从他，因此被认为应被剥夺继承权，永不允许他的名字出现在家族中——那么，离弃天主的人，在祂身上主和父这两个同样值得敬畏的名字融为一体，更是多么严重？因为买奴隶的人给予他的恩惠，除了为自身利益供给他的食物外，还有什么？而生儿子的人并不能使他受孕、出生或活着；由此可见，他不是父亲，只是生育的工具。因此，离弃那位既是真正的主又是真父的人，难道不该受天主亲自规定的惩罚吗？祂为邪恶的灵准备了永火，并藉先知向不虔敬和悖逆的人亲口威胁这火。\par

% structure: p0111 source=h2 target=subsection reason=relative-heading-rank; base=h2

\subsection{第二十章　论虚妄与罪行、不虔敬的迷信及基督徒所受的酷刑。}

因此，让那些毁灭自己和他人灵魂的人明白他们犯下了何等不可宽恕的罪过：首先，因为他们事奉最邪恶的邪灵而被天主判受永罚，从而自招死亡；其次，因为他们不允许别人事奉天主，却竭力将人引向致命的仪式，并竭尽全力不让任何仰望天堂且得保障其状态的生命在地上不受伤害。我还能称他们什么，只能称他们为可怜的人，他们服从那些他们以为是神明的掠夺者的煽动？他们对这些掠夺者的状况、起源、名字和本性一无所知；却固守民众的信念，甘心犯错，偏爱自己的愚昧。如果你问他们信念的根据，他们无法给出任何理由，只诉诸祖先的判断，说祖先有智慧，祖先认可了这些，祖先知道什么是最好的；这样他们就剥夺了自己的一切感知能力：他们抛弃理性，却信赖他人的错误。因此，对一切事物一无所知，他们既不认识自己，也不认识他们的神。但愿他们甘愿独自犯错，独自无知！但他们也把别人拉入他们邪恶的同伴中，仿佛要从多数人的毁灭中得到安慰。正是这种无知使他们在逼迫智慧者时如此残酷；他们假装是在促进智慧者的利益，想把他们唤回好的心智。\par

他们是否试图通过交谈或提出某种理由来达到这一点？绝非如此；而是试图用暴力和酷刑来达成。哦，奇妙而盲目的迷妄！他们认为在那些努力保持信德的人心中有恶意，而在刽子手心中有善意。那么，难道在那些违反一切人道法则、违反一切正义原则而受酷刑的人心中有恶意，倒不如说，在那些对无辜者的身体施加连最残忍的强盗、最狂怒的敌人、最野蛮的蛮族也从未施行过的手段的人心中有恶意？他们如此自欺，以至于相互转嫁并混淆善恶之名吗？那么，他们为何不把白昼称作黑夜，把太阳称作黑暗？再者，将善名加于恶、智者之名加于愚者、义人之名加于不虔者，同样是厚颜无耻。此外，如果他们对自己的哲学或辩才有任何自信，就让他们武装起来，驳倒我们的这些论证，如果他们能的话；让他们与我们短兵相接，逐一审视每一点。他们应当为自己的神辩护，免得如果我们的信仰壮大（正如它日益壮大），他们的信仰连同其神龛和虚妄的戏弄被抛弃；既然他们用暴力无法达成任何效果（因为天主的宗教越是受压迫就越增长），那就让他们转而通过理性和劝诫来行动。\par

让他们的司祭——无论是低级的还是最高级的——站出来；他们的\OriginalTerm{弗拉门祭司}{flamens}、\OriginalTerm{占卜官}{augurs}，以及祭祀之王，还有他们迷信的司祭和侍从。让他们召集我们开会；让他们劝我们崇拜他们的神；让他们说服我们存在许多神祇，其神性和眷顾掌管万物；让他们展示他们的圣礼和诸神的起源与开端是如何传给人类的；让他们解释其源头和原理；让他们陈明崇拜他们的回报是什么，忽视的惩罚又是什么；为什么他们希望被人类崇拜；如果他们是幸福的，人的虔诚对他们有何贡献；并且让他们用神圣的见证来证实这一切，而不是靠他们自己的断言（因为凡人的权威没有分量），就像我们做的那样。没有暴力和伤害的余地，因为宗教不能靠强制施加；必须用言语而非拳脚来进行，以便影响意志。让他们拔出智力的武器；如果他们的体系是真实的，就让它被陈述出来。我们准备好倾听，只要他们教导；他们保持沉默时，我们当然不相信他们，正如即使他们发怒我们也不屈服。让他们效法我们，陈述整个事情的体系：因为我们并不像他们所说的那样诱骗；而是教导、证明、展示。因此，没有人被我们强行扣留，因为缺乏信德和忠诚的人对天主毫无用处；然而也没有人离开我们，因为真理本身留住了他。如果他们对真理有任何自信，就让他们这样教导；让他们说话，让他们开口；我敢说，让他们冒险与我们讨论这类问题；那么，他们的错误和愚昧必定会被他们所轻视的老妇人和我们的孩童所嘲笑。因为他们尤其聪明，从书本上知道诸神的世系、事迹、命令、死亡和坟墓；他们也可能知道那些他们被引入的礼仪本身，其起源要么是人类的行为，要么是偶然事件，要么是死亡。难以置信的疯狂才会想象他们是神，而无法否认他们曾是凡人；或者如果他们无耻到否认这一点，他们自己人的著作都会反驳他们；简而言之，礼仪本身的起源就会定他们的罪。因此，他们甚至可以从这件事本身知道真理与谬误之间有多大的差别；因为他们自己尽管有辩才也无法说服人，而无学识、未受教育的人却能说服，因为事情本身和真理在说话。\par

那么，他们为何发怒，以致在试图减轻愚昧时反而加重了它？酷刑与虔诚相距甚远；真理不可能与暴力结合，正义也不可能与残忍结合。但他们有充分的理由不敢教导任何关于神圣事物的事，以免既被我们的人嘲笑，又被自己的人抛弃。因为大多数平民，如果查明这些奥秘是为纪念死者而设立，就会谴责它们，并寻求更真实的崇拜对象。\par

\Quote{由此而来的是神秘的敬畏之礼}\par

它们是由狡猾之人设立的，为的是使百姓不知道他们所敬拜的是什么。但既然我们熟悉他们的体系，为什么他们既不信任我们这些熟悉两者的人，又嫉妒我们，因为我们选择了真理而非谎言？但他们说：\Quote{公共宗教礼仪是必须被捍卫的}。唉，这些可怜的人怀着何等可敬的热忱迷失了方向！因为他们意识到，在人中间没有什么比宗教更高尚，并且宗教应当以我们全部的力量来捍卫；然而，正如他们在宗教本身的事情上受骗，他们在捍卫的方式上也同样受骗。因为宗教是应当通过死亡而非杀戮来捍卫；通过忍耐而非残忍；通过忠信而非罪行：因为前者属于邪恶，后者属于良善；并且良善必须在宗教中有其地位，而非邪恶。因为如果你希望通过流血、酷刑和罪行来捍卫宗教，它就不再被捍卫，而是被玷污和亵渎。因为没有什么像宗教那样属于自由意志；在其中，如果敬拜者的心灵对它不倾向，宗教便立刻被夺去，不再存在。因此正确的方法是，你通过忍耐或死亡来捍卫宗教；在此，信德的保持既悦乐天主自身，又为宗教增添权威。因为如果在这地上的争战中，有人以某种卓越的行动向他的君王保持信德，若他继续活着，便更加受爱戴和悦纳；若他倒下，便获得至高荣耀，因为他为领袖而死；那么，对万有之主、天主的信德更当如何保持呢？祂不仅能为活着的人，也能为死者酬报德行的奖赏！因此，对天主的敬拜，既然属于属天的争战，就需要最大的献身和忠信。因为如果天主自身不被敬拜者所爱，祂如何爱敬拜者？如果祈祷者心怀不诚或不敬来祈祷，祂如何赐予他所求的一切？但这些人，当他们前来献祭时，向他们的神祇奉献的没有一样是来自内心的，没有一样是他们自己的——没有正直的心灵，没有敬畏或恐惧。因此，当这些无价值的祭献完成后，他们便将宗教完全留在圣殿里，连同圣殿，正如他们所发现的那样；既不带任何东西来，也不带走任何东西。因此，这类的宗教仪式既不能使人变善，也不能坚定而不改变。如此，人们很容易被引离它们，因为在其中学不到与生命有关的事，学不到与智慧有关的事，学不到与信德有关的事。因为那些神的宗教是什么？它的能力是什么？它的纪律是什么？它的起源是什么？它的原则是什么？它的基础是什么？它的实质是什么？它的趋向是什么？或者它应许了什么，使人可以忠信地持守并勇敢地捍卫？我在其中看到的无非是仅关乎手指的仪式。但我们的宗教因此是坚定、稳固、不变的，因为它教导正义，因为它常与我们同在，因为它完全存在于敬拜者的灵魂中，因为它以心灵本身作为祭献。在那宗教中，所需要的不过是动物的血、香烟的气息和无意义的奠祭倾倒；但在我们的宗教中，所需要的却是良好的心灵、纯洁的胸怀和无辜的生命：那些仪式常被不贞的淫妇、无耻的老鸨、污秽的娼妓所光顾；它们常被角斗士、强盗、窃贼和巫师所光顾，他们祈祷的不过是能无罪地行恶。因为强盗献祭时所求的，或角斗士，岂非他们能杀戮？投毒者，岂非他能逃过注意？娼妓，岂非她能罪至极致？淫妇，岂非要么丈夫死亡，要么她的不贞被掩盖？老鸨，岂非她能剥夺许多人的财产？窃贼，岂非他能犯下更多的盗窃？但在我们的宗教中，即使是轻微的寻常过错也无立足之地；若有人没有健全的良心前来献祭，他必听见天主向他发出的威吓：我是指那位洞察内心隐秘之处、始终与罪恶为敌、要求正义、要求忠信的天主。这里哪里有邪恶心灵或邪恶祈祷的位置？但那些不幸的人，既不明白他们自己的罪行所显示的敬拜是何等邪恶，因为他们被一切罪行玷污，却前来祈祷；他们以为只要洗净皮肤就是献上虔诚的祭献；仿佛任何河流能洗去，任何海洋能净化，那被封闭在他们胸中的私欲。那么，清洁被邪恶欲望玷污的心灵，借由德行与信德的唯一洗涤驱散一切恶习，岂不是好得多！因为凡这样做的人，虽然带着污秽肮脏的身体，却已是足够纯洁了。\par

% structure: p0118 source=h2 target=subsection reason=relative-heading-rank; base=h2

\subsection{第21章 论异教神与真神的敬拜，以及埃及人所敬拜的动物。}

但他们，因为他们不知道敬拜的对象或方式，盲目而无意识地陷入相反的做法。因此他们敬拜他们的仇敌，用祭品安抚他们的强盗和谋杀者，他们将他们的灵魂连同乳香本身放置在可憎的祭台上焚烧。可怜之人也愤怒，因为其他人没有以同样方式灭亡，带着难以置信的心灵盲目。因为他们若看不见太阳，又能看见什么？仿佛，如果他们是神祇，他们需要人对付蔑视他们的人的帮助。为什么，因此，他们对我们发怒，如果他们没有任何能力做任何事？除非他们毁灭他们权力所不信任的神祇，他们比那些根本不敬拜他们的人更不虔敬。\Person{西塞罗}在他的\WorkTitle{论法律}中，吩咐人怀着圣洁走近祭献，说：\Quote{让他们穿上虔敬，让他们放下财富；若有人行事相反，天主自己将作复仇者。}这话说得很好；因为对天主绝望是不对的，你敬拜祂正是因为你认为祂有能力。因为如果祂不能为祂自己的敬拜者报仇，祂又怎能为祂自己的受屈报仇？所以我愿意问他们，他们认为强迫人违背意愿献祭特别对谁有好处？是对于他们强迫的人吗？但这不是对拒绝的人的恩惠。但我们必须顾及他们的利益，即使违背他们的意愿，因为他们不知道什么是好的。那么，为什么他们如此残忍地骚扰、折磨、削弱他们，如果他们希望他们得救？或者，虔敬从何而来如此不虔敬，以至于他们要么以这种悲惨的方式毁灭，要么使那些他们希望促进其福利的人变得无用？或者他们为神祇服务？但这不是从人违背其意愿强取来的祭献。因为除非它是自愿且从灵魂中献上的，它就是诅咒；当人们因放逐、伤害、监禁、酷刑而被迫献祭时。如果以这种方式被敬拜的是神祇，如果仅因为这个原因，他们就不应被敬拜，因为他们希望以这种方式被敬拜：他们无疑值得人的憎恶，因为奠祭伴随着眼泪、呻吟和从所有肢体流出的血来进行。\par

但相反，我们不要求任何人，无论愿意与否，被迫敬拜我们的天主，祂是所有人的天主；如果有人不敬拜祂，我们也不愤怒。因为我们信靠祂的威严，祂有能力报复对祂的蔑视，如同祂也有能力报复加在祂仆人身上的灾难和伤害。因此，当我们遭受如此不虔敬的事情时，我们甚至不抵抗言语；而是将报复交托给天主，不像那些行动的人，他们想要显得是自己神祇的捍卫者，并且毫无约束地对那些不敬拜他们的人发怒。由此可以理解，敬拜他们的神祇如何不好，因为人本应被善引向善，而不是被恶；但因为这是恶的，甚至它的职责也缺乏善。但毁灭宗教体系的人必须受罚。我们比埃及民族更恶劣地毁灭了它们吗？埃及人敬拜最可耻的野兽和牲畜形象，并崇拜一些甚至羞于提及的东西为神祇。我们比那些人做得更恶劣吗？他们声称敬拜神祇，却公开而可耻地嘲笑他们——因为他们甚至允许以笑声和欢乐演出的关于他们的哑剧。这是什么宗教，或者说，这威严必须被看作多大，在寺庙中被敬拜而在剧院中被嘲弄？做了这些事的人并没有遭受被侵犯神祇的报复，反而受到尊敬和赞扬。我们比某些哲学家更恶劣地毁灭了它们吗？他们说根本没有神祇，而是所有事物自发产生，所发生的一切都是偶然发生的。我们比伊壁鸠鲁学派更恶劣地毁灭了它们吗？他们承认神祇存在，但否认他们关心任何事情，并说他们既不愤怒也不受宠惠影响。通过这些言辞，他们清楚地劝人根本不要敬拜神祇，因为他们既不关心敬拜者，也不对不敬拜他们的人发怒。此外，当他们反对恐惧时，他们努力达成的无非是没有人应该害怕神祇。然而，这些事被人乐意地听取，并被不受惩罚地讨论。\par

% structure: p0121 source=h2 target=subsection reason=relative-heading-rank; base=h2

\subsection{第二十二章 论魔鬼对基督徒的愤怒，以及不信者的错误。}

他们如此迫害我们，并非因为我们的不敬拜他们的诸神，而是因为真理在我们这一边；而真理（\Quote{正如最有真知灼见者所言}）招致仇恨。那么，我们还能作何感想，只能认为他们不知道自己所受的苦吗？因为他们所行的是一种盲目而失去理性的狂怒，我们看得见，而他们自己却浑然不觉。因为迫害者并非这些人本身，——他们对无辜者并无愤怒的理由——而是那些被污染和被遗弃的邪灵，它们明知真理又憎恨真理，便潜入他们的心智，驱策无知的人陷入狂怒。这些邪灵，只要天主子民之间存有和平，便躲避义人，惧怕他们；而一旦它们夺取人的身体、折磨他们的灵魂，义人便以誓言驱逐它们，一听到真天主的名号，它们就逃遁。因为它们听到这名号便战栗、喊叫，声称自己受烙印、被鞭打；当被问及它们是谁、从何而来、又是如何潜入人身时，它们便招认。于是，被天主之名的力量折磨和煎熬，它们便从那人身上出来。由于这些打击和威胁，它们始终仇恨圣洁公义的人；因为自己无力伤害他们，便借助公众的仇恨迫害那些它们视为可憎的人，并以所能使出的一切暴力施行残忍之事：或是以痛苦削弱他们的信德，或是若无法做到，就干脆将他们从地上消灭，好让无人能遏制它们的邪恶。我并非不知对方可能会作何答辩。那么，为何那位全能的天主，那位你们承认统治万有、为万主之主的大能者，竟容许这些事发生，既不报复也不保护祂的敬拜者？简言之，为何那些不敬拜祂的人反而富有、有权势、有福？为何他们享受尊荣与王权，并且使这些敬拜者屈从于他们的权力与统治之下？\par

对此我们也须给出理由，好叫疑惑不再留存。因为人们之所以认为宗教没有天主的大能，主要的原因正在于此：他们被现世和眼前之福的表象所影响，而这些福乐与心灵的关怀毫无关系；他们又看见义人没有这些福乐，而不义之人却应有尽有，于是断定敬拜天主是无用的，因为其中没有包含这些可见之物；同时他们又认为敬拜其他神祇的礼仪是真实的，因为其敬拜者享有财富、尊荣和王国。然而，持此见解的人并未留心考察人的能力与本质，它完全在于心灵，而不在于身体。因为他们所见的无非是可见之物，即身体；而这身体是被看见、被触摸的，它软弱、易朽、必死；他们所渴望和赞叹的一切福乐——财富、尊荣、政权——都属于身体，因此它们与身体本身一样易朽。但灵魂，它是人的唯一构成，既然不为肉眼所见，其福乐也就不可见，因为只在于德行；因此灵魂必须像德行本身一样坚定、恒久、长存，而灵魂之善就在于德行。\par

% structure: p0124 source=h2 target=subsection reason=relative-heading-rank; base=h2

\subsection{第二十三章 论基督徒的正义与忍耐}

逐一举出各种德行的表现，并就每一种表明：一个智慧而公义的人多么需要远离那些福乐——不义之人因享有这些福乐而使人以为敬拜他们的神祇是真实有效的——这将是一项冗长的工作。就我们当前的探讨而言，从一个德行的例子来证明我们的观点便已足够。例如，忍耐是一项伟大而首要的德行，被民众的呼声、哲学家和演说家一致给予最高的赞扬。倘若不能否认这是一项至高无上的德行，那么，为了获得忍耐，公义智慧之人就必须处于不义之人的权下；因为忍耐就是以平静的心态承受所遭受或临到我们身上的恶事。因此，公义智慧之人，因他操练德行，便在自己身上有忍耐；但若他从未遭遇逆境，他便会完全没有忍耐。反之，生活在顺境中的人是没有忍耐的，并且缺少那最大的德行。我称他为没有忍耐，因为他没有遭受任何事。他也无法保持纯洁无瑕，而纯洁无瑕是公义智慧之人特有的德行。他也常行不义，贪求他人之物，并因不义而夺取所贪求之物，因为他没有德行，受制于恶习与罪过；又忘记了自己的脆弱，被傲慢之心所膨胀。\par

由此，不义之人及不认识天主的人，财势尊荣丰盈。因为这一切都是不义的赏报，既然它们不能持久，且凭着贪欲与暴力获得。然而义人与智者，因视这一切为人间之物（正如\Person{莱利乌斯}所说），且视自己的善为属天之善，故不觊觎他人之物，以免违反人性之律而伤害任何人；也不渴求权势或尊荣，以免加害于人。因他知道众人皆由同一天主所造，同处一境，并因兄弟之谊而相连。他满足于自己所有，而且仅止于微薄，因他念及自身的脆弱，不求超出生活所需；甚至从他所有中，他也分施给穷乏之人，因他是虔诚的；而虔诚是一种极大的德行。此外，他轻视脆弱邪恶的欢愉——这些正是财富所追求的——因他节制而克制情欲。他亦无骄傲与傲慢，不自大，不昂首妄自尊大；而是平静温和，谦卑有礼，因他自知其境况。因此，既然他不伤害任何人，不贪恋他人财产，甚至如有人强夺其物他也不护卫，因他懂得如何以温和忍受所受的伤害——因他具有德行——故义人必然受制于不义之人，智者必受愚者凌辱，使不义者因不义而犯罪，而义者因自身具有美德。\par

若有人愿更详细了解为何天主容许邪恶不义者成为有权势、幸运且富有，而另一方面却容许虔诚者谦卑、困苦且贫穷，可取\Person{塞内卡}的那部题为 \WorkTitle{善人虽遇厄运，然有上智安排} 的书，其中他谈了许多，绝非以这世界的愚昧，而是以智慧且近乎神感的言语。他说：\Quote{天主视人类如己子，却容许腐化堕落者生活奢侈安逸，因祂认为他们不值得祂的纠正。但祂常惩戒祂所爱的善人，并以不断的劳苦磨炼他们修德。祂不容许他们因脆弱易逝的财货而腐化堕落。}因此，若我们因自己的过失常受天主惩戒，无人应感奇怪。反而，当我们受困受压时，我们尤其应感谢最慈爱的天父，因祂不容我们的腐败更进一步，而是以鞭笞责打纠正它。由此我们明白，我们是天主所关顾的对象，因祂在我们犯罪时发怒。因为祂本可将财富与王国赐予祂的子民，如从前赐予犹太人那样——我们正是他们的继承者和后裔——但祂却有意让他们生活在他人权柄与统治之下，免得他们因顺境的幸福而腐化，滑入奢侈并轻视天主的诫命，正如我们那些常常被这些属地易逝的财富所削弱的前辈，他们偏离了纪律，挣断了法律的束缚。因此，祂预见了祂将给予敬拜祂的人的安息有多大，如果他们遵守诫命；但同时如果他们没有听从，祂也会纠正他们。因此，为避免他们像他们的父辈因放纵而腐化，祂愿意他们受那些祂置于他们权下者的压迫，以便祂能在他们动摇时坚固他们，在腐化时更新他们使之刚毅，在忠信时试炼并考验他们。一个将军若不面对敌人，怎能证明其士兵的勇力？然而，敌人虽违背他的意愿出现，因为他是凡人，且可能被战胜；但既然天主不可对抗，祂便自己兴起对抗祂名号的敌人——不是为了与天主本身争战，而是对抗祂的士兵——以便祂或证明祂仆人的热忱与忠信，或坚强他们，直到以苦难的鞭笞纠正他们松弛的纪律。\par

还有另一个原因，使祂允许迫害加在我们身上，好使天主的子民得以增加。说明这事为何及如何发生并不困难。首先，许多人因憎恨残忍而被驱离假神的敬礼。因为谁不厌恶这样的祭献呢？其次，有些人因美德和信德本身而感到喜悦。有些人怀疑，这么多人视对神的敬礼为邪恶并非没有理由，以致他们宁愿死也不愿做别人为保全生命所行的事。有人想知道，那甚至受到以死捍卫的善是什么，它超越今生一切令人愉悦和珍爱的事物，财物、光明、身体的痛苦、脏腑的折磨都不能阻止他们。这些事影响巨大；但这些原因一直特别增加了我们追随者的人数。围观的人们听见他们在这些折磨之中说，他们不是向人手所造的石头献祭，而是向那在天上的永生天主献祭：许多人明白这是真的，并将它接纳到心中。其次，正如在不确定之事中常有发生，当他们彼此询问这毅力的原因时，许多与宗教有关的事因谣言在彼此间传播并被细心观察而得以学习；因为这些事是美好的，故不能不令人喜悦。此外，如同经常发生的，随后的报复大大促使人相信。确实，还有一个不小的原因是：魔鬼的不洁之灵得了许可，进入许多人的身体；后来当这些魔鬼被驱出时，被治愈的人便皈依了他们所经验过的力量之宗教。这些众多原因汇集起来，奇妙地赢得一大群人归向天主。\par

% structure: p0129 source=h2 target=subsection reason=relative-heading-rank; base=h2

\subsection{第24章 论天主对基督徒的施刑者所施的报复}

因此，无论邪恶的君王对我们策划什么，天主自己都允许发生。然而，那些极不义的迫害者——对他们而言，天主的名是羞辱和嘲弄的对象——切莫以为他们可以免于刑罚，因为他们仿佛是祂对我们义怒的仆役。他们必受天主的审判，因为他们领受了权柄，却滥用至不人道的程度，甚至傲慢地侮辱天主，将祂永恒的名放在脚下，加以不虔不义的践踏。为此，祂应许要快快报复他们，从地上铲除这些邪恶的怪物。不过，尽管祂惯于甚至在今世报复对祂子民的迫害，祂却命令我们耐心等待那天上审判的日子，在那日祂将按各人的功过亲自尊荣或惩罚。因此，让亵渎者的灵魂不要以为那些被他们如此践踏的人会被轻视而无人报复。那些贪婪凶残的豺狼，折磨了无辜和正义的灵魂，却毫无罪行，他们必定会得到报应。只愿我们努力，使人间惩罚我们的唯一原因只是正义；愿我们竭尽全力，好使我们同时从天主那里获得对我们受苦的报复和赏报。\par

\SourceRecord{Translated by William Fletcher.；Ante-Nicene Fathers；Vol. 7.；Edited by Alexander Roberts, James Donaldson, and A. Cleveland Coxe.；Buffalo, NY: Christian Literature Publishing Co.,；1886.；<http://www.newadvent.org/fathers/07015.htm>.}

% structure: p0001 source=h1 target=section reason=page-title

\section{\WorkTitle{神圣训导}，第六卷（论真实敬礼）}

% structure: p0002 source=h2 target=subsection reason=relative-heading-rank; base=h2

\subsection{第一章 论真实天主的敬礼、纯真，以及假神的敬礼}

我们藉着圣神的教导和真理本身的帮助，完成了我们目标所向的工作；宣讲和解释真理的使命，是由良心、信德以及我们的主自己托付给我的，没有祂，凡事都不能知晓或清晰陈述。我现在来从事这部作品中最主要、最伟大的部分——教导必须以何种方式或何种祭献来敬礼天主。因为这是人的本分，一切的总和与幸福生活的全部过程都包含在这一个目标之中，因为我们由祂所塑造并接受生命的气息，正是为此，并非如\Person{阿那克萨戈拉}所认为的那样，是要我们观看天空和太阳，而是要以纯洁、无玷污的心灵敬礼那造了太阳和天空的天主。不过，虽然在前几卷中，按照我平庸的才能所允许的，我维护了真理，但敬礼的方式本身尤其可以阐明真理。因为那神圣、超越的尊威从人那里什么都不要求，唯独纯真；若有人将此献给天主，他便已经献上了足够虔诚的祭献。但是人们，忽略了正义，尽管被各种罪行和暴行玷污，却自以为虔诚，如果他们在庙宇和祭台上沾染了牺牲的血，如果他们用大量芬芳和陈酒浸湿了炉床。此外，他们还准备神圣的筵席和精选的宴席，好像他们向那些将要品尝其中一些东西的人奉献一样。任何稀罕见到的，任何在工艺或芬芳上珍贵的东西，他们都判断为令他们的神祇喜悦，并非依据他们对神性的任何认知——他们对此一无所知——而是依据他们自己的欲望；他们也不明白天主并不需要世间的资源。\par

因为他们除了大地一无所知，他们只凭身体的感知和快乐来估量善与恶的事物。正如他们按照其快乐来判断宗教，他们也同样安排他们整个生活的行为。既然他们一度彻底从对天空的默观中转向，并使那属天的官能成为身体的奴隶，他们便放纵自己的情欲，好像他们将要把快乐与自己一同带走，他们时时刻刻急于享受快乐；然而，灵魂应当使用身体的服役，而不是身体使用灵魂的服役。同样这些人判断财富是最大的善。如果他们不能通过善行获得财富，他们就努力通过恶行获得财富；他们欺骗、强夺、抢劫、埋伏、起假誓；总而言之，他们不考虑也不顾念任何事物，只要他们能用黄金闪烁，用银器、珠宝和衣服显眼，能把财富花费在他们贪得无餍的食欲上，并总是由成群的奴隶簇拥着穿过被迫让路的人群。因此，他们将自己献给快乐的服役，熄灭了心灵的力量和活力；当他们特别以为自己活着的时候，他们正以最快的速度冲向死亡。因为，正如我们在第二卷中所指出的，灵魂关系到天上，身体关系到大地。那些忽视灵魂的善而追求身体的善的人，被黑暗和死亡所缠累，这属于大地和身体，因为生命和光明来自天上；而缺乏这些的人，由于服侍身体，远离了对神圣之事的理解。同样的盲目处处压迫着这些可怜的人；因为他们既不知道谁是真实的天主，也不知道什么是真实的敬礼。\par

% structure: p0005 source=h2 target=subsection reason=relative-heading-rank; base=h2

\subsection{第二章 论假神与真天主的敬礼}

因此他们将肥美健壮的牲畜祭献给天主，仿佛祂饥饿；他们向祂奠酒，仿佛祂口渴；他们为祂点燃灯火，仿佛祂在黑暗之中。然而，倘若他们能揣测或想像那些天上的美物——我们仍被尘世身体束缚时无法想象其伟大——他们立刻就会明白，用这些空洞的供奉是何等愚昧。或者，倘若他们默观那称为太阳的天上光明，他们便会立刻察觉，天主自己赐予人如此明亮灿烂的光，祂何需他们的烛火？在那小小的圆盘上——因距离遥远，其大小看似不过人头一般——尚且如此光辉璀璨，以致肉眼不能直视，若你注视片刻，昏花之眼便会被云雾黑暗笼罩；那么请问，在祂那里没有黑夜的天主，我们当设想祂有怎样的光明、怎样的辉煌呢？因为祂调节这光本身，使其既不以极度明亮或炽热伤害生物，又赐予它恰如血肉之躯所能承受、五谷成熟所需的分量。那么，向那位光之创始与赐予者献上烛火火炬的人，岂能算是有理智吗？祂向我们索求的是另一种光，且是不带烟尘的，而是（如诗人所言）清澈明亮；我指的是心灵之光，诗人因此称我们为\Emph{\OriginalTerm{人}{\GreekText{φῶτες}}}，这光除非认识天主，无人能献上。然而，他们的神祇，因为属于尘世，需要灯火以免陷于黑暗；而他们的敬拜者，因不尝天上之事，即使在其虔诚的宗教仪式中，也被拉回尘世。因为大地需要光明，其体制与本性是黑暗的。因此他们不赋予神祇天上的知觉，而是属人的。为此他们相信，那些对我们而言必要且悦乐的事，对神祇也一样：我们饥饿时需要食物，口渴时需要饮水，寒冷时需要衣裳，太阳隐没时需要灯光才能看见。\par

因此，再没有比从他们的敬拜本身——这敬拜完全属于尘世——更清楚明白地证明并理解，那些神祇既然曾存活，如今已死了。因为牲畜的血溅在祭台上，这能有什么天上的影响呢？除非他们幻想神祇以人不敢触碰之物为食。无论谁向神祇献上这等食物，即使他是凶手、奸夫、巫师或弑亲者，他仍是幸福亨通的。神祇喜爱他、保护他、赐予他所欲求的一切。因此\Person{Persius}以其风格恰当地嘲讽了这类迷信：\Quote{你用什么贿赂，}他说，\Quote{赢得神祇的垂听？是用肺和肠吗？}他清楚地看出，平息上天威严不需要肉，而需要纯洁的心、正义的灵，以及如他所说，因天生荣誉感而慷慨的胸怀。这才是天上的宗教——不在于腐败之物，而在于灵魂的德性，灵魂源自天上；这才是真敬拜，敬拜者的心灵将自己作为无玷的祭品献给天主。然而，如何才能获得这敬拜，如何才能奉献它，本书的讨论将会显明；因为再没有比引导人行义更辉煌、更合于人的事了。\par

在\Person{Cicero}的著作中，\Person{Catulus}在\Emph{\WorkTitle{Hortensius}}里，当他把哲学置于一切之上时，他说他宁愿要一篇关于义务的短论，也不要一篇为煽动者\Person{Cornelius}辩护的冗长演说。这显然不应视为\Person{Catulus}的意见——他或许并未说出此话——而是\Person{Cicero}的意见，是他写下的。我相信他写这话是为了推荐他即将撰写的\WorkTitle{论义务}诸书，在\WorkTitle{论义务}中他见证哲学里没有比给出生活准则更好、更有益的了。然而，如果那些不认识真理的人这样做，我们这蒙天主教导和光照、能给出真正准则的人，岂不更应当做？不过，我们教导时，并非像传授美德的基础课程——那将是无穷尽的工作——而是如同指导一个在他们看来似乎已经成全的人。因为当他们仍正确教导正直的准则时，我们将把那些他们所未知的事添加进去，以完成和完善他们所不具有的义。但我要略过那些与我们共同的事，以免显得借用了那些我决意揭发其错误之人的东西。\par

% structure: p0009 source=h2 target=subsection reason=relative-heading-rank; base=h2

\subsection{第三章 论道路、恶习与德行，以及天堂的赏报与地狱的惩罚}

有两条道路，哦，君士坦丁皇帝，人类生命必须沿着它们前进——一条通向天堂，另一条沉入地狱；这些道路诗人们在他们的诗歌中，哲学家们在他们的辩论中都曾介绍过。确实，哲学家们将一条描绘为属于美德，另一条属于恶习；他们描绘那条属于美德的在起初入口处陡峭崎岖，如果有人克服困难攀上顶峰，他们说此后他便有一条平坦的道路，一片明亮愉快的平原，并且享受着丰盛而令人愉悦的劳动果实；但那些被最初接近的困难所吓退的人，滑落并转向恶习的道路，这道路起初看来愉快且更为平坦，但当他们稍许前进之后，那愉快的外观就消失了，出现了一条陡峭的道路，时而石子嶙峋，时而荆棘丛生，时而被深水阻断，时而被急流冲击，以至于他们必定感到困难、犹豫、滑倒并跌倒。所有这些都是为了表明，承担美德有极大的辛苦，但一旦获得它们，就有最大的益处和坚定不朽的快乐；而恶习则以某种自然的甜言蜜语诱捕人的心智，并以空虚快乐的假象俘虏他们，引向痛苦的悲伤和惨境——这实在是明智的讨论，如果他们知道美德本身的形态和界限的话。因为他们既没有学到美德是什么，也没有学到从\OriginalTerm{天主}{God}那里有何奖赏等待它们：但我们将在这两卷书中展示这一点。\par

但这些人，因为他们不知道或怀疑人的灵魂是不朽的，便以人间的荣誉或惩罚来评价美德和恶习。因此，有关这两条道路的整个讨论都关乎节俭和奢侈。因为他们说，人类生命的历程类似于字母Y，因为每一个人，当他到达青春早期的门槛，并到达\Quote{道路一分为二的地方}时，便犹豫不决，不知道应当转向哪一边。如果他遇到一位向导，能引导他摇摆不定地走向更好的事物——即，如果他学习哲学或雄辩术，或某些可敬的艺术，以此他能转向良好的行为，这若不付出巨大辛劳是不可能的——他们说他会过一种光荣而丰富的生活；但如果他没有遇到一位节制的教师，他就落入左手边的道路，这条道路呈现出更好的一面——即，他沉溺于懒惰、懈怠和奢侈，这些对于一个不知真正善好的人而言，暂时显得令人愉快，但之后，失去了一切尊严和财产，他将在极其悲惨和耻辱中生活。因此，他们将那些道路的终点归于身体，归于我们在世上所过的此生。诗人或许做得更好，他们主张这双重的道路在阴间；但他们在这点上受骗了，即他们将这些道路归于死者。因此，两者都说了真话，但两者都不正确；因为道路本身应当归于生命，它们的终点归于死亡。因此我们说得更好、更真实，我们说这两条道路属于天堂和地狱，因为义人被许诺了不朽，而不义的人受到永罚的威胁。\par

但我要解释这些道路如何要么升上天，要么沉入地狱，并将阐明哲学家们所不知道的这些美德是什么；然后我要展示它们的奖赏是什么，以及恶习是什么，它们的惩罚是什么。因为或许有人期望我将分别谈论恶习和美德；然而，当我们讨论善恶的主题时，相反之物也可以被理解。因为，无论你引入美德，恶习会自动离开；或者如果你除去恶习，美德会自动接替。善恶之事的本性是如此确定，以致它们总是彼此对抗并驱逐对方：因此，没有美德就不能除去恶习，没有除去恶习就不能引入美德。因此，我们以一种与哲学家们惯常呈现方式非常不同的方式来提出这些道路：首先，因为我们说每条道路都有一个向导，而且每个向导都是不朽的：但一位受到尊荣，他掌管美德和善质；另一位受到谴责，他掌管恶习和邪恶。但他们只在右侧放一个向导，而且并非唯一，也非持久的；因为他们引入任何一门好技艺的教师，他可能把人从懒惰中召回，教他们节制。但他们除了男孩和年轻人之外，不认为任何人进入那条道路；因为艺术是在这些年龄学习的。另一方面，我们引导每种性别、每个年龄和每个民族的人进入这条天路，因为作为那条道路向导的\OriginalTerm{天主}{God}，没有拒绝任何人的不朽。这些道路本身的形状也并非如他们所想。因为，在彼此不同且对立的物事中，Y字母有何必要呢？但较好的一条朝向日出，即较坏的一条朝向日落：因为追随真理和正义的人，得到不朽的奖赏，将享受永恒的光明；但被那邪恶向导诱骗，宁愿恶习而非美德、谎言而非真理的人，必被带向日落和黑暗。因此，我将描述每一条，并指出它们的特性和习性。\par

% structure: p0013 source=h2 target=subsection reason=relative-heading-rank; base=h2

\subsection{第四章 论生命之路、快乐，以及基督徒的艰辛。}

因此，有一条美德与善的道路，它并不像诗人们所说，通往埃律西昂平原，而是通往世界的最高要塞：——\par

\Quote{左手的道路将罪人交付痛苦，\newline{}引向塔耳塔洛斯的罪孽统治。}\par

因为这事属于那控告者，他发明了虚假的宗教，将人从属天的道路上引开，领他们走向丧亡之路。这路的样貌和形状在眼前如此呈现：似乎平坦开阔，满布各种花卉与果实，令人愉悦。因为地上被视为美物的一切，都安放在其中——我指的是财富、尊荣、安逸、快乐，以及各种诱惑；但与之相伴的还有不义、残忍、骄傲、背信、情欲、贪婪、纷争、无知、谎言、愚昧及其他恶行。然而这路的结局如下：当他们到达一个再无归路之处时，这路连同它的一切美丽突然消失，无人能在跌入深渊之前预见这欺诈。因为凡被眼前之物的外表所迷惑，并沉溺于追求和享受它们的人，未能预见死后将要发生之事，背离了天主；他必被投入地狱，受永罚。\par

但那属天的道路却显得艰难多山，或布满可怕的荆棘，或被突出的石块绊阻；每个人都必须以最大的劳苦和足底磨损，并加倍小心以防跌倒，才能行走其上。在这条路上，他安放了正义、节制、忍耐、信德、贞洁、克己、和睦、知识、真理、智慧及其他德行；但与之相伴的还有贫穷、耻辱、劳苦、痛苦及各种艰难。因为凡将希望延伸到现世之外，选择了更美之物的人，必会没有这些地上的东西，以便轻装上阵，毫无阻碍地克服道路的艰难。因为一个身披王室排场或满载财富的人，不可能走上这些艰难道路，更不可能坚持到底。由此可知，恶人和不义之人更容易满足欲望，因为他们的路是向下倾斜的；而善人达成愿望则是困难的，因为他们行走在艰难陡峭的路上。因此，义人既踏上艰难崎岖的道路，就必成为轻视、嘲笑和仇恨的对象。因为凡被欲望或快乐拖入深渊的人，都嫉妒那能达到德行的人，并恼恨有人拥有他们自己所没有的东西。因此，他将贫穷、卑微、卑贱、易受伤害，却忍耐一切痛苦；若他持续忍耐，毫不松懈地直到最后一步和终点，德行的荣冠将赐予他，天主必因他在世上为正义所受的劳苦，以不朽来赏报他。这些就是天主为人生所指定的两条道路，在每条路上，他都显示了善恶之事，但次序颠倒。在一条路上，他首先指出了暂时的凶恶，随后是永恒的美善，这是更好的次序；在另一条路上，先是暂时的美善，随后是永恒的凶恶，这是更坏的次序。如此，凡选择了现世的凶恶与正义相伴的人，必将获得比他藐视的那些更大更确定的美善；但凡宁取现世的美善而不取正义的人，必将陷入比他逃避的那些更大更持久的凶恶。因为此肉身生命是短暂的，因此它的美善与凶恶也必然是短暂的；但由于那与这属世生命相反的属灵生命是永恒的，因此它的美善与凶恶也是永恒的。由此导致，短促的美善之后继以永恒的凶恶，短促的凶恶之后继以永恒的美善。\par

因此，既然善恶之事同时摆在人面前，每个人都应思量：以永恒的美善补偿短暂的凶恶，比为短暂可朽的美善而忍受永恒的凶恶，好得多少。因为在此生，当一场与敌人的争战摆在你面前时，你必须先劳苦，然后才能享受安息；你必须忍受饥渴，必须忍受酷热和严寒，必须卧于地上，必须警醒并经历危险，以便保全你的子女、房屋和财产，享受和平与胜利的一切福分。但若你宁取眼前的安逸而不愿劳苦，你必给自己造成最大的伤害：因为敌人将趁你不抵抗而突袭，你的土地将被荒废，房屋被抢掠，妻子儿女成为掠物，你自己将被杀或被俘。为避免这些事发生，必须放弃眼前的利益，以获得更大更持久的利益。同样，在整个人生中，因为天主为我们预备了一个仇敌，使我们能获得德行，所以必须放下眼前的满足，以免仇敌制服我们。我们必须警醒、设岗、出征、流血；总之，必须耐心忍受一切令人不快和痛苦的事，更当甘心乐意，因为天主的元帅为我们的劳苦定下了永恒的赏报。既然在此属世的争战中，人们付出如此多的劳苦去获取那些会朽坏的东西——它们如何被获得，也必如何朽坏，那么我们绝不应拒绝任何劳苦，因为我们所获得的是绝不会失落的。\par

因为创造人类参与这场战争的天主，愿意他们严阵以待，以敏锐的心志警惕我们唯一敌人的诡计或公开攻击。这敌人像有技巧、有经验的将军一样，用各种伎俩企图陷害我们，按其本性及禀性施以怒火。他向某些人注入无法满足的贪婪，使他们被财富如锁链般束缚，从而驱离真理之路。他用愤怒的刺激点燃另一些人，使他们在专注于伤害他人时，转而离开对天主的默观。他将另一些人投入无度的情欲，使他们沉溺于肉体的享乐，无法注视德行。他向另一些人灌输嫉妒，使他们为自己的痛苦所占据，除了憎恨之人的幸福外，一无所思。他使另一些人因野心欲望而膨胀。这些人将全部精力与关怀用于担任官职，以便在年鉴上留下印记，以名字标记年份。另一些人的欲望更高，不是以暂时之剑统治行省，而是以无限且永恒的权力，希望被称为全人类的君主。此外，他看见那些虔诚的人，就用各种迷信缠住他们，使他们成为不虔诚的人。对于那些寻求智慧的人，他将哲学投到他们眼前，用以光明之相使他们盲目，免得有人把握并持守真理。这样，他堵塞了一切接近人的通道，占据了道路，以普遍的错谬为乐。但为使我们能驱散这些错谬，并战胜罪恶的根源本身，天主光照了我们，以真实而天上的德能武装了我们。关于这德能，我现在必须讲述。\par

% structure: p0020 source=h2 target=subsection reason=relative-heading-rank; base=h2

\subsection{第五章　论假德性与真德性；并论知识}

然而在我开始阐述各项德行之前，我必须先勾勒德性本身的性质。哲学家们并未正确定义其性质，或它由何构成；我还要描述其运作与职责。因为他们只保留了名称，却失去了其力量、本质和效果。但他们在德性定义中通常所言，\Person{卢齐利乌斯}用几句诗概括并表达出来。我宁愿引入这几句诗，以免在反驳许多人的意见时显得冗长：\par

\Quote{德性，\Person{阿尔比努斯}啊，是付出适当的代价，\newline{} 关注我们所从事和赖以生存的事务。\newline{} 德性对于人，是认识万物的本性。\newline{} 德性对于人，是知道何为正当、有益而可敬，\newline{} 何者为善，何者为恶，\newline{} 何者为无用、卑劣而不名誉。\newline{} 德性是知道所求对象的目的及获取之道；\newline{} 德性是能赋予财富以价值；\newline{} 德性是给予荣誉真正应得之物；\newline{} 成为恶人与恶习的敌人和对头，\newline{} 另一方面成为善人与善习的捍卫者；\newline{} 看重他们，祝福他们，与他们结为挚友；\newline{} 此外，首先考虑国家的利益，\newline{} 其次是父母的利益，最后把我们自己的利益放在第三位。}\par

诗人简要总结的这些定义，\Person{马库斯·图利乌斯}遵循斯多葛派的\Person{帕奈提乌斯}，从中引申出生活的职责，并写成了三卷书。\par

但我们将很快看到这些是多么虚假，从而显示天主的垂顾赐予我们开启真理是何等伟大。他说，德性在于知道什么是善与恶，什么是卑劣，什么是可敬，什么是有益，什么是无用。他本可缩短论述，如果他只谈善与恶——因为没有任何有益或可敬之物不是同时也为善，也没有任何无用或卑劣之物不是同时也为恶。这一点对哲学家们看来也是如此，\Person{西塞罗}在上述著作的第三卷中也显示了这一点。但知识不能是德性，因为它不在我们内，而是从外面进入我们。能够从一人传到另一人的东西不是德性，因为德性是个人的固有物。因此，知识在于来自他人的恩惠，因为它依赖于听闻。德性完全是我们自己的，因为它依赖于行善的意志。所以，如同在开始旅程时，知道道路并无益处，除非我们也有行走的努力和力量；同样，如果我们的德性欠缺，知识确实无用。因为通常，即使是犯罪的人也感知什么是善与恶，尽管不完全；每当他们行为不当，他们知道自己犯罪，因此设法隐藏自己的行为。尽管善与恶的本性并未逃脱他们的注意，但他们被邪恶的犯罪欲望所压倒，因为他们缺乏德性，即行正直可敬之事的欲望。因此，善与恶的知识是一回事，德性是另一回事，这一点显而易见：知识可以没有德性而存在，正如许多哲学家的情况那样；既然明知正确之行而不为，理当受谴责，那么邪恶的意志和邪恶的心灵，既然无知不能成为借口，将受到公正的惩罚。因此，正如善与恶的知识不是德性，同样行善避恶才是德性。不过知识如此与德性结合，以致知识先于德性，德性跟随知识，因为知识若无行动跟进则毫无用处。因此\Person{贺拉斯}说得更好一些：\Quote{德性是逃离恶行，首要的智慧是摆脱愚昧。}但他言辞不当，因为他用相反者来定义德性，如同说：善即非恶。因为我若不知什么是德性，就不知什么是恶。因此两者都需要定义，因为事理如此：两者必须同时被理解否则都不能被理解。\par

但我们当行他所当做之事。克制愤怒、控制欲望、抑制情欲是一种德行；因为这就是躲避邪恶。几乎所有不义和不诚实的行为都源于这些情感。因为若这种被称为愤怒的情感力量被削弱，人类一切的恶争都将平息；无人会图谋，无人会冲出去伤害他人。同样，若欲望被抑制，无人会通过陆路或海路使用暴力，无人会率领军队去掠夺和毁坏他人的财物。此外，若情欲的烈焰被压制，每一年龄段和性别都将保持其圣洁；无人会遭受或做出任何可耻之事。因此，若这些情感被德行平息和镇定，一切罪行和可耻行为都将从人的生活和品性中除去。而情感的平定和镇定意味着我们做一切合宜之事。德行的全部职责在于不犯罪。确实，不认识天主的人无法履行这职责，因为对从他而来之善的无知必使人不知不觉堕入恶习。因此，为更简洁明确地固定每项主题的职责：知识是认识天主，德行是敬拜他；前者包含智慧，后者包含义德。\par

% structure: p0026 source=h2 target=subsection reason=relative-heading-rank; base=h2

\subsection{第六章 论至善与德行，以及知识与义德。}

我已说过首要的一点，即对善的认识并非德行；其次，我已指明德行是什么，以及它在于何处。接下来还应指明哲学家们对善与恶无知；这一点简述即可，因为在第三卷讨论至善问题时已基本阐明。由于他们不知至善是什么，他们必然在其他非至善的善恶上犯错；因为人若没有善恶所出之源本身，就不能以正确的判断来权衡它们。善之源是天主；恶之源则是那位常与神圣之名仇敌者，我们多次论及。善与恶即源于此二者。从天主而来的善旨在使人获得不朽，这乃至善；而从另一者而来的恶则旨在使人远离天上之事，沉溺于尘世之事，从而使其受永死之罚，此乃至恶。因此，那些既不认识天主也不认识天主的仇敌的人，岂非对善恶一无所知？故他们将善的归宿归于身体和这短暂、必朽坏的生命，而不更进一步。但他们的一切规诫和所引为善的事物都附着于尘土，躺卧在地，因为它们与必死的身体一同消亡；它们并非旨在为人获得生命，而是获取或增加财富、荣誉、光荣和权势，这些都是完全必朽的，正如那劳苦追求它们的人一样。因此有言：\Quote{德行在于知晓所求之物的目的及获取它的方式}，因为他们教导以何种方式和实践去寻求财产，因为他们看到财产常被不义地寻求。但这种德行并非为智者所设；因为寻求财富并非德行，其获得和拥有都不在我们的能力之内；故此歹人比善人更容易获取和保有财产。因此，德行不能在于寻求那些事物——德行的力量和意义正在于轻视它们；它也不会求助于那些它以崇高之心所欲践踏和踩在脚下的事物；一个专注于天上善事的灵魂，不得因获得这些脆弱之物而偏离其不朽的追求。但德行的进程尤其在于获得那些无人、甚至死亡本身也不能从我们手中夺走的事物。既然如此，以下所言为真：\Quote{德行在于能为财富估定其价值}。这诗句与前两句含义几乎相同。但这位诗人以及他所追随的一切人，都未能知晓这价值本身的性质为何；因为诗人以及他所追随的一切人认为，这意味着正确使用财富——即生活有节制、不设奢华宴席、不随意浪费、不将财产用于多余或可耻之物。\par

或许有人会说：\Quote{你说什么？你否认这是德行吗？}我确实不否认；因为如果我否认，我似乎就要证明相反的一面。但我否认这是真正的德行；因为它不是那属天的原则，而是完全属于尘世的，因为它所产生的效果仅存留于尘世。至于如何善用财富，以及从财富中寻求什么益处，我将在我开始谈论虔诚的职责时更公开地阐述。现在接着的其他东西绝非真实的；因为向恶者宣告敌意，或为善者辩护，可能对恶者也是共同的。因为有些人借着伪装的善行为自己铺平通往权力的道路，并且做许多善人习惯做的事，而且由于他们为欺骗而做，就更为轻易；我希望在实际中行善能像伪装行善一样容易。但是，当他们开始达到其目的和愿望，登上权力的最高阶梯时，那时，真正抛开伪装，这些人便显露出他们的本性；他们攫取一切，施以暴力，劫掠破坏；他们甚至压迫那些他们曾为之辩护的善人自身；他们砍掉自己攀登的阶梯，使无人能够模仿他们来反对他们自己。然而，让我们假设这为善者辩护的职责只属于善人。但承担它是容易的，履行它却是困难的；因为当你投身于一场争战和交锋时，胜利是由天主支配的，而非在你自己的能力之内。而且，恶者大多在数量和联合上比善者更强有力，以致战胜他们与其说需要德行，不如说需要好运。有谁不知道更好、更正义的一方常常被击败？因此，严酷的暴政总是针对公民爆发。全部历史充满了例子，但我们只满足于一个。\Person{格奈乌斯·庞培}愿意成为善者的辩护者，因为他拿起武器保卫共和国、保卫元老院、保卫自由；然而同一个人，被击败后，与自由本身一同灭亡，并被埃及宦官阉割后，未得埋葬而被抛弃。\par

因此，无论是成为恶者的敌人还是善者的辩护者，都不是德行，因为德行不能受制于不确定的机遇。\par

\Quote{此外，要首先考虑祖国的利益。}\par

当人类的共识被取消时，德行便根本不存在；因为祖国的利益若不是另一国家或民族的损失，又是什么呢？——即扩展从他人那里暴力夺取的边界，增加国家的权力，改善财政收入——所有这些都不是德行，而是德行的颠覆：因为首先，人类社会的联合被取消，无辜被取消，不侵占他人财物被取消；最后，正义本身也被取消，它无法忍受人类的分裂，凡武器闪烁之处，正义必定被驱逐和消灭。\Person{西塞罗}的这句话是真实的：\Quote{但是那些说应当照顾公民，而不应当照顾外邦人的人，他们摧毁了人类共同的团体；而当这一团体被移除时，仁慈、慷慨、友善和正义就完全被取消了。}因为一个人若伤害、憎恨、掠夺、杀害，怎能是正义的？而竭力为祖国服务的人做了所有这些事：因为他们不知道这种服务是什么，认为只有用手才能抓住的东西才是有用的、有益的；而这唯一的东西却不能被抓住，因为它可能被夺走。\par

因此，凡是为祖国获得了这些善——如他们自己所称——的人，即通过推翻城市、毁灭民族，用金钱充实国库，夺取土地并使同胞富裕——他就被赞颂到天上：他被认为拥有最伟大、最完美的德行。这不仅是一般百姓和无知者的错误，也是哲学家的错误，他们甚至为不正义制定规则，以免愚行和邪恶缺乏纪律和权威。因此，当他们谈论与战争相关的职责时，那整个论述既不适应正义，也不适应真正的德行，而是适应此生和公民制度；而这并非正义，事情本身已经表明，\Person{西塞罗}也已证明。\Quote{但是，}他说，\Quote{我们并不拥有真正法律和真正正义的真实而活生生的形象，我们只有描绘和草图；我甚至希望我们遵循这些，因为它们是从自然和真理的优秀原本中取来的。}那么，他们所认为的正义不过是一种描绘和草图。但智慧呢？难道同一个人不承认它在哲学家中并不存在吗？\Quote{而且，}他说，\Quote{当\Person{法布里奇乌斯}或\Person{阿里斯提德斯}被称为正义时，并非从这些人身上作为智慧人寻找正义的榜样；因为这些人中没有一个是我们在真正意义上的智慧人。而那些被尊崇并称为智慧人的，如\Person{马尔库斯·加图}和\Person{盖乌斯·莱利乌斯}，以及那著名的七贤，实际上也不是智慧人；而是因为他们恒常实践\InnerQuote{中间职责}，才具有某种智慧人的相似和外表。}因此，如果智慧据哲学家们自己的承认被从他们那里取走，正义也从那些被视为正义的人那里取走，那么随之而来的是，所有那些关于德行的描述必定是虚假的，因为没有人能知道真正的德行是什么，除非他是正义且智慧的人。但没有人是正义和智慧的，除非是天主以天上的诫命教导了他。\par

% structure: p0033 source=h2 target=subsection reason=relative-heading-rank; base=h2

\subsection{第七章　论谬误与真理之路：此路单一、狭窄、陡峭，并以天主为其向导}

所有那些因他人的公认愚蠢而被视为有智慧的人，披着美德的外表，追逐着阴影和轮廓，却没有真实的东西。之所以发生这种情况，是因为那条通往西方的欺骗性道路有许多路径，由于人类生活中各种追求和体系的不同与多样。因为正如那条智慧之路包含着某种类似于愚蠢的东西，正如我们在前书中所展示的，同样，这条完全属于愚蠢的道路也包含着某种类似于智慧的东西，而他们察觉到了人类普遍的愚蠢并抓住这一点；并且正如它有明显的恶习，它也有某种看起来像美德的东西；正如它有公开的邪恶，它也有正义的类似和外表。因为那条路线的先驱，其力量和能力完全在于欺骗，除非他向人们展示一些类似于真理的东西，否则他如何能完全将人们引向欺诈呢？因为为了隐藏祂不朽的奥秘，天主在这条路上放置了一些人们会视之为邪恶和可耻的事物，以便他们远离智慧和真理（他们正在无任何向导的情况下寻求这些），反而陷入他们想要避免和逃避的事物本身。因此，他指出了那条毁灭和死亡之路，它有许多弯曲，要么是因为生命有许多种类，要么是因为有许多神祇被崇拜。\par

这条路的欺骗性和背信弃义的向导，为了使真理与虚假、善与恶之间看起来有些区别，将放纵的人引向一个方向，将所谓有节制的人引向另一个方向；将无知的人引向一个方向，将有学问的人引向另一个方向；将懒惰的人引向一个方向，将积极的人引向另一个方向；将愚蠢的人引向一个方向，将哲学家引向另一个方向，而且就连这些人也不是在同一条路上。因为那些不避开快乐或财富的人，他将其从这条公共且常走的路上稍稍拉开；而那些要么希望追随美德，要么声称鄙视世事的人，他将其拖过一些崎岖的悬崖。尽管如此，所有那些展示出荣誉表象的道路并非不同的道路，而是岔路和小径，它们看起来确实与那条公共的路分开并向右分支，但最终还是回到同一条路，并且最终都通向同一个结局。因为那位向导将它们全部联合起来，在正是有必要将善与恶、强与弱、智与愚分开的地方；即在对诸神的敬拜中，他用同一把剑将他们全部杀死，因为他们毫无区别地都是愚蠢的，并将他们投入死亡。但这条路——即真理、智慧、美德和正义之路，所有这些都只有一个源泉、一个力量之源、一个居所——是单纯的，因为怀着相似的意念和最大的和谐，我们跟随并敬拜唯一的天主；它是狭窄的，因为美德赐予较少数的人；它是陡峭的，因为良善，极其崇高和卓越，若不经历最大的困难和劳苦，就无法达到。\par

% structure: p0036 source=h2 target=subsection reason=relative-heading-rank; base=h2

\subsection{第八章 哲学家的错误，以及法律的多变性。}

这是哲学家们寻求的道路，但他们之所以找不到，是因为他们宁愿在地上寻求它，而它不可能在那里显现。因此，他们就像在大海上漂泊，不明白自己正被带往何方，因为他们既看不清道路，也不跟随任何向导。因为这条生命之路应当以与船只在大海中寻找航向相同的方式来寻求：除非他们观察天上的某种光，否则他们将漫无目的地漂流。但凡是努力持守正确生命之路的人，都不应注视大地，而应注视上天；说得更明白些，他不应跟随人，而应跟随天主；不应事奉这些尘世的偶像，而应事奉天上的天主；不应以身体为参照来衡量一切，而应以灵魂为参照；不应关注此生，而应关注永生。因此，若你总是将目光投向天上，观察太阳从何处升起，并以此作为你生命的向导，如同航海一般，你的脚将自然地导向那条路；并且那天上的光，对于健全的心智而言，是比我们在血肉之躯中所见的这个太阳明亮得多的太阳，它将如此统治和引导你，使你毫无错误地抵达智慧和美德的最佳港湾。\par

因此，我们必须承担天主的法律，它引领我们走上这条路；这部神圣的、天上的法律，\Person{马尔库斯·图利乌斯}在其\WorkTitle{论共和国}的第三卷中，几乎以神圣的口吻进行了描述；我将附上他的话语，以免我过多赘述：\Quote{确实有一条真正的法律，正确的理性，与本性相合，遍及众人，不变、永恒，它通过命令召唤我们尽本分，通过禁止阻止我们犯错；但它对善人既不徒然命令或禁止，也不通过命令或禁止影响恶人。不允许改变这条法律的条款，也不允许我们修改它，更不能完全废除它。确实，无论是元老院还是人民，都不能使我们免除这条法律；也不需另寻他人来解释或诠释它。不会在\Place{罗马}有一套法律，在\Place{雅典}有另一套；不会有现在的一套法律，以后又有另一套；而是同一条法律，永恒不变，将在所有时代束缚所有民族；将有一位共同的、万有的主宰和统治者，就是天主，他是这条法律的制定者、仲裁者和提出者；凡不服从这条法律的，就是逃避自己，轻视人的本性，将因这件事本身而遭受最严厉的惩罚，即使他逃脱了其他被认为存在的惩罚。}有谁认识天主的奥秘，能如此深刻地论述天主的法律，如同一个远离真理知识的人所阐述的那样？但我认为，那些无意识地说出真理的人，应当被视为如同受到某种神灵的感动而说预言。但倘若他也知晓并解释了这部法律本身包含哪些诫命——因为他清楚地看到了神圣法律的力量和要旨——那么他就不是履行哲学家的职责，而是先知了。正因为他未能做到这一点，就必须由我们来完成，因为这部法律本身已由那位唯一伟大的、万有的主宰和统治者——天主——交付给了我们。\par

% structure: p0039 source=h2 target=subsection reason=relative-heading-rank; base=h2

\subsection{第九章　论天主的法律与诫命；论仁慈与哲学家的错误。}

这条法律的第一个要点是：认识天主本身，惟独服从祂，惟独敬拜祂。因为一个人若不认识自己灵魂之父的天主，就不能保持人的品格：这是最大的不敬。因为这种无知导致他事奉别的神，没有比这更大的罪过了。因此，由于对真理和至善的无知，人便如此轻易地步入邪恶；因为天主，即他避而不认识的那一位，本身就是美善的源泉。或者，如果他愿意追随天主的正义，但若对神圣法律无知，他便会将本国的法律当作真正的正义来接受，尽管这些法律显然不是由正义，而是由功利所制定。为何在万民之中有种种不同的法律？正是因为每个国家都为自己制定了它认为对其事务有益的法律。然而，功利与正义有多大差异，\Place{罗马}人民自己就教导了我们：他们通过\OriginalTerm{费提亚祭司团}{Fecials}宣告战争，依法行事，总是渴望并夺取他人的财产，从而获得了对整个世界的占有。但这些人若不做任何违反自己法律的事，便自以为正义；这甚至可以归因于恐惧，如果他们因惧怕现世的惩罚而避免犯罪。但让我们承认，他们出于本性，或如哲学家所说，出于自愿，做了法律所迫使的事。那么，他们会因此而正义吗，仅仅因为他们服从了人的制度？而这些人本身可能犯了错，或曾是不义的——正如\OriginalTerm{十二铜表法}{Twelve Tables}的制定者，他们无疑根据时代状况促进了公共利益。民法是一回事，它随各地风俗而变化；但正义是另一回事，是天主为所有人设立的统一而简单的标准；不认识天主的人，也必然不认识正义。\par

但是，让我们假设有人可能凭借天生的、内在的良善获得真正的德行，就像我们听说过的\Place{雅典}的\Person{客蒙}那样，他既向有需要的人施舍，又款待穷人，还给赤身露体的人衣服穿；然而，当那最最重要的一件事——认识天主——缺乏时，那么所有这些好事都是多余而空洞的，以致在追求它们时，他徒然劳苦。因为他所有的正义都将像一个没有头的身体，尽管所有肢体都处于正确的位置、形状和比例，但由于那最重要的事物缺失，它既无生命也无感觉。因此，那些肢体只有肢体的形状，却没有任何用处，正如一个没有身体的头；而他也类似这样：一个人认识天主，却生活在不义之中。因为他拥有那最重要的事物，但拥有它是徒然的，因为他缺乏德行，如同缺乏肢体一样。\par

因此，为使身体存活且具有感觉，认识天主如同头一样是必需的，而诸德行如同身体一样。如此，便有一个完美而活着的人；然而，整个实体在于头；虽然头不能完全脱离其他部分而存在，但可以在缺少某些部分的情况下存在。这样的人将是一个不完美且有缺陷的活物，却仍然是活的，如同认识天主却在某些方面犯罪的人。因为天主赦免罪过。因此，没有某些肢体仍可以存活，但绝不可能没有头而存活。这就是为何哲学家们，尽管他们可能生性善良，却毫无知识和智慧。他们的一切学识和德行都没有头，因为他们不认识天主——德行之首和知识之首；不认识祂的人，虽看却是瞎的，虽听却是聋的，虽说话却是哑的。但当他一认识万物的造物主和父亲时，他就会看、听和说话。因为他开始有了一个头，其中安置了所有的感官，即眼、耳和舌。因为他确实看见了：他以心灵之眼目睹了真理——天主在其中，或天主——真理在其中；他听见了：他将神圣的言语和赋予生命的诫命铭刻于心；他说话了：在谈论天上的事物时，他讲述了超越万有的天主的德能和威严。因此，不承认天主的人无疑是邪恶的；他所有自以为拥有或具备的德行，都归于那完全属于黑暗的死亡之路。所以，任何人若因获得了这些空洞的德行而自贺，是毫无理由的，因为那不单是失去现世福乐的人有祸，而且他也必定是愚蠢的，因为他一生承担了最大的劳苦，却毫无目的。因为如果挪去了永生的希望——天主向那些持守祂宗教的人所应许的，为了获得永生而应当追求德行，并忍受一切临到的灾祸——那么，愿意服从那些徒然给人带来灾祸和劳苦的德行，必然是最大的愚蠢。因为如果德行在于刚毅地忍受和经受穷乏、流放、痛苦和死亡——这些是他人所惧怕的——那么我请问，它本身有什么好处，以致哲学家们说它应当为其自身而被追求？他们确实以那些多余而无用的惩罚为乐，而他们原本可以平安生活。\par

因为如果我们灵魂是可朽的，如果德行在身体解散后将不复存在，那么我们为何要回避赐予我们的福乐，如同忘恩负义或不配享受神圣恩赐一般？因为要享受这些恩赐，我们必须活在邪恶和不虔敬中，因为德行——即正义——带来贫穷。因此，一个没有更大盼望摆在面前，却宁愿选择劳苦、折磨和痛苦，而不选择他人生活中所享受的福乐的人，是心智不健全的。但如果德行是应当采纳的——正如这些人所正确说过的，因为显然人天生就是为了德行——那么它应当包含某种更大的盼望，能为德行所必须忍受的灾祸和劳苦提供伟大而显赫的慰藉。德行本身是艰难的，除非以其艰难被最大的福乐所补偿，否则无法被视为一种善。我们无法以其他方式同样地戒除这些现世的福乐，除非有其他的、更大的福乐，为了它们而值得放弃对快乐的追求，并忍受一切邪恶。但这些福乐，如我在第三卷中所表明的，没有别的，就是永生之福。谁能赐予这些呢？除了天主，祂亲自将德行摆在我们面前。因此，一切的总和与实质都包含在认识并敬拜天主之中；人的一切盼望与安全都集中于此；这是智慧的第一步：知道谁是我们的真正圣父，并以当有的虔敬唯独敬拜祂、服从祂、以最深的献身将自己交付于祂的侍奉；让我们的全部行动、关怀和注意力都用于赢得祂的恩宠。\par

% structure: p0044 source=h2 target=subsection reason=relative-heading-rank; base=h2

\subsection{第十章　论对天主的宗教与对人的仁慈，以及论世界的起始。}

我已经说了对天主当尽的义务，现在我要说对人所当给予的；然而，你所给予人的，其实是给予天主，因为人是天主的肖像。但是，正义的首要职责是与天主联合，其次是与人和好。前者称为宗教，后者称为仁慈或慈善；这德行是义人和敬拜天主者所特有的，因为唯独它包含了共同生活的原则。天主没有将智慧赐予其他动物，却用天然的防御使它们在危险中更安全。但祂使人赤身露体、毫无防御，为要更赋予他智慧，便赐给他除此以外的仁慈之情，使人保护、爱惜并珍爱他人，并在一切危险中接受和提供援助。因此，仁慈是人类社会最大的纽带；破坏这纽带的人应被视为不虔敬和杀亲者。因为我们若都出于同一个人——天主所创造的那一位，那么我们显然同属一脉；因此，憎恨一个人，即使是有罪的人，也必被视为最大的邪恶。为此，天主命令我们永远不可结成仇恨，而要常常消除仇恨，以致我们安抚那些仇敌，提醒他们彼此的血缘关系。同样，如果我们都由同一天主所激发和赋予生命，我们不是兄弟，又是什么呢？而且，确实更为紧密，因为我们是在灵魂上而非身体上联合。所以，卢克莱修说得好：\Quote{简言之，我们都出自同一神圣种子；我们都同有一位父亲。}因此，那些伤害人的人——他们违反人类本性的一切法律和权利，掠夺、折磨、杀害和放逐——应被视为野兽。\par

由于这种兄弟情谊的关系，天主教导我们永远不要作恶，而要常行善。祂也规定了行善在于什么：给受压迫和困难的人提供帮助，给贫困的人施舍食物。因为天主是仁慈的，祂愿意我们成为社会性的动物。因此，对于其他人，我们应当想到自己。如果我们不援助他人，就不配在自己危险中被解救；如果我们拒绝给予帮助，就不配得到帮助。哲学家们没有这样的训导，因为他们被虚假德行的外表所迷惑，把怜悯从人身上夺走了，他们想要医治，却反而败坏。虽然他们通常承认人类社会的相互参与应当保持，但他们却通过非人德行的严酷将自己完全与之分离。因此，也必须驳斥那些认为不应施舍给任何人的错误。他们不仅提出了建立城市的唯一起源和原因，而且有些人叙述说，最初从大地出生的人，在森林和平原间过着流浪生活，没有被任何言语或正义的相互纽带联结，而是以树叶和草为床，以洞穴和岩洞为居所，成为野兽和更强壮动物的猎物。然后，那些逃脱或被撕裂的人，或者看到邻居被撕裂的人，被自己的危险所提醒，求助于他人，恳求保护，起初通过点头表达愿望；然后尝试言语的初步，通过给每个事物命名，逐渐完成了言语系统。但当他们看到他们自己的人数也不足以安全对抗野兽时，也开始建造城镇，要么使他们夜间安息安全，要么通过设置障碍而非战斗来抵御野兽的入侵和攻击。\par

啊，不配为人者的头脑，竟产生这些愚蠢的琐事！可怜又可悲的人，将这些愚蠢的记录写下并传给后世记忆；当他们看到聚集在一起的设计、或相互交往、或避免危险、或防范邪恶、或为自己准备床铺和洞穴，甚至对哑巴动物来说也是自然的，然而却认为人类除非通过榜样，否则不可能被提醒和学习什么应该恐惧、什么应该避免、什么应该做，或者认为除非野兽吞噬他们，否则他们永远不会聚集或发现言语的方法！这些事情在其他人看来是无意义的，正如它们实际一样；他们说聚集的原因不是野兽的撕裂，而是人性本身的感觉；因此他们聚集起来，因为人的本性逃避孤独，渴望交往和群居。他们之间的分歧不大；因为原因不同，事实却相同。每种说法都可能为真，因为没有直接对立。但是，无论哪种都不是真实的，因为人类并非如同诗人所述那样，仿佛从某条龙的牙齿中生长出来，分布在整个大地上；而是一个人由天主形成，从那个人，全地充满了人类，如同洪水后又同样发生一样，他们当然无法否认。因此，起初并没有发生这种聚集；而且，地球上从未有过不能说话的人，除了婴儿，每一个有理智的人都明白。然而，让我们假设这些懒惰而愚昧的老人徒然所说的事情是真的，以便我们可以特别通过他们自己的情感和论证来反驳他们。\par

如果人们聚集在一起是为了通过相互帮助来保护自己的软弱，那么我们必须援助需要帮助的人。因为，既然人们为了保护而与他人进入并缔结合约，那么要么违反、要么不遵守从起源之初就在人们之间订立的这一契约，应被视为最大的不敬。因为那从提供帮助中撤回自己的人，必然也要从接受帮助中撤回；因为拒绝帮助他人的人认为自己不需要任何人的帮助。但是那从整体中退出并分离自己的人，必须不按人的习俗生活，而按野兽的方式生活。但如果这做不到，那么人类社会的纽带必须无论如何保留，因为人没有他人就完全无法生活。但是社会的保存是善意的相互分享；那就是提供帮助，以便我们能够接受帮助。但是，如果如那些人所断言，人类的聚集是出于人性本身，那么人无疑应当承认人。但如果那些无知且尚未开化的人，在言语实践尚未建立时就做了这些，那么我们应该认为，那些文雅的、通过交谈和一切事务相互联系、习惯于人类社会而不能忍受孤独的人，应当做什么呢？\par

% structure: p0049 source=h2 target=subsection reason=relative-heading-rank; base=h2

\subsection{第十一章 论施惠的对象}

因此，若我们愿意正当地被称为人，就必须保持人道。但保持人道若非爱人，只因他是人且与我们相同，又是什么呢？因此，纷争与不睦与人的本性不相符；\Person{西塞罗}的那句话也是真的，他说人若顺从本性，便不能伤害人。所以，若伤害人违反本性，则有益于人必定符合本性；凡不这样做的人，便剥夺了自己作为人的称号，因为人道之责在于援救人的急需与危难。因此，我要问那些不认为智者应当被说服而产生怜悯的人：若一人被野兽攫住，向一个武装者求援，他们是否认为他应当被援助？他们不至于无耻到否认那合乎人道要求的事应当去做。同样，若有人被火围困，被倒塌的建筑物压住，陷入海中，或被河流冲走，他们是否认为不援助他是人的本分？他们若这样想，便不是人；因为无人能免于这类危险。是的，他们必会说，拯救一个濒死之人是人的本分，甚至是勇敢者的本分。因此，若在危及人命的此类灾祸中，他们承认施以援手是人道之责，那么，如果一个人遭受饥饿、口渴或寒冷，他们为何认为不应施以援手？然而，这些事与那些意外情况在性质上相同，都需要同一种人道，他们却加以区别，因为他们不以真理本身衡量一切，而以眼前的利益。因为他们希望那些被他们从危险中救出的人会回报他们。但因他们不指望穷人有回报，便认为施予这类人的一切都属白费。因此，\Person{普劳图斯}那句可憎的话出现了：——\par

\Quote{施食予乞丐者，其行不善；\newline{}因其所施乃白费，且\newline{}延长他人之生命以增其苦。}\par

但也许那诗人是为剧中人而说的。\par

\Person{马尔库斯·图利乌斯}在他的\WorkTitle{论义务}一书中说了什么？他不是也劝告人根本不要施用慷慨吗？因为他这样说：\Quote{从我们产业中支出的厚赐，抽干了我们慷慨的源头；于是慷慨被慷慨毁掉；因为你越是向更多的人施行慷慨，你越不能向许多人施行慷慨。}随后他又说：\Quote{还有什么事比那样做更愚蠢，即你这样做却无法继续做你自愿做的事？}这位智慧教授显然在阻止人行善，劝他们谨慎保护财产，安全保存钱柜，而不是追求正义。当他意识到这是不人道和邪恶时，不久后，仿佛出于悔恨，在另一章中他这样说：\Quote{然而，有时我们必须施行厚赐；这种慷慨也不应完全摒弃；并且，当合适的人需要援助时，我们应从自己的财产中给予他们。}\Quote{合适}是什么意思？当然是那些能够偿还并回报恩惠的人。如果\Person{西塞罗}现在活着，我必会喊道：这里，这里，\Person{马尔库斯·图利乌斯}，你背离了真正的义；你用一个词就把它取消了，因为你用利益衡量了敬虔和人道的义务。因为我们不应将厚赐施与合适的人，而应尽可能施与不合适的人。因为那不带任何回报希望而做的事，才是凭着义、敬虔和人道做的。\par

这就是真实而纯正的义，你说你对其没有真实和活生生的形象。你自己在许多地方声称美德不是有酬劳的；你在\WorkTitle{法篇}一书中承认慷慨是无偿的，用这些话：\Quote{毋庸置疑，那被称为慷慨宽宏的人，是受责任感而非利益驱动。}那么，为什么你把厚赐施与合适的人，除非是为了以后得到回报？因此，以你作为义的倡导者和教师，凡是不合适的人，将因赤身、口渴、饥饿而憔悴；而富有且物资充裕甚至奢侈的人，将不会援助他的绝境。如果美德不要求报酬；如果，如你所说，它应当为其自身而被追求，那么请估量义——它是美德之母和首要者——以其本身的价值，而不是按你的利益：尤其要施与那你不指望任何回报的人。你为什么挑选人？你为什么看外表？那向你哀求的人，无论他是谁，都应当被你当作一个受尊敬的人，因为他视你为一个人。抛弃那些义的轮廓和草图，紧握义本身，真实而栩栩如生。对盲人、弱者、跛子、赤贫者慷慨吧，他们若不接受你的慷慨就会死。他们对人无用，但对天主有用，天主存留他们的生命，赋予他们气息，恩赐他们光明。尽力珍惜并以仁慈支持人的生命，使之不被熄灭。能够援救濒死之人的人，若不做，便是杀了他。但他们，因既不持守本性，也不知晓其中的赏报，害怕损失，反而损失了，陷入他们主要防备的事中；以致他们施予的，要么完全失去，要么只能得益极短时间。因为他们拒绝给可怜人一点施舍，希望不损失自己而保持人道，却浪费了财产，以致要么为自己获得脆弱易逝之物，要么肯定因自己巨大损失而一无所获。\par

因为对于那等人，他们被民众赞美的虚荣所驱使，将甚至足以供应大城市的财富耗费在表演的展示上，我们又该说什么呢？我们岂不是要说，那些将既失去于自己，又为任何受赠者所不获的东西施予民众的人，是愚昧和疯狂的吗？因此，既然一切快乐都是短暂且易逝的，尤其是眼目和耳朵的快乐，人们要么忘记且不感激他人所花费的代价，要么如果民众的任性得不到满足，他们甚至会感到愤怒：以致最愚蠢的人甚至以恶为自己招致了恶；或者，如果他们因此成功讨好了人，他们所获得的不过是空洞的青睐和几日的谈论。这样，每天都有最轻浮之人的田产被浪费在无谓之事上。那么，那些向同胞展示更有用且更持久的礼物的人，行事是否更为明智呢？例如，那些通过建造公共工程来为自己名字寻求长久记忆的人？他们将自己的财产埋入地下，也并非行事正确；因为对他们的纪念既不带给死者任何东西，他们的工程也不是永恒的，因为它们要么被一次地震倾覆和摧毁，要么被偶然的火烧毁，要么被敌人的攻击推翻，或者至少因年久失修而朽坏倒塌。正如演说家所说，人的手所造的任何东西，没有一样是时间不能削弱和毁灭的。但我们所说的这正义和慈悲，却日日更加兴盛。因此，那些将恩惠施予同族人和门客的人行事更好，因为他们确实将一些东西施予人，并使人受益；但那不是真实和公正的施恩，因为没有需要之处便没有施惠。因此，凡为求名声而给予无需要之人的东西，都是被浪费的；或者它连本带利被偿还，这样它便不是施惠。虽然它使受赠者喜悦，但它仍不是正义的，因为如果不这样做，也不会带来任何恶。因此，施舍的唯一可靠且真实的职责是扶助有需要和无力自助的人。\par

% structure: p0056 source=h2 target=subsection reason=relative-heading-rank; base=h2

\subsection{第十二章 论施舍的种类与慈悲善工}

这就是那保护人类社会的完美正义，哲学家们论及它。这是财富首要且最真实的用处：不将财富用于个人特殊的享乐，而用于许多人的福利；不为一己的即时享受，而为正义，因为只有正义永不朽坏。因此，我们必须绝对牢记，在施行慈悲的职责中，必须完全排除得到回报的希望：因为这项工作和职责的赏报只应从天主那里期待；如果你从人那里期待它，那么那便不是仁慈，而是借贷以求利；而且，一个人不是为别人，而是为自己所给予的，便不能算是积了功德。然而，事情归结为这一点：一个人若对他人施予而不指望从他那里得到任何好处，他实际上是将恩惠施予了自己，因为他将从天主那里获得赏报。天主也命令，我们若有时设宴，应当邀请那些不能回请我们、不能回报我们的人，这样我们生活的每一个行动都不可缺乏慈悲之举。然而，不要让人以为他被禁止与朋友交往或对邻人友善。但天主已向我们指明什么是我们真实且正义的工作：我们应当如此与邻人相处，只要我们知道一种生活方式关乎人，另一种关乎天主。\par

因此，款待是一项主要的德行，正如哲学家们所说；但他们却将其从真正的正义中引开，并强行应用于私利。\Person{西塞罗}说：\Quote{\Person{泰奥弗拉斯托斯}正确地赞扬了款待。因为（在我看来）显赫之人的家宅向显赫的宾客敞开，是非常得体的。}他在这里犯了与他当时同样的错误，即他说我们必须将我们的慷慨施予\Quote{合适}的人。因为一个正义而智慧之人的家宅，不应对显赫者开放，而应对卑微和低下的人开放。因为那些显赫而有权势的人不可能缺乏任何东西，因为他们自己的财富足以保护并荣耀他们。但一个正义之人只应做有益之事。但如果善举得到回报，它就被毁掉并终止了；因为我们无法完全拥有已经付出代价的东西。因此，正义的原则适用于那些保持完好无损的善举；但除非将它们给予那些完全无法使我们受益的人，否则它们无法如此保持。但在接待显赫者时，他只着眼于效用；这位聪慧的人也没有掩饰他希望从中得到的好处。因为他声称，这样做的人将因领袖们的恩宠而成为外邦人中有权势的人，这些领袖将因款待和友谊的权利而与他绑定。哦，假如这是我的目标，\Person{西塞罗}的矛盾可以用多少论证来证明！他不会被我的话定罪，更多的是被自己的话。因为他还说，一个人越是把自己的所有行为都指向自己的私利，他就越不是一个好人。他还说，一个简单坦诚的人不应讨好他人、假装和声称什么、做一套而显另一套、假装将自己给别人的东西留给自己；这反而是一个有计谋、狡猾、欺骗和背信弃义之人的作为。但他怎能坚持那种野心勃勃的款待不是恶意的呢？\Quote{你是否跑遍所有的城门，以便邀请各邦国和城市中的首领前来你家，借此你可以通过他们获得其公民的影响力；并希望被称为正义、善良、好客，尽管你是在谋求促进自己的利益？}但他这样说岂不是太不谨慎了？还有什么比这更不适合\Person{西塞罗}的呢？但由于对真正正义的无知，他明知且有预谋地落入了这个陷阱。为了使自己得到宽恕，他声明他所给出的训导并非基于真正的正义——他并不拥有真正的正义——而是基于正义的草图和轮廓。因此，我们必须宽恕这位使用草图和轮廓的导师，也不应要求一个承认对真理无知的人给出真理。\par

赎回俘虏是一项伟大而高尚的正义实践，同一位\Person{图利乌斯}也赞同此事。\Quote{而这种慷慨，}他说，\Quote{甚至对国家也有用，即应将俘虏从奴役中赎回，并为那些资源微薄的人提供生计。我如此偏好这种慷慨的行为，远胜于在表演上挥霍。这是伟大而杰出之人的作为。}因此，供养穷人和赎回俘虏是正义之人的合宜工作，因为在非正义之人中，如果有人做这些事，他们就被称为伟大和杰出。因为对于那些施予善举而无人预期他们如此行的人，是值得极大赞扬的。因为那些对亲属、邻人或朋友行善的人，要么不值得赞扬，或者当然不值得大加赞扬，因为他有义务去做，如果他不行那自然本身和关系所要求之事，他便是邪恶可憎的；而如果他做了，他这样做与其说是为了获得荣耀，不如说是避免指责。但若是对陌生人和不认识的人这样做，他确实值得赞扬，因为他是出于纯粹的善意而行。因此，正义存在于无需任何义务而施予善举的地方。因此，他不应该将这项慷慨的义务优先于看戏的花费；因为这是进行比较，从两种好处中选择更好的一方。因为那种将财产抛入大海的挥霍是徒劳而微不足道的，且远离一切正义。因此，它们甚至不能被称为礼物，其中无人有所领受，除了那不配领受的人。\par

保护并捍卫那些孤苦无依、亟需援助的孤儿寡妇，同样是一件伟大的义德工作；因此，天主的法律将这诫命颁给众人，因为所有善心的法官都认为，以本性仁慈善待他们，并竭力使他们受益，乃是自己的职责。但这些工作尤其属于我们，因为我们领受了法律，也领受了天主亲自训诲的圣言。因为外邦人虽能察觉到保护需要保护者乃本性之正义，却不明白其中的缘由。原来，那属于永恒慈悲的天主，之所以命令我们捍卫并照料孤儿寡妇，是为了使任何人都不至于因顾念和怜惜自己的亲骨肉，而不敢为义德和信德赴死，相反，他应满怀快慰与勇气去面对死亡，因为他知道自己将所爱之人托付于天主的眷顾，他们绝不会缺少保护。此外，承担照顾和支持那些需要人帮助的病患，这是最伟大的仁慈之举，也是极大的善行；凡行此善举者，不仅向天主献上活的祭献，而且他暂时给予他人的，将从天主那里获得永恒的赏报。虔诚的最后一项也是最伟大的职分，是埋葬陌生人和贫穷者；那些德性与义德的导师对此却绝口不提。因为他们以效益衡量一切义务，故无法洞见此事。在以上所提及的其他事上，虽然他们未能遵循正道，但由于在这些事中发现了某种益处，恰如被真理的一点暗示所牵引，所以他们偏离得较近；但这一项，因他们看不出任何益处，便放弃了。\par

不仅如此，有些人甚至认为埋葬是多余的，并说无人掩埋、无人过问并非灾祸；但他们的不虔敬的智慧遭到全体人类的摒弃，也受到命令执行此礼仪的天主圣言的谴责。他们却不敢说此事不该做，而是说，倘若被忽略也不会带来什么不便。因此，在这件事上，他们行使的职责与其说是传授诫命的人，不如说是提供安慰的人，为的是让一位贤人若偶然遭遇此情形，不因而自认为不幸。但我们所说的，不是一位贤人应当忍受什么，而是他本人应当做什么。因此，我们现在并非探究埋葬的全部仪式是否有益；而是说，即使它像他们所想象的那样毫无益处，也必须履行，单凭这一点——它在众人眼中被视为正当而仁慈之举——就够了。因为所审察的是心意，所衡量的乃是意图。所以，我们决不容许天主的肖像和造物沦为野兽和飞禽的猎物，而要将其归还给大地，它原本就是从那里来的；即便是一个素不相识的人，我们也要尽到亲属的职责；既然亲属缺席，就让仁慈取而代之；凡有需要人之处，我们就当认为我们的职责在那里被要求。然而，义德的本质，岂不更在于我们出于仁慈对陌生人施行那些我们出于亲情对自家亲属所做的事吗？而且，当这仁慈不再施与那毫无知觉的人，而唯独施与天主时，它便更为确定、更为义德，因为一件义德的工作对天主而言是最中悦的祭献。或许有人会说：\Quote{如果我做这一切，我将一无所有。因为假设有许多人匮乏、受冻、被掳、死亡，而行此善举者甚至会在一天之内耗尽自己的财产，难道我要把自己或祖先辛勤挣来的家产都抛出去，以致此后我自己也要靠别人的怜悯过活吗？}\par

你们为何如此怯懦地惧怕贫穷呢？连你们的哲学家都赞美贫穷，并见证说，没有比贫穷更安全、更宁静的了。你们所惧怕的，正是忧患的避风港。你们岂不知道，凭借这些邪恶的财富，你们会暴露于多少危险、多少意外之中吗？这些财富若能在不流血的情况下从你们手中溜走，那就算善待你们了。但你们却满载战利品而行，带着那些连你们的亲属都会为之动心的赃物。那么，你们为何迟疑不肯善用那或许一次抢劫就会夺走、或骤然的流放令、或敌人的掠夺就会掠去的东西呢？你们为何害怕将脆弱易逝的财物变成永恒的，或将你们的宝藏托付给天主作为保管者呢？若如此，你们便无需惧怕盗贼、强盗、锈蚀或暴君。凡在天主面前富有的人，永不会贫穷。如果你们如此珍视\OriginalTerm{义德}{justice}，那就卸下压迫你们的重担，追随它吧；挣脱锁链和桎梏，好能毫无阻碍地奔向天主。轻视并践踏尘世之事，乃是伟大而崇高灵魂的本分。但若你们不明白这德行，无法将你们的财富献于天主的祭台，以便为自己预备比这些脆弱之物更坚固的产业，那我就要解除你们的恐惧。所有这些诫命并非只赐给你们，而是赐给所有同心合意、如同一人团结在一起的全体子民。若你们单独无力成就伟大的事业，就要尽力培养义德，只是须在善工上超越他人，如同你们在财富上超越他们一样。不要以为你们被劝告要减少或耗尽你们的财产；而是要将你们原本会花费在多余事物上的东西，转向更好的用途。用那些你们原本购买野兽的钱财去赎回俘虏；用那些你们原本喂养野兽的钱财去供养穷人；用那些你们原本供养人上剑的钱财去埋葬无辜的死者。资助那些与野兽搏斗的败坏邪恶之人，并装备他们去犯罪，这有什么益处呢？要将那些将要被可怜地丢弃的东西，转移到伟大的祭献上，好使你们因这些真实的礼物，从天主那里获得永恒的恩赐。\OriginalTerm{慈悲}{mercy}有极大的赏报；因为天主应许，祂将赦免一切罪。祂说：\Quote{如果你们俯听你们哀求者的祈祷，我也必俯听你们的祈祷；如果你们怜悯困苦的人，我也必在你们困苦时怜悯你们。但如果你们不理会、不帮助他们，我也必以同样的心意对待你们，并按你们自己的律法审判你们。}\par

% structure: p0063 source=h2 target=subsection reason=relative-heading-rank; base=h2

\subsection{第十三章  论\OriginalTerm{悔改}{Repentance}、\OriginalTerm{慈悲}{Mercy}与\OriginalTerm{罪的赦免}{Forgiveness of Sins}。}

因此，每当你们被请求援助时，要相信这是天主在考验你们，看看你们是否配得蒙垂听。要省察你们自己的良心，并尽力医治你们的创伤。然而，不要因为过犯能因慷慨而消除，就以为你们获得了犯罪的许可。因为如果你们因犯罪而对天主慷慨，过犯就会被消除；但如果你们因依赖自己的慷慨而犯罪，过犯就不会被消除。因为天主特别渴望人能从他们的罪中被洗净，因此祂命令他们悔改。但悔改不是别的，正是承认并声明自己不再犯罪。因此，那些无意中、不小心滑入罪中的人得蒙赦免；故意犯罪的人不得赦免。然而，如果有人已从一切罪的污秽中被洗净，不要认为他可以因没有过犯需要抹去而停止慷慨的善工。其实，当他已成为义人时，他更应履行义德，好使他先前为医治创伤所行的，以后能为德行的赞美和荣耀而行。此外，没有人能在身披血肉之躯时毫无过犯，因为这血肉之躯的软弱在三方面——行为、言语和思想——都受罪的辖制。\par

凭这些阶梯，\OriginalTerm{义德}{justice}达至最高峰。德行的第一步是戒除恶行；第二步是戒除恶言；第三步是戒除恶念。登上第一步的人已足够正义；登上第二步的人现已拥有完美的德行，因为他既不在行为上犯罪，也不在言语上犯罪；登上第三步的人似乎真正达到了\OriginalTerm{天主的肖像}{likeness of God}。因为即使不将那些在行动上邪恶或在言语上不当的事纳入思想，也几乎超出了人的尺度。因此，即使是义人，虽能止于一切不义的行为，但有时也会因自身的脆弱而跌倒，以致在愤怒中说恶言，或在见到悦目之物时心中默默渴求它们。既然必死的境况不容许人完全洁净无瑕，肉身的过犯便应藉着不断的慷慨而被消除。因为聪明、正义、堪当永生的人的唯一工作，就是将他的财富只用在义德上；确实，缺少义德的人，即使他的财富超过\Person{克罗伊斯}{Crœsus}或\Person{克拉苏}{Crassus}，也应被视为贫穷、赤身、如同乞丐。因此，我们必须努力，好使我们穿上\OriginalTerm{义德之服}{garment of justice}和\OriginalTerm{虔敬}{piety}，这衣裳无人能夺去，且能为我们提供永远的装饰。因为如果崇拜偶像的人敬拜无知的塑像，并将他们一切贵重之物献给它们，尽管它们既不能使用这些礼物，也不能因领受而感恩，那么，敬拜天主活生生的肖像，以赢得永生天主的宠幸，岂不更加公义和真实？因为这些肖像会使用他们所领受的，并感恩；同样，天主——在祂眼前你们行了善——必会悦纳并赏报你们的虔敬。\par

% structure: p0066 source=h2 target=subsection reason=relative-heading-rank; base=h2

\subsection{第十四章  论\OriginalTerm{情感}{Affections}，以及\OriginalTerm{斯多葛派}{Stoics}关于它们的意见；并论\OriginalTerm{德行}{Virtue}、\OriginalTerm{恶习}{Vices}与\OriginalTerm{慈悲}{Mercy}。}

因此，若怜悯在人身上是一种卓越而出众的恩赐，且被善人恶人一致视为至善，那么哲学家们显然远离了人之善，因为他们既未命令也未实践任何此类事，反而一直将那几乎在人中居首位的德行视为恶习。在此，我乐于援引哲学的一个主题，以便更充分地驳斥那些把怜悯、欲望和恐惧称为灵魂疾病之人的错误。他们确实试图区分德行与恶习，这实则是一件非常容易的事。因为谁能无法像他们那样区分大方之人和奢侈之人，或节俭之人和吝啬之人，或平静之人和懒惰之人，或谨慎之人和怯懦之人呢？因为这些善的事物有其限度，一旦超越限度，便堕落为恶习；同样，恒毅若不为真理而承担，便成为无耻。同理，勇敢若在没有必要逼迫或不为荣耀原因的情况下贸然涉险，便化为鲁莽。直言不讳若攻击他人而非反击攻击者，便是固执。严厉若非将其自身限制在罪有应得的惩罚内，便成为残忍。\par

因此，他们说，那些看似邪恶的人并非出于自愿犯罪，也不是偏好选择恶事，而是因着善的表象而犯错误，在不知善恶区别的情况下落入邪恶。这些话并非不真实，但一切都指向身体。因为节俭、恒毅、谨慎、平静、庄重、严厉确实是德行，但这些都是关联于短暂生命的德行。但我们这些轻视此世生命的人，面前摆着其他德行，关于这些德行，哲学家们甚至丝毫无法揣测。因此，他们将某些德行视为恶习，将某些恶习视为德行。因为斯多葛派从人身上除去所有情感，灵魂正是因着这些情感的冲动而动——欲望、喜乐、恐惧、悲伤：前二者源于未来的或现在的善事，后者源于恶事。同样，他们称这四种（如我所说）为疾病，与其说是天性植入我们内的，不如说是因着扭曲的见解而承担的；因此他们认为，若能除去对善恶的虚假观念，这些情感便可被根除。因为若智者视无物为善或恶，他便不会被欲望点燃，不会被喜乐激动，不会被恐惧惊扰，也不会因悲伤而垂头丧气。我们即将看到，他们是否实现了他们所愿的，或他们实际达成了什么；与此同时，他们的目的狂妄得近乎疯狂，竟以为自己在施以疗法，并能与天性的力量与法则相抗衡。\par

% structure: p0069 source=h2 target=subsection reason=relative-heading-rank; base=h2

\subsection{第十五章 论情感与逍遥学派对它们的看法}

因为，这些情感是自然的而非自愿的，这由所有受造物之本性显明，它们皆被所有这些情感所动。因此，逍遥学派做得更好，他们说这些情感不可能从我们身上去掉，因为它们与我们一起出生；他们努力显示，天主或本性（他们如此称呼）是多么明智且必要地用这些情感装备了我们；然而，这些情感若过度泛滥，往往会变成恶习，人可以借由施加限制来有益地调节它们，以便给人留下足以满足本性的部分。这并非不明智的讨论，如果，如我所说，一切都未指向此生。因此，斯多葛派是疯狂的，他们不是调节而是切除这些情感，试图以某种方式剥夺人那些由本性植入的能力。这无异于想要除去鹿的怯懦、蛇的毒液、野兽的愤怒或牛群的温顺。因为那些分别赋予哑巴动物的品质，一齐赋予了人。然而，若如医生所言，喜乐的情感居于脾脏，愤怒居于胆囊，欲望居于肝脏，恐惧居于心脏，那么杀死动物本身比从身体中撕下任何东西更容易；因为这是要改变活物的本性。但这些有技艺的人不明白，当他们从人身上除去恶习时，他们也除去了德行，而他们正是在为德行腾出地方。因为若在愤怒的冲动中节制和克制自己是德行（这是他们无法否认的），那么没有愤怒的人也就没有德行。若控制身体的欲望是德行，那么没有可调节的欲望的人必然缺乏德行。若抑制想要贪图他人之物的欲望是德行，那么没有这种需要克制的事物的人，当然不可能有德行。因此，哪里没有恶习，哪里甚至没有德行的位置，正如没有对手就没有胜利的位置。如此便得出，此生中不可能有善而无恶。因此，情感是心灵力量的一种自然丰饶。因为正如一片天然肥沃的田地会产出丰盛的荆棘，未经耕耘的心灵也会因其自发繁茂的恶习而蔓生如荆棘。但当真正的耕耘者出现时，恶习立即退却，德行的果实便萌发出来。\par

因此，天主起初造人时，以奇妙的远见先在人心中植入了这些情感，使人能够接受德行，如同大地能够接受耕作；祂将恶习的素材置于情感中，将德行的素材置于恶习中。因为确实，如果那些使德能得以显示或存在的事物缺乏，德行将不存在，或不在运作。现在让我们看看那些完全消除情感的人达到了什么效果。关于那四种他们认为产生于对善恶事物的看法的情感，他们以为通过根除这些情感，智者的心灵就能得到医治；由于他们明白这些情感是天生植入的，没有这些情感，什么也不能被推动，什么也不能做成，他们便将某些其他事物放在它们的位置和空缺上：他们用倾向代替欲望，仿佛渴望善比倾向善好不到哪里去；他们同样用欢喜代替喜乐，用谨慎代替恐惧。但在第四种情感上，他们无法替换名称。因此，他们完全消除了悲伤，即灵魂的悲哀和痛苦，这是不可能做到的。因为如果瘟疫蹂躏了他的国家，或敌人推翻了它，或暴君压碎了它的自由，谁能不悲伤？如果他目睹了自由的颠覆，或邻人、朋友或善人的放逐或最残酷的屠杀，谁能不悲伤？——除非某人的心灵受到如此打击，以致所有感觉都离他而去。因此，他们要么应当完全消除所有情感，要么应当完成这个有缺陷且薄弱的讨论；也就是说，应当将某种东西置于悲伤的位置，因为既然前面那些情感已如此安排，这一种自然会随之而来。\par

因为正如我们因现时的善事而欢喜，我们也因恶事而烦恼和悲伤。因此，如果他们因认为喜乐是恶的而给它另一个名称，那么同样，因认为悲伤也是恶的，理当给悲伤另一个名称。由此可见，他们所缺乏的不是对象本身，而是一个词语；由于缺乏这个词，他们违反自然的允许，想要消除那种最伟大的情感。因为我本可以更详细地驳斥这些名称的更改，并表明许多名称附着于同一对象，是为了美化文体和增加其丰富性，或者至少它们彼此之间并没有很大差别。因为欲望起源于倾向，谨慎来源于恐惧，喜乐不过是欢喜的表达。但让我们假设它们是不同的，正如他们自己所主张的那样。因此他们会说，欲望是持续不断的倾向，喜乐是过度表现自己的欢喜；恐惧是过度的谨慎，超出了节制的界限。于是，他们认为应当消除的那些事物，实际上并没有被消除，而是被调节了，因为名称改变了，事物本身却保留了。因此，他们不知不觉地回到了逍遥学派通过论证所达到的那一点：即恶习既然不能被消除，就应当通过节制来调节。因此，他们错了，因为他们未能达到他们所追求的目标，而是绕了一条漫长而崎岖的弯路，回到了同一条路径上。\par

% structure: p0073 source=h2 target=subsection reason=relative-heading-rank; base=h2

\subsection{第十六章 论情感，并驳斥逍遥学派关于情感的意见；情感的正确使用与滥用。}

但我认为，逍遥学派甚至没有接近真理，他们承认情感是恶习，但通过节制来调节它们。因为我们必须摆脱即使是适度的恶习；是的，更确切地说，应当从一开始就做到没有恶习。因为没有任何东西天生是恶的；但是，如果我们滥用情感，它们就成为恶习；如果我们善用它们，它们就成为德行。然后必须说明，应当被调节的是情感的原因，而不是情感本身。他们说，我们不应过度欢喜，而应适度、有节制地欢喜。这就像他们说，我们不应快速奔跑，而应从容行走。但行走的人可能出错，奔跑的人可能走在正路上。如果我证明有一种情况，其中不仅欢欢喜喜适度，甚至最小程度的欢喜也是恶的；而相反，另有一种情况，其中甚至狂喜也毫无过错，那该怎么办呢？那么，请问，这种中庸对我们有何益处？我问，他们认为智者是否应该在他看到敌人遭遇不幸时欢喜；或者，如果通过征服敌人或推翻暴君，他的同胞获得了自由和安全，他是否应该抑制他的欢喜？\par

没有人怀疑，在前一种情况下，稍微高兴一点，在后一种情况下，高兴得太少，都是极大的罪行。对于其他\OriginalTerm{情感}{拉丁文 affectus}，我们也可以这样说。但是，正如我所说，智慧的对象不在于调节这些情感本身，而在于调节它们的原因，因为它们是从外部被激发的；也不应该约束它们本身，因为它们即使程度很小也可能伴随着最大的罪恶，而程度很大也可能毫无罪恶。相反，它们应当被限定在特定的时间、环境和地点，这样当我们被允许正确使用它们时，它们就不会是恶习。因为正如行走在正确的道路上是对的，而偏离它是恶的，同样，被情感驱动去做正确的事是善的，而被驱动去做败坏的事是恶的。因为\OriginalTerm{感官欲望}{拉丁文 concupiscentia}，如果它不偏离其合法的对象，尽管它很强烈，却仍是无罪的。但如果它渴望非法的对象，尽管它很温和，却仍是极大的恶习。因此，愤怒、欲望或被情欲激动本身并不是疾病；但暴怒、贪婪或放纵才是疾病。因为暴怒的人即使对于不应该发怒的人或在不应该发怒的时候也发怒。贪婪的人甚至渴望那些不必要的东西。放纵的人甚至追求那些被法律禁止的东西。全部问题应当围绕这一点：既然这些\OriginalTerm{冲动}{拉丁文 impetus}的猛烈无法被约束，也不应当被约束，因为它们为了维持生活的职责而被必然植入，那么它们应当被引导到正确的道路上，在那里甚至可以在不绊倒和没有危险的情况下奔跑。\par

% structure: p0076 source=h2 target=subsection reason=relative-heading-rank; base=h2

\subsection{第十七章 论情感及其运用；论忍耐以及基督徒的至善。}

但是我在反驳他们的欲望中被带得太远了；因为我目的是要表明，那些哲学家认为是恶习的东西，非但不是恶习，甚至是伟大的美德。关于其他，为了教导起见，我将选取那些我认为与主题最为密切的。他们把恐惧或畏惧视为极大的恶习，并认为它是心灵的极大软弱；其对立面是勇敢：如果一个人拥有勇敢，他们说就没有恐惧的位置。那么有谁相信，正是这种恐惧可以是最高度的\OriginalTerm{刚毅}{拉丁文 fortitudo}呢？绝不。因为本性似乎不允许任何事物回归到其对立面。但是，我并非像\Person{苏格拉底}在\Person{柏拉图}著作中所做的那样，通过巧妙的结论，迫使与他争辩的人承认他们所否认的事物，而是以简单的方式，表明最大的恐惧就是最大的美德。没有人怀疑，恐惧痛苦、贫困、流放、监禁或死亡是胆小和软弱心灵的一部分；如果有谁不惧怕这一切，他就被判断为最勇敢的人。但是\OriginalTerm{敬畏天主}{拉丁文 timor Dei}的人不受这一切事物的恐惧。证明这一点无需论证：因为对天主的敬拜者所施加的惩罚，在各时代都已目睹，并且至今仍在全世界被目睹，在折磨他们时，人们发明了新的和不寻常的酷刑。因为回忆起各种死亡方式，当野蛮怪兽的屠杀甚至超越了死亡本身时，心灵会退缩。但是一种幸福且不可战胜的\OriginalTerm{忍耐}{拉丁文 patientia}，在毫无呻吟中忍受了这些可憎的肢体撕裂。这种美德使所有人、各省份、甚至使酷刑者自己（当残忍被忍耐征服时）感到极大的惊异。但这种美德不是由别的，而是由对天主的敬畏所引起。因此（如我所说），恐惧不应被连根拔除，如斯多葛派所主张，也不应被约束，如逍遥派所希望，而应被引导到正确的道路上；并且应当除去忧虑，只留下这一个：因为既然这是唯一合法和真实的畏惧，只有它能做到使其他一切都不被畏惧。欲望也被算作恶习；但如果它渴望地上的事物，它就是恶习；反之，如果它渴望天上的事物，它就是美德。因为渴望获得正义、天主、永生、永远的光明以及天主应许给人类的一切事物的人，将会藐视这些财富、荣誉、权力和统治本身。\par

斯多葛派或许会说，对于获得这些事物，倾向是必要的，而不是欲望；但事实上，倾向并不足够。因为许多人虽有倾向，但一旦痛苦侵入要害，倾向便退让，而欲望却坚持；如果欲望使别人所追求的一切事物在他眼中都成为可鄙的，那便是最大的德行，因为它是自律之母。因此，我们更应当努力做到这一点：使情感得到正确引导，对情感的错误使用便是恶习。因为这些心灵的激动好比一辆套好马匹的战车，驾驭者的首要职责是认路；如果他坚持正道，无论走得多快，都不会撞上障碍。但如果他偏离了路线，即使走得平稳温和，也会在崎岖处颠簸，或滑下悬崖，或至少被带到不需要去的地方。因此，那由情感如同快马牵引的生命之车，若坚持正道，便尽到其职。因此，畏惧和欲望，若倾注于尘世事物，便成为恶习；但若指向天主之事，便成为德行。相反，他们视吝啬为德行；吝啬若是贪求拥有，便不能是德行，因为它完全致力于增加或保存地上的财物。但我们不将至善归于身体，而是仅以灵魂的保存来衡量每一项义务。但是，若如前所述，我们绝不可吝惜财产以保全仁慈与正义，那么节俭便不是德行；这名称在德行的外表下迷惑欺骗。因为节俭诚然是戒绝享乐；但在这一点上它是恶习，因为它源于拥有的爱，而我们应该既戒绝享乐，又绝不可吝惜金钱。因为节俭地使用金钱，即适度使用，是一种心灵软弱的表现，或出于害怕匮乏，或由于无法收回而绝望，或无力轻视尘世之物。但另一方面，他们称不吝惜财产的人为挥霍者。因为他们这样区分慷慨之人和挥霍者：慷慨之人施予配得的人、在适当的时候、施予足够的数量；而挥霍者挥霍给不配的人、在不需要的时候、并毫不顾及自己的财产。\par

那么，我们该当称那因怜悯而给穷人食物的人为挥霍者吗？但是有很大区别：你是因情欲而将金钱花在妓女身上，还是因仁爱而给予悲惨的人；是浪荡子、赌徒和皮条客挥霍你的钱，还是你将其用于虔诚和天主；你是将它花在自己的口腹之欲上，还是存入正义的宝库。因此，正如胡乱花费是恶习，妥善花费便是德行。如果不吝惜财富——这些可被替代——以维持人的生命——这不可替代——是德行，那么节俭便是恶习。因此，我除了称他们为疯狂之人外，别无他名；他们剥夺了人——一种温和合群的动物——其名称；他们根除了情感——人性完全在于情感——并希望使人达到一种心灵不可动摇的麻木状态，同时他们渴望使灵魂摆脱扰动，并如他们自己所说，使之平静安宁；这不仅是不可能的，因为灵魂的力量和本性在于运动，而且甚至不应该如此。因为水如果总是静止不动，便是有害的、更加浑浊；同理，灵魂如果不动且迟钝，连对自己也无用；甚至无法维持生命本身；因为它将既不行事也不思想，因为思想本身无非是心灵的激动。总之，主张这种灵魂不动的人想要剥夺灵魂的生命；因为生命充满活动，而死亡是静止。他们正确地将某些事物视为德行，但他们没有维持其应有的比例。\par

坚毅是一种德行；不在于我们抵抗那些伤害我们的人，因为我们应当向他们让步；至于为何应当如此，我稍后将说明；而在于：当人们命令我们违反天主的法律、违反正义行事时，没有任何威胁或惩罚能阻止我们优先选择服从天主的命令而非人的命令。同样，轻视死亡是一种德行；不是指我们追求死亡，并自愿加害于己，正如许多杰出的哲学家常做的那样——那是邪恶和不敬的行为——而是指当我们被迫背弃天主、背叛我们的信仰时，我们应当选择承受死亡，并以刚毅的精神向那些不能自制之人的愚蠢和盲目的暴力捍卫我们的自由，以不屈的精神挑战世间的一切威胁和恐怖。这样，我们以崇高无敌的思想践踏他人所畏惧的事物——痛苦和死亡。这便是德行；这便是真正的坚毅——唯独在这件事上坚持并保守：没有任何恐怖和暴力能使我们远离天主。因此，西塞罗的这句话是真实的：\Quote{没有人，}他说，\Quote{能是公正的，如果他害怕死亡、痛苦、流放或匮乏。} 塞内卡在他的道德哲学书籍中也说：\Quote{这就是那位有德之人，不以王冠、紫袍或扈从的拥簇著称，却也毫不逊色；他看见死亡临近时，并不像见到新鲜事物那样惊慌；无论是全身受折磨，还是口中烈焰，或是双手被钉在十字架上，他不问自己受什么苦，只问如何受得高尚。} 然而，敬拜天主之人无不恐惧地承受这些事。因此他是公正的。由此可知，凡与独一天主的宗教隔绝的人，根本不能认识或持守德行及其精确界线。\par

% structure: p0081 source=h2 target=subsection reason=relative-heading-rank; base=h2

\subsection{第十八章 论天主的一些诫命，以及论忍耐}

然而，让我们撇开那些哲学家：他们要么一无所知，却把这无知本身当作最大的知识；要么因为他们以为自己知道那些他们实际无知的事，于是荒谬而傲慢地愚昧。因此，让我们（为了回到我们的目的）——唯有我们得到了天主所启示的真理，并从天上获得了智慧——践行那光照我们的天主所命令的事：让我们通过彼此互助，承担并忍受生命的劳苦。然而，如果我们做了任何善工，不要为此追求荣耀。因为天主告诫我们，行义者不应自夸，以免他表现出履行仁爱义务，并非出于遵守天主诫命的渴望，而是为了取悦人，并且已经获得了他所追求的光荣的赏报，而不能再得到那属天和神圣报偿的酬劳。当这些德行被领会时，天主崇拜者应当遵守的其他事情就容易了：即，任何人都不应为了欺骗或伤害而说谎。因为培育真理的人，在任何事上都不可以诡诈，也不可背离他所追随的真理本身。在这正义及一切德行的道路上，谎言没有位置。因此，真正而正直的旅人不会使用\Person{卢奇里乌斯}的话：\par

\Quote{我不应该对一位朋友和熟人撒谎；}\par

但他会认为，即使对敌人和陌生人也不应该说谎；他绝不会让他的舌头——他心灵的译员——与他的情感和思想相冲突。如果他借出任何钱财，他不会收取利息，以使那援助急需的恩惠保持纯粹，并且他可以完全远离他人的财产。因为在这种义务中，他应当满足于自己的所有；既然他的责任在其他方面是不吝惜自己的财产以便行善，但收取多于他所给予的，是不义的。这样做的人，在某种程度上是埋伏着，想从他人的急需中获取掠夺物。\par

但义人不会错过任何仁慈行事的机会：他不会以这种利益玷污自己；反而会如此行动，使得他所借出的，在没有任何损失的情况下，被算作他的善工。他不可接受穷人的馈赠；这样，如果他提供了什么，那便是好的，因为是免费的。如果有人辱骂，他必须以祝福回应；他自己绝不可辱骂，好叫敬畏\OriginalTerm{良善圣言}{the good Word}的人口里不出恶言。此外，他还必须勤勉小心，免得因自己的过失而树敌；如果有人竟如此无耻，对良善义人施加伤害，他必须以平静和节制承受，不为自己复仇，而将此事保留给天主的审判。\BibleRef{罗 12:19} 他必须在任何时候、任何地方保持无罪。这条诫命不仅限于他不应亲自施加伤害，还包括当自己受到伤害时不应报复。因为有一位最伟大而公正的审判者坐在审判座上，祂是万事的观察者和见证者。让他宁愿选择祂来判决他的案件，因为祂的判决无人能因任何人的辩护或恩惠而逃脱。因此，义人成为众人轻视的对象；并且因为人们认为他无力自卫，他将被视为懒惰和无所作为；但如果有人向敌人报了仇，他就被认为是有精神、有作为的人——众人都尊敬他。虽然善人有能力帮助许多人，但他们却仰慕那能伤害人的，而非那能帮助人的。然而，人的邪恶不能败坏义人，使他不再努力服从天主；他宁愿被轻视，只要他始终履行善人的职责，而从不作恶人。\Person{西塞罗}在他那部\WorkTitle{论义务}中说：\Quote{但如果有人想要解开他灵魂中这种模糊的概念，让他立刻教导自己：善人是那帮助他能帮助的人，并且除非被挑衅，不伤害任何人。}\par

唉，他用两个词毁了一个简单而真诚的意见！何必添加这些词：\Quote{除非因受伤而激怒？}他竟可以将恶习如同最可耻的尾巴附在善人身上，并形容他缺乏忍耐——这乃是所有美德中最大的。他说，善人若被激怒就会施加伤害；但正因如此，他若施加伤害，就必然失去善人之名。因为报复伤害与施加伤害同样属于恶人的行为。因为争斗、打架和争执从何而来？岂不正是由于不耐与不义对立，常常掀起巨大风暴吗？但若你用忍耐（没有比这更真实、更配得人的美德）面对不义，它便立即熄灭，如同将水浇在火上。但若那激起对抗的不义遇到了与它相等的急躁，如同浇了油一般，便会燃起如此大火，以致没有水流能将其扑灭，唯有流血才能止息。因此，忍耐的益处何其大，而这位智者却剥夺了善人的这一益处。因为唯有忍耐能使邪恶不发生；若它被赐予众人，人间便再无恶行和欺诈。因此，对善人来说，还有什么比放纵愤怒的缰绳更灾难深重、更违背其品性呢？愤怒不仅剥夺他善人之名，甚至剥夺他作为人的称号，因为正如他自己最真实所说，伤害他人不合乎人的本性。因为若你激怒牛马，它们会用蹄或角反击；蛇和野兽，除非你追捕它们要杀死它们，否则它们不会制造麻烦。再回到人的例子，就连那些没有经验和愚拙的人，一旦受到伤害，就会被盲目无理性的暴怒驱使，竭力报复那些伤害他们的人。那么，智慧善人与邪恶愚者有何区别？若不是前者拥有不可战胜的忍耐（而愚者缺乏此忍耐），若不是前者懂得管理自己、平息自己的愤怒（而后者因无德而不能克制）？但这一情况显然欺骗了他，因为在探究美德时，他认为美德在于各种争胜中得胜。他丝毫无法看见：一个向忧愁和愤怒屈服、沉溺于这些本应与之抗争的情感、并且冲向不义所召唤之处的，并未履行美德的责任。因为努力报复伤害的人，就是想要模仿那伤害他的人。因此，模仿恶人的人绝不能是善人。\par

因此，他用两个词从善人与智者身上夺走了两个最大的美德：纯真与忍耐。但是，正如\Person{萨路斯特}所记\Person{阿庇乌斯}之言：因为他自己行那犬类般的辩才，便希望人也像狗一样生活，以致被攻击时就反咬一口。为了显示这种报复伤害何等有害、它惯常产生何等杀戮，还有什么比这位老师自身的悲惨灾难更合适的例子呢？他本想遵从哲人的这些箴言，却毁灭了自己。因为如果他在受到伤害时保持了忍耐——如果他学会了善人的本分是掩饰和容忍侮辱，而他的急躁、虚荣和疯狂没有倾泻出那些以另一种来源之名标题的高贵演说——他就绝不会因自己的头颅被钉在其上，而玷污了他昔日曾赢得荣耀的讲坛，那场公敌宣告也不会彻底毁灭国家。因此，智慧善人不应渴望争斗、置身危险，因为得胜不在我们掌握之中，每一场争战都是不确定的；而智慧卓越之人的本分是——不渴望除掉对手（这不能没有罪过和危险），而是终结争战本身（这可以既有利又正义地做到）。因此，忍耐应被视为一个极大的美德；而为了使义人能获得此美德，天主——如前所述——愿意他被视为迟钝而受轻视。因为除非他受到侮辱，便无从知道他在克制自己方面有何等的刚毅。现在，若他在受伤激怒之际开始用暴力追击攻击者，他便被战胜了。但若他用理性抑制了那份情绪，他便完全掌握了自我：他能管束自己。这种自我抑制恰当地被称为忍耐，这独一的美德对抗所有的恶习与情感。它将不安而动摇的心灵召回宁静；它缓和、它使人回归自我。因此，既然全然不激起情感既不可能也无益；但在情感爆发为伤害之前，应尽可能及时加以平息。天主命令我们不可让太阳在我们含怒时落下\BibleRef{弗 4:26}，免得祂离去，作我们疯狂的见证。最后，\Person{马尔库斯·图利乌斯}在违背他自己的训诫（我最近曾论及）方面，对忘记伤害给予了最高赞美。\Quote{我怀有希望，}他说，\Quote{啊，\Person{凯撒}，你惯于遗忘一切，唯独不遗忘伤害。}但若他——一个不仅远离天上的正义，也远离公共与民事正义的人——如此行，我们这些可说是堪当永生的人，更当如何行呢？\par

% structure: p0088 source=h2 target=subsection reason=relative-heading-rank; base=h2

\subsection{第十九章 论情感及其运用；以及论三位复仇女神}

当\OriginalTerm{斯多葛学派}{Stoics}试图从人身上根除情感如同疾病时，\OriginalTerm{逍遥学派}{Peripatetics}反对他们，后者不仅保留情感，而且为之辩护，并认为人身上没有一样东西不是出于伟大的理性和远见而产生的。他们如此说确实正确，如果他们知道每一事物真正的限度。因此他们说，愤怒这种情感本身是美德的磨刀石，仿佛没有人能在不激怒的情况下勇敢地与敌人作战；他们由此清楚地表明，他们既不知道什么是美德，也不知道天主为何将愤怒赐给人。如果愤怒赐给我们是为了这个目的，即我们可以用它来杀人，那么还有什么比人更残忍，还有什么比天主为共融和纯洁所造的这个动物更接近野兽的呢？因此，有三种情感驱使人鲁莽地陷入一切罪行：（1）愤怒，（2）欲望，（3）情欲。为此，诗人说有三位\OriginalTerm{复仇女神}{furies}折磨人的心灵：愤怒渴求报复，欲望渴求财富，情欲渴求享乐。但天主为这一切设定了固定的界限；如果它们越过这些界限并开始变得过大，它们必然扭曲本性，变成疾病和恶习。指出这些界限是什么并不需要太大的努力。\OriginalTerm{贪欲}{Cupidity}赐给我们是为了提供生活必需之物；\OriginalTerm{私欲偏情}{concupiscence}是为了生育后代；\OriginalTerm{义愤之情}{indignation}是为了约束在我们权下之人的过错，也就是说，为了使稚嫩的年龄通过更严厉的管教形成正直和公正：因为如果这段生命时期不受恐惧约束，放纵会产生胆大妄为，而这将爆发为各种可耻和大胆的行为。因此，对年轻人运用愤怒既是正义的也是必要的，而对同龄人运用愤怒则既有害又不敬。不敬，因为伤害了人性；有害，因为如果他们反抗，则必须要么毁灭他们，要么自己灭亡。但我所说愤怒之情赐给人的原因，可以从天主自己的诫命中理解，他命令我们不要对辱骂和伤害我们的人发怒，而应始终将手放在年轻人身上；也就是说，当他们犯错时，我们要用不断的鞭打纠正他们，免得因无用的爱和过分的纵容使他们被训练成邪恶并滋养出恶习。但那些不谙世事、不明理性的人，已将那些赐给人用于善途的情感驱逐出去，他们比理性要求的走得更远。因此他们生活得不义和不敬。他们对同龄人运用愤怒：由此产生分歧、放逐和战争，违反正义。他们运用欲望聚敛财富：由此产生欺诈、抢劫和所有种类的罪行。他们只运用情欲享受快乐：由此产生放荡、奸淫和一切腐败。所以，凡是将这些情感限制在适当范围内的人（不认识天主的人做不到），他就是忍耐的，他就是勇敢的，他就是正义的。\par

% structure: p0090 source=h2 target=subsection reason=relative-heading-rank; base=h2

\subsection{第二十章 论感官及其在禽兽与人中的快乐；以及眼目的快乐与表演。}

余下的是要反对五种感官的快乐，并要简短地讨论，因为本书的篇幅要求适度；所有这些快乐，由于邪恶而致命，应当被美德克服和征服，或者，如我之前在论及情感时所说，应召回它们的适当职责。其他动物没有快乐，除了唯一与生殖相关的快乐。因此，它们为自身本性的需要而使用感官：它们观看，以便寻求维持生命所必需的东西；它们彼此听闻，互相辨别，以便能够聚集在一起；它们或通过嗅觉发现，或通过味觉感知，哪些食物有益；它们拒绝并抛弃无用的东西，它们以胃的饱足来衡量饮食之事。然而，最巧夺天工的天主的远见赐给人无限的快乐，且易陷入恶习，因为祂在人面前放置了美德，这美德常与快乐为敌，如同家中的敌人。西塞罗在\WorkTitle{大加图}中说：\Quote{的确，放荡、奸淫和可耻的行为，除了快乐的诱惑之外，别无其他诱因。既然自然或某位神赐给人最卓越的心智，那么没有什么比快乐更与这一神圣的恩惠和礼物为敌了。因为当情欲主导时，就没有节制的余地；当快乐至高无上时，美德也无法存在。}但另一方面，天主为此赐予美德，使它能制服和征服快乐，并且当快乐越过赋予它的界限时，能将其约束在规定的界限内，免得它用享乐安抚和迷惑人，使人受其控制，并以永死惩罚人。\par

源于眼睛的愉悦各式各样，来自于观看那些在人与人交往中、或在自然界或手工制品中令人愉悦的物体。哲学家们正确地夺走了这种愉悦。因为他们说，观看天堂远比观看雕刻的作品更加卓越、更配得上人，去欣赏那件极其美丽的作品——被闪烁的星光如花朵般点缀——胜过欣赏那些绘画、塑造的、用珠宝装饰的东西。但当他们雄辩地劝诫我们轻视尘世之物，并敦促我们仰望天堂时，他们却并不轻视这些公共的表演。因此，他们既以这些为乐，也乐意出席其中；然而，既然它们是恶习最大的诱因，拥有极强地败坏人心的倾向，就应当从我们中间被除去；因为它们不仅对幸福生活毫无贡献，反而造成最大的伤害。因为一个人若以目睹一个虽然被公正定罪的人当众被杀为乐，他的良心所受的玷污，不亚于他成为一场秘密杀人的旁观者和同谋。然而，他们竟称那些流人血的表演为游戏。人性已从这些人身上消失得如此之远，以至于当他们毁灭人的生命时，竟以为自己在以游戏取乐，比所有那些他们以流血为乐的人更为有罪。我现在要问，那些看到人们被置于死亡的刀下、恳求怜悯时，不仅允许他们被杀，而且要求处死他们，并以残酷而不人道的投票赞成他们死亡，不满足于伤口，不以流血为满足的人，还能是正义和虔诚的吗？此外，他们甚至命令，即使受伤倒下了，也要再次攻击，用打击耗尽他们的生命，以防有人以假死欺骗他们。他们甚至对格斗者发怒，除非两人中有一人迅速被杀；好似他们渴求人血，憎恨延迟。他们要求再给他们新一批的格斗者，以便尽快满足他们的眼睛。沉浸在这种实践中，他们丧失了人性。因此，他们甚至不放过无辜者，而是将在杀戮恶人中学到的东西施加于所有人。因此，那些努力持守正义之路的人，不适合成为这种公共杀人的同伴和同谋。因为当天主禁止我们杀人时，祂不仅禁止公开的暴力（这甚至不为公共法律所允许），而且警告我们不要去做那些在人间被视为合法的事。因此，一个义人既不可投身战争（因为他的战争就是正义本身），也不可指控任何人犯下死罪，因为无论你是用言语还是用刀剑杀人，都没有区别，因为被禁止的正是杀人行为本身。因此，关于天主的这条诫命，不应有任何例外；但杀害天主所愿视为神圣动物的人，总是非法的。\par

因此，没有人可以想象，连绞杀新生婴儿也是允许的，这是最大的不敬；因为天主将灵魂吹入他们体内是为了生命，而不是死亡。但人，为了不使任何罪行玷污他们的手，竟剥夺了那些尚属无辜单纯之灵魂的光明，而这光明并非他们自己所赐。事实上，一个连自己的血都不肯放过的人，怎么能期望他克制不去流别人的血呢？但这些人毫无疑问是邪恶不义的。那些被虚假的虔诚迫使遗弃自己孩子的人又怎样呢？他们能被视为无辜吗？他们将自己的亲骨肉扔给狗作为猎物，并且凭自己所能，比绞死他们更加残忍地杀害了他们。谁能否认，一个人若给别人提供怜悯的机会，他本身就是不敬的？因为，尽管他所愿的——孩子被收养——可能发生，但他无疑将自己的骨肉交给了奴役或妓院。但谁不明白、谁不知道，在各性别的情况下，甚至因着错误，会发生或惯常发生什么事呢？因为仅凭一个带有双重罪疚的\Person{俄狄浦斯}的例子就足以说明。因此，遗弃如同杀害一样邪恶。但那些杀害亲子的人确抱怨自己资源匮乏，声称没有足够的能力抚养更多的孩子；仿佛事实上，他们的资源并非掌握在拥有者手中，或者天主没有每天使富人变穷、穷人变富。因此，若有人因贫穷而无法抚养孩子，与其用邪恶之手破坏天主的工程，不如不结婚。\par

如果无论如何都不允许杀人，那么我们也根本不能参与其中，以免任何流血玷污良心，因为这血是为了取悦民众而献的。而且我认为，舞台的败坏影响更为污染。因为喜剧的主题是侮辱处女，或妓女的爱情；而且那些撰写这些可耻事迹的人越是雄辩，就越能通过他们观点的优雅来说服人；和谐而精炼的诗句更容易留在听众的记忆中。同样，悲剧作家的故事将邪恶君王的弑亲和乱伦置于眼前，并表现悲惨的罪行。演员那些不端庄的姿态，除了教导和激发情欲之外，还能产生什么效果呢？他们身体软弱，按妇女的步态和服装变得柔弱，用可耻的姿态模仿淫荡的女人。我何必提哑剧演员呢？他们传授败坏的影响，在假装通奸时教导通奸，通过假装的行动训练人去做真正的事。青年男女看到这些事情毫无羞耻地演出，并被所有人乐意观看时，他们会做什么呢？他们明确地被告知自己能做什么，并被情欲点燃，这种情欲尤其通过观看而激发；每个人根据自己的性别在这些表演中塑造自己。他们嘲笑这些事，同时也赞同它们，罪恶紧随着他们，他们更加败坏地回到自己的房间；不仅是男孩，他们本不应过早地沾染恶习，而且还有老人，他们在这样的年纪犯罪是不合适的。\par

赛车会的活动除了轻浮、虚妄和疯狂之外，还有什么呢？因为他们的灵魂被狂热的激动所冲昏，其冲动与那里进行的战车竞赛一样剧烈；以至于那些为观看表演而来的人，现在他们自己更成为一种表演，当他们开始呼喊、狂喜、从座位上跳起来时。因此，应当避免所有的表演，不仅是为了不让任何恶习扎根于我们应当平静安宁的胸怀中，而且是为了不让任何愉快的习惯性放纵安抚和俘虏我们，使我们偏离天主和善工。因为游戏庆典是敬神的节日，因为它们是因神祇的生日或新庙宇的奉献而设立的。起初，狩猎（即所谓的表演）是为\OriginalTerm{萨图尔努斯}{Saturnus}的荣誉，舞台游戏是为\OriginalTerm{利伯尔}{Liber}的荣誉，而赛车会是为\OriginalTerm{尼普顿}{Neptune}的荣誉。然而，渐渐地，同样的荣誉也开始献给其他神祇，并以其名字设立了专门的游戏，正如\Person{西西尼乌斯·卡皮托}在他的《论游戏》一书中所教导的那样。因此，如果有人出席那些人们为宗教目的聚集的表演，他就离开了对天主的敬拜，转而投向了那些其所庆祝其生日和节日的神祇。\par

% structure: p0096 source=h2 target=subsection reason=relative-heading-rank; base=h2

\subsection{第二十一章 论耳的快乐与神圣文学}

耳的快乐来自声音和曲调的甜美，这确实与我们所说的视觉快乐同样产生恶习。因为谁会不认为那些在家中设有舞台艺术的人是奢侈和卑劣的呢？但独自在家奢侈放纵，还是与民众一起在剧院中，并无区别。然而我们已经谈论了表演；还剩下一样东西必须被我们克服，以免被那些穿透到最内在知觉的事物所俘获。因为所有那些与言语无关的事物，即空中的悦耳声音和琴弦的乐音，可以轻易地被忽视，因为它们不附着于我们，也无法被写下来。但是一首精心创作的诗歌和一段以其甜美诱人的演讲，会俘虏人们的心灵，并随意推动它们。因此，当有学问的人投身于天主的宗教时，除非他们受到某位娴熟导师的指导，否则他们不信。因为习惯于优美和精练的演讲或诗歌，他们轻视圣经简单朴实的语言，认为它粗俗。因为他们寻求能抚慰感官的东西。但是任何悦耳的东西都会说服人，并在愉悦时深深扎根于胸中。那么，天主，既是心灵、声音和舌头的创造者，难道不能雄辩地说话吗？恰恰相反，祂以最大的远见，希望神圣的事物不加修饰，以便所有人都能理解祂向所有人所说的话。\par

因此，渴望真理、不愿自欺的人，必须摒弃有害和致伤享乐，这些享乐会像美味的食物束缚身体一样束缚心灵；真实的事物必须优先于虚假的，永恒的事物优先于短暂的事物，有用的事物优先于愉快的事物。除了你看到那些出于虔诚和正义而行的事之外，不要让任何东西取悦于视觉；除了那些滋养灵魂、使你成为更好的人的事之外，不要让任何东西取悦于听觉。尤其是这个感官不应被扭曲为恶习，因为它被赐予我们正是为了使我们认识天主。因此，如果听旋律和歌曲是愉快的，那么唱和听天主的赞美诗也应是愉快的。这才是真正的快乐，它是德性的伴随者和同伴。它不是脆弱和短暂的，像那些像牲畜一样受身体奴役的人所渴望的那样；而是持久的，毫无间断地提供愉悦。如果有人逾越了它的界限，除了快乐本身之外不再从快乐中寻求任何东西，那么他就是为自己设计死亡；因为正如在德性中有永生，在快乐中也有死亡。因为选择暂时事物的人将失去永恒的事物；喜欢地上事物的人将不会拥有天上的事物。\par

% structure: p0099 source=h2 target=subsection reason=relative-heading-rank; base=h2

\subsection{第二十二章 论味觉与嗅觉的快乐}

但关于味觉和嗅觉的快乐——这两种感官仅与身体相关——我们无需讨论；除非有人要求我们说，一个明智而善良的人如果成为食欲的奴隶，如果浑身涂满香膏、头戴花环而行，那是不光彩的：而做这些事的人显然是愚蠢而无知的，是无价值的，甚至没有一丝德性的观念触及他。或许有人会说：\Quote{那么，这些东西被造出来，除了供我们享用之外，还有什么目的呢？} 然而，人们常常说过，如果没有可以被克服的事物，就不会有德性。因此，天主创造万物，是为了在两者之间提供一场竞赛。那么，这些快乐的诱惑是那一位的工具，它的唯一任务是征服德性，并从人那里驱逐正义。它以这些抚慰和享乐俘虏人的灵魂；因为它知道快乐是死亡的策划者。因为正如天主只通过德性和劳苦召唤人走向生命，同样，那一位却以欢愉和快乐召唤我们走向死亡；正如人通过骗人的邪恶达到真正的善，同样，他们通过骗人的善达到真正的恶。因此，必须警惕那些享乐，如同陷阱或网罗，免得被享乐的温柔所俘虏，我们连同我们所奴役的身体本身一起被置于死亡的权下。\par

% structure: p0101 source=h2 target=subsection reason=relative-heading-rank; base=h2

\subsection{第二十三章 论触觉的快乐与情欲，以及论婚姻与节制}

现在我来讨论那种通过触觉感受到的快乐：这一感官确实遍及全身。但我认为，我应当谈论的不是装饰或衣服，而仅仅是\OriginalTerm{情欲}{libido}，它最需要被约束，因为它为害最烈。当天主设计了两个性别的道理时，祂赋予它们彼此渴慕，并以结合为乐。因此，祂在所有生物的身体中混合了最炽热的\OriginalTerm{欲望}{cupiditas}，使它们极其贪婪地投入这些情感中，并借此使各族类得以繁衍增殖。这种欲望和渴求在人类中更为强烈和炽热；或许是因为天主愿意有更多的人群，又或许是因为祂唯独赐予了人\OriginalTerm{德性}{virtus}，以便在克制享乐和节制自我方面有赞美和荣耀。因此，我们的那位\OriginalTerm{仇敌}{adversarius}知道这种欲望的力量有多大——有些人宁愿称其为必然性——并将它从正当和善的事物转向邪恶和扭曲的事物。因为它灌输非法的欲望，使人玷污他人的（配偶），而他们拥有自己的（配偶）本可以无罪。它用刺激性的形象呈现在眼前，提供燃料，并供给恶习的养料；然后它搅动和激发内心深处的所有刺激，点燃那天然的火焰，直到迷惑并缠绕那人，使他陷入网罗。而且，为了没有人因惧怕惩罚而远离他人的（配偶），它也设立了妓院，并公开了不幸妇女的羞耻，以嘲弄那些行此事的人和必须忍受此事的人。\par

它将这些淫秽之事——那些为圣洁而生的灵魂——如同淹没在泥沼的深渊中，熄灭了羞耻，摧毁了贞洁。它也混淆男性与男性，并图谋了违背天性和天主之命的可憎结合：如此它浸染了人，并武装他们去行一切恶事。对于那些将软弱无力、需要保护的年龄置于自己的\OriginalTerm{情欲}{libido}之下加以摧残和玷污的人，还有什么能是圣洁的呢？这件事实在无法按其罪行之重大来叙述。我只能称这些人为不虔敬者和弑亲者，他们不满足于天主所赐的性别，除非他们亵渎而放肆地嘲弄自己的性别。然而，在他们眼中，这些事是轻微的，几乎是可敬的。我该说什么关于那些实行可憎的——不是情欲，而是\OriginalTerm{疯狂}{insania}——的人呢？我不愿说：但如果我们认为那些行此事的人不羞愧，他们又会怎样呢？然而，因为这事在发生，我必须说。我指的是那些连头也不放过的最可怕的情欲和可咒诅的狂怒的人。我该用什么言辞或什么义愤来追述如此重大的恶行呢？罪行的巨大超过了言语的能力。因此，既然情欲产生这些行为并犯下这些罪行，我们必须以最大的德性武装自己来对抗它。凡不能克制那些情感的人，就当将他们约束在合法婚姻的界限内，以便既满足他急切渴望的，又不至于陷入罪恶。因为那些堕落的人想要什么呢？当然，正当的行为伴随着快乐：如果他们追求快乐本身，他们可以享受正当和合法的快乐。\par

如果某种需要禁止这一点，那么应当运用最大的\OriginalTerm{德性}{virtus}，使\OriginalTerm{节制}{continentia}与\OriginalTerm{欲望}{cupiditas}相抗衡。天主不仅命令我们应当远离他人的身体（这些不可触摸），也远离公开和普通的身体；祂教导我们，当两个身体结合时，就成为一个身体。因此，谁若陷入泥沼，就必然被泥沼玷污；身体固然可以很快洗净，但灵魂被不洁身体的接触所污染，除非经过长时间和许多善工，才能从所沾染的污秽中净化。因此，每个人都应当给自己立下这样的准则：两性的结合是为生育而赐给生物的，这一法则建立在这样的情感上，以使后代延续。正如天主给我们眼睛，不是为观看和获得快乐，而是为看见那些与生活需要相关的行为，同样，我们接受身体的生殖部分——其名称本身也表明——除了为了繁衍后代之外，没有其他原因。我们应当以最大的虔敬服从这一神圣律法。所有自称是天主门徒的人，都应当如此生活和受教，以便能管束自己。因为那些纵情享乐、随从\OriginalTerm{情欲}{libido}的人，是将自己的灵魂交给身体，并判处死亡；因为他们把自己卖给了身体，而死亡对身体有权力。因此，每个人都当尽其所能，养成谦逊，培养羞耻心，以\OriginalTerm{良心}{conscientia}和心智守护贞洁；不仅要遵守公共法律，而且跟随天主法律的人应当超越一切法律。如果他习惯了这些美德，他就会耻于堕落到更坏的事：只要正直和可敬的事物令他喜悦——这些事物对善良之人比邪恶和可耻的事物对恶人更为甘甜。\par

我尚未完全论述贞洁之职；天主不仅将其限于私室之内，且规定于床笫之间；即凡有妻者，不得再另拥女仆或自由女，而应保守婚姻之忠诚。因为，与世俗法律之规定不同，唯有女子，若另有他夫，方为淫妇；而丈夫，即使有多妻，亦不被视为犯奸淫罪。但神圣法律如此将二人结合于婚姻（即成为一体），以同等权利相连，以致凡破坏此身体之合一而分裂之者，皆被视为奸淫者。天主使其他动物在受孕后抗拒雄性，唯独使女子忍受男子，别无他故；以免女子抗拒时，情欲迫使男子另求他女，如此则不能持守贞洁之荣耀。但女子亦不能获得端庄之德，若她不能犯罪。因为，谁能称那受孕后抗拒雄性的哑畜为端庄呢？它之所以如此，是因为若允许，则必陷入痛苦与危险。所以，不做你不能做之事，并无赞美可言。端庄之所以在人身上受赞美，是因为它不是天生，而是出于自愿。因此，双方应彼此持守忠诚：更应以节欲之榜样教导妻子，使她行为端庄。因为，你若要求他人履行你自己不能做到的事，便是不义。这种不义确实导致了奸淫，因为女子难以忍受自己守信而对方不报以相互的爱德。总之，没有哪个如此丧失羞耻心的淫妇，不以这理由为自己的恶行辩护，说她犯罪并非加害，而是报复。\Person{昆体良}对此表达得极好：\Quote{人，既不节制他人的婚姻，又不守护自己的婚姻，这两者在本质上是相连的。因为丈夫若专注于败坏他人的妻子，便无暇顾及家中的神圣；而妻子若落入这样的婚姻，便会被榜样所刺激，要么以为自己是在效仿，要么以为自己是在报复。}\par

因此，我们必须谨慎，以免因我们的放纵而给恶行提供机会；而应使二人的性情彼此适应，以平等之心负轭。我们应在对方身上思想自己。因为正义之首要大体在于：你不愿别人对你做的事，也不对别人做。这是天主关于节欲所吩咐的。但为免有人以为可以规避天主的诫命，又加上这些：凡娶被休之妻者，为奸淫；凡除奸淫之故外休妻另娶者，也为奸淫；因为天主不愿身体分离拆散。此外，不仅应避免奸淫，也应避免思想；不可注视他人之妻而心中贪恋；因为若心灵自己描绘了快乐的图像，心灵便成了淫妇。因为犯罪的确在于心；心以思想拥抱无节制的情欲之果；罪在于此，一切过犯在于此。因为即使身体未受任何玷污，若心思不洁，端庄之理仍不成立；若情欲已玷污了良心，便不能视为清白贞洁。也不可有人认为，给情欲套上缰绳，将其纳入贞洁端庄的界限是困难的，因为人甚至能战胜它，且许多人保持了幸福无玷的身体完整，也有许多人极其快乐地享受着这种属天的生活方式。天主吩咐如此行，并非像强制规定（因为人应当生育），而是像允许。他知道他给这些情感加上了多大的必然性。\Quote{若有人能行此事，他说，他将获得非凡无比的赏报。}这种节欲之德仿佛是顶峰，是一切德行的圆满。若有人能奋力攀登至此，主必认他为仆人，师傅必认他为门徒；他将在世上得胜，他将相似于天主，因为他得了天主的德能。这些事看似艰难，但我们所谈论的人，是踏灭一切尘世之物、预备登天之路的人。因为德行在于认识天主，所有事在你不知时都是沉重的；一旦你知道，便变得容易：我们这些趋向至善的人，必须通过这些困难本身走出去。\par

% structure: p0107 source=h2 target=subsection reason=relative-heading-rank; base=h2

\subsection{论忏悔、宽恕及天主的诫命}

然而，不要让人因着被情欲征服、被贪欲驱使、被错误欺骗、或被强力强迫而转向不义的途径，就灰心失望，或对自己绝望。因为他能够被带回，并被释放，只要他悔改自己的行为，转向更好的事物，并向天主做补赎。诚然，西塞罗认为这是不可能的，他在《学园派》第三卷中的话是：\Quote{但是，如果在那些在旅途中迷失了道路的人的情况下，允许那些走上歧途的人通过悔改来纠正他们的错误，那么纠正轻率就会更容易。} 这完全是允许的。因为如果我们认为我们的孩子在我们察觉他们悔改自己的过失时得到纠正，并且尽管我们曾剥夺他们的继承权、抛弃他们，我们却再次接纳、抚爱、拥抱他们，为什么我们应当怀疑我们天父的仁慈能再次因悔改而被平息呢？因此，那位既是主又是最慈爱的父亲的天主应许，祂将赦免悔改者的罪，并涂抹那重新开始实践正义者的一切过犯。因为正如过去生活的正直无益于那生活败坏的人，因为后来的邪恶摧毁了他的正义行为；同样，先前的罪也不会阻碍那改善生活的人，因为后来的正义已抹去了先前生活的污点。因为那对自己所做的事悔改的人，明白他先前的错误；因此希腊人更好且更有意义地称之为\OriginalTerm{悔改}{\GreekText{μετάνοια}}，这在拉丁语中我们可以说成是回归正确的领悟。因为那为他的错误而忧伤的人，回归了正确的领悟，并仿佛从疯狂中恢复了他的心智；他责备自己的疯狂，并将他的心智锚定在更好的生活方式上；那时他特别警惕此事，以免再陷入同样的陷阱。简言之，甚至哑巴动物，当它们被诡计诱捕时，如果它们设法解脱出来而逃脱，就会变得对未来更加谨慎，并且总是避免它们曾察觉其中有诡计和陷阱的一切事物。悔改使人谨慎勤勉，以避免那他曾因受骗而一度陷入的过错。\par

因为没有人能如此明智且谨慎，以至于从不会失足；因此天主，知道我们的软弱，出于祂的怜悯，为人类打开了一个避难的港口，使得悔改的良药能帮助我们的脆弱所面临的这种需要。所以，若有人犯了错，让他回头，并尽快恢复和改正自己。\par

\Quote{但要重新向上攀登，\newline{} 进入光明的白日，\newline{} 那时劳苦的重担便临到。}\par

因为当人们品尝了导致毁灭的甜蜜快乐时，他们很难脱离它们；如果他们从未品尝过它们的诱惑，他们本会更容易追随正确的事物。但如果他们从这种有害的奴役中挣脱出来，他们所有的错误都将被宽恕，只要他们以更好的生活纠正了自己的错误。且不要让人以为，若他的过错没有证人，他就占了便宜：因为万事都知晓于那位我们在他面前生活的那一位；即使我们能向所有人隐藏某事，我们也无法向天主隐藏它，对祂而言，没有什么能被隐藏，没有什么能是秘密。塞内卡以一句卓越的格言结束了他的劝诫：\Quote{存在，} 他说，\Quote{某位伟大的神祇，比所能想象的更伟大；我们正是为他而努力生活。让我们向他证明自己。因为良心得到坚定并无益处；我们暴露在天主的目光之下。} 还有什么比一个不认识真宗教的人所说的，更真实地由认识天主的人说出的话呢？因为他表达了天主的威严，说它太大，人的心智无法容纳；他也触及了真理的源泉，看出人的生命并非如伊壁鸠鲁派所主张的那样是多余的，而是他们努力为天主而生活——如果他们确实正义且虔诚地生活。如果曾有人向他指明天主，他本可成为天主的真崇拜者；如果他获得了一位真正的智慧导师，他肯定已蔑视了芝诺和他的老师索提翁。让我们向他证明自己，他说。这话实在像天国的话，只是前面已有无知的承认。良心被约束并无益处；我们暴露在天主的目光之下。因此，没有谎言的余地，没有伪装的余地；因为人的眼睛被墙壁阻隔，但天主的全能力量不能被体内的部分阻隔，以致无法看透并认识整个人。同一位作家在同一部作品的第一卷中说：你在做什么？你在筹划什么？你在隐藏什么？你的守护者跟随你；一个人因旅行与你分离，另一个因死亡，另一个因健康不佳；这个紧紧跟随着你，你永不能没有他。你为什么选择一个隐秘的地方，并移开证人？假设你已成功地避开了所有人的注意，愚蠢的人啊！没有一个证人，对你有什么益处，如果你有自己良心的见证呢？\par

而且图利以同样非凡的方式论及良心和天主说：\Quote{让他记住，}他说，\Quote{他有一位见证者，就是天主，即依我判断，是他自己的心，天主赐予人者，再没有比这更神圣的了。}同样，在论及公正良善之人时，他说：\Quote{因此，这样的人不仅不敢做出，甚至不敢想任何不敢公开宣称的事。}所以，让我们洁净我们的良心，它向天主的眼目是敞开的；并且，正如同一位作者所说，\Quote{让我们常如此生活，以至于记得我们将要交账；}并让我们时时刻刻想到，我们被注视着，并非如他所说，在世界的某个舞台上被人注视，而是从上头被那位将要做审判者也是见证者注视，祂在要求交我们生命的账时，不许任何人否认自己的行为。因此，要么逃避良心，要么我们自己自愿敞开我们的心，撕裂我们的伤口，倾倒出毁灭；这些伤口除了那一位，无人能医治，祂曾使瘸子行走，使盲者复明，洁净了污秽的肢体，并使死者复活。祂将熄灭欲望的炽热，祂将根除情欲，祂将移除嫉妒，祂将平息愤怒。祂将赐予真正而持久的健康。所有人都当寻求此治疗，因为灵魂比身体面临更大的危险，应当尽快对隐疾施治。因为若有人眼睛清亮，肢体健全，身体极其健康，我仍不称他为健全，若他被愤怒驱使，被骄傲膨胀鼓胀，是情欲的奴隶，被欲望燃烧；但我更应称他为健全，若他不抬眼羡慕他人的顺遂，不羡慕财富，以贞洁的眼看待他人之妻，一无所贪，不觊觎他人之物，不嫉妒任何人，不轻视任何人；他谦卑、慈悲、慷慨、温和、有礼：和平永恒地居留于他的心灵中。\par

那人是健全的，他是正义的，他是完美的。因此，凡遵守这一切天上诫命者，便是对真天主的敬拜者，其祭献是心灵的温和、无罪的生命和善行。而凡彰显这一切美德者，他在每次履行任何良善虔敬的行为时，都献上祭献。因为天主并不要求畜牲的祭献，也不要求死亡与流血，而是要求人与生命。此祭献既不需要神圣的枝丫，也不需要净化仪式，更不需要草皮块，这些东西显然极其虚妄，只需要那些从内心深处发出的东西。因此，在天主的祭台上——它真是极其伟大，安放在人心之中，不能被血玷污——放置着正义、忍耐、信德、无罪、贞洁和节制。这是最真实的礼仪，这是天主的法律，如西塞罗所称，辉煌而神圣，它常命令正直和可敬之事，禁止错谬和可耻之事；凡遵守此至圣而确实的法律者，必能公正合法地生活。我已陈述了此法律的几个主要点，因为我曾许诺只讲那些构成德行与正义特质的事。若有人愿意包含其他所有部分，让他从源头本身——那溪流流向我们的源头——去寻求吧。\par

% structure: p0114 source=h2 target=subsection reason=relative-heading-rank; base=h2

\subsection{论祭献及堪当献于天主的供物，并论赞美天主的仪式。}

现在让我们简要论一下祭献本身。\Quote{象牙，}柏拉图说，\Quote{并非献给天主的洁净供物。}那么什么呢？绣花和昂贵的织物？不，凡可朽坏或可被偷走的，都非献给天主的洁净供物。但既然他看到这一点，即从尸体取下的东西不应当献给活人，他为何没有看到，有形的供物不应当献给无形者？塞内加说得多么好、多么真实：\Quote{你当视天主为伟大、温和、以温柔尊严可敬、且常临在的朋友；祂不该以宰杀牺牲和大量流血来敬拜——因为屠杀无辜动物有何喜乐？——而应以纯洁的心灵和良善可敬的意图。不应用高高堆起的石头为祂建造殿宇；祂应由各人在自己心中祝圣。}因此，若有人认为衣裳、珠宝及其他被视为珍贵之物为天主所看重，他便完全不知天主为何，因为他以为那些东西能取悦祂，连人因其轻视都会受赞扬。那么，什么是纯洁、什么堪当于天主，除了祂自己在神圣法律中所要求的？\par

有两件事应当奉献，即礼品与祭献：礼品作为常献，祭献则作为临时之献。但在那些丝毫不理解神性本质的人看来，凡以金银制成的，以及以紫布和丝绸织成的，均可称为礼品；祭献则是牺牲以及一切在祭台上焚烧之物。但天主既不使用前者也不使用后者，因为祂不受朽坏的影响，而这些都是完全可朽坏的。因此，在每种情况下，都应将无形之物献给天主，因为祂接受这样的奉献。祂的礼品是灵魂的无辜；祂的祭献是赞美与颂歌。因为既然天主是不可见的，就应当以不可见的事物来敬拜祂。所以，除了那由德行与义德构成的宗教之外，没有其他真正的宗教。但天主如何对待人的义德是容易理解的。因为如果人将成为义的，并领受不朽性，他将永远服事天主。然而，人若非为义德而生，这一点古代的哲学家以及连西塞罗也隐约感到怀疑。因为他在讨论法律时说：\Quote{但在有识之士所讨论的一切事物中，最重大的无疑莫过于完全理解我们乃是为义德而生的。}因此，我们应当向天主献上并且摆上的，正是祂生出我们为要我们领受的那唯一之物。然而，这双重祭献如何真实，\Person{赫尔墨斯·特里斯墨吉斯忒斯}是一位合适的见证者，他在事实上和在言语上都同意我们，即同意我们所追随的先知们。他这样论及义德说：\Quote{儿子啊，你要敬拜和尊崇这道。}但天主的敬拜在于一件事：不作恶。同样，在那完美的讲论中，当他听到\Person{阿斯克勒庇俄斯}由他的儿子询问是否要把馨香和其他供神祭献的香料献给其父时，他喊道：\Quote{讲吉言啊，阿斯克勒庇俄斯。因为对于那至善者怀有这样的想法是最为不敬的。因为这些东西以及类似的东西并不适合于祂。祂充满万有，凡存在的一切，祂一无所缺。然而，让我们感谢祂，并敬拜祂。因为祂的祭献唯独在于祝福。}他说得对。\par

因为我们应当以言语祭献天主；既然天主就是圣言，正如祂自己所承认的。因此，敬拜天主的主要礼仪是出自义人之口的对天主的赞美。然而，为使这赞美蒙天主悦纳，需要极大的谦卑、敬畏和虔诚，以免有人自恃于自己的正直与无辜，从而陷入骄傲与傲慢的过失，并因此行为丧失其德行的赏报。但为获得天主的宠幸，并免于一切污点，让他不断呼求天主的仁慈，不求别的，只求罪过的赦免，即使他没有罪。若他渴望别的事物，无需用言语向那位知道我们心愿的祂表达；若有什么好事临到他，让他感恩；若有灾祸，则让他补赎，并承认灾祸是因他的过失而临到他；即使在灾祸中，他也应照样感恩，在好事中补赎，好使他在任何时候都保持一致，坚定、不变、不动摇。他不要认为这仅应在圣殿中实行，也应在家里，甚至在他的床上。总之，让他常与天主同在，在心中尊奉祂为圣，因为他自己就是天主的殿。但如果他以这样的勤勉、顺从和虔诚服事了天主，他的父和主，义德就圆满和完善了；凡持守这义德的，正如我们先前所见证的，他已服从了天主，并履行了宗教和他自身职责的义务。\par

\SourceRecord{Translated by William Fletcher.；Ante-Nicene Fathers；Vol. 7.；Edited by Alexander Roberts, James Donaldson, and A. Cleveland Coxe.；Buffalo, NY: Christian Literature Publishing Co.,；1886.；<http://www.newadvent.org/fathers/07016.htm>.}

% structure: p0001 source=h1 target=section reason=page-title

\section{\WorkTitle{神圣训谕，第七卷（论幸福生活）}}

% structure: p0002 source=h2 target=subsection reason=relative-heading-rank; base=h2

\subsection{第一章 论世界，以及即将相信和不信的人；并在此对无信者的谴责。}

很好：根基已经奠定，正如这位杰出的演说家所说的。但我们不仅奠定了根基，这些根基坚固而适于支撑这工程；而且我们建造了整座大厦，以巨大而坚固的建筑物，几乎达到了顶端。剩下的是一件更为容易的事情，要么覆盖它，要么装饰它；然而，没有它，之前的工程既无用又令人不快。因为，无论是摆脱假宗教而认识真宗教，有何益处？无论是看见虚假智慧的虚妄，还是知道什么是真的，有何益处？我说，捍卫那属天的正义有何益处？以巨大的艰难持守对天主的崇拜，而这是最大的德行，除非永远福乐的天主赏报随之而来，又有何益处？我们必须在本书中谈论这个主题，免得之前的一切显得徒劳无益：如果我们把为之而着手这些事的那个目标留在不确定之中，恐怕有人偶然认为如此大的劳苦是徒然而行；他怀疑天上的赏报，这赏报是天主为那轻视现世甘甜享乐、宁愿选择孤独无酬的德行的人所指定的。让我们也通过圣经的见证和可能的论证来满足我们主题的这一部分，使得同样清楚地显明，将来的事应当优先于现在的事；天上的事优先于地上的事；永恒的事优先于暂时的事：因为恶行的赏报是暂时的，德行的赏报是永恒的。\par

因此，我将阐述世界的体系，以便容易理解它是什么时候以及如何被天主创造的；\Person{柏拉图}，他曾谈论世界的创造，却既不能知道也不能解释，因为他不知晓那神圣的奥秘，这奥秘只有通过先知和天主的教导才能学到；因此他说世界是为永恒而造的。然而，情况大不相同，因为凡是具有坚实沉重之形体的事物，既然在一定时间获得开端，也必然有终结。因为\Person{亚里士多德}，当他看不出如此巨大的事物如何能消亡，并希望逃避这一反驳时，说世界始终存在，并将永远存在。他完全没有看到，凡存在的物质事物必定在某个时候有过开端，并且没有任何事物可以存在除非它有一个开端。因为当我们看到土、水、火消亡、消耗、熄灭时，它们显然是世界的部分，就可以理解，其部分是可朽的，那整体也就是可朽的。因此，任何易于朽坏的事物必然是产生的。但一切眼睛所见之物都必然是物质的，并且能够分解。因此，唯有\Person{伊壁鸠鲁}跟随\Person{德谟克利特}的权威，在这件事上说了真话，他说世界在某时有过开端，也将在某时灭亡。然而，他也不能给出任何理由，无论是通过什么原因或什么时间，如此巨大的工程会被毁灭。但是，自从天主向我们启示了这一点，而我们不是通过猜测，而是通过来自天上的教导达到认识，我们将仔细教导它，以便最终向那些渴望真理的人显明，哲学家们并未看见也未能理解真理；他们对真理的认识是如此肤浅，以至于他们完全不知道那如此令人愉快和悦意的智慧芬芳是从何处吹向他们的。\par

与此同时，我认为有必要提醒那些将要阅读本书的人：堕落和邪恶的心智，由于他们心智的敏锐被世俗的情欲所迟钝，这些情欲压垮一切感知并使之软弱，他们将要么完全不能理解我们所叙述的这些事情，要么即使理解了，也会假装并不愿意它们为真；因为他们被恶行所吸引，并且明知故犯地偏向自己的邪恶，被其愉悦所俘获，而他们离弃了德行的道路，因其苦涩而恼怒。因为那些被贪婪和某种对财富的不知餍足的渴望所燃炽的人——当他们卖掉或挥霍了他们所喜爱的东西之后，他们无法过简朴的生活——无疑更偏好那迫使他们放弃急切欲望的东西。同样，那些被情欲的刺激所驱使的人，正如诗人所说，\par

\Quote{冲入疯狂与烈火，}\par

说我们提出的是明显不可信的事情；因为关于自我节制的诫命刺伤了他们的耳朵，这些诫命禁止他们享受欢乐，而他们已将灵魂连同身体都交给了欢乐。但是，那些因野心膨胀或因权力之爱而炽热，将全部精力用于获取荣誉的人，即使我们将太阳本身捧在手中，也不会相信那命令他们藐视一切权力和荣誉、并生活在谦卑中，且如此谦卑以至于能忍受伤害，若受了伤也不愿报复的教导。就是这些人闭着眼睛，以任何方式大声反对真理。但那些心智健全的人，即没有如此沉溺于恶行以至于无可救药的人，将会相信这些事，并乐意接近它们；我们所说的一切，在他们看来都将是敞开、明了、简单的，并且最重要的是，真实而不可反驳。\par

无人赞助德行，除非他能追随德行；但并非所有人都容易追随它：只有那些被贫穷和匮乏锻炼过、并能胜任德行的人才能如此。因为如果忍受苦难是德行，那么那些始终享受美好事物的人就不能胜任德行；因为他们从未经历过苦难，也无法忍受苦难，因为他们长期习惯和渴望那些他们唯一知晓的美好事物。因此，贫穷和卑微的人，他们无所牵挂，比那些被众多障碍缠累的富人更容易相信天主；更确切地说，他们被束缚和镣铐所奴役，听从情欲——这个女主人——的示意，情欲用解不开的绳索诱捕了他们；他们无法仰望天堂，因为他们的心思低垂向地，固定在地上。但德行之路上不容许背负沉重负担的人。正义引领人通往天堂的小径非常狭窄；除非他无牵挂且轻装上阵，否则无人能坚持此路。因为那些富足的人，背负着众多且沉重的负担，行走在死亡之路上，那路非常宽阔，因为毁灭统治着广大的领域。天主为正义所赐的诫命，以及我们在天主的教导下关于德行与真理所提出的内容，对他们来说是苦涩的，如同毒药。如果他们敢于反对这些，他们必须承认自己是德行和正义的敌人。现在我要进入主题的剩余部分，以便结束这部著作。但剩下的是，我们应当论述天主的审判，这审判将在我们的主返回大地、按各人的功过给予每个人赏报或惩罚时确立。因此，正如我们在第四卷中论述了祂的首次来临，同样在本卷中我们将叙述祂的第二次来临，这也是犹太人所承认并盼望的；但那是徒劳的，因为祂必须为了那些祂曾为召唤而来的人的困惑而返回。因为那些在祂卑微时以不敬对待祂并施以暴力的人，将在祂以征服者的身份显能时经历祂；并且，由于天主的报应，他们将遭受他们所读却不理解的一切；因为，他们被各种罪恶玷污，且被圣者的血所洒，被那同一位置他们于不义之手的人注定要受永罚。但我们将有单独针对犹太人的主题，在其中我们将证明他们的错误和罪责。\par

% structure: p0009 source=h2 target=subsection reason=relative-heading-rank; base=h2

\subsection{第二章　论哲学家的错误、天主的智慧以及黄金时代。}

现在让我们教导那些对真理无知的人。至高天主的安排已如此决定：这个不义的时代，在完成其指定时限的历程后，应当结束；并且，一切邪恶立刻被消灭，义人的灵魂被召回幸福的生命，一个宁静、安宁、和平的，简言之，如同诗人所称的黄金时代，将在天主亲自统治下繁荣。这正是哲学家所有错误的特别原因：他们没有理解世界的体系，而这体系包含全部智慧。但仅凭我们自身的感知和天生的智力，无法理解它，而他们希望在没有老师的情况下自行做到。因此他们陷入了各种且常常矛盾的意见中，无法从中脱身，\par

并且他们停留在同样的泥沼中，\par

正如喜剧作家所说，因为他们的结论与假设不符；因为他们假设为真实的事物，在缺乏真理和属天事物的知识下，是无法被肯定和证明的。而这种知识，正如我多次说过的，除非来自天主的教导，否则无法存在于人里面。因为如果一个人能理解属天事物，他也能行出来；因为理解就如同跟随它们的踪迹。但他不能做天主所做的事，因为他披戴着可朽的身体；因此他甚至不能理解天主所做的事。是否可能做到这一点，每个人都可以从天主的行为和工作的宏大中轻易衡量。因为如果你默观世界及其包含的一切，你必定会明白，天主的工程远超人的工程。因此，正如天主工程与人工程之间的差距有多大，天主的智慧与人的智慧之间的距离也必定有多大。因为天主是不朽不坏的，是永生的，因此是完美的，因为祂是永恒的，祂的智慧也如同祂自身一样是完美的；没有任何事物能反对它，因为天主自身不受任何事物的支配。\par

但因为人受制于情欲，他的智慧也容易出错；正如许多事物阻碍人的生命，使其不能持久，同样，他的智慧也必须被许多事物所阻碍：以至于它在完全感知真理方面并不完美。因此，没有人的智慧，如果它试图凭借自身达到对真理的概念和认识；因为人的心智，与脆弱的身体相连，并被封闭在黑暗的居所中，既不能自由漫游，也不能清楚地感知真理，而真理的知识属于天主本性。因为祂的工程只有天主知道。但人不能通过反思或辩论获得这种知识，而是通过学习和聆听那唯一能够知道和教导的那一位。因此，\Person{马尔库斯·图利乌斯}从\Person{柏拉图}那里借用了\Person{苏格拉底}的思想，苏格拉底曾说他自己离开生命的时刻已到，但那些在他面前辩护的人仍然活着，他说：哪样更好，只有不死的神明知道；但我认为没有人知道。因此，所有哲学家的派别都必须远离真理，因为建立这些派别的是人；那些没有任何神圣声音的宣告所支持的东西，不可能有任何基础和稳固性。\par

% structure: p0014 source=h2 target=subsection reason=relative-heading-rank; base=h2

\subsection{第三章　论自然与世界；以及对斯多葛派和伊壁鸠鲁派的谴责。}

既然我们正在谈论哲学家的错误，斯多葛派将本性分为两部分：一部分是作用者，另一部分是提供自身可作用于其上的。他们说前一部分含有全部知觉能力，后一部分是质料，且两者不能独立作用。如何能说作用者与被作用者是同一回事呢？若有人说陶工就是泥土，或泥土就是陶工，他岂不是显然疯了？但这些人在本性的同一名称下包含了两种截然不同的事物：天主与世界，造物主与作品；并说两者缺一不可，仿佛天主与世界在本性中混合在一起。因为他们有时如此混淆它们，以至天主自身便是世界的心灵，世界便是天主的身体；仿佛世界与天主同时存在，而天主未曾亲自创造世界。他们自己在别的时候也承认这一点，当他们说世界是为人的缘故而被造的，并且天主若愿意，可以没有世界而存在，因为天主是神圣而永恒的心灵，分离且脱离身体。既然他们无法理解他的能力和威严，便将他与世界，即与他自己的作品混为一谈。\Person{维吉尔}的这句诗便由此而来：——\par

\Quote{一种精神，其天火\newline{}在框架每一部分中闪耀，\newline{}并搅动伟大的整体。}\par

那么，他们自己的说法——世界是由天主眷顾所创造并治理的——又成了什么呢？因为若他创造了世界，则他在世界之先就已存在；若他治理世界，则显然不是像心灵治理身体那样，而是像主人管理房屋，舵手管理船只，御者管理战车。然而，他们并未与所治理的事物相混合。因为如果我们所见的一切都是天主的肢体，那么天主就被它们弄得毫无感觉，因为肢体没有感觉；而且可朽，因为我们看见肢体是可朽的。\par

我可以列举多少土地因突然震动而开裂或骤然下陷；多少城市和岛屿被海浪吞没而沉入深渊；沼泽淹没肥沃的平原，河流和池塘干涸；山峰要么陡峭崩塌，要么夷为平地。许多地区和多座山的地基被暗中藏匿的火焰毁灭。而这还不够，如果天主不宽恕自己的肢体，除非人也被允许对天主的身体拥有某种权力。大海被修建，高山被削平，大地的深处被挖掘以提取财富。为何我还说，我们甚至耕地时也无不在撕裂天主的身体？以致我们对天主的肢体施暴时，同时既是邪恶的又是不敬的。那么，天主会忍受自己的身体被折磨，甘心削弱自己，或允许人这样做吗？除非那与世界及世界各部分混合的神圣理智放弃了地球的最初表面，隐入了最深处，以便从持续的撕裂中感觉不到痛苦。但如果这是琐屑而荒谬的，那么他们自己也像那些没有领悟到神圣精神无所不在、万物靠它维系的人一样缺乏理智——但并非以那种方式，即不可朽坏的天主自己与沉重而可朽的质料混合。因此，他们从\Person{柏拉图}那里得来的说法更为正确：世界是由天主创造的，也由天主眷顾治理。所以，\Person{柏拉图}以及持同一意见的人应当教导并解释如此伟大作品的设计之原因和理由：他为什么或为了谁而创造它。\par

但斯多葛派也说世界是为了人的缘故而被造的。我听到了。但\Person{伊壁鸠鲁}却不知道人本身是为何或由谁创造的。因为\Person{卢克莱修}说世界不是由诸神创造的，他如此说道：\par

\Quote{再说，为了人的缘故，他们愿意安排世界光荣的本性}——\par

然后他引述：——\par

\Quote{这纯粹是愚蠢。因为我们的感激能给予不朽有福者什么好处，以至于他们为了我们的缘故要着手安排任何事？}\par

这有充分的理由。因为他们没有提出任何理由说明人类为何或由天主所创造与建立。阐明世界和人的奥秘是我们的职责；而他们由于缺乏这些，既不能到达也不能看见真理的圣所。因此，正如我刚才所说，当他们假定了一个真实的前提——即世界是由天主创造的，且是为了人的缘故而创造——却又因随后的推理失败，最终无法维护他们所假定的。总之，\Person{柏拉图}为使天主的作品不软弱易毁，说它将永远存留。如果世界是为了人的缘故被造，且被造成永恒，那么为何因它而被造的人本身不是永恒的呢？如果因之被造的人是可朽的，那么世界本身也必是可朽的、会瓦解的，因为它并不比那些为它而被造的人更有价值。但如果他的推理是一贯的，他就会明白，世界必会朽坏，因为它是被造的；除了那不能被触摸的，没有什么能永远存留。\par

但是，如果有人说世界并非为人的缘故而被创造，他是没有论据的。因为如果他说造物主为自己而设计了如此宏伟的作品，那么我们为何被造？我们为何享受世界本身？人类以及其他生物的创造有何意义？我们为何中断他人的利益？简而言之，我们为何成长、衰退和死亡？我们自身的生产意味着什么？永久的延续又意味着什么？无疑天主愿意我们被看见，并塑造出带有祂自身各种形象印记的受造物，使祂能以此自娱。尽管如此，如果真是这样，祂会将生灵视为祂的眷顾对象，尤其是人——祂使万物都服从于人的统治。至于那些说世界永恒存在的人：我暂且不提世界没有某个起始就不可能存在这一点——他们无法从中摆脱；但我这样说：如果世界永恒存在，它就不可能有系统性的安排。因为在任何事物被做成或安排之前，需要先有谋划以决定如何做成；没有既定计划的预见，什么也不能做成。因此，计划先于每一项工作。所以，未曾被造的事物就没有计划。但世界拥有计划，它据此而存在并受治理；因此它也是被造的：如果它被造，也将会被毁灭。那么，如果可能，就让他们指出一个理由，说明它起初为何被造，或者将来为何被毁灭。\par

由于\Person{伊壁鸠鲁}或\Person{德谟克利特}未能教导这一点，他们便说世界是自发产生的，种子从四面八方汇聚；当这些种子再次分解时，就会出现混乱和毁灭。于是他们扭曲了自己正确看到的东西，并且由于对体系的无知，完全推翻了整个体系，把世界和其中所发生的一切都降格为一场最微不足道的梦幻——如果人类事务中没有计划存在的话。但是，既然我们看到世界及其所有部分都按奇妙的计划治理；既然天穹的架构、星辰与天体运行的和谐（即使在多样性中也保持和谐）、四季恒常而奇妙的秩序、大地的丰饶多样、广阔的平原、高山防御与堆积、森林的翠绿与多产、泉水的极有益处的喷涌、河水的适时泛滥、大海丰富而充足的涌入、相反且有益的风的吹拂，以及万物，都以最大的规律性固定下来：谁还能如此盲目，以至于认为它们是没有原因而被造的呢？在这一切中闪耀着最富有远见的天主眷顾安排之奇妙秩序。因此，如果没有任何东西是无因而存在或发生的；如果至上天主的天主眷顾从事物的秩序中显明，祂的伟大从它们的宏大中显明，祂的权能从它们的治理中显明：那么声称没有天主眷顾的人就是迟钝和疯狂的。如果他们的目的是为了肯定一位神的存在而否认诸神的存在，我不会不认同；但当他们这样做的意图是声称没有神时，凡不认为他们是无知的人，他自己就是无知的。\par

% structure: p0026 source=h2 target=subsection reason=relative-heading-rank; base=h2

\subsection{第四章　万物都被造为有用，甚至那些看似邪恶的事物也是如此；人是为何在如此软弱的身体中享有理性的？}

但是我们在第一卷中已经充分讨论了天主眷顾的主题。因为如果天主眷顾存在——正如从它奇妙的工作性质所显示的那样——那么必然是同一天主眷顾创造了人和其它动物。因此，让我们看看人类被造的原因是什么，既然正如斯多葛学派所说，世界是为人的缘故而被造的——尽管他们在这个问题上本身就有不小的错误，说世界不是为某一个人的缘故，而是为许多人的缘故而被造。因为单个个体的名称已经包含了全体人类。但这源于他们不知道只有一个人是由天主创造的，并且他们认为人们在所有土地和田地里像蘑菇一样生长出来。但\Person{赫耳墨斯}并非不知道人是由天主所造，且是按照天主的肖像造的。但我回到我的主题上来。我想，没有任何事物是为其自身而被造的；但凡被造的事物，必然是为某个目的而被造的。因为谁会如此无知或如此漠不关心，以至于试图随意做某件事，而从其中不期望任何效用、任何益处呢？建造房屋的人，并非仅仅为了房屋是房屋而建造，而是为了让人居住。建造船只的人，并非仅仅为了船只是可见的而付出劳动，而是为了让人能在其中航行。同样，设计并制作任何器皿的人，并非仅仅为了显示他做了此事，而是为了使制成的器皿能够容纳某种必需之物供使用。同样，其他一切被造之物，显然不是多余地，而是为了某些有用的目的而被造的。\par

因此，显然世界是由天主创造的，并非为了世界本身；因为世界没有知觉，它既不需要太阳的温暖、光、风的呼吸、雨水的湿润，也不需果实的滋养。甚至不能说天主为了自己的缘故创造了世界，因为天主没有世界也可以存在，正如世界受造之前祂已存在；而且天主本身并不使用其中所包含和产生的一切事物。所以，显然世界是为了生物而建造的，因为生物享用其所构成的一切；为了使它们生存和存在，一切必需之物都在固定的时节供应。再者，其他生物是为了人而被造的，这一点很清楚：它们服从于人，被赐予人作为保护和服役；因为无论陆地还是水中的生物，都不像人那样感知世界的系统。对此我们必须答复哲学家，尤其是\Person{西塞罗}，他说：\Quote{天主既然为我们创造了万物，为什么还要造那么多的毒蛇和蝮蛇？为什么要在陆地和海洋散布那么多有害的东西？}这是一个非常广泛的讨论主题，但必须简要地触及，如同顺便提及。既然人由不同且对立的元素构成，即灵魂与身体，也就是天与地、轻灵与可感知、永恒与暂时、有知觉与无感觉、光明与黑暗，理性本身和必然性要求将善与恶都摆在人面前——善事供他使用，恶事供他防备和避免。\par

因为智慧被赐予人，正是为此：认识了善与恶的本质，他就可以运用理性的力量追求善而避免恶。因为其他动物没有赐予智慧，它们既有天然的外衣保护，也有爪牙武装；但代替这一切，天主赐予了人最卓越的礼物，仅仅是理性。因此，祂造人赤身裸体、毫无武装，好使智慧成为他的防御和遮盖。祂将防御和装饰安置在内里而非外面，不在肉体而在内心。因此，如果没有恶事供他防备，并使他与善的和有用的事物区分，智慧对他就没有必要。所以，让\Person{西塞罗}知道：理性之所以赐予人，要么是为了让人捕鱼以供己用，并且为了自身安全而避开毒蛇和蝮蛇；要么是因为善与恶被摆在他面前，因为他已领受了智慧，而智慧的全部力量就在于分辨善恶。因此，伟大、正确且令人赞叹的是人的力量、理性及能力，为了人，天主创造了世界本身和一切存在之物，并赐予他如此大的尊荣，使他凌驾于万物之上，因为唯有人能欣赏天主的化工。因此，我们的\Person{阿斯克勒庇俄斯}在他写给我的那本书中讨论至高天主的\OriginalTerm{天主眷顾}{providence}时，说得极好：\Quote{正因如此，任何人都可以有充分的理由相信，\OriginalTerm{天主眷顾}{providence}将最接近自己的位置赐给了那能够理解其安排的人。那太阳就是如此：谁观看它以致明白它为何是太阳，以及它对系统其他部分有何影响？这是天，谁仰望它？这是地，谁居住其上？这是海，谁航行其上？这是火，谁使用它？}因此，至高天主安排这些并非为了自己，因为祂一无所缺，而是为了人，使人能恰当地使用它们。\par

% structure: p0030 source=h2 target=subsection reason=relative-heading-rank; base=h2

\subsection{第5章 论人的创造、世界的安排与至善}

现在我们说明天主为何造人本身。因为如果哲学家们知道这一点，他们或者会坚持那些他们发现为真实的事情，或者不会陷入最大的错误。因为这是首要之事；这是一切的关键。如果有人不拥有此点，真理就完全从他那里溜走。简言之，正是这一点使他们与理性相悖；因为如果这一点曾光照他们，如果他们知道了人之全部奥秘，学园就绝不会与他们的辩论和全部哲学完全对立。因此，正如天主不是为了自己而创造世界，因为祂并不需要世界的益处，而是为了使用世界的人，同样，祂造人本身也是为了祂自己。\Person{伊壁鸠鲁}说，人在天主身上有什么益处，使得祂为自己而造人呢？确实，为的是有人能理解祂的工程；能以其理智欣赏，并以言语表达在其安排中所显示的上智，其受造之物秩序，其完成过程中所施展的能力。这一切的总和就是，人应当崇拜天主。因为理解这些事的人就崇拜祂；他以应有的敬畏追随祂，视祂为万物的创造者，祂是他真正的父亲，祂按其工程的发明、起始和完成来衡量祂威严的卓越。还有什么更明显的证据可以证明天主既为了人而创造世界，也为自己而造人，比这一点更清楚呢：在一切生物中，唯有人被如此塑造，以致他的眼睛朝向天空，他的面容朝向天主，他的面容与他的创造者共融，以至于天主似乎以张开的手将人从地面举起，并将他提升到默观祂本身。\Quote{那么，}他说，\Quote{人所献上的崇拜，对那已受祝福、一无所缺的天主有什么裨益呢？或者，如果天主赐予人如此荣誉，以至为他创造世界，赋予他智慧，使他成为一切活物之主，并爱他如同儿子，为什么又使他屈从于死亡与腐朽？为什么将祂所爱的对象暴露于一切邪恶？当时人应当幸福，如同与天主紧密相连，并如天主一样永恒，因为人正是为崇拜和默观祂而受造的。}\par

虽然我们在前几卷中已零散地传授了这些道理的大部分，但因主题现在特别要求如此（因为我们承担了讨论幸福生活的主题），所以我们要更仔细更充分地解释这些事，使天主的安排、祂的工程和旨意可得知晓。尽管祂总能够凭祂自己的不朽之灵产生无数的灵魂（就如祂产生天使一样，天使拥有不朽而无任何邪恶的危险和恐惧），但祂设计了一项不可言喻的工程：祂如何可能创造无限数量的灵魂，这些灵魂起初与脆弱而软弱的身体结合，祂将他们置于善恶之间，祂将德性摆在他们面前（他们由两种本性组成），好使他们不能通过娇嫩而轻易的生活道路达到不朽，而是要以极大的困难和辛劳达到永生的不可言喻赏报。因此，为了用沉重且易受损伤的肢体包裹他们（因为他们无法在中间虚空中存在，身体沉重向下沉降），祂决定首先为他们建造一个住所和居处。如此，以不可言喻的能力和权能，祂创造了世界超绝的工程；将轻的元素悬于高处，将重的元素压于深处，坚固了天上的事物，奠定了地上的事物。目前无需逐一细述每一点，因为我们在第二卷中已一起讨论过了。\par

因此祂在天上安放了光体，它们的规律、光亮和运行，都极适当地配合众生的益处。此外，祂赐给大地（祂设计为其居所）孕育和产生万物的丰饶，好按各物之性质与需求供给养分。然后，当祂完成了属于世界状态的一切事物时，祂用地本身造了人，这地祂从起初就为人预备作居所；即，祂用属地的身体包裹并覆盖了人的灵，使之由不同且对立的质料组成，人可能易于感受善恶；正如大地本身孕育五谷，从地取来的人体也得到了生育后代的能力。因为他由脆弱之物形成，不能永远存在，当他属世生命的空间过去时，他可能离去，并通过永恒的后裔更替他所承载的脆弱而软弱的身体。那么，既然祂为那人建造了世界，为何使他脆弱而必死呢？首先，为产生无数的生灵，并使祂以众多充满全地；其次，为将德性（即对邪恶和劳苦的忍耐）摆在人面前，人以此能够赢得不朽的赏报。因为人由两部分组成：身体和灵魂，其中一为属地，一为属天，所以两重生命被赐予人：一为暂时的，为身体而设；另一为永恒的，属于灵魂。我们在出生时接受了前者；我们通过奋斗达到后者，为使人不朽不能毫无困难地存在。那属地的生命如同身体，因此有终结；但这属天的生命如同灵魂，因此没有限制。我们在无知时接受了第一重生命；第二重生命则是明知地接受；因为它是赐给德性而非本性，因为天主愿意我们在生命中为自己获得生命。\par

为此，祂赐予我们现世生命，好让我们或借恶行失去那真正而永恒的生命，或借德行赢得它。至善并不包含于这形体生命中，因为正如这生命是出于天主的必然而被赐予我们，同样也将出于天主的必然而毁灭。因此，有终结者不包含至善。但至善包含于我们凭自身获得的精神生命中，因为它不能包含恶，也不会有终结；对此，自然和身体的结构提供了一个论证。因为其它动物俯向地面，因为它们是属土的，且不能享有来自天上的不朽；但人是直立的，仰望上天，因为不朽已为他预备；然而，除非由天主赐予人，这不会到来。否则，义人与不义之人就没有区别了，因为每个出生的人都将成为不朽的。所以，不朽不是天性的结果，而是德行的赏报和酬劳。最后，人并非出生时立即直立行走，而是起初四肢爬行，因为他的身体和现世生命的性质是我们与哑兽所共有的；之后，当他的力量增强时，他直起身来，他的舌头得以舒展，使他说话清晰，他不再是哑兽。这个论证教导我们，人出生是可朽的；但他后来成为不朽的，当他开始按照天主的旨意生活时，即追随义德，这包含于恭敬天主中，因为天主使人仰望上天和祂自己。这发生在当人藉天上的圣洗而净化，将他的幼稚连同过去生活的一切污秽一并除去，并接受天主的德能增加，成为完美成全的人。\par

因此，因为天主在人面前提出了德行，尽管灵魂与身体相连，却是相反且彼此对抗的。对灵魂有益的事物，对身体有害，即远离财富、禁止享乐、蔑视痛苦和死亡。同样，对身体有益的事物，对灵魂有害，即贪欲和情欲，由此财富被贪求，各种享乐被纵容，灵魂被削弱和毁灭。因此，义人和智慧之人必然要经历一切恶事，因为勇敢战胜邪恶；而不义之人则在财富、荣誉、权力中。因为这些好处属于身体，且是属地的；这些人也过着属地的生活，无法获得不朽，因为他们已将自己交付给了享乐，而享乐是德行的敌人。因此，这暂时生命应服从那永恒生命，如同身体服从灵魂。所以，谁若偏好灵魂的生命，就必须轻视身体的生命；除非他轻视最低的事物，否则他无法以任何其他方式追求那最高的事物。但谁若拥抱身体的生命，并将其欲望转向属地，就无法达到那更高的生命。但谁若宁愿为永恒而善生，就将在时间上过恶的生活，并在世上经历一切烦恼和劳苦，以便获得神圣和天上的安慰。而谁若宁愿在时间上过善的生活，就将永远过恶的生活；因为他因偏爱属地的好处胜过天上的好处，而被天主的判决判定受永罚。因此，天主希望受人恭敬，并被当作圣父光荣，好使人拥有德行和智慧，而只有德行和智慧产生不朽。因为除祂自己外，无人能赐予那不朽，因唯独祂拥有不朽，祂将赐予人的虔诚——人以之恭敬天主——这份赏报：永远蒙福，并永远在天主面前，与天主共融。\par

\EditorialNote{以下章节的段落等在许多手稿中缺失，非常不确定是否由\Person{拉克唐修}撰写。}\par

也没有人能借故归咎于那位既造善又造恶的那一位。因为为什么祂愿意恶存在，如果祂憎恶它？为什么祂不只造善，使无人犯罪、无人作恶？虽然我在几乎前几卷书中都解释过这一点，并已在上面有所触及，尽管简略，但仍须反复提及，因为整个问题取决于这一点。因为除非祂造了相反的事物，否则不可能有德行；善的能力除非通过与恶的对比，否则根本无法显现。因此，恶无非是善的解释。所以，如果恶被除去，善也必须被除去。如果你砍掉左手或脚，你的身体将不完整，生命本身也将不复存在。因此，为了身体结构的适当调整，左肢与右肢最适当地结合。同样，如果你使棋子全都相同，没有人会玩。如果你只给竞技场一种颜色，没有人会认为值得观看，一切竞技赛的乐趣都被剥夺了。因为最初设立比赛的人是偏爱一种颜色的；但他引入另一种颜色作为对手，以产生竞争和某种偏袒性。同样，天主在确定何者为善并赐予德行时，也指定了与之争斗的相反者。如果缺少敌人和战斗，就没有胜利。除去竞争，德行本身也一无所有。人与人之间的相互竞争何其多，而且以各种不同的技艺进行！然而，如果他没有可以与之竞争的对手，就没有人会被认为在勇敢、速度或卓越上超群。而如果缺乏胜利，则胜利的光荣和赏报也必须随之而去。因此，为了使德行本身通过不断的练习而得到加强，并使之在与恶的斗争中臻于完美，祂两者都赐予，因为两者之中缺少任何一个都无法保持其力量。因此，存在差异，整个真理体系依赖于差异。\par

我不忽视更熟练者在这里可能提出的反对意见。如果善不能没有恶而存在，你怎能说，在第一个人得罪天主之前，他仅生活在善的实践中，或者今后他将仅生活在善的实践中？这个问题应当由我们来审查，因为在之前的几卷中我略过了它，以便在此处补全题目。我们上面已说过，人的本性由对立元素构成；因为身体出于尘土，是可把握的、暂时的、无感觉的、黑暗的；但灵魂出于天上，是非物质的、永恒的、有感觉的、充满光明的；因为这些品质彼此对立，人必然受制于善与恶。善归于灵魂，因为灵魂不能被分解；恶归于身体，因为身体脆弱。因此，既然身体与灵魂相连结合，善与恶必然结合在一起；除非它们（身体与灵魂）分离，否则彼此不能分开。最后，善与恶的知识同时赐给了第一个人；他一旦明白了这一点，就立刻被逐出那没有恶的圣地；因为当他只与善交往时，他并不知道这本身就是善。但在他获得了善恶知识之后，他就不许再留在那个福乐之地，而被放逐到这个普通的世界，以便他立刻体验到那两种事物的本性，因为他已经同时认识了它们。因此显然，智慧赐给了人，使他能辨别善恶，区分有利与不利、有用与无用，使他能判断并思虑应当防备什么、渴望什么、避免什么、追随什么。因此，智慧不能没有恶而存在；而人类的第一位始祖，当他只与善交往时，如同婴儿一般生活，对善恶一无所知。然而，确实，将来人必须既智慧又幸福，毫无罪恶；但这不能在灵魂仍被身体的居所包裹时发生。\par

但当身体与灵魂分离时，恶也就与善分离；正如身体消逝而灵魂长存，恶将消逝而善将永存。那时，人接受了不朽的衣袍，将如同天主一样，智慧而远离邪恶。因此，那希望我们只与善交往的人，特别渴望我们脱离身体而活，因为恶存在于身体中。但如果恶被除去，那么如我所言，或智慧、或身体将从人那里被夺去：智慧，使他不知恶；身体，使他不感觉恶。然而如今，既然人既有智慧以认识，又有身体以感知，天主愿意两者在此生中并存，以便德行与智慧相一致。于是，他将人置于两者之间，使人有自由选择追随善或恶。但他又在恶中掺入一些看似善的东西，即各种令人愉悦的享受，以便通过这些诱惑将人引向隐藏的恶。同样，他在善中掺入一些看似恶的东西，即艰难、痛苦和辛劳，借其严酷和不悦，使灵魂受挫而退缩，远离隐藏的善。但此时需要智慧的作用，使我们更多地用心灵而非身体去看，这是极少数人能做到的；因为德行艰难罕见，而享乐却普遍常见。因此，必然发生的是，智人被看作愚人：他追求看不见的善事，却让看得见的从手中溜走；他躲避看不见的恶事，却陷入眼前的恶事。当我们为了信德而不拒绝酷刑或死亡时，此事就发生在我们身上，因为我们被逼到最大的邪恶，以致背叛信德、否认真天主，并向死亡和带来死亡的诸神献祭。这就是为什么天主使人成为可朽的，并使他遭受各种苦难，虽然天主曾为人创造了世界：为了使人能够有德行，并且他的德行能为他赢得不朽的赏报。如今，如我们已指出的，德行就是崇拜真天主。\par

% structure: p0040 source=h2 target=subsection reason=relative-heading-rank; base=h2

\subsection{第六章 为何创造世界与人。崇拜假神何等无益。}

现在，让我们用一个简短的定义来概括整个论证。世界被创造是为了这个目的：我们可以出生；我们出生是为了这个目的：我们可以认识世界和我们的创造者——天主；我们认识他是为了这个目的：我们可以崇拜他；我们崇拜他是为了这个目的：我们可以获得不朽作为我们劳作的赏报，因为崇拜天主包含最大的劳苦；为了这个目的，我们被赏以不朽，以便与天使相似，永远服事至高之父和主，并成为天主的国度直到永远。这就是万物的总和，这就是天主的奥秘，这就是世界的奥秘；那些追随当下满足、投身于追求地上和脆弱之物、并通过致命的享乐将他们那本为追求天上事物而生的灵魂沉入泥潭和淤泥的人，与此奥秘疏远了。\par

现在，让我们接着探究：敬礼这些神祇是否合理？因为如果神很多，如果人们之所以敬拜他们，只是为了让他们赐予财富、胜利、尊荣以及一切只在今世有益的事物；如果我们被造毫无缘由——如果人的被造没有天主的眷顾——如果我们出于偶然为自己而生、为了自己的享乐而生——如果死后我们归于无有——那么，有什么比人事乃至世界本身更多余、更空虚、更虚妄呢？尽管世界具有难以置信的宏大，并以如此奇妙的秩序构建而成，却仍被琐事占据。因为风的吹动为何要使云移动？为何要闪现闪电、轰鸣雷声、降下雨露，使大地结出果实并滋养各种产物？总之，为何整个大自然都要辛劳，使维持人生命的事物一无所缺，如果这一切都是虚妄的，如果我们彻底消亡，如果我们身上没有对天主更大的益处？但若说这些话是说不出口、不可想象的：你所见极符合理性的事，却不是出于某个重要理由而建立的——那么，在这些堕落宗教的错误和哲学家的见解（他们想象灵魂会消亡）中，又能有什么理由呢？确实没有。因为他们能说什么理由，说明神祇为什么如此准时地按时供给人一切？难道是为了让我们向他们献上谷物、酒、香料和牲畜的血？这些事不能为不死者所悦纳，因为它们是会朽坏的；也不能对没有身体的存在有用，因为这些是用来给有身体者使用的。然而，如果他们需要这些东西，他们本可以随时赐予自己。因此，无论灵魂消亡还是永远存在，敬礼神祇涉及什么原则？世界是由谁建立的？人为何、何时、多久、多远被造？或出于什么原因？他们为什么出生、死亡、接续、更新？神祇从那些死后不再存在的人的敬拜中得到什么？他们做了什么事、承诺了什么、威胁了什么，配得上人或神？或者，如果灵魂死后仍存，他们（神）对他们做什么或将要做什么？他们需要灵魂的宝库做什么？他们自己从何而来？他们如何、为何、从何处而来这么多？因此，如果你离开我们上面所包含的总体事物，一切体系都被摧毁，万物归于无有。\par

% structure: p0043 source=h2 target=subsection reason=relative-heading-rank; base=h2

\subsection{第七章 论哲学家及其真理的多样性。}

由于哲学家们不理解这一关键点，他们便无法理解真理，尽管他们大多看见并解释了该关键点本身所包含的那些事物。但不同的人以不同的方式提出了所有这些见解，并未将原因、结果和理由联系起来，以便他们能够联合并完成包含整体内容的关键点。但要表明哲学家和学派几乎分有全部的真理，却是容易的。因为我们并不像学园派惯常所做的那样推翻哲学；他们的做法是反驳一切，这更像是诽谤和嘲笑；但我们表明，没有一个学派如此偏离正轨，也没有一个哲学家如此虚妄，以至于看不到一些真理。然而，当他们疯狂地渴望反驳时，当他们即使虚假也要辩护自己的论点、即使真实也要推翻他人的论点时，真理不仅从他们那里逃走了，而且他们自己主要因自己的过错失去了它。但若有人能将分散在个人中、散落在各学派中的真理收集起来，并归纳成一个整体，他肯定与我们一致。然而，除非他对真理有经验和知识，无人能做到这一点。但认识真理只属于那受天主教导的人。因为他无法以其他方式拒绝虚假的事物、选择和认可真实的事物；但即使他偶然做到了这一点，他也最肯定地扮演了哲学家的角色；尽管他不能以神圣的见证来辩护这些东西，但真理会凭自身的光解释自己。因此，那些人的错误令人难以置信：他们一旦认可某个学派并投身其中，就谴责所有其他学派为虚假和虚妄，并武装起来准备争战，既不知道应该辩护什么，也不知道应该反驳什么；他们不加区分地攻击那些与他们意见不同的人所提出的一切。\par

由于他们这些极其顽固的争论，没有任何哲学更加接近真理，因为全部真理已被这些人分别包含。柏拉图说世界是由天主创造的；先知们也这样说；西彼拉的诗句也表明同样的事。因此，那些说万物是自发产生的或来自原子集合的人是错误的，因为如此伟大、如此装饰且如此巨大的世界，若非有一位最精巧的作者，既不能被造成，也不能被安排和整理，而且那种使万物得以维系和治理的秩序本身，也表明了一位具有最精巧心智的工匠。斯多葛派说世界及其中的一切都是为了人的缘故而造的；圣经也教导我们同样的事。因此，德谟克利特是错的，他认为万物像虫子一样从地里涌出，没有任何作者或计划。因为人被造的理由属于一个神圣的奥秘；由于他无法知道这一点，他把人的生命贬低为虚无。阿里斯通断言人生来是为了实践德行；我们也从先知那里得到提醒并学到同样的事。因此，亚里斯提卜受骗了，他使人受制于快乐，即邪恶，如同禽兽。斐瑞塞德斯和柏拉图主张灵魂是不朽的；但这正是我们宗教中的独特教义。因此，狄凯亚科斯连同德谟克利特都错了，他们论证灵魂与身体一同消亡并消散。斯多葛派的芝诺教导有阴间，善人的居所与恶人分离；前者享有宁静愉快的区域，后者在黑暗之地和可怕的泥沼深渊中受罚；先知们也显示同样的事。因此，伊壁鸠鲁错了，他认为那是诗人的虚构，并把所说的那些阴间惩罚解释为发生在今生。因此，哲学家们触及了全部真理和我们神圣宗教的每个奥秘；但当别人否认时，他们无法捍卫他们发现的东西，因为其体系与具体细节不一致；他们也未能像我们上面所做的那样，将他们认为真实的东西归纳为一个概要。\par

% structure: p0046 source=h2 target=subsection reason=relative-heading-rank; base=h2

\subsection{论灵魂不朽}

因此，唯一的至善是不朽，为了领受不朽我们最初被造和出生。我们向此目标前进；人性关注此目标；德行将我们提升至此目标。既然我们发现了这个善，我们还应谈论不朽本身。柏拉图的论证虽对此主题大有贡献，但对于证明和充实真理力量薄弱，因为他既没有总结并汇集这整个伟大奥秘的计划，也没有把握至善。虽然他感知到灵魂不朽的真理，但他并未将其作为至善来谈论。因此，我们能以更确定的标记引出真理；因为我们不是通过可疑的猜测收集它，而是通过神圣的教导认识它。柏拉图如此推理：凡凭自身感知且永远运动的，是不朽的；因为那没有运动开端的东西，也不会有终结，因为它不能被自身遗弃。但这个论证甚至会使哑巴动物也获得永恒存在，除非他通过增加智慧来区分。因此，为了避免这种普遍联系，他补充说：人的灵魂不能不成为不朽的，因为其发明之奇妙技巧、思考之敏捷、感知和学习之迅速、对过去的记忆和对未来的预见，以及对无数技艺和科目的知识——这些其他生物不具备——显得神圣和属天；因为能够构想如此伟大之物并包含如此伟大之物的灵魂，在地上找不到起源，因为它没有与任何地上的混合物结合。人身上沉重且易朽的部分必须归于尘土；而轻巧精微的部分是不可分的，当它从身体的居所（如同监狱）中释放后，它飞向天堂及其本性。这是柏拉图学说的简要总结，在他自己的著作中被广泛而详尽地阐述。\par

毕达哥拉斯此前也持同样见解，他的老师斐瑞塞德斯也是，西塞罗报告说他是第一个谈论灵魂不朽的人。虽然所有这些人以口才出众，但至少在这场争论中，那些反对此意见的人也有不少权威：首先是狄凯亚科斯，然后是德谟克利特，最后是伊壁鸠鲁；以至于他们所争论的事情本身变得可疑。最后，图利乌斯在陈述了所有这些关于不朽和死亡的见解后，宣称他不知道真理是什么。他说：\Quote{这些意见中哪一个是真的，或许某位神能看见。} 他在另一处又说：\Quote{既然每种见解都有最博学的捍卫者，就无法推测出确定性。} 但我们不需要推测，因为神性本身已向我们揭示了真理。\par

% structure: p0049 source=h2 target=subsection reason=relative-heading-rank; base=h2

\subsection{论灵魂不朽与德行}

因此，藉着这些论证——既非\Person{柏拉图}，也非其他任何人所发明——灵魂的不死得以证明并被认知；我们将简要汇集这些论证，因为我的论述急于讲述天主的伟大审判，这审判将在世界临近终结时在地上举行。首先，由于天主不能被人类看见，以免有人因而想像天主不存在，因为祂未被肉眼看见，在其它奇妙的安排中，祂也造了许多事物，其力量显明可见，但实体却看不见，如声音、气味、风，为的是藉这些事物的标记和例子，我们可从祂的能力、运作和工程中感知天主，尽管祂不落于我们眼前。有什么比声音更清晰，比风更强劲，比气味更强烈呢？然而，这些，当它们经空气传播而达于我们的感官，并以它们的能力推动感官时，并不为视觉所辨别，而是由身体的其他部分所感知。同样，天主不是通过视觉或其他脆弱的感官被我们所感知，而是由心灵之眼来注视，因为我们看见祂辉煌奇妙的工程。至于那些完全否认天主存在的人，我不仅不应称他们为哲学家，甚至不应称他们为人，因为他们与哑巴动物极为相似，只有身体，心灵毫无辨别力，把一切归于肉体感官，认为只有他们亲眼所见之物才存在。又因他们看见恶人遭遇不幸，或善人得享昌盛，便相信万事皆由命运主宰，世界由自然建立，而非出于上智安排。\par

由此他们立刻陷入与此类观点必然伴生的荒谬之中。但如果有一位无形无体、不可见而永恒的天主，那么灵魂既不可见，在离开身体后，便不会消亡，这是可信的；因为显然存在某种东西，它感知并充满活力，却不显现于视觉。但有人说，灵魂若没有包含感知机能的身体部分，如何能保持感知，这是难以用心灵把握的。那又怎样说天主呢？祂没有身体，又如何充满活力，难道容易把握吗？但如果他们相信有神灵存在——如果真有，他们显然是没有身体的——那么人类的灵魂也必然以同样方式存在，因为从理性本身和辨别力可以感知，人类与天主有某种相似性。最后，甚至\Person{马尔库斯·图利乌斯}所看到的那个论证，也具有足够的力量：灵魂的不死可以从以下事实辨别出来：没有别的动物对天主有任何认识；而宗教几乎是唯一将人与哑巴动物区别开来的事物。既然这仅属人类，它无疑证实了我们所追寻、渴望并培养之事，是即将与我们亲近且十分密切的。\par

当有人思考其他动物的本性——至高天主的眷顾使它们卑下，身体弯向地面、俯伏于地，由此可知它们与上天没有交往——难道他还不明白所有动物中唯独人是属天和神圣的。人的身体从地面升起，面容高昂，直立姿势，追寻其根源，仿佛藐视尘世之低微，伸向高处，因为他认识到至善应在最高之处寻求，并记住天主使他荣耀的状态，仰望造物主。\Person{特里斯墨吉斯忒斯}极其正确地称这仰望为对天主的默观，这在哑巴动物中是不存在的。既然智慧，独赐予人的智慧，不外是对天主的认识，那么灵魂显然不会消亡，也不会分解，而是永远存留，因为它凭借其本性的冲动寻求并热爱永恒的天主，感知它从何而来，以及它将要回到何处。此外，以下事实也是不死的一个不轻的证据：唯独人类使用天上的元素。因为世界的本性由两种相互对立的元素组成——火与水——其中一种归于天，另一种归于地；其他生物，因为是地上的、可朽的，便使用尘世而沉重的元素：唯独人类使用火，这是一种轻的、向上上升的、天上的元素。但沉重的事物压向死亡，而轻的事物提升向生命；因为生命在高处，死亡在低处。正如没有火就没有光，同样没有光就没有生命。因此，火是光与生命的元素；由此显然，使用它的人，分享不朽的境界，因为造成生命的事物与他亲近。\par

惟独人蒙受德行之恩赐，也是灵魂不朽的一大明证。因为若灵魂消灭，恩赐便不合乎本性；因它有害于现世的生命。那与畜生共有的尘世生命，既藉着各种悦目而可口的果实追求快乐，又避免痛苦；痛苦那令人不快的感触损伤了生灵的本性，并试图引领它们走向瓦解生命的死亡。因此，若德行禁止人获取本性所向往的美善，并促使人承受本性所回避的祸患，则德行便是邪恶，且与本性对立；追求德行的人必被视为愚昧，因为他既回避现世的善，又同样追求恶，而无更大利益的希望。当我们可以享受最甜美的快乐时，若不选择过寒微、匮乏、受轻蔑和羞辱的生活，或根本不再活着，却受痛苦折磨而死，而由此等祸患中我们得不到任何补偿以弥补放弃的快乐，岂不显得毫无理智？但若德行并非邪恶，且因其轻视邪恶可耻的快乐而行高尚之事，并因其为尽本分而不惧痛苦或死亡而行勇敢之事，则它必获得比它所轻视的那些事物更大的善。然而，当死亡已被经受时，除了永生之外，还能期盼什么更大的善呢？\par

% structure: p0054 source=h2 target=subsection reason=relative-heading-rank; base=h2

\subsection{第十章 论恶习与德行，以及生命与死亡。}

现在让我们依次转向那些与德行相反的事物，以便也从这些事推论出灵魂的不朽。所有恶习都是暂时的，因为它们因现时而激动。愤怒的冲动在复仇完成后平息；身体的快乐止息情欲；欲望或因完全享受所追求的对象而消灭，或因其他情感的激动而消失；野心一旦获得渴望的荣耀便失去力量；其他恶习同样不能站稳并持久，反因它们所渴望的享受本身而终结。因此它们隐退又返回。但德行是永久的，毫无间断；凡一旦接受德行的人不能离开它。因为若它有任何中断，若我们有时能没有它，那常与德行对立的恶习便会返回。因此，若它擅离职守，若它有时退隐，它便未被掌握。但德行一旦为自己建立稳固的居所，就必在每项行动中从事；除非它以永恒的守卫坚固其居住的胸怀，否则它不能忠实地驱赶和击退恶习。因此，德行本身的不同断持续表明，人的灵魂——若接受了德行——便持久存留，因为德行是永恒的，且惟独人的心智接受德行。所以，既然恶习与德行相反，两者的整个体系必彼此不同且相反。因为恶习是灵魂的骚动与扰乱；德行反之是心灵的温和与安宁。因为恶习是暂时而短暂的；德行是永恒而不变的，且始终与其自身一致。因为恶习的果实，即快乐，与恶习本身同样短暂而暂时，因此德行的果实与赏报是永恒的。因为恶习的利益是即刻的，故德行的利益是未来的。\par

由此，在此生中没有德行的赏报，因为德行本身仍存在。因为正如恶习在完成其行为之后，随之而来的是快乐和它们的赏报；同样，德行终结之后，其赏报随之而来。但德行唯在死亡中才终结，因为其最高职分在于经受死亡：因此德行之赏报在死后。最后，\Person{西塞罗}在\WorkTitle{图斯库卢姆论辩集}中虽带疑惑地认识到，至善除死后外并不临到人。他说：\Quote{若情况如此，人将怀着自信的心走向死亡，在死亡中我们确知要么有至善，要么无恶。}因此，死亡并不消灭人，而是使他进入德行之赏报。但如同\Person{西塞罗}所说，那以恶习和罪行污染自己并做了快乐的奴隶的人，确实受定谳遭受永远的惩罚，圣经称之为\InnerQuote{第二次死亡}，它既是永远的，又充满最严厉的折磨。因为如同向人提出两种生命，一种属于灵魂，另一种属于身体；同样也提出两种死亡——一种关乎身体，所有人依本性必须经历；另一种关乎灵魂，它因邪恶而获得，因德行而避免。正如这生命是暂时的并有固定界限，因它属于身体；同样，死亡也是暂时的并有固定终结，因它影响身体。\par

% structure: p0057 source=h2 target=subsection reason=relative-heading-rank; base=h2

\subsection{第十一章 论末世，以及灵魂与身体。}

因此，当天主为死亡所定的时期结束时，死亡本身也将终结。既然暂时的死亡紧随暂时的生命而来，那么灵魂复活而获得永生就是必然的，因为暂时的死亡已经终结。再者，正如灵魂的生命是永恒的，在其中它领受其不朽性的神圣而不可言喻的果实；同样，它的死亡也必是永恒的，在其中它因自己的过犯而承受永久的惩罚和无限的痛苦。因此，事情处于这样的状态：在此世身体和尘世方面幸福的人，将永远悲惨，因为他们已经享受了他们所偏爱的美好事物，这发生在那些崇拜假神而忽视真天主的人身上。其次，那些追随正义、在此世悲惨、被藐视、贫穷，并因正义本身而常受侮辱和伤害的人（因为美德无法以其他方式获得），将永远幸福，既然他们已经忍受了恶，也当享受善。这显然发生在那些蔑视尘世之神和易朽善物、追随天主天上宗教的人身上，天主的善物是永恒的，如同赐予它们的天主自身一样。关于身体和灵魂的作为，我还能说什么呢？它们岂不表明灵魂不受死亡支配吗？因为就身体而言，既然它本身是脆弱和必朽的，它所设计的一切作为也同样易于消亡。因为\Person{图利乌斯}说，凡人手所做的一切，没有一样不会在某个时候归于毁灭，或由于人的伤害，或由于时间——万物的毁灭者。\par

但我们确实看到心灵的产物是不朽的。因为有多少人，他们献身于蔑视现世事物，将他们的天才和伟大事迹的纪念留给记忆，显然借此为他们的心灵和美德赢得了不朽的名声。因此，如果身体的作为是必朽的，因为身体本身是必朽的，那么由此可知，灵魂被证明是不朽的，因为它的产物并非必朽。同样，身体和灵魂的欲望也表明一个是必朽的，另一个是永恒的。因为身体只渴望暂时的事物，即食物、饮料、衣服、休息和享乐；若没有灵魂的同意和协助，它甚至不能渴望或获得这些事物。但灵魂本身渴望许多不涉及身体职责或享受的事物；这些事物并非易朽，而是永恒的，如美德的声誉、名声的纪念。因为灵魂甚至与身体相对，渴望对天主的敬拜，这在于克制欲望和情欲、忍受痛苦、蔑视死亡。由此可以相信灵魂不会消亡，而是与身体分离，因为身体没有灵魂便一无所能，但灵魂没有身体却能做许多伟大的事。我何必提及那些肉眼可见、手可触摸的事物不能永恒，因为它们承受外在的暴力；而那些既不可触摸也不可见、仅在其力量、方式和效果中显明的事物是永恒的，因为它们不受外在的暴力？但如果身体因同样可见可触而必朽，那么灵魂因不可触摸也不可见而永恒就是合理的。\par

% structure: p0060 source=h2 target=subsection reason=relative-heading-rank; base=h2

\subsection{第十二章　论灵魂与身体，及其结合、分离与回归。}

现在让我们驳斥那些持相反论点者的论据，\Person{卢克莱修}在其第三卷书中叙述了这些论点。他说，既然灵魂与身体一同诞生，它必然与身体一同死亡。但两种情况并不相似。因为身体是坚固的，能被眼睛和手抓住；而灵魂是细微的，避开了触摸和视线。身体由泥土形成，变得坚实；灵魂中没有任何凝固的、尘世重量的东西，如\Person{柏拉图}所主张的。因为它不可能有如此巨大的力量、如此高超的技能、如此迅速的行动，除非它源自天上。因此，身体既然由沉重和易朽的元素构成，且是可触可见的，它就腐烂并死亡；它也无法抵抗暴力，因为它落入视线和触觉之下；但灵魂，因其细微而避开了一切接触，任何攻击都不能使它分解。因此，尽管它们从出生起就结合在一起并相连，一个由尘世材料构成，可以说是另一个的容器，另一个由天上的精微物抽出，当任何暴力使二者分离——这种分离称为死亡——时，各自便返回自己的本性：来自尘世的归于尘土，来自天上气息的则保持不变，永远繁荣，因为神圣的精神是永恒的。最后，同一\Person{卢克莱修}忘记了他所主张的和他所捍卫的教义，写下了这些诗句：——\par

\Quote{那先前来自土中的，重归土中；那从以太边界送来的，又被天穹区域带走。}\par

但这样的语言并非他该使用的，他主张灵魂与身体一同消灭；然而他被真理所征服，真实的体系不知不觉地潜入他心中。此外，他所提出的推论——即灵魂解体，也就是与身体一同消灭，因为它们是同时产生的——既是错误的，也可以转向相反的方向。因为身体并不与灵魂一同消灭；当灵魂离去后，身体仍完整地存留多日，并且常常通过医药处理，身体会长时间保持完整。因为如果它们一同消灭（如同它们一同产生），灵魂就不会匆忙离去并抛弃身体，而是会在同一时刻同样消散；而且身体，当气息仍在其中时，也会像灵魂离去那样迅速地分解和消灭：是的，确实，身体一旦分解，灵魂就会消失，如同从破碎的器皿中倾泻而出的湿气。因为如果尘世而脆弱的身体在灵魂离去后并不立即流散并归于它来源的尘土，那么灵魂（它并不脆弱）将永存于世，因为它的起源是永恒的。他说，既然理智在孩童时期增长，在青年时期强健，在老年时期衰退，那么显然它是可朽的。首先，灵魂与心智并非同一回事；因为我们活着是一回事，思考是另一回事。因为在睡眠者身上，安静下来的是心智，而非灵魂；在疯狂者身上，心智熄灭，灵魂却存留；因此他们并非被称为没有灵魂，而是被剥夺了心智。所以心智（即理智）随着年龄增长或衰退。灵魂始终处于自己的状态；从它获得呼吸能力那一刻起，它保持不变直到终点，直到它从身体的囚禁中被释放，飞回自己的居所。其次，灵魂虽然由天主所嘘，但由于被幽禁在尘世血肉的黑暗居所中，它并不拥有属于神性的知识。因此它通过学习和听闻来听到和学到一切，并藉着学习和听闻获得智慧；老年并不减少智慧，反而增加智慧，如果青年时期是在德行中度过的；倘若过度衰老使肢体衰弱，那么视力消失、舌头麻木、听力变聋，这并非心智的过错，而是身体的过错。但有人说记忆力衰退。这有什么奇怪呢？如果心智被坍塌房屋的废墟所压迫，并忘记过去，那么它只有在逃脱了幽禁它的牢狱之后，才能成为神圣的。\par

但他又说，灵魂也会遭受痛苦和悲伤，并因醉酒而失去知觉，由此可见它是脆弱而可朽的。因此，德行和智慧是必要的，这样，因遭遇和见到不配之事而产生的悲伤，可以被刚毅所击退；而快乐不仅可以通过戒酒，还可以通过戒除其他事物来克服。因为如果灵魂缺乏德行，如果它沉溺于快乐并因此变得软弱，它就会沦为死亡的奴仆，因为正如我们所指出的，德行是不朽的缔造者，而快乐则是死亡的缔造者。但死亡，如我所阐述的，并非完全消灭和摧毁，而是以永恒的折磨来惩罚。因为灵魂不能完全毁灭，因为它从永恒的天主之神那里获得了起源。他说，灵魂也能感受到身体的疾病，并遭受自我的遗忘；正如它会生病，也常得医治。因此，这就是为什么特别要运用德行：为的是使心智（而非灵魂）不被身体的任何痛苦所困扰，或遭受自我的遗忘。由于心智位于身体的某一部位，当疾病的猛烈攻击损坏了那一部位时，心智就会移动位置；如同被摇动一般，它离开自己的位置，等到治疗和健康重新塑造了它的居所时，它才会返回。因为，既然灵魂与身体结合，如果它缺乏德行，它就会因身体的感染而生病，并且因分享其脆弱性，这种软弱会影响心智。但当它脱离身体后，它将独立兴盛；此时任何脆弱的情况都不会再攻击它，因为它已经卸下了脆弱的覆盖物。他说，如同眼睛被挖出并脱离身体后什么也看不见，同样，灵魂脱离后也什么都感知不到，因为它本身就是身体的一部分。这是错误的，与所设想的情况并不相似；因为灵魂不是身体的一部分，而是存在于身体之中。如同器皿中的东西不是器皿的一部分，房子里的东西不能说成是房子的一部分；同样，心智不是身体的一部分，因为身体要么是灵魂的器皿，要么是灵魂的容器。\par

现在，有一种更为空洞的论据，它说灵魂似乎是可朽的，因为它不是迅速脱离身体，而是从各肢体逐渐展开，从脚趾开始；仿佛如果它是永恒的，它就会在一瞬间迸发出来，就像那些死于刀剑下的人所发生的那样。但那些因疾病而死的人呼出气息的时间更长，以致当肢体变冷时，灵魂被呼出。因为，既然灵魂蕴藏在血液的物质中，正如光蕴藏在油中，当那物质被热病的热量消耗时，肢体的末端必然会变冷；因为更纤细的静脉延伸到身体的末端，当源头干涸时，极端和更小的溪流就会枯竭。然而，不能因此就认为，因为身体的知觉消失了，灵魂的感觉就熄灭并消亡了。因为当身体衰败时，不是灵魂变得无知觉，而是当灵魂离开时，身体变得无知觉，因为它带着一切知觉离去。但由于灵魂因其临在而赋予身体知觉，并使其存活，灵魂自身不可能没有生命和知觉，因为它本身就是意识和生命。至于那说：\par

\Quote{但如果我们的心灵是不朽的，它在临终时就不会如此抱怨它的解体，反而会像蛇一样欣喜于出游和脱去它的旧衣，}\par

我从未见过有人在死亡中抱怨自己的解体；但他或许见过某个伊壁鸠鲁派的人甚至在死亡时还在哲思，用最后一口气谈论着自己的解体。\par

怎能知道他是感到自己处于解体状态，还是正在脱离身体，当他的舌头在离世时变得哑默？因为只要他还能知觉并有说话的能力，他就尚未解体；当他已经经历了解体，他就不再能知觉或说话，所以他要么还不能抱怨自己的解体，要么已经不能再抱怨了。但据说，他在经历解体之前就明白自己必须经历它。我何必提及我们看到许多临终者，并非抱怨他们在经历解体，而是见证他们在离世，启程并行走？他们用手势示意，或者如果他们仍能说话，也用声音表达。由此显明，所发生的是分离而非解体；这表明灵魂继续存在。伊壁鸠鲁派体系的其他论点与\Person{Pythagoras}相对立，他主张灵魂从因年老和死亡而衰败的身体中迁移，进入新生而刚出生的身体中；并且相同的灵魂始终被再生，有时在人身，有时在绵羊，有时在野兽，有时在飞鸟；它们因此是不朽的，因为它们经常更换住所，由各种不同的身体构成。而这个无知之人的意见，既然荒谬，更值得舞台演员而非哲学学派，甚至不值得被认真反驳；因为这样做的人似乎害怕有人会相信它。因此，我们必须略过那些为虚假反对虚假而讨论的内容；反驳那些反对真理的论据就足够了。\par

% structure: p0069 source=h2 target=subsection reason=relative-heading-rank; base=h2

\subsection{第十三章 论灵魂及其永生的见证。}

我认为，我已经证明灵魂不是可分解的。现在我要提出见证人，他们的权威可以确认我的论点。我不再列举先知们的证言，因为他们的体系和占卜只在于这一点：教导人被创造是为了敬拜天主，并从祂那里获得永生；而我要提出那些拒绝真理者不能不相信的人。\Person{Hermes}在描述人的本性时，为了表明人如何由天主创造，提出了这样的说法：\Quote{并且，那同一个以两种本性——不朽的与可朽的——构成了一个本性，即人的本性，使同一部分不朽，部分可朽；并将此置于那神圣而不朽的本性与那可朽而可变的本性之间，让他看见一切，从而赞叹一切。}但或许有人会把他算在哲学家之列，尽管他被列入众神中，并被埃及人以\Person{Mercury}之名尊敬，并且对他的权威不致赋予比对\Person{Plato}或\Person{Pythagoras}更多的尊重。因此，让我们寻求更大的见证。有一个名叫\Person{Polites}的人向\Place{Miletus}的\Person{Apollo}询问：灵魂死后是继续存在还是解体？他用以下的诗句回答：——\par

\Quote{只要灵魂被锁链束缚于身体，感知可朽的苦难，它就屈服于必死的痛苦；但当身体耗竭之后，它找到了极其迅速的死亡解体，它就被完全带到空中，永不衰老，永远保持无损；因为天主首生的眷顾作了这样的安排。}\par

西比尔诗歌说了什么？它们不是宣告了这一点吗，当它们说时候将到，天主将审判活人死人？——我们将以后引证它们的权威。因此，德谟克利特、伊壁鸠鲁和狄凯阿克斯关于灵魂消散的见解是错误的；他们绝不敢在任何一个巫师面前谈论灵魂的毁灭，因为巫师知道，灵魂可以通过某些咒语从冥界被召唤出来，它们就在眼前，让自己被肉眼看见，说话，预言未来的事；如果他们胆敢如此，他们就会被事实本身和他们所看到的证据所压倒。但由于他们不理解灵魂的本性——它如此精微，以至于逃脱了人类心智的眼睛——他们便说灵魂会消逝。至于阿里斯托克森，他否认甚至在灵魂活于身体内时存在着任何灵魂，这又怎么说呢？但是，正如通过琴弦的绷紧而产生和谐的声音和音乐家所称的和声一样，他认为感知能力存在于身体之中，是由于内脏的连接和四肢的活力；再没有比这更无稽的言论了。他的眼睛确实完好无损，但他的心是盲目的，他未能看到自己活着，并且拥有那使他构想出这一想法的理智。但这在许多哲学家身上都发生过：他们不相信任何不是肉眼可见的事物存在；然而，对于感知那些其力量和性质更多是被感受到而不是被看见的事物，心智的视力应当比身体的视力清晰得多。\par

% structure: p0073 source=h2 target=subsection reason=relative-heading-rank; base=h2

\subsection{第十四章　论世界的起初与终末}

既然我们已经谈论了灵魂的不朽，接下来我们要教导如何以及何时将它赐予人；这样，他们也能看到那些人的乖戾和愚昧的错误，那些人妄想某些凡人因凡人的法令和教条而成为了神，或因他们发明了技艺，或因他们传授了某些大地出产之物的用途，或因他们发现了有益于人类生活的事物，或因他们杀死了野兽。这些事何等不配获得不朽，我们在前几卷中已经表明，现在我们将表明，以便清楚显明：唯有正义才能为人获得永生，唯有天主才能赐予永生的奖赏。因为那些据说因自己的功绩而成为不朽的人，既然既不拥有正义，也不拥有任何真正的德行，他们为自己获得的不是不朽，而是因自己的罪恶和情欲而来的死亡；他们也不配获得天堂的奖赏，而应受地狱的惩罚，这刑罚连同他们所有的崇拜者都悬在他们头上。而我指出，这审判的时候临近了，好使正义者得到应得的赏报，恶人受到应受的惩罚。\par

柏拉图及许多其他哲学家，由于对万物的起源以及世界被造之初的那段原始时期一无所知，便说在这美丽的世界秩序完成之后，已经过去了成千上万个世纪；在这方面，他们或许追随了迦勒底人，正如西塞罗在其\WorkTitle{论占卜}第一卷中所记载的，迦勒底人愚蠢地声称他们的记载中包含了四十七万年。在这一点上，因为他们认为自己不会被揭穿，便以为可以随意撒谎。但是我们，蒙圣经教导认识真理的人，知道世界的起初和终末；关于终末，我们现在将在本卷的结尾处讲述，因为我们已在第二卷中解释了起初。因此，让那些从世界起初列举成千上百个世纪的哲学家们知道，第六千年还没有结束，当这个数目完成时，终结必然发生，人类事务的状况将被重塑得更好；其证明必须先叙述，以便事情本身变得清晰。天主在六天之内完成了世界和这项令人赞叹的自然工程，正如圣经的奥义中所记载的，并祝圣了第七天，因为祂在那天歇了祂的工程。这就是安息日，它在希伯来语中以其数字得名，因此七是合法且完全的数字。因为有七天，它们按顺序的循环构成了年的周期；又有七颗不落之星，还有七个被称为行星的光体，它们不同且不等的运动被认为导致了环境和时间的各种变化。\par

因此，既然天主的全部工程都在六天内完成，世界也必须经过六个时期，即六千年，保持其当前状态。因为天主的大日以千年为周期，正如先知所显示的，他说道：\BibleQuote{主啊，在祢眼前，千年如同一日。} 正如天主在那六天劳动以创造如此伟大的工程，同样，祂的宗教和真理也必须在六千年中劳动，在此期间邪恶盛行并掌权。再者，既然天主在完成祂的工程后，在第七天休息并祝圣了它，那么在第六千年结束时，一切邪恶都必从地上被消灭，正义将统治一千年；并且必有一段宁静和安息，免除世界长久以来承受的劳苦。至于这事如何发生，我将按次序说明。我们常说起，次要的和微小的事物是伟大事物的预像和预先的投影；正如我们这以日出日落为界限的一天，是那以千年之圆周为界限的大日的写照。\par

同样，地上人的塑造也预表了将来天上百姓的形成。因为正如当为人使用而设计的一切都完成后，最后在第六天，天主也创造了人，并将他引入这世界，如同进入一个精心预备的家园；同样，如今在伟大的第六天，真正的\Person{人}正藉天主的圣言而被塑造，即是说，一个圣洁的子民正藉天主的训诲和诫命被塑造成义。并且，正如当时一个必朽的、不完全的人从尘土中被造，使他能在今世活一千年；同样，如今从这地上的年代中，一个完美的人被塑造，使他被天主赐予生命后，能在这个世界上统治一千年。但终结将以何种方式发生，人类的命运将有何结局，若有人查考圣经，便会明白。此外，世上的先知之声也与天上的预言一致，宣告不久之后万物的终结和毁灭，描述这疲惫耗损的世界仿佛到了最后的晚年。至于先知和先见所说，在最后的终结临到世界之前将要发生的事，我将从各处搜集汇编，附列如下。\par

% structure: p0078 source=h2 target=subsection reason=relative-heading-rank; base=h2

\subsection{第十五章 论世界的荒芜与帝国的更替}

在圣经的奥秘中记载着，一位希伯来人的首领因缺粮所迫，带着他所有的家人和亲属进入了\Place{埃及}。当他的后裔长久留在\Place{埃及}，增长成为一个大民族，并被沉重而无法忍受的奴役枷锁压迫时，天主以无法治愈的打击击打了\Place{埃及}，并解救祂的子民，带领他们穿过大海，那时海浪被劈开，向两边分开，百姓在干地上走过。\Place{埃及}王企图追赶逃跑的他们，海水复归原处，他和他的全体百姓都灭亡了。这一如此辉煌而奇妙的事迹，虽然当时向人显示了天主的能力，也是同一天主将在时代的最后终结时所行的一件更大事迹的预兆和预象，因为祂将解救祂的子民脱离世界的压迫性束缚。但在当时，天主的子民是一个民族，只在一个国家里，所以只击打了\Place{埃及}。但如今，天主的子民从各种语言中被聚集，居住在万民之中，并被统治他们的人压迫，因此，必须使所有民族，即整个世界，遭受天上的打击，以便正义的子民，即崇拜天主的人，得以解放。正如当时给了\Place{埃及}人预兆，显示即将到来的毁灭，同样，在最后的时期，遍及世界一切元素将有奇异的预兆发生，使万民明白即将来临的毁灭。\par

因此，随着这世界末日的临近，人类事务的状况必然发生变化，并因邪恶的盛行而变得更糟；以致于我们现今这些不义和不敬已增至极点的时代，相比之下，将可被判定为幸福而几乎是黄金般的。因为正义将如此减少，而不敬、贪婪、私欲和情欲将大大增加，以至于如果那时有善人存在，他们将沦为恶人的猎物，并在各方面被不义者骚扰；只有恶人将富裕，而善人将在一切诽谤和匮乏中受苦。一切正义将被混淆，法律将被废除。那时除了用手夺取或护卫的东西外，无人拥有任何事物；狂妄和暴力将主宰一切。人与人之间没有信实、和平、仁慈、羞耻或诚实；因此也没有安全、治理或任何脱离罪恶的安息。因为全地将陷入骚动，战争将到处肆虐，所有民族都将武装起来彼此敌对，邻国将彼此冲突；首先，\Place{埃及}将为她愚昧的迷信付出代价，被血覆盖如同河流。然后刀剑将横越世界，割倒一切，像庄稼一样打倒一切。我心不敢叙述，但我仍要叙述，因为这事将要发生——这荒芜和混乱的原因将是：罗马之名，世界现今由它统治，将从地上被除去，政权将回到\Place{亚洲}；\Place{东方}将再次统治，\Place{西方}将沦为奴役。任何人都不应感到惊奇，如果一个如此庞大、由如此多伟人长期增强、且由如此巨大资源所巩固的王国，仍然会在某一时刻倒下。凡是人力所建立的一切，同样也能被人力所摧毁，因为凡人的工程都是必朽的。同样，古代的其他王国，尽管长久繁荣，也都被摧毁了。因为据记载，\Person{埃及人}、\Person{波斯人}、\Person{希腊人}和\Person{亚述人}曾统治世界；在他们全部灭亡之后，统治权也传到了\Place{罗马}人手中。正因为他们比所有其他王国更大，所以也将以更大的覆灭而倾倒，因为那些高于其他建筑的建筑物，坠落时具有更大的力量。\par

因此，塞内加并非不熟练地将罗马城的时代按年龄划分。他说，最初是它的婴儿期，在国王罗慕路斯治下，罗马由此诞生，并如同被养育；然后是它的童年，在其他国王治下，它被增长，并以更多的训诲制度和机构得以成型；但最终，在塔尔奎尼乌斯统治时期，当它已如同开始长大时，它不再忍受奴役；并且，摆脱了傲慢暴政的轭，它宁愿服从法律而非国王；当它的青年期随着布匿战争的结束而终结时，然后，以确认的力量，它开始成为成年。因为当迦太基——它长期以来的权力对手——被除掉后，它通过陆地和海洋向全世界伸展双手，直到制服了所有国王和民族，当战争资源已匮乏时，它滥用了导致自我毁灭的力量。这是它的第一次老年，当时，被内战撕裂，并被内在的邪恶压迫，它再次退回到单一统治者的治理之下，如同回复到第二次婴儿期。因为，失去了它在布鲁图斯的指导和权威下所捍卫的自由，它变得如此衰老，好像没有能力支撑自己，除非依赖其统治者的帮助。但如果事情如此，剩下的还有什么，除了老年将跟随死亡？而且这将如此发生，先知们的预言以其他名号的掩盖简要地宣告，以致无人能轻易理解它们。然而，西彼拉们公开说，罗马注定要灭亡，而且确实是因天主的审判，因为它仇恨祂的名；并且，作为正义的敌人，它摧毁了持守真理的民众。还有，希斯塔斯佩斯——一位非常古老的米底亚国王，一条现在称为希达斯佩斯的河由此得名——通过一个男孩的阐释传下了一个奇妙的梦给后人的记忆，那男孩发出预言，早在特洛伊民族建立之前就已宣告，罗马的帝国和名号将从世上被除去。\par

% structure: p0082 source=h2 target=subsection reason=relative-heading-rank; base=h2

\subsection{第十六章 论世界的浩劫及其预言征兆。}

但恐有人以为此不可信，我将显示它将如何发生。首先，王国将被扩大，而最高权力分散于多人并分裂，将被削弱。然后，内乱将不断被播下；不会有任何止息脱离致命战争，直至同时有十王兴起，他们将瓜分世界，不是为治理，而是为消耗。这些王，将其军队扩至极大，并丢弃田地的耕作——这乃是倾覆与灾难的开端——将荒废、粉碎并消耗一切。然后，一个最强大的敌人将从北方区域的极边界突然兴起攻击他，此敌将摧毁那亚细亚属地中三名在位者，被其余者接纳为盟友，并被立为众王子之首。他将以不可容忍的统治骚扰世界；将混杂神界与人间之事；将图谋不可言说与可憎之事；将心中筹算新计谋，以便为自己建立政府：他将更改法律，制定自己的法律；他将玷污、掠夺、劫洗、杀害。最终，名字被更改，统治的座位被转移，混乱与人类的骚动将随之而来。那时，真正可憎与可恶的时刻来临，在其中，对任何人而言，生命都不会愉快。\par

城市将彻底倾覆并灭亡；不仅因火与剑，也因持续的地震、水泛滥、频繁的疾病和反复的饥荒。因为大气将被污染，变得腐败而有害——时而因不合时节的雨，时而因不毛的干旱，时而因寒冷，时而因酷热。大地也不会给人果实；没有田地、树木或葡萄树会出产任何东西；但当它们在开花时给予最大希望之后，却在果实上失败。泉源也必连同江河一同干涸，以致没有足够的水供饮用；水将变为血或苦涩。因这些事，地上的走兽、空中的飞鸟、海中的鱼都将消亡。天上奇异的异象也以最大的恐怖使人心惊惶：彗星的轨道、太阳的黑暗、月亮的颜色、流星的滑落。然而，这些事不会以惯常的方式发生；而是将有星星忽然显现，未知且肉眼未曾见过；太阳将永久昏暗，以致昼夜之间几乎没有分别；月亮现在不只缺光三小时，而是被永恒的血覆盖，运行失常，以致人不易确定天体的运行或时间的系统；因要么冬天有夏天，要么夏天有冬天。然后，年将缩短，月将减少，日将收缩为短暂时光；群星将大量坠落，以致整个天空无光而显黑暗。最高的山也将倒塌，与平原齐平；海将变得不可航行。\par

为使人类与大地的灾祸无所不备，将从天闻号角声，西彼拉以下列方式预言此事：——\par

\Quote{从天上发出的号角将发出哀号之声。}\par

然后，众人都要因那悲惨的声音而战栗发抖。但那时，因天主对不认识正义之人的忿怒，刀剑、火灾、饥荒和疾病将要肆虐；而恐惧则将永远悬在众人头顶。那时，他们呼求天主，但天主不俯听他们；他们渴望死亡，但死亡不来；黑夜也不能使他们的恐惧得安息，睡眠也不临到他们的眼睑，而是焦虑和警醒耗尽人的灵魂；他们哀哭、悲叹，咬牙切齿；他们庆幸已死的人，哀哭活着的人。通过这些和许多其他的灾祸，大地将荒凉，世界将变得面目全非、荒无人烟，正如\OriginalTerm{西比勒}{Sibyl}的诗句所表达的：—\par

\Quote{世界将因人的毁灭而失去美丽。}\par

人类将如此减少，以至于几乎只剩下十分之一的人；曾经有千人出去的地方，几乎只有百人能出去。在敬拜天主的人中，三分之二将要灭亡；而那经受试炼的三分之一将要存留。\par

% structure: p0090 source=h2 target=subsection reason=relative-heading-rank; base=h2

\subsection{第十七章 论假先知、义人的苦难及对他的毁灭。}

但我将更清楚地说明这事发生的方式。当时代的终结临近时，一位伟大的先知将从天主那里被派遣，使人们转向认识天主，他将获得行奇事的能力。凡不听从他的人，他将封闭天，使之不降雨；他将使他们的水变为血，用干渴和饥饿折磨他们；如果有人试图伤害他，火将从他的口中发出，烧死那人。借着这些奇事和权能，他将使许多人转向敬拜天主；当他的工作完成时，另有一个王将从叙利亚兴起，由邪灵所生，是人类族类的倾覆者和毁灭者，他将把先前邪恶者所剩下的连同他自己一起毁灭。他将与天主的先知作战，并要战胜他、杀死他，任他暴尸不葬；但三天后，他将复活；在众目睽睽下，众人观看并惊奇，他将被提到天上。但那王不仅本身最可耻，而且也是谎言的先知；他将自立并自称为神，命令人敬拜他如同天主子；他将被赐予能力行征兆和奇事，借这些景象引诱人崇拜他。他将命令火从天降下，命令太阳停住并偏离轨道，命令偶像说话；这一切都将因他的话而成就——凭着这些奇迹，甚至许多有智慧的人也将被他诱惑。然后，他将企图摧毁天主的圣殿，逼迫义人；那时将有困苦和灾难，如同从世界开始以来从未有过的那样。\par

凡相信他并与他联合的人，都要被他如羊一般烙印；凡拒绝他印记的人，或逃到山中，或被捕获，在精心设计的酷刑下被杀。他还要用先知的书卷包裹义人，将他们焚烧；他被赐予能力使全地荒凉四十二个月。那将是正义被抛弃、纯洁被仇恨的时候；恶人像敌人一样掠取善人；没有法律、秩序或军纪被遵守；没有人尊敬白发，没有人认识敬虔的义务，也没有人怜悯性别或幼年；万物都变得混乱不堪，违背公义和自然法则。大地就像被一场共同的掠夺所荒废。这些事发生时，正义者和真理的追随者将离开恶人，逃往荒野。那邪恶的王听到这话，怒火中烧，率大军前来，调集所有兵力，包围义人所处的山，要捉拿他们。但他们见自己被四面包围、围困，就大声呼求天主，请求天上的援助；天主必垂听他们，从天上派遣一位伟大的王来拯救并释放他们，用火与剑消灭所有恶人。\par

% structure: p0093 source=h2 target=subsection reason=relative-heading-rank; base=h2

\subsection{第十八章 论末时世界的命运及占卜者所预言的事。}

这一切将如此发生，是众先知在圣神感动下所宣告的，也是占卜者在魔鬼的唆使下所预言的。因为我前面提到的\OriginalTerm{希斯塔斯普}{Hystaspes}，描述了这末时的邪恶后，说虔诚和忠信的人与恶人分离，将流泪哀号向天伸手，恳求\OriginalTerm{朱庇特}{Jupiter}的保护；朱庇特将垂顾大地，倾听人的声音，并消灭恶人。这一切都是真的，除了一点：他把天主将行的事归给了朱庇特。但那也是从记载中被魔鬼诡诈地删去了，即那时天主子将被派遣，消灭所有恶人，释放义人。然而，\OriginalTerm{赫尔墨斯}{Hermes}并没有隐瞒这事。因为在那本题为\WorkTitle{完备论著}的书中，列举了我们所谈论的那些邪恶之后，他补充了这些话：\Quote{但当这些事这样发生时，那位主、父、天主，以及第一位和唯一的天主的创造者，垂顾所发生的事，以自己的旨意（即良善）对抗混乱，召回迷失者并洁净邪恶，部分用多水淹没，部分用极快的火焚烧，有时用战争和瘟疫压迫，使他的世界恢复到古代的状态，并重建它。}\OriginalTerm{西比勒们}{Sibyls}也表示，若不如此，天主子就不能由祂至高的父派遣，来从恶人手中释放义人，并消灭不义者及其残忍的暴君。其中一位西比勒这样写道：——\par

\Quote{他要求，想要毁灭幸福之城；一位从众神那里派来对付他的王要杀死所有伟大的国王和首领；然后审判就这样从永生者临到世人。}\par

还有另一位女预言者：—\par

\Quote{然后天主将从太阳那里派来一位王，他要使整个大地止息灾难性的战争。}\par

又有另一位说：—\par

\Quote{他将除去加在我们颈上的不堪忍受的奴役之轭，废除不敬神的律法和暴力的锁链。}\par

% structure: p0100 source=h2 target=subsection reason=relative-heading-rank; base=h2

\subsection{第十九章 论基督来临审判与击败假先知。}

因此，当世界遭受压迫，人的力量不足以推翻如此强大暴政之时，因为它要以大群盗贼军队压迫被俘的世界，如此巨大的灾难将需要天主的援助。因此，天主被这危险的不确定性以及义人的悲惨哀号所惊醒，将立刻派遣一位拯救者。然后在深夜的死寂和黑暗中，天穹将被打开，使降临的天主之光如闪电般显现在全世界：关于此事，女预言者用这些话说道：—\par

\Quote{当祂来临，必有火与黑暗在漆黑的夜中。}\par

这就是我们因我们的君王和天主的来临而守夜庆祝的那一夜：这一夜有双重含义；因为在这夜里，祂当时受难而得生命，日后将接受世界的国度。因为祂是拯救者、审判者、复仇者、君王和天主，我们称祂为基督。在祂降临之前，祂将给出这个记号：一把剑将突然从天而降，使义人知道圣战的统帅即将降临；祂将带着天使队伍降临到大地中央，一把不灭的火在祂前面行，天使的权能将把那围困山丘的群众交在义人手中，他们将被从第三时辰杀到黄昏，血如激流涌流；所有军队被摧毁后，那恶者独自逃脱，他的权能从他那里灭亡。\par

现在这就是那被称为假基督的；但他将妄称自己为基督，与真理争战，被击败后逃跑；他将多次重启战争，多次被征服，直到在第四次战斗中，所有恶人被杀死、制伏并被俘，他终于要为自己的罪行受惩罚。但其他也曾折磨世界的王公和暴君也将与他一起被锁链带到王面前；祂将斥责他们，责备他们，以他们的罪行谴责他们，定他们的罪，并将他们交付应得的酷刑。这样，邪恶被消灭，不虔敬被压制，世界就会获得安息，这世界曾长期受错误和邪恶的支配，忍受可怕的奴役。人手所造的神明将不再被敬拜；相反，偶像将被逐出他们的庙宇和床榻，连同他们奇妙的礼物一起被投到火中烧毁：女预言者也按照先知所说的事预告了这些将要发生的事：—\par

\Quote{但凡人将要打碎偶像和一切财富。}\par

厄立特里亚女预言者也作了同样的许诺：—\par

\Quote{并且人手所造的神明之作要被烧毁。}\par

% structure: p0108 source=h2 target=subsection reason=relative-heading-rank; base=h2

\subsection{第二十章 论基督的审判、基督徒与灵魂。}

此后，冥府将被打开，死人将复活，那同一位君王和天主要审判他们，至高的圣父要赐给祂审判和统治的大权。关于这审判和统治，在厄立特里亚女预言者中这样记载：—\par

\Quote{当这事命定完成，永生天主的审判现在临到世人时，伟大的审判和开端就要临到人。}\par

然后在另一处说：—\par

\Quote{然后裂开的大地要显出一个阴间的深渊；所有君王都要来到天主的审判座前。}\par

在同一作品中的另一处说：—\par

\Quote{滚动诸天，我要打开大地的洞穴；然后我要复活死人，解除命运和死亡的毒刺；之后我要召他们来审判，审判虔诚和不虔敬之人的生命。}\par

然而，并非所有人那时都要被天主审判，只有那些在天主的宗教中受过操练的人。因为那些不认识天主的人，既然不能为他们作出无罪判决，他们已经被审判和定罪了，因为圣经见证恶人不会复活受审。因此，认识天主的人将要被审判，他们的行为，也就是他们的恶行，将与其善行相比较、被衡量；这样，如果善行和义行更多更重，他们就被赐予幸福的生命；但如果恶行超过，他们就被定罪受罚。这里，也许有人会说：如果灵魂是不死的，它怎能受苦并感受惩罚呢？因为如果它要因罪受罚，显然它将感受到痛苦，甚至死亡。如果它不可能死亡，甚至不可能痛苦，那它就不可能受苦。\par

斯多葛派如此回应这个问题或论证：人的灵魂继续存在，不会被死亡干预而消灭；而且，义人的灵魂，纯洁、无苦、有福，返回它们起源的天上居所，或被带到某个幸福的平原，在那里享受奇妙的快乐；但恶人的灵魂，因以邪恶的情欲玷污自己，具有一种介于不朽和必朽之间的中间本性，并因肉体的传染而有些软弱；且受其欲望和贪念的奴役，染上无法抹去的污点和尘世的玷污；当这污点因时间长久而完全内化时，灵魂便屈从其本性，以致虽然它们来自天主而不能完全被消灭，然而透过身体的污点而容易受折磨，这污点因罪而被烧入，产生痛苦的感觉。诗人如此表达这一情感：——\par

\Quote{不，当生命终于逃遁，\newline{} 留下尸身冰冷亡殒，\newline{} 即便那时也不消逝\newline{} 泥土的痛苦遗产：\newline{} 许多长期积存的污点\newline{} 必然深留纹理之间。\newline{} 所以他们忍受惩罚的痛苦\newline{} 为古时的罪孽，使之纯净。}\par

这些事情接近真理。因为灵魂脱离身体后，正如同一诗人所说，是\par

\Quote{没有沉睡之夜的幻象，\newline{} 没有气流一半轻盈，}\par

因为它是一个精神，由于其极其轻盈而无法被感知，但对我们这些有形体的（人）而言，却能被天主感知，因为属于祂能够做一切事。\par

% structure: p0121 source=h2 target=subsection reason=relative-heading-rank; base=h2

\subsection{第二十一章 论灵魂的折磨与惩罚。}

首先，因此，我们说天主的权能如此伟大，连无形体的事物祂也能感知，并按祂的旨意管理它们。因为连天使也敬畏天主，因祂能以某种不可言喻的方式惩罚他们；魔鬼也畏惧祂，因他们受祂折磨和惩罚。所以，若灵魂虽是不朽的，却仍能受天主的苦，这有何奇怪呢？因为灵魂自身没有坚固和可触摸的东西，就不能受坚固有形物体的暴力；但因它们只生活在自己的精神中，只有天主能处置它们，天主的能量和本质是精神的。然而，圣经告知我们恶人将如何受罚。因为他们曾在身体内犯罪，他们将再次穿上肉体，以便在身体内补偿；但这肉体将不是天主用来装扮人的那肉体，像我们这地上的肉体，而是不可毁灭、永存的，以便能忍受折磨和永火，这火的本质不同于我们这为生活必需品使用的火，后者若没有某种材料的燃料支撑就会熄灭。但那神圣之火总是靠自己活着，没有任何滋养而繁荣；它没有烟混杂，而是纯净、液态、流动，如水的样子。因它不被任何力量向上推，如我们的火，被地上身体的污点和混杂的烟所束缚，被迫跃起并颤动地飞向天的性质。\par

因此，同一神圣之火以同一力量和权能，既焚烧恶人，又将他们重新成形，并将消耗他们身体的部分补回，还为自己供应永恒的营养：诗人将其转移到\Person{提提俄斯}的兀鹰身上。因此，身体毫无消耗，反而恢复物质，它只焚烧并使他们感受痛苦。但当祂审判了义人，祂也将以火试炼他们。那时，那些罪过在重量或数量上超出的，将被火炙烤焚烧；但那些被完全的正义和成熟的德行浸透的，将不会觉察那火；因为他们自身有天主的东西，能击退和拒绝火焰的暴力。纯洁的力量如此之大，以至于火焰无害地退缩；这火从天主获得这权能，即焚烧恶人，服从义人的命令。然而，不要让人以为灵魂死后立即受审判。因为所有人被拘禁在一个共同的监禁之所，直到大审判者调查他们功过的时候到来。那时，那些虔诚被认可的人将获得不朽的赏报；但那些罪恶和罪行被揭露的人将不再复活，而将与恶人一同被藏在同样的黑暗中，注定承受确定的惩罚。\par

% structure: p0124 source=h2 target=subsection reason=relative-heading-rank; base=h2

\subsection{第二十二章 论诗人的错误与灵魂从下界的返回。}

有些人以为这些是诗人的虚构，不知诗人从何处得到它们，并说这些事是不可能的；这对他们来说显得如此并不奇怪。因为诗人以不同于真理的方式叙述此事；虽然他们比历史学家、演说家和其他类型的作家古老得多，但因他们不知道神圣奥秘的秘密，且未来的复活的传闻以模糊的形式传到他们那里，他们漫不经心地轻轻听到后，以虚构故事的方式传下来。然而，他们也证明自己并非遵循可靠的权威，而只是意见，如\Person{马罗}所说，\par

\Quote{耳朵听到的，让舌头传扬。}\par

因此，虽然他们部分地败坏了真理的奥秘，但事情本身却显得更为真实，因为它部分地与先知们相符：这足以作为我们对此事的证据。然而，他们的错误中也包含了一些道理。因为当先知们不断宣告天主之子将要审判死者时，这一宣告并未逃脱他们的注意；既然他们认为天上没有别的统治者，只有\Person{朱庇特}，于是他们便宣称\Person{朱庇特}之子是冥界的君王，但不是\Person{阿波罗}、\Person{利伯尔}或\Person{墨丘利}（这些被认为是天上的神祇），而是一位既是有死的又是公正的人物，或者是\Person{米诺斯}、\Person{埃阿科斯}，或者是\Person{拉达曼提斯}。因此，他们以诗人的自由曲解了他们所领受的；或者，由于意见通过不同的口舌和各种言谈散布开来，改变了真理。因为他们预言，在冥界度过一千年之后，他们将再次获得生命，正如\Person{维吉尔}所说的：——\par

\Quote{所有这些人，当十个世纪的时光轮转，\newline{}命运之轮已滚过，\newline{}神的声音从远方远处\newline{}召唤到忘川河畔，\newline{}以便他们可以再次返回人间，\newline{}不再记起从前的事，\newline{}并带着盲目的倾向渴望\newline{}回到有肉身的身体：}\par

但这一点却逃脱了他们的注意，即死者复活不是在死后一千年，而是当他们再次获得生命后，要与天主一同统治一千年。因为天主将要来临，清除世界一切污秽，使义人的灵魂归回他们更新的身体，并提升他们至永福。因此，其他事情都是真实的，除了那遗忘之河——他们虚构此河正是为了无人能提出异议：为什么他们不记得自己曾经活过，或是谁，或是完成了什么事？但无论如何，此事被认为不可信，整件事被拒绝，仿佛是无拘无束、凭空杜撰的。但当我们坚持复活教义，并教导灵魂将回到另一个生命，不是忘记自我，而是拥有同样的知觉和形貌时，我们遇到这样的反对：这许多时代已经过去，有哪一个个体曾从死里复活，使我们能通过他的例子相信这是可能的？但复活不可能在邪恶仍然盛行时发生。因为在这个世界上，人们被暴力、刀剑、暗算、毒药所杀害，遭受伤害、贫困、监禁、酷刑和流放。此外，正义被憎恨，所有愿意跟随天主的人不仅被憎恨，而且遭受各种侮辱，被多种惩罚折磨，被驱赶去敬拜人手所造的不敬虔之神，不是凭借理性或真理，而是通过可怕的身体折磨。\par

那么，人应当再次复活到同样的事物中吗？或者回到一个无法安全生活的生命中？既然义人如此被轻视，如此容易被带走，我们可以设想，如果有任何人从死者中返回，通过恢复他原来的状态而重获生命，会发生什么呢？他必定会从人的眼前被带走，免得他被看见或听见，所有人都一齐离开诸神，转而敬拜和信奉独一真天主。因此，必须在邪恶被除去时仅发生一次复活，因为那些复活的人既不再死亡，也不再受任何伤害，这样他们才能度过幸福的一生，因为死亡已被废止。但是诗人们知道今世充满一切邪恶，于是引入了遗忘之河，免得灵魂记住他们的劳苦和邪恶，拒绝返回上界；因此\Person{维吉尔}说：——\par

\Quote{啊，父亲！思想能否设想，\newline{}这些幸福的灵魂会离开这领域，\newline{}去寻求上方的天空，\newline{}与迟钝的泥土重新结合？\newline{}这可怕的对光明的渴望，\newline{}从何而来，请说，为何？}\par

因为他们不知道它何时或必须如何发生；因此他们以为灵魂是重新出生的，并且重新回到子宫，回到婴儿状态。因此，\Person{柏拉图}在讨论灵魂的本性时也说，由此可知灵魂是不朽和神圣的，因为孩童的心智是可塑的，易于感知，而且他们如此迅速地理解所学的科目，以至于他们似乎不是第一次学习，而是重新回忆和记起它们：在这件事上，这位智者最愚蠢地相信了诗人们。\par

% structure: p0133 source=h2 target=subsection reason=relative-heading-rank; base=h2

\subsection{第二十三章 论灵魂的复活及此事实的证据}

因此，他们不会重生（这是不可能的），但会复活，并由天主穿上身体，记得他们从前的生活和一切行为；进入天上财宝的拥有，享受无数资源的愉悦，他们将立即在天主面前感谢天主，因为祂摧毁了一切邪恶，并将他们提升到祂的国度和永生。关于这个复活，哲学家们也试图像诗人一样败坏地谈论。因为\Person{毕达哥拉斯}声称灵魂进入新的身体；但愚蠢的是，它们从人进入牲畜，从牲畜进入人；并且他本人是从\Person{欧福尔布斯}恢复而来的。\Person{克里西普}说得更好，\Person{西塞罗}说他支撑着斯多葛派的柱廊，他在关于天主眷顾所写的书中，当谈到世界的更新时，引入了这些话：\Quote{但既然是这样，显然没有不可能的事，并且我们死后，当某些时间周期再次来临时，我们恢复到我们现在所处的状态。} 但是让我们从人间事物回到神圣的事物。\Person{西彼拉}这样说：——\par

\Quote{因为整个人类难以相信；但当世界和凡人的审判即将来临时，即天主亲自设立审判，同时审判不敬虔的人和圣善的人，那时他终将把恶人投入火中的黑暗。但所有圣善的人将在地上复活，天主同时赐予他们精神、尊荣和生命。}\par

但不仅先知，甚至吟游诗人、诗人和哲学家都同意将有死人的复活，那么谁也不要问我们这如何可能：因为不能为天主的工作指派任何理由；但如果起初天主以一种不可言喻的方式塑造了人，我们就可以相信那位造了新人的天主能够恢复旧人。\par

% structure: p0137 source=h2 target=subsection reason=relative-heading-rank; base=h2

\subsection{第24章 论更新的世界。}

现在我将续述其余的事。因此，至高全能天主的圣子要来审判活人死人，正如西彼拉所见证并说：—\par

\Quote{因为那时全地必有凡人的混乱，当全能者亲自来到他的审判座上，审判活人死人的灵魂和全世界。}\par

但他，当他摧毁了不义，执行了他的大审判，并召回了从起初就活着的义人之后，将与人同住一千年，以最公正的命令统治他们。这是西彼拉在另一处，当她发出受默感的预言时宣告的：—\par

\Quote{听啊，你们凡人，一位永生的君王统治着。}\par

那时，仍活在身体里的人将不会死亡，反而在那千年中要生出无数的后代，他们的后裔将是圣洁的，为天主所爱；但从死里复活的人将作为审判官管理活人。但万民不会完全消灭，有些将被留下作为天主的胜利品，成为义人得胜的缘由，并受制于永久的奴役。大约同时，万恶的创造者，魔鬼的首领，将被锁链捆绑，在千年属天统治期间被囚禁，那时正义要在世上掌权，使他不能对天主的子民行恶。在他来临之后，义人将从全地聚集，审判完成后，圣城将建立在大地中央，天主亲自作为建造者要与义人同住，并在其中统治。西彼拉指出这座城时说道：—\par

\Quote{天主所造的那城，他使它比星辰、太阳和月亮更加灿烂。}\par

那时，那笼罩并遮蔽天空的黑暗将从世界除去，月亮将接受太阳的光辉，不再亏损；太阳将比现在明亮七倍；大地将敞开它的丰产，自行结出最丰盛的果实；岩石山将滴下蜂蜜；酒河奔流，奶川涌流；总之，世界本身要欢喜，所有自然要欢跃，从邪恶、不敬、罪愆和错误的统治中被解救和释放。在此期间，野兽不再靠血喂养，飞鸟不再靠猎物存活；一切都要和平安宁。狮子和牛犊一同站在马槽旁，狼不叼羊，猎犬不追猎物；鹰和雕不伤害；婴儿与蛇嬉戏。总之，那时将要发生诗人们所说的在萨图尔努斯统治下发生的事。他们的错误源于此——先知们把许多未来的事当作已经完成的来提出和讲述。因为异象由圣神带到他们眼前，他们看见这些事，仿佛就在他们眼前完成和实现。当他们的预言逐渐传开，那些未受宗教奥义教导的人不知道为什么说这些话，就以为所有这些事在远古时代已经实现，而这些事显然不能在人的统治下完成和实现。但当不敬神的宗教被摧毁、罪孽被压制后，大地将服从天主——\par

\Quote{水手自己也要放弃大海，松木船不再\newline{}交换货物；所有土地将生产一切。\newline{}大地不再忍受耙犁，葡萄园不再忍受修剪钩；\newline{}强壮的农夫也要从轭下松开公牛。\newline{}平原将渐渐被柔软的麦穗染黄，\newline{}未修剪的荆棘上挂着绯红的葡萄，\newline{}坚硬的橡树滴下带露的蜂蜜。\newline{}羊毛也不再学会仿制各种颜色；\newline{}但公羊自己在草地上将改变它的毛色，\newline{}时而为甜美的绯红，时而为藏红花染料；\newline{}猩红色自行覆盖吃草的羔羊。\newline{}山羊自己带着膨胀的乳房回家；\newline{}畜群不再惧怕巨大的狮子。}\par

这些事是诗人根据库迈西彼拉的诗歌预言的。但厄立特里亚西彼拉这样说：—\par

\Quote{但狼在山上不与羔羊争斗，山猫与山羊羔同吃草；野猪与牛犊同食，且与一切羊群同食；食肉的狮子在马槽里吃糠，蛇与失母的婴儿同睡。}\par

在另一处，谈到万物的丰饶时：—\par

\Quote{那时天主将赐予人类极大的喜乐；因为大地、树木和无数的地上羊群，将向人提供真正的葡萄果实、甜蜂蜜、白奶和谷物，这是对凡人最好的东西。}\par

另一位同样说：—\par

\Quote{唯有敬虔者的圣地要出产这一切，从岩石和泉源涌出蜜的河流，以及神享的奶将流向所有义人。}\par

因此，人们将过着最宁静的生活，资源丰足，并与天主一同为王；列国的君王要从地极带着礼物和供物前来，朝拜和尊崇那位伟大的君王，祂的名将在天下万国和世上统治的列王中受到赞颂和敬拜。\par

% structure: p0153 source=h2 target=subsection reason=relative-heading-rank; base=h2

\subsection{第25章 论末后时期，及罗马城}

这些是先知所说关于日后将要发生的事。但我认为没有必要引述他们的证言和言辞，因为那将是无穷的工作；我的书卷也无法容纳如此众多的话题，因为有如此多的人异口同声说着类似的事。同时，也是为了避免读者因我堆积从各处收集和转述的内容而感到厌倦。此外，我这样做是为了证实我所说的一切，不是用我自己的著作，而是特别借助他人的著作，并表明不仅在我们中间，即使在那些诋毁我们的人中间，真理也被保存着，只是他们拒绝承认。但若有人想更准确地知道这些事，可以从源头汲取，他将知道比我在这几卷书中包含的更令人惊叹的事。也许现在有人会问，我们所说的这些事何时会发生？我上面已经指出，当六千年完成时，这个改变必定发生，而最终结局的最后一日如今已临近。我们被允许知道关于先知所说的记号，因为他们预言了记号，我们可据此逐日期待时代的终结，并心存敬畏。然而，当这个数字完成时，那些关于时代著作的人，从圣经和各类历史中收集资料，教导我们从世界开始以来年数的多少。虽然他们意见不一，他们所计算的数目也相差甚大，但所有预期都不超过二百年。事情本身表明，世界的倾覆和毁灭即将发生；然而，只要罗马城还存在，似乎就无需惧怕这类事。但当那世界首都倾覆，开始变为一条街——正如西彼拉所说将要发生的事——谁还能怀疑人类和整个世界的结局已经到来呢？正是那座城，唯独那座城，仍支撑着一切。我们必须祈求天上的天主——如果祂的安排和旨意果真能被延缓——免得在我们想之先，那可憎的暴君就来了，他必行如此大事，挖出那只眼睛，因它的毁灭，世界本身也将倾覆。现在让我们返回，陈述那时将要发生的其他事。\par

% structure: p0155 source=h2 target=subsection reason=relative-heading-rank; base=h2

\subsection{第26章 论魔鬼的被释放，及第二次和最大的审判}

我们之前曾说过，在神圣统治开始时，魔鬼的首领将被天主捆绑。但当王国的千年，即世界的七千年开始结束时，他也将被重新释放，从监牢中出来，去聚集所有那时仍在义人统治下的民族，使他们攻打圣城；从世界各地将聚集起无数民族，他们围困和包围该城。那时，天主的最后愤怒将降临于这些民族，完全消灭他们；首先，祂将极其猛烈地震动大地，因这震动，叙利亚的群山将崩裂，山丘将骤然下陷，所有城墙将倒塌；天主将使太阳停止，三日不落，并将其点燃；极大的炎热和烈火将降于敌对的、不虔敬的人民，还有硫磺雨、冰雹和火滴；他们的灵魂会因酷热而熔化，身体会被冰雹击碎，他们也会彼此用刀剑击杀。群山将被尸体填满，平原被骸骨覆盖；但天主的子民在那三日期间将隐藏在地穴中，直到天主对列国的愤怒和最后的审判结束。\par

然后，义人将从他们的藏身之处出来，发现一切都被尸体和骸骨覆盖。但所有恶人的种族将完全灭亡；世上不再有任何民族，只有天主的民族。然后，连续七年，森林将不受惊扰，也不从山上砍伐木材，但列国的武器将被焚烧；从此不再有战争，只有和平与永远的安息。但当千年完成时，世界将被天主更新，天将被卷起，地也要改变，天主将人转化为天使的样式，他们白如雪；他们将永远在全能者的眼前事奉，向他们的主献祭，永远服事祂。同时，将发生第二次公开的全体复活，其中不义的人将被复活承受永罚。这些人曾崇拜他们手所作的，或对世界的主宰和父亲一无所知，或否认祂。但他们的主连同他的仆役将被捕并判罚，所有恶人的队伍也将按他们的行为，在天使和义人面前被永火焚烧，永远受刑。\par

这是圣先知们的教义，我们基督徒所追随的；这是我们的智慧，那些崇拜脆弱之物或坚持空洞哲学的人，嘲笑它为愚昧和虚妄，因为我们不习惯于公开辩护和主张它，因为天主命令我们在安静和沉默中隐藏祂的奥迹，并将其保存在我们自己的良心中；不要固执地与那些对真理无知、并非为学习而是为指责和嘲笑而猛烈攻击天主及其宗教的人争辩。因为奥迹应当被最忠实地隐藏和遮盖，尤其是由我们这些拥有信德之名的人。但他们指责我们的沉默，仿佛这是出于邪恶的良心；因此他们也捏造一些关于圣洁无辜之人的可憎之事，并欣然相信他们自己的捏造。\par

\EditorialNote{献给君士坦丁的致辞在某些抄本和版本中缺失，但米涅根据一些重要手稿将其插入正文，符合拉克唐修的风格和精神。}\par

但是，最圣洁的皇帝啊，自从伟大的天主为了重建正义之家和保护人类而兴起你以来，所有虚构之事都已沉寂；因为当你统治罗马国家时，我们这些崇拜天主的人不再被视为可咒骂和不敬神的。既然真理现在从隐晦中出现并被带到光中，我们这些努力实践正义之工的人不再被指责为不义。没有人再以天主之名责备我们。我们这些在所有人群中唯一虔诚的人，不再被称为不虔诚；因为我们藐视死者的偶像，崇拜永生的真天主。至高天主的眷顾已提升你到皇帝尊严，使你能够以真正的虔诚废除他人的有害法令，纠正错误，以父亲的仁慈为人类的安全提供保障——简言之，从国家中除掉恶人，他们被卓越的虔诚所击倒，天主将他们交在你手中，以便向所有人显明真正的威严何在。\par

因为那些想要废除对至高无上天主的崇拜以便捍卫不敬神的迷信的人，已倒在废墟中。但你，你捍卫并热爱祂的名，因德行和繁荣而卓越，以最大的幸福享受你不朽的荣耀。他们遭受并已遭受他们罪过的惩罚。天主强有力的右手保护你免受一切危险；祂赐给你一个安宁和平的统治，伴随着所有人的最高祝贺。世界的主和统治者并非毫无理由地选择你超过众人，借着你祂可以更新祂神圣的宗教，因为在所有人中只有你存在，可以作为德行和圣洁的卓越榜样；在其中你不仅可以等同，而且可以超越古代君王的荣耀，这是一个非常伟大的事，虽然名声将他们列为善人。他们确实或许只是凭本性相似于义人。因为那不认识宇宙统治者的天主的人，可以达到义人的相似，但不能达到义人本身。但你，既因你性格中固有的圣洁，又因你在每一行动中对真理和天主的承认，完全实践了正义之工。因此，在安排人类状况时，天主理应运用你的权威和服务。我们每日祈祷，恳求祂特别保护你——祂愿成为世界的守护者；然后求祂赐予你一种性情，使你能够永远继续热爱天主的圣名。因为这对所有人都是有益的：对你而言是幸福，对他人而言是安宁。\par

% structure: p0162 source=h2 target=subsection reason=relative-heading-rank; base=h2

\subsection{对虔诚者的鼓励和确认}

既然我们已经完成了所承担工作的七道课程，并且已朝着目标前进，剩下的是要劝勉所有人将智慧与真正的宗教一同领受，其力量和职责在于：藐视世俗之物，放下我们从前服务脆弱之物、渴望脆弱之物时所持的错误，使我们能转向天上宝藏的永恒赏报。为了获得这些，必须尽快放下今生诱人的快乐，它们以有害的甜蜜抚慰人的灵魂。从这些尘世的污秽中被抽离，去到那位最公义的审判者与慈父面前，祂以辛劳换安息，以死亡换生命，以黑暗换光明，以短暂属世之物换永恒属天之物——我们为成就正义之工而在这个世界忍受的艰苦与苦难，与这样的赏报根本无法相比或相提并论。因此，如果我们希望成为有智慧和幸福的人，不仅必须深思并牢记\Person{泰伦提乌斯}的那些话：\par

\Quote{我们必须永远在磨坊里碾磨，我们必须被鞭打，被戴上镣铐；}\par

而且必须忍受比这些更可怕的事，即监牢、锁链和酷刑：必须经历痛苦，总之，必须接受并忍受死亡本身，当我们的良心清楚那暂时的快乐不会没有惩罚，而德行不会没有神圣的赏报时。因此，所有人都应努力尽快使自己转向正路，或者，在领受并实践德行、耐心完成今生的劳苦之后，配得天主作为他们的安慰者。因为我们的父和主，祂建造并坚固了天，在其中安置了太阳及其他天体，以祂的权能称量大地，以群山围绕它，以海洋环绕它，以河流划分它，并从无中创造并完成了这世界工造中一切存在之物；祂观察了人类的错误，派遣了一位引导者，为我们开启正义之路：让我们都跟随祂，听从祂，以最大的忠诚服从祂，因为唯独祂，如\Person{卢克莱修}所说：\par

\Quote{以说真话的训诫洁净了人的胸怀，为欲望和恐惧设定了界限，解释了什么是我们所有人努力要达到的至善，并指明了道路，沿着一条窄路，我们可径直抵达那里。}\par

祂不但指出了这条路，而且亲自走在我们前面，好叫没有人因德行之艰难而畏惧。愿那布满欺骗的毁灭之路，若有可能，被人弃绝，因为死亡藏在其中，被快乐的诱惑所掩盖。\par

当一个人的年岁渐趋衰老，越来越近地看到那必须离开此生的日子时，就应思忖如何纯洁地离开，如何清白地去见审判者；而不是像那些人一样，他们的黑暗心智被剥夺了光明，当身体的力气衰竭时，反被提醒这最后的迫切处境，应当更热切而炽烈地投身于满足自己的私欲。趁现在还来得及，机会尚存，愿每个人都从这个深渊中挣脱出来，以全心转向天主，好能无忧无虑地期盼那一天，那时天主宰、统治万物之主将审判各人的言行和思想。凡此世所贪求之物，他不但要轻视，更要避开它们，并认为他的灵魂比那些骗人的财富更为宝贵，因为这些财富的拥有是不确定且短暂的；它们日渐消逝，离去远较进入时更为迅速，即使我们能享受它们直到最后一刻，也必然要留给他人。除了善度纯洁的一生外，我们什么也不能带走。那人将带着丰厚的资财出现在天主面前，那人将富足地出现，因他将拥有克己、慈悲、忍耐、爱德和信德。这便是我们的产业，既不能从任何人手中夺去，也不能转让给他人。又有谁不愿为自己预备和获取这些美善呢？\par

让饥饿的人来，以天上的食粮为食，好戒除长久的饥饿；让口渴的人来，以便从永流不竭的泉源中满口汲取救恩之水。藉此天主的饮食，盲者将看见，聋者将听见，哑者将说话，跛者将行走，愚者将变得智慧，病者将得强壮，死者将复活。因为凡以德性践踏世上腐败的人，那至高而真实的仲裁者将使他复活，进入永恒的光明。不要有人倚靠财富，不要有人倚靠权杖，甚至不要倚靠王权：这些都不能使人不朽。因为凡抛弃合乎人性的行为，追逐现世之物，匍匐于地的人，必受惩罚，如同背离其主、其统帅、其圣父的逃兵。因此，让我们致力于正义，它能如不可分离的同伴，唯独引领我们到天主那里；并且\Quote{当灵魂统治这四肢时}，让我们以不懈的服事服事天主，坚守岗位与警醒，勇敢地与我们所认识的敌人交战，好使我们战胜仇敌，凯旋得胜，从主那里获得祂亲口应许的勇毅之赏报。\par

\SourceRecord{Translated by William Fletcher.；Ante-Nicene Fathers；Vol. 7.；Edited by Alexander Roberts, James Donaldson, and A. Cleveland Coxe.；Buffalo, NY: Christian Literature Publishing Co.,；1886.；<http://www.newadvent.org/fathers/07017.htm>.}
